\documentclass{report}
\usepackage[a4paper,left=3.5cm,right=3cm,top=2.5cm,bottom=2.5cm]{geometry}


\usepackage[spanish]{babel} %seleccionar español como idioma principal
\usepackage{graphicx} % Required for inserting images
\usepackage[utf8 ]{inputenc} %español
\usepackage[hidelinks]{hyperref} %ocultar colores de enlaces si se prefiere
\usepackage[style=ieee]{biblatex}
\addbibresource{referencias.bib} % Asegúrate de que este path es correcto
\usepackage{amssymb}
\usepackage{amsmath}
\usepackage{bm}
\usepackage{listings}
\usepackage{minted}
\usepackage{tikz} % Required for TikZ diagrams
\usepackage{float}
\usepackage{csquotes}
\usepackage[dvipsnames,svgnames,table]{xcolor}
\usepackage{lscape}
\usepackage{booktabs}
\usepackage{tabularx}
\usepackage{ragged2e}
\usepackage{adjustbox}
\usepackage{siunitx}
\usepackage{array, booktabs, tabularx}
\usepackage{multirow}
\usepackage{subcaption}
\usepackage[automake]{glossaries} % para el glosario con automake
\makeglossaries % inicializar el glosario
\usepackage[spanish]{babel}
\usepackage{svg}
\usepackage{pgfgantt}

\usetikzlibrary{arrows.meta, positioning, shapes, calc, babel}

\usepackage{threeparttable}% para notas dentro de la tabla
\sisetup{
  detect-mode,
  table-number-alignment = center,
}


%\title{Desarrollo de una Aplicación Móvil para diaDiabéticos} %% cambiar mas adelante
\lstset{
        tabsize=2, % tab = 2 espacios
        backgroundcolor=\color[HTML]{F0F0F0}, % color de fondo
        captionpos=b, % posición de pie de código, b=debajo
        basicstyle=\ttfamily, % estilo de letra general
        columns=fixed, % columnas alineadas
        extendedchars=true, % ASCII extendido
        breaklines=true, % partir líneas
        prebreak = \raisebox{0ex}[0ex][0ex]{\ensuremath{\hookleftarrow}}, % marcar final de línea con flecha
        showtabs=false, % no marcar tabulación
        showspaces=false, % no marcar espacios
        keywordstyle=\bfseries\color[HTML]{007020}, % estilo de palabras clave
        commentstyle=\itshape\color[HTML]{60A0B0}, % estilo de comentarios
        stringstyle=\color[HTML]{4070A0}, % estilo de strings
}
\author{alumno  Adriano García-Giralda Milena\\ \\Director1:María José Rodríguez Fortiz \\ Director2:María del Campo Bermúdez Edo}
\date{January 2026}

\begin{document}

% Incluye el contenido de prefacio.tex aquí
\input{portada/prefacio} 


%&\maketitle

\tableofcontents %indice de contenidos automatico


\listoffigures %indice de figuras automatico
% 3. Definir cada entrada del glosario

\input{glosario}

%%%%%%%%%%%%%%%%%%%%%%%%%%%%%%%%%%%%%%%%%%%%%%%%%%

\chapter{Introducción}
\label{chap:introduccion}
% Breve descripción del problema y la solución propuesta (la app móvil).

En este capítulo se presentan el contexto, la motivación y los objetivos consecuentes del proyecto. 

%Es importante citar ideas que no sean propias con \cite{sample1} y todas las referencias ponerlas en formato .bib en "referencias.bib"
%\cite{sample1}

\section{Contexto}
En las últimas décadas, la proliferación de dispositivos conectados (sensores, actuadores, sistemas embebidos, \textit{wearables}, smartphones, \textit{gateways}, etc.) ha transformado de forma radical la interacción con nuestro entorno. Este fenómeno, conocido como Internet de las Cosas (\gls{iot}), puede entenderse como la interconexión de \emph{``cosas''} físicas dotadas de capacidades de sensado, comunicación y, en muchos casos, actuación, cuyo propósito último es \textbf{convertir el mundo real en datos} para \textbf{monitorizar, modelar y extraer inferencias} sobre procesos físicos, biológicos o sociales (y, cuando procede, cerrar el bucle con acciones de control) \cite{atzori2010iot,gubbi2013iot,choudhary2024iotoverview}. 
En otras palabras, el valor del IoT no reside únicamente en ``conectar dispositivos'', sino en \textbf{capturar señales del entorno y analizarlas} para comprender cómo se comporta un sistema real, detectar eventos, anticipar estados o tomar decisiones basadas en evidencia.

Este ecosistema ha experimentado un crecimiento sostenido y a gran escala. Según el informe más reciente de IoT Analytics, el número de dispositivos IoT conectados alcanzó los \textbf{18,5\,mil millones} a finales de 2024, y se estima que crecerá un \textbf{14\,\%} durante 2025 hasta superar los \textbf{21,1\,mil millones} de conexiones. Las previsiones a largo plazo sitúan esta cifra en torno a los \textbf{39\,mil millones de dispositivos conectados en 2030} \cite{iotanalyticsconnecteddevices2025}. 



Este crecimiento plantea un reto fundamental: \textbf{cómo extraer conocimiento útil a partir de datos generados de forma distribuida, sin centralizar información sensible ni sobrecargar la infraestructura}. Sin embargo, en IoT muchos dispositivos finales (sensores, wearables o móviles) presentan limitaciones de CPU, memoria y energía, además de conectividad intermitente \cite{shi2016edge}. Por ello, durante años el enfoque dominante ha sido \textbf{centralizar} el procesamiento en servidores remotos, trasladando los datos hacia la nube.

En aplicaciones con datos personales o biométricos, esta centralización conlleva inconvenientes relevantes: incrementa el tráfico de red y la latencia y, además, introduce riesgos de privacidad al transferir y concentrar información sensible. Estas limitaciones motivan arquitecturas donde el procesamiento se distribuye más cerca del origen del dato y donde el aprendizaje se realiza sin necesidad de mover información cruda.

Para reducir latencia, ancho de banda y dependencia de la conectividad, han surgido paradigmas que desplazan parte del procesamiento hacia el borde (\textit{edge}) y niveles intermedios (\textit{fog}), conformando un continuo \textit{edge--fog--cloud} \cite{satyanarayanan2017emergence,naha2018fog}. En este marco, algunas tareas pueden ejecutarse cerca del dispositivo, mientras que otras se delegan a nodos con mayor capacidad, buscando un equilibrio entre rendimiento, disponibilidad de recursos y privacidad.

En este contexto surge el \textbf{Aprendizaje Federado} (\gls{fl}), un paradigma de entrenamiento distribuido en el que \textbf{los datos permanecen en el dispositivo} y lo que se comparte son \textbf{actualizaciones del modelo} \cite{mcmahan2017communication,hard2018federated}. Su idea central puede resumirse de forma sencilla: \textit{en lugar de mover los datos hacia el modelo, se traslada el modelo hacia los datos}. Además en el contexto europeo y español, esta federación viene determinada principalmente por el \gls{gdpr}\cite{gdpr2016} y la \gls{lopdgdd}\cite{lopdgdd2018}, que establecen obligaciones estrictas para el tratamiento de datos personales y, en particular, de categorías sensibles como datos biométricos y de salud. En consecuencia, resulta especialmente relevante adoptar arquitecturas que reduzcan la transferencia y centralización de datos, favoreciendo enfoques donde el procesamiento se realice lo más cerca posible del origen.

Respecto al caso de estudio, vamos a escoger la detección y evaluación de estrés, específicamente en personas mayores. La literatura científica confirma que el estrés no es un problema menor: está sólidamente asociado con un aumento del riesgo de enfermedad coronaria incidente de aproximadamente un 27\% en individuos con altos niveles de estrés percibido frente a aquellos con niveles bajos \cite{Richardson2012}. De forma similar, en adultos con enfermedades crónicas, la depresión y el estrés reducen de manera significativa su calidad de vida relacionada con la salud, con coeficientes de regresión de hasta $\beta \approx -0.24$ en el caso de la depresión, lo que refleja un impacto clínicamente relevante \cite{Hwang2024}. En personas mayores con responsabilidades intensas de cuidado, estudios recientes han evidenciado que estos cuidadores presentan un 53\% más de probabilidades de desarrollar síntomas depresivos y un 26\% más de riesgo de deterioro funcional en comparación con personas mayores no cuidadoras \cite{Litzelman2021}. Además, se ha observado un incremento del 63\% en el riesgo de mortalidad a lo largo de seis años entre cuidadores mayores que reportan niveles altos de sobrecarga emocional \cite{Thomas2018}.

Además, estudios recientes en grandes muestras de población de edad avanzada muestran que hay experiencias que afectan a la salud mental, como la pérdida del cónyuge, que se asocia con un aumento del 69\% en el riesgo de mortalidad \cite{Nilaweera2023} y con un incremento del 72\% en la probabilidad de desarrollar demencia, mientras que las dificultades financieras graves elevan este riesgo en un 53\% \cite{Nilaweera2023b}.  



\section{Motivación}

En cuanto a la motivación técnica, de los retos detectados se derivan desafíos estrechamente relacionados: \textbf{(i) entrenar modelos a partir de múltiples fuentes distribuidas}, \textbf{(ii) decidir dónde ejecutar el procesamiento a lo largo del continuo \textit{edge--fog--cloud}}, y \textbf{(iii) coordinar la agregación del aprendizaje de forma eficiente en entornos IoT}. Este Trabajo Fin de Máster se sitúa en la intersección de estos problemas, proponiendo el \textbf{aprendizaje federado (y, en particular, su variante jerárquica)} como eje principal de la arquitectura desarrollada para \textbf{analizar datos biométricos orientados a la detección de estrés en personas mayores}.

Por lo tanto, este Trabajo Fin de Máster se sustenta en una motivación técnica clara: \textbf{el análisis de datos biométricos mediante técnicas modernas de aprendizaje automático requiere cómputo y coordinación}, pero en escenarios IoT dicho procesamiento no puede depender exclusivamente de arquitecturas centralizadas. Sin embargo, pasar ``del cloud al edge'' (o al \textit{computing continuum}) no resuelve automáticamente el problema: aparecen retos abiertos como \textbf{cómo distribuir el entrenamiento de modelos}, \textbf{cómo coordinar la agregación de conocimiento entre nodos heterogéneos}, \textbf{cómo gestionar recursos limitados y conectividad intermitente}, y \textbf{cómo diseñar sistemas escalables y reproducibles} que puedan evaluarse con garantías en entornos realistas. Es precisamente en este punto donde se sitúa la aportación principal de este TFM.

Este trabajo propone y valida una \textbf{arquitectura modular jerárquica} para el entrenamiento distribuido de modelos de aprendizaje automático dentro del continuo \textit{edge--fog--cloud}. La idea central es \textbf{organizar el aprendizaje por niveles}, de modo que los dispositivos cercanos al usuario contribuyan al entrenamiento con datos locales, mientras que nodos intermedios coordinan y agregan el aprendizaje antes de integrarlo a nivel global. Este enfoque permite aproximar el modelo al usuario, reducir la necesidad de transmitir datos sensibles y aprovechar el cómputo distribuido disponible, manteniendo al mismo tiempo un esquema de coordinación compatible con entornos IoT.


En cuanto a la motivación personal, tras la realización de mi TFG disfruté enormemente de la experiencia de investigar y desarrollar un proyecto propio. Tanto fue así que lo comenté con mi tutor, el profesor Francisco García Moreno, quien me aconsejó ponerme en contacto con el grupo de investigación MYDASS (Modelling \& Development of Advanced Software Systems) de la Universidad de Granada, perteneciente al Departamento de Lenguajes y Sistemas Informáticos, del cual él forma parte. De esta manera, tendría la oportunidad de experimentar de primera mano cómo se trabaja en el ámbito de la investigación y comprobar si realmente me interesa continuar mi formación con un futuro doctorado.

Como consecuencia de ello, me matriculé en este Máster y conocí a las que hoy son mis tutoras, María Bermúdez Edo y María José Rodríguez Fortiz, quienes me introdujeron en profundidad en el proyecto nacional que están liderando en la actualidad: PID2023-149185OB-I00: “Una plataforma de salud con servicios de computación continua para detectar el estrés y la ansiedad en las personas mayores”, concedido en la convocatoria de Proyectos de Generación de Conocimiento 2023.

Este proyecto, que lleva como nombre OnTheEdge(juego de palabras por llevar a una persona al límite y el tipo de computación que se busca investigar), intenta crear una plataforma de Salud y computación continua para medir el estrés de personas mayores en el computing continuum.

A nivel personal, fue un proyecto que desde el primer momento suscitó en mí un gran interés. No solo por la posibilidad de aportar mi granito de arena a un problema de gran relevancia social, sino también por la oportunidad de adentrarme en un campo de investigación innovador y con un impacto directo en la calidad de vida de las personas. Explorar nuevas formas de aplicar el aprendizaje automático para mejorar el paradigma clásico (centralizado y costoso) fue lo que inicialmente despertó mi interés. Pero la posibilidad de contribuir a una iniciativa con impacto social real, que combina el rigor científico, la innovación tecnológica y el compromiso ético con el bienestar de las personas mayores, fue lo que finalmente me hizo decantarme por este Trabajo de Fin de Máster.



\section{Objetivos del proyecto}
\label{Objetivos}
% Descripción de los objetivos principales y secundarios del TFG.
\subsection{Objetivo general}
El objetivo principal de este trabajo es \textbf{diseñar, desarrollar e implementar un sistema de aprendizaje federado} basado en una arquitectura distribuida conforme al paradigma del \textit{Computing Continuum}, incorporando una capa de abstracción en el nivel \textit{fog} que permita superar el modelo tradicional cliente-servidor.  
El sistema servirá como plataforma experimental para explorar la escalabilidad, la privacidad y la interoperabilidad del aprendizaje federado en entornos \textit{edge–fog–cloud}.

\subsection{Objetivos específicos}

Para alcanzar este propósito general, se establecen los siguientes objetivos específicos:

\begin{enumerate}
    \item \textbf{Revisar y analizar el estado del arte} en materia de aprendizaje federado, \textit{computing continuum} y detección de estrés mediante sensores portátiles, identificando los principales avances, limitaciones y oportunidades de integración entre estos campos.
    
    \item \textbf{Profundizar en el conocimiento y manejo de las tecnologías involucradas}, incluyendo frameworks de aprendizaje federado (como \textit{Flower} o \textit{TensorFlow Federated}), contenedores y orquestadores (\textit{Docker}, \textit{Docker Compose}) y sistemas de mensajería para coordinación distribuida (por ejemplo, \textit{MQTT}).
    
    \item \textbf{Diseñar e implementar una arquitectura modular jerárquica} que incorpore una capa de agregación intermedia en el nivel \textit{fog}, facilitando la comunicación, la agregación parcial de modelos y la escalabilidad del sistema a múltiples clientes y agregadores.
    
    \item \textbf{Desarrollar un prototipo funcional de un sistema federado} capaz de soportar distintas técnicas de \textit{machine learning}, tanto clásicas como basadas en redes neuronales, permitiendo realizar experimentos reproducibles y controlados bajo diferentes esquemas de partición de datos (horizontal y por transferencia).
    
    \item \textbf{Integrar mecanismos de registro y análisis de métricas experimentales}, que posibiliten la comparación de resultados, la trazabilidad de los entrenamientos y la validación de las estrategias de agregación empleadas.
    
    \item \textbf{Evaluar la escalabilidad, eficiencia y privacidad del sistema} mediante pruebas controladas con múltiples nodos distribuidos, verificando su comportamiento al incrementar el número de clientes y agregadores.

    \item \textbf{Evaluar experimentalmente el impacto de la incorporación de nodos Fog} en el comportamiento global del sistema federado, analizando cómo afectan a: la latencia, la escalabilidad, la convergencia del modelo, la robustez frente a fallos y la complejidad operativa del sistema.
    
\item \textbf{Aplicar y evaluar el sistema en un caso de uso representativo}, en nuestro caso centrado en la detección de estrés a partir de señales fisiológicas y contextuales, con el fin de validar la idoneidad y robustez del sistema para el tratamiento de datos sensibles en entornos reales.

\end{enumerate}


\section{Estructura de la Memoria}
La presente memoria se organiza en una serie de capítulos diseñados para guiar al lector de forma progresiva desde los fundamentos teóricos del proyecto hasta su implementación práctica y los resultados obtenidos. A continuación, se describe brevemente el contenido de cada uno de ellos.

\begin{itemize}

    \item \textbf{Capítulo~\ref{chap:introduccion}. Introducción}\\
    En este capítulo se contextualiza el trabajo, presentando las bases conceptuales necesarias para comprender el proyecto. Se incluyen nociones fundamentales sobre \textit{Machine Learning}, los paradigmas distribuidos modernos (Cloud, Edge y \textit{Computing Continuum}). Además, se detallan la motivación personal del proyecto y los objetivos generales y específicos que guían su desarrollo.

    \item \textbf{Capítulo~\ref{chap:estado-del-arte}. Estado del Arte}\\
    Se presenta una revisión exhaustiva de la literatura relacionada con el Aprendizaje Federado, abarcando su origen, funcionamiento, variantes y tendencias actuales. También se analiza el concepto de \textit{Computing Continuum} como evolución de los modelos de computación distribuidos. Finalmente, se estudian los trabajos existentes en detección de estrés mediante dispositivos \textit{wearables} y datasets públicos, e incluye un apartado específico de tecnologías y herramientas empleadas.

    \item \textbf{Capítulo~\ref{chap:stakeholders}. Análisis de \textit{Stakeholders}}\\
    Aquí se identifican y analizan todos los actores implicados directa o indirectamente en el proyecto, evaluando sus necesidades, intereses y el impacto que el sistema desarrollado puede tener sobre ellos.

    \item \textbf{Capítulo~\ref{chap:metodologia}. Metodología de Desarrollo y Planificación}\\
    Se expone la metodología seleccionada para organizar el trabajo, basada en un enfoque ágil Scrum. Se detalla la planificación temporal, las fases del proyecto, la estimación de recursos y costes, así como la definición de las iteraciones previstas. También se incluye el listado de tareas, requisitos e historias de usuario modeladas en el proyecto.

    \item \textbf{Capítulo~\ref{chap:desarrollo}. Desarrollo del Proyecto}\\
    Este capítulo constituye el núcleo práctico de la memoria. Se documenta el desarrollo iterativo mediante \textit{sprints} o iteraciones, incluyendo el \textit{Sprint 0} (análisis inicial, \textit{backlog} y diseño preliminar), y las sucesivas iteraciones. Para cada \textit{sprint} se documenta el trabajo obteniendo como producto final del análisis una revisión del \textit{backlog}.

    \item \textbf{Capítulo~\ref{chap:costes}. Estimación de Costes}\\
    Se estima la dedicación total del proyecto y se detallan los recursos empleados, incluyendo criterios de cálculo y equipamiento considerado en la planificación económica.

    \item \textbf{Capítulo~\ref{chap:conclusiones}. Conclusiones}\\
    Se revisan los objetivos alcanzados, se presenta la valoración personal y el aprendizaje obtenido, y se plantean líneas de trabajos futuros.

\end{itemize}


\chapter{Estado del Arte}
\label{chap:estado-del-arte}


\section{Introducción}

Este capítulo tiene como objetivo situar este Trabajo Fin de Máster en el contexto actual de la investigación en aprendizaje federado, arquitecturas distribuidas y computación en el continuo \textit{Edge--Fog--Cloud}. Para ello, se ha llevado a cabo una revisión de la literatura científica relevante, así como del uso de conjuntos de datos ampliamente aceptados en la comunidad.

La búsqueda bibliográfica se ha realizado principalmente a través de \textbf{Google Scholar} y el repositorio de preprints \textbf{arXiv}, por ser dos de las plataformas más utilizadas y reconocidas en el ámbito de la investigación en aprendizaje automático y sistemas distribuidos. Google Scholar permite acceder de forma unificada a artículos revisados por pares publicados en revistas y conferencias de alto impacto, mientras que arXiv facilita el acceso temprano a trabajos recientes y propuestas emergentes que, aunque no siempre revisadas formalmente, suelen marcar tendencias y líneas de investigación activas.

La selección de las fuentes se ha realizado atendiendo a criterios de \textbf{actualidad, relevancia científica e impacto}. En particular, se han priorizado trabajos publicados entre \textbf{2020 y 2025}, un periodo que recoge la consolidación y rápida expansión del aprendizaje federado, incluyendo avances en arquitecturas jerárquicas, eficiencia en comunicaciones, privacidad y despliegues en entornos reales. Este intervalo temporal permite abarcar tanto contribuciones fundacionales recientes como propuestas maduras que abordan retos abiertos en escenarios no centralizados.

Las fuentes proceden principalmente de \textbf{revistas indexadas y actas de conferencias internacionales de referencia}, complementadas con preprints ampliamente citados cuando aportan resultados relevantes o enfoques novedosos.

En cuanto a los conjuntos de datos considerados, se han incluido principalmente aquellos publicados desde \textbf{2010 hasta la actualidad}, motivado en gran medida por la \textbf{escasez de datasets disponibles en el ámbito del estrés}. Para áreas no centradas específicamente en aprendizaje federado o en el uso de datasets, la revisión se ha apoyado en los trabajos \textbf{más citados y reconocidos} según Google Scholar.


\section{Machine Learning}

El \textit{Machine Learning} (\gls{ml}) constituye una rama de la inteligencia artificial cuyo objetivo es desarrollar modelos capaces de aprender patrones a partir de datos. En lugar de programar explícitamente cada regla de decisión, un modelo de \gls{ml} ajusta sus parámetros internos para resolver una tarea específica utilizando ejemplos previos. Estas tareas incluyen clasificación, regresión, detección de anomalías o predicción de valores futuros. En el caso concreto de este TFM, el interés se centra en la \textbf{clasificación y/o predicción de datos biométricos} orientados a la \textbf{detección de estrés en personas mayores}.

Un modelo de \gls{ml} se entrena utilizando un conjunto de datos que contiene ejemplos del fenómeno que se desea estudiar. Cada ejemplo está formado por:
\begin{itemize}
    \item un conjunto de \textbf{características} (\textit{features}) que describen el estado o señales registradas (en nuestro caso, variables fisiológicas o contextuales);
    \item y una \textbf{etiqueta} o valor objetivo, que representa la categoría o resultado que se desea aprender (por ejemplo, niveles de estrés).
\end{itemize}

Para garantizar que el modelo aprende correctamente, generalice y no memorice los datos, es habitual dividir el conjunto de datos en tres subconjuntos bien diferenciados:
\begin{itemize}
    \item \textbf{Conjunto de entrenamiento (train):} utilizado para ajustar los parámetros del modelo mediante optimización iterativa.
    \item \textbf{Conjunto de validación (validation):} empleado durante el entrenamiento para ajustar hiperparámetros, detectar sobreajuste y seleccionar el mejor modelo.
    \item \textbf{Conjunto de prueba (test):} reservado exclusivamente para evaluar el rendimiento final y estimar la capacidad real de generalización del modelo.
\end{itemize}

Esta separación es fundamental, ya que permite medir el rendimiento del modelo en datos nunca vistos y garantiza un comportamiento robusto antes de su implantación en un sistema real. En este trabajo, este aspecto cobra especial relevancia debido a la naturaleza altamente variable de las señales fisiológicas que se utilizarán para detectar estrés y a la necesidad de detectarlo en condiciones reales de uso.

El entrenamiento consiste en ajustar los parámetros del modelo para minimizar el error entre sus predicciones y las etiquetas reales. Este ajuste se realiza habitualmente mediante un método de optimización conocido como \textbf{descenso por gradiente} (\textit{gradient descent}) \cite{lecun2015deep,goodfellow2016deep}.

De forma intuitiva, el error del modelo puede imaginarse como una superficie con valles y montañas, donde cada punto representa una posible configuración de los parámetros. El objetivo del entrenamiento es encontrar una configuración situada en un valle, donde el error sea mínimo. En este contexto, el \textbf{gradiente} indica la dirección en la que el error aumenta más rápidamente; por tanto, para reducirlo, los parámetros deben modificarse en la dirección opuesta al gradiente.

Este proceso puede expresarse de forma simplificada mediante la actualización iterativa:
\[
w \leftarrow w - \eta \, \nabla L(w),
\]
donde \(w\) representa los parámetros del modelo, \(L(w)\) es la función de pérdida que mide el error, \(\nabla L(w)\) es el gradiente de dicha función y \(\eta\) es la \textbf{tasa de aprendizaje}, que controla el tamaño de los ajustes realizados en cada iteración.

Estas actualizaciones se repiten muchas veces recorriendo el conjunto de datos en múltiples iteraciones, denominadas \textit{épocas}, hasta que el error deja de disminuir de forma apreciable. En modelos modernos —especialmente redes neuronales profundas— este entrenamiento implica realizar un gran número de operaciones matemáticas repetidas, lo que conlleva un coste elevado en términos de cómputo, memoria y consumo energético. Una vez entrenado, el modelo puede aplicarse en la fase de \textbf{inferencia}, donde genera predicciones a partir de nuevos datos.


Este flujo completo es comúnmente denominado \textbf{pipeline} y suele incluir:
\begin{enumerate}
    \item \textbf{Recopilación de datos} (por ejemplo, señales fisiológicas de usuarios).
    \item \textbf{Preprocesamiento y extracción de características}.
    \item \textbf{División de datos en train, validation y test}.
    \item \textbf{Entrenamiento} del modelo con datos etiquetados.
    \item \textbf{Validación} y ajuste de hiperparámetros.
    \item \textbf{Evaluación final} en el conjunto de test.
    \item \textbf{Inferencia en tiempo real} sobre dispositivos o servidores.
\end{enumerate}

Sin embargo, estos pipelines presentan limitaciones significativas en entornos IoT. Tanto el entrenamiento como la inferencia —especialmente con modelos modernos requieren:
\begin{itemize}
    \item una \textbf{alta capacidad de cómputo} (CPU/GPU),
    \item \textbf{memoria} suficiente para almacenar modelos y lotes de datos,
    \item y \textbf{energía} disponible en los dispositivos durante un tiempo prolongado.
\end{itemize}

 necesidades contrastan con las restricciones típicas de los dispositivos IoT, que suelen disponer de pocos recursos computacionales, baterías limitadas y conectividad variable \cite{shi2016edge}. Como consecuencia, ejecutar el pipeline completo de \gls{ml} directamente en un dispositivo pequeño, como un teléfono móvil por ejemplo, no es viable en la mayoría de los casos; tradicionalmente, se entrenan y ejecutan estos modelos sobre hardware dedicado en servidores en la nube, como veremos a continuación.
 
\section{Computing Continuum: Integración entre Cloud, Fog y Edge Computing}
\label{sec:continuum}

El paradigma del \textit{computing continuum} surge como una evolución natural de las arquitecturas distribuidas, proponiendo una integración fluida entre la \textit{computación en la nube} (\textit{cloud}), la \textit{computación en el fog} (\textit{fog}) y la \textit{computación en el edge} (\textit{edge}). Su objetivo principal es ofrecer una infraestructura unificada capaz de desplegar tareas de procesamiento y almacenamiento de forma dinámica a lo largo de las distintas capas, optimizando recursos, reduciendo la latencia y preservando la privacidad de los datos \cite{gkonis2023survey,bittencourt2025continuum,moreschini2022cloudcontinuum}.

Para entender cómo se pueden desplegar de forma eficiente sistemas distribuidos complejos, resulta conveniente introducir el concepto de \textit{computing continuum}. En las siguientes subsecciones se revisan los fundamentos de las capas \textit{cloud}, \textit{fog} y \textit{edge}, así como los objetivos que persigue su integración cooperativa. A partir de esta base, se analizan las ventajas que este paradigma ofrece para el aprendizaje federado, permitiendo situar la arquitectura propuesta en un contexto adecuado y justificar el uso de un enfoque jerárquico Edge–Fog–Cloud para el entrenamiento distribuido de modelos en aplicaciones biomédicas sensibles.

\subsection{Conceptos fundamentales}

\textbf{Cloud computing.}  
La computación en la nube constituye el modelo tradicional de provisión de recursos informáticos mediante grandes centros de datos centralizados. Ofrece elasticidad, escalabilidad y administración simplificada, permitiendo ejecutar aplicaciones complejas sin necesidad de infraestructura local \cite{mouradian2017comprehensive,naha2018fog}.

\textbf{Edge computing.}  
El \textit{edge computing} traslada el procesamiento de datos hacia la periferia de la red, es decir, más cerca del origen de los datos (sensores, cámaras, dispositivos IoT). Esto permite respuestas más rápidas, ahorro de ancho de banda y un mayor control de la acosta de un incremento de la complejidad del sistema \cite{varshney2017characterizing}.  

\textbf{Fog computing.}  
El término \textit{fog computing} fue introducido por Cisco como una extensión intermedia entre el edge y la nube. Los nodos fog (gateways, routers o servidores locales) permiten el procesamiento y almacenamiento en proximidad, actuando como una capa de agregación y orquestación entre los dispositivos finales y el \textit{cloud}. Este enfoque busca mitigar las limitaciones de latencia y congestión de red propias del modelo puramente en la nube \cite{mouradian2017comprehensive,naha2018fog}.

\subsection{Objetivos del computing continuum}

El \textit{computing continuum} no se limita a la coexistencia de estas tres capas, sino que promueve su integración cooperativa para lograr una distribución dinámica de cargas de trabajo y datos a lo largo de la infraestructura. Su propósito es que el procesamiento se realice en el punto más adecuado según las condiciones del sistema (latencia, energía, disponibilidad, privacidad o coste), garantizando una experiencia continua y transparente al usuario \cite{gkonis2023survey,bittencourt2025continuum}.  

Entre los objetivos más destacados se incluyen:
\begin{itemize}
    \item Minimizar la latencia mediante el procesamiento local o en capas intermedias.
    \item Reducir el tráfico de red hacia el \textit{cloud}, optimizando el uso del ancho de banda.
    \item Mejorar la resiliencia ante desconexiones o fallos mediante capacidades de cómputo distribuidas.
    \item Facilitar la movilidad y la interoperabilidad entre dominios heterogéneos.
    \item Preservar la privacidad al mantener los datos lo más cerca posible de su origen.
\end{itemize}

Trabajos recientes, como el de Gkonis et al. \cite{gkonis2023survey} y Bittencourt et al. \cite{bittencourt2025continuum}, enfatizan que esta integración continua representa la base de las infraestructuras digitales del futuro, donde la frontera entre nube, fog y edge tiende a desaparecer.
%\subsection{Paradigmas distribuidos: Cloud Computing , Edge Computing y Computing Continuum}


\subsection{Limitaciones del Cloud Computing centralizado en entornos IoT}

Tradicionalmente, los sistemas de \gls{ml} y de análisis de datos han seguido una arquitectura \textbf{centralizada}, donde la mayor parte del procesamiento se realiza en servidores remotos en la nube. En este modelo, los dispositivos IoT actúan principalmente como generadores de datos: capturan señales, las transmiten a través de la red y delegan las tareas de almacenamiento, entrenamiento e inferencia a servicios especializados en el \textit{cloud} \cite{shi2016edge}.

Este enfoque resulta eficaz para cargas de trabajo intensivas, dado que los centros de datos en la nube disponen de recursos computacionales abundantes, escalabilidad elástica y hardware especializado (p.\,ej., GPU/TPU). Sin embargo, cuando se aplica a entornos IoT, este modelo presenta una serie de \textbf{limitaciones estructurales} que hacen que la solución centralizada sea inviable o ineficiente. La Tabla \ref{tab:cloud_limitaciones} resume las principales limitaciones identificadas en la literatura, categorizadas por su impacto operativo, técnico y normativo.

\begin{table}[H]
\centering
\small
\caption{Limitaciones del \textit{cloud computing} centralizado en entornos IoT: análisis de impacto técnico y operativo.}
\label{tab:cloud_limitaciones}
\renewcommand{\arraystretch}{1.3}
\begin{tabular}{p{3.2cm} p{9.2cm}}
\toprule
\textbf{Limitación} & \textbf{Descripción e impacto en IoT} \\
\midrule
\textbf{Latencia} & 
La distancia a centros de datos remotos introduce retardos incompatibles con aplicaciones sensibles al tiempo, como analítica en tiempo real, sistemas biomédicos o AR/VR, donde se requieren respuestas rápidas a eventos locales \cite{shi2016edge,satyanarayanan2017emergence}. \\
\midrule
\textbf{Consumo de ancho de banda} & 
Sensores de alto volumen (vídeo, señales biomédicas, multimodalidad) generan tráfico continuo; enviar datos crudos a la nube incrementa costes operativos y puede saturar enlaces en entornos con conectividad limitada \cite{shi2016edge,naha2018fog,varghese2016challenges}. \\
\midrule
\textbf{Dependencia de conectividad} & 
La prestación del servicio queda condicionada a la disponibilidad de una red estable; fallos temporales o escenarios móviles comprometen la continuidad operativa \cite{satyanarayanan2017emergence,naha2018fog}. \\
\midrule
\textbf{Riesgos de privacidad} & 
Centralizar datos sensibles (biomédicos, de localización) aumenta la superficie de exposición, dificulta el cumplimiento normativo (GDPR, LOPDGDD) y reduce la confianza del usuario \cite{shi2016edge,varghese2016challenges}. \\
\midrule
\textbf{Coste computacional} & 
El entrenamiento de modelos de aprendizaje automático y profundo requiere infraestructura especializada costosa; el enfoque cloud-céntrico concentra esta inversión en grandes centros de datos \cite{varghese2016challenges}. \\
\midrule
\textbf{Desaprovechamiento de recursos} & 
Numerosos dispositivos y nodos cercanos (móviles, gateways, routers, estaciones base) disponen de capacidad de cómputo infrautilizada que podría reducir latencia y tráfico si se explota localmente \cite{shi2016edge}. \\
\bottomrule
\end{tabular}
\end{table}


Estas limitaciones han impulsado la búsqueda de \textbf{alternativas distribuidas} que acerquen el procesamiento a la fuente de datos. La literatura ha propuesto diversos enfoques (Edge Computing, Fog Computing, aprendizaje distribuido clásico, split learning, gossip learning) para abordar estos problemas. Sin embargo, cada enfoque presenta compromisos distintos entre latencia, privacidad, escalabilidad y complejidad operativa.

La Tabla \ref{tab:comparativa_paradigmas_fl} compara cómo diferentes paradigmas responden a las limitaciones del cloud centralizado, evidenciando que el \textbf{aprendizaje federado} ofrece la combinación más equilibrada de ventajas para entornos IoT sensibles: mantiene los datos locales sin centralización, coordina el entrenamiento de forma escalable y está diseñado explícitamente para privacidad.

\begin{table}[H]
\centering
\small
\caption{Comparativa de paradigmas distribuidos para ML en IoT frente a limitaciones del cloud centralizado.}
\label{tab:comparativa_paradigmas_fl}
\renewcommand{\arraystretch}{1.25}
\setlength{\tabcolsep}{4pt}

\begin{tabular}{p{3.3cm} p{2.1cm} p{2.1cm} p{2.1cm} p{2.1cm} p{2.1cm}}
\toprule
\textbf{Criterio} & \textbf{Cloud} & \textbf{Edge/Fog} & \textbf{DL dist.} & \textbf{Split} & \textbf{FL} \\
\midrule
\textbf{Latencia} & Alta & Baja & Media & Baja & \textbf{Baja} \\
\midrule
\textbf{Ancho de banda} & Alto & Medio & Alto & Medio & \textbf{Bajo} \\
\midrule
\textbf{Privacidad} & Baja & Media & Baja & Media & \textbf{Alta} \\
\midrule
\textbf{Escalabilidad en IoT} & Media & Media & Media & Media & \textbf{Alta} \\
\midrule
\textbf{Conectividad} & Alta & Media & Alta & Media & \textbf{Baja} \\
\midrule
\textbf{Complejidad} & Baja & Media & Alta & Alta & \textbf{Alta} \\
\bottomrule
\end{tabular}
\caption*{\footnotesize
La evaluación cualitativa presentada se basa en revisiones y trabajos de referencia sobre computación distribuida y aprendizaje federado en IoT, incluyendo \cite{shi2016edge,satyanarayanan2017emergence,naha2018fog} para \textit{cloud/edge/fog}, \cite{dean2012distbelief,li2014parameterserver} para aprendizaje distribuido clásico, \cite{vepakomma2018splitnn,gupta2018split} para \textit{split learning}, y \cite{mcmahan2017communication,kairouz2021advances} para aprendizaje federado.
}
\end{table}

La Tabla~\ref{tab:comparativa_paradigmas_fl} resume cómo distintos paradigmas de computación y aprendizaje distribuido abordan las limitaciones del \textit{cloud computing} centralizado en entornos IoT, atendiendo a criterios clave para la viabilidad operativa y regulatoria de sistemas basados en datos sensibles, como la latencia, el consumo de ancho de banda, la preservación de la privacidad, la escalabilidad en escenarios heterogéneos y la dependencia de la conectividad \cite{shi2016edge,satyanarayanan2017emergence,naha2018fog}.

Las arquitecturas \textit{edge/fog} mitigan eficazmente la latencia y el tráfico de red al acercar el procesamiento a la fuente de datos, pero no resuelven por sí solas los problemas de privacidad durante el entrenamiento de modelos \cite{shi2016edge,satyanarayanan2017emergence}. De forma similar, los enfoques de aprendizaje distribuido clásico y \textit{split learning} reducen la centralización de datos crudos, aunque mantienen dependencias de infraestructuras centrales o de intercambios de activaciones y gradientes que pueden resultar problemáticos en dominios regulados \cite{dean2012distbelief,li2014parameterserver,vepakomma2018splitnn,gupta2018split}.

En contraste, el aprendizaje federado permite mantener los datos locales en los dispositivos y posibilita un entrenamiento colaborativo y escalable, reduciendo de forma explícita la exposición de información sensible. Aunque introduce una mayor complejidad operativa asociada a la coordinación y agregación de modelos, esta resulta asumible y necesaria en entornos IoT sensibles, donde la privacidad, la robustez y el cumplimiento normativo son requisitos prioritarios \cite{mcmahan2017communication,kairouz2021advances}.


\section{Aprendizaje Federado (Federated Learning)}
\label{sec:fl}
\subsection{Orígenes y evolución del Aprendizaje Federado}
Como se introdujo en el contexto, el \gls{fl} es un paradigma de aprendizaje automático distribuido en el que múltiples clientes (p.\,ej., dispositivos móviles u organizaciones) entrenan de forma colaborativa un modelo bajo la orquestación de un servidor, manteniendo los datos de entrenamiento \emph{descentralizados} y evitando la centralización de información sensible \cite{mcmahan2017communication,kairouz2021advances}. 

Este es un paradigma de aprendizaje automático distribuido introducido por Google en 2016, con el objetivo de entrenar modelos colaborativos sin necesidad de centralizar los datos de los usuarios \cite{mcmahan2017communication}.  
En lugar de transferir los datos al servidor, cada dispositivo entrena localmente un modelo y comunica únicamente las actualizaciones de los parámetros al servidor central, donde se realiza una agregación global \cite{google2017federated}.  

El algoritmo propuesto por McMahan et al. \cite{mcmahan2017communication}\cite{pmlr-v54-mcmahan17a}, denominado \textbf{Federated Averaging (FedAvg)}, marcó el inicio de este campo. La clave de \textit{FedAvg} reside en desacoplar la frecuencia de entrenamiento local de la frecuencia de comunicación: en lugar de enviar gradientes tras cada iteración. Cada cliente ejecuta múltiples pasos locales de descenso de gradiente sobre sus propios datos y posteriormente comunica únicamente el modelo resultante al final de cada ronda \cite{mcmahan2017communication,kairouz2021advances}.

Posteriormente, el servidor central agrega estas actualizaciones mediante una media ponderada por el tamaño de los conjuntos de datos locales, aproximando el descenso de gradiente global sin necesidad de centralizar los datos.  
De hecho, McMahan et al.\cite{mcmahan2017communication} reportan reducciones de 10--100$\times$ en rondas de comunicación frente a SGD sincronizado, manteniendo un comportamiento robusto en el escenario federado.


Desde entonces, el aprendizaje federado ha evolucionado desde un esquema centralizado (basado en la interacción directa entre un servidor global y múltiples clientes) hacia arquitecturas más complejas, que incorporan agregaciones jerárquicas, esquemas asíncronos y estrategias adaptativas de entrenamiento. Estas extensiones buscan dar respuesta a los principales retos del FL: la heterogeneidad de los datos, los costes de comunicación, las limitaciones computacionales y la preservación de la privacidad \cite{kairouz2021advances,li2020federated}.

\subsection{Definición general y principio de funcionamiento}

El aprendizaje federado permite entrenar un modelo global \( w \) a partir de un conjunto de clientes \( \{1, 2, ..., K\} \), cada uno con su propio conjunto de datos \( D_k \), sin transferir la información sensible fuera de su entorno local.  
Durante cada ronda \( t \), el servidor central distribuye el modelo actual \( w_t \) a un subconjunto de clientes. Estos realizan varias iteraciones de entrenamiento local, ajustando los parámetros del modelo en función de sus propios datos y obteniendo los pesos actualizados \( w_t^k \) o sus gradientes asociados.  

El proceso de actualización global se define mediante una agregación ponderada\cite{pmlr-v54-mcmahan17a}:
\[
w_{t+1} = \sum_{k=1}^{K} \frac{n_k}{n} w_t^k,
\]
donde \( n_k \) representa el número de muestras del cliente \( k \) y \( n = \sum_k n_k \).  
De esta forma, el modelo global se ajusta en función de las contribuciones locales, manteniendo la privacidad y cumpliendo con las normativas de protección de datos (como el GDPR)\cite{rgpd2016}.

\subsection{Agregación de modelos, pesos y gradientes}

En el núcleo del proceso federado se encuentra el componente denominado \textbf{agregador}, encargado de fusionar las actualizaciones locales provenientes de los clientes \cite{bonawitz2019towards}.  
Cada nodo local calcula gradientes o pesos actualizados del modelo a partir de la función de pérdida definida para su conjunto de datos \( D_k \):

\[
w_t^k = w_t - \eta \nabla F_k(w_t),
\]

donde \( \eta \) es la tasa de aprendizaje y \( \nabla F_k(w_t) \) el gradiente local.  
El agregador combina estas actualizaciones y genera el nuevo modelo global \( w_{t+1} \), de acuerdo con la ecuación anterior \cite{pmlr-v54-mcmahan17a}.  

En los enfoques tradicionales, esta operación se realiza de forma centralizada en el servidor principal. Sin embargo, en entornos de gran escala o redes heterogéneas (como los sistemas IoT o móviles), la comunicación directa con el servidor puede resultar ineficiente o costosa.  
Por ello, los estudios recientes han introducido el concepto de \textbf{agregación jerárquica}, donde se emplean niveles intermedios de agregación situados entre los clientes y el servidor global.  

Estos niveles, implementados en nodos intermedios denominados \textit{Fog nodes}, permiten agrupar y combinar las actualizaciones de un subconjunto de clientes antes de enviarlas al nivel superior.  
La formulación general de una agregación jerárquica puede expresarse como:

\[
\underbrace{w_{t+1}^{(f)}}_{\text{Fog}} = \sum_{k \in \mathcal{C}_f} \frac{n_k}{n_f} w_t^k,
\quad
\underbrace{w_{t+1}^{(\text{global})}}_{\text{Cloud}} = \sum_{f=1}^{F} \frac{n_f}{n} w_{t+1}^{(f)},
\]

donde \( \mathcal{C}_f \) representa el conjunto de clientes asociados al nodo Fog \( f \), y \( n_f = \sum_{k \in \mathcal{C}_f} n_k \).  

Este enfoque, conocido como \textbf{Hierarchical Federated Learning (HFL)} o \textbf{Fog Federated Learning}, busca mejorar la escalabilidad y reducir la latencia de comunicación, además de aumentar la resiliencia del sistema frente a fallos y variaciones de red.  
Trabajos como los de Shi et al. \cite{shi2020communication} y Liu et al. \cite{liu2022hierarchical} demuestran que la agregación multinivel reduce significativamente el tiempo de convergencia y el consumo de ancho de banda, sin degradar el rendimiento del modelo.  
Más recientemente, Prigent et al. \cite{prigent2024flcontinuum} y Moussa et al. \cite{moussa2023decentralised} han extendido este paradigma dentro del contexto del \textit{computing continuum}, integrando de forma cooperativa las capas \textit{Edge}, \textit{Fog} y \textit{Cloud}.

En este tipo de arquitecturas, los nodos \textit{Fog} no solo actúan como agregadores intermedios, sino que también pueden asumir un conjunto de \textbf{tareas complementarias} que resultan críticas en despliegues reales\cite{bonomi2012fog,hu2017fog_survey}:

\begin{itemize}
    \item \textbf{Filtrado de ruido en gradientes o pesos:} las actualizaciones locales pueden estar afectadas por ruido estadístico, datos atípicos o comportamientos anómalos de algunos clientes. Los nodos Fog pueden aplicar técnicas de suavizado, detección de valores extremos o agregación robusta para mitigar estas perturbaciones antes de enviar la información al nivel superior \cite{qi2024aggregation_survey}. 
    \item \textbf{Almacenamiento temporal de modelos y versiones intermedias:} el Fog puede actuar como un punto de \textit{cache}, almacenando modelos parciales o versiones recientes del modelo global. Esto reduce la necesidad de acceder continuamente a la nube y permite recuperar el entrenamiento tras desconexiones o fallos temporales de red \cite{bonomi2012fog,hu2017fog_survey,maia2024cloud_edge_3c}.

    \item \textbf{Gestión del balance de carga y coordinación local:} dado que los dispositivos Edge presentan capacidades muy heterogéneas, el nodo Fog puede gestionar qué clientes participan en cada ronda, ajustar la frecuencia de actualización o redistribuir la carga de entrenamiento para evitar cuellos de botella y mejorar la eficiencia global \cite{trindade2024resource_mgmt_fl,cilic2024reactive_hfl}.

    \item \textbf{Inserción de mecanismos de privacidad diferencial:} en lugar de aplicar ruido directamente en dispositivos con recursos limitados, el nodo Fog puede encargarse de añadir mecanismos de privacidad diferencial o de agregación segura, reduciendo el coste computacional en el Edge y manteniendo garantías de privacidad antes de la comunicación con la nube  \cite{islam2025pphffl,qiu2025hierarchical_aggregation}.
\end{itemize}

La incorporación de estas funciones permite acercar la inteligencia al edge, reducir la dependencia del servidor central y desacoplar la dinámica altamente variable de los dispositivos Edge del proceso de agregación global. Como resultado, se obtiene un aprendizaje federado más escalable, eficiente y adaptado a entornos IoT heterogéneos y sensibles.



\subsection{Arquitecturas federadas en el computing continuum: del estado del arte hacia el diseño.}
Antes de entrar en la propuesta de este trabajo, resulta útil revisar algunos enfoques recientes en los que se plantean \textbf{arquitecturas de aprendizaje federado desplegadas a lo largo del \textit{computing continuum}} (Edge--Fog--Cloud y variantes). El objetivo de este repaso no es agotar el estado del arte, sino \textbf{identificar patrones de diseño recurrentes} (p.ej., agregación jerárquica, colaboración \textit{edge-to-edge}, servicios por capa y mecanismos de coordinación) y \textbf{explicitar compromisos de diseño} (latencia, coste de comunicación, privacidad, heterogeneidad y complejidad operativa), especialmente relevantes en escenarios biomédicos.

El aprendizaje federado distribuido en arquitecturas Edge-Fog-Cloud no es un concepto abstracto, sino una materialización concreta de un principio más amplio: el \textbf{Distributed Computing Continuum} (DCC). A continuación, revisamos cuatro perspectivas complementarias que ilustran cómo diseñar sistemas federados reales y qué patrones y compromisos emergen en cada nivel.

\subsubsection{Perspectiva 1: Evolución histórica hacia el continuum integrado}

\begin{figure}[H]
\centering
\includegraphics[width=0.96\linewidth]{Plantilla-Latex-TFG/imagenes/Evolution.jpg}
\caption{\textbf{De mainframes a computing continuum.} Evolución de paradigmas de computación distribuida: mainframe → grid → cluster → cloud → fog/edge → serverless → continuum integrado (DCCS). \textit{El eje horizontal representa crecimiento en dispersión de recursos; el eje vertical, dinamismo en asignación de carga.} Adaptado de Donta et al. (2023).}
\label{fig:dccs_evolution_donta2023}
\end{figure}

La Figura~\ref{fig:dccs_evolution_donta2023} presenta la evolución natural que está siguiendo el paradigma Cloud el \textbf{computing continuum no es esporádico}, sino la convergencia natural de décadas de computación distribuida. Cada paradigma anterior (mainframe centralizado, grids heterogéneos, clusters homogéneos, clouds públicos) ha dejado lecciones: escalabilidad requiere distribución, pero distribución pura es difícil de coordinar. El continuum resuelve esto integrando el espectro completo: desde dispositivos cercanos a datos hasta centros de datos remotos, \textbf{decidiendo dinámicamente dónde ejecutar cada tarea} según \textbf{capacidad local, latencia disponible, coste energético y sensibilidad del dato} \cite{donta2023dccs}.

Para \textbf{aprendizaje federado biomédico}, esto significa: \textit{no ejecutar todo en la nube}, sino elegir dónde entrenar, agregar e inferir. Si el dato de estrés es sensible, entrenar localmente en el wearable. Si el volumen es alto, agregar en una puerta de entrada local (Fog). Si se necesita coordinación global, usar la nube. Esta elasticidad es lo que hace FL viable en IoT sensible.

\subsubsection{Perspectiva 2: Arquitectura técnica de tres capas}

\begin{figure}[H]
\centering
\includegraphics[width=0.96\linewidth]{Plantilla-Latex-TFG/imagenes/computers-12-00198-g007.png}
\caption{\textbf{Arquitectura DCCS concreta: capas, interfaces, y flujos de datos.} Cloud (centro de datos con cómputo masivo), Fog/Edge (gateways, servidores locales, estaciones base), IoT Devices (sensores, wearables, actuadores). \textit{Nótese la doble infraestructura: plano de red (conectividad) y plano de cómputo (orquestación de servicios).} Adaptado de Donta et al. (2023).}
\label{fig:dccs_arch_donta2023}
\end{figure}

La Figura~\ref{fig:dccs_arch_donta2023} presenta una arquitectura DCCS, la cual  \textbf{requiere orquestar dos planos simultáneamente}:

\begin{itemize}
    \item \textbf{Plano de red}: garantizar conectividad entre capas (estaciones base, Internet core, ruteamiento). No es trivial en entornos móviles o con fallos parciales.
    
    \item \textbf{Plano de cómputo}: decidir qué servicios (entrenamiento, inferencia, almacenamiento, monitorización) corren en cada nodo. Aquí es donde \textbf{FL jerárquico} toma decisiones concretas.
\end{itemize}

Para \textbf{entornos pervasivos como en IoT}, la Figura~\ref{fig:dccs_arch_donta2023} ilustra un punto crítico: \textbf{la capa Fog}. Con ella, los dispositivos IoT se ya no se ven sobrecargados intentando hablar directamente con la nube, y la latencia/energía se reducen significativamente. Con ella, la capa Fog actúa como \textbf{amortiguador} (buffering), \textbf{agregador local} (combining updates) y \textbf{punto de decisión} (qué procesar localmente, qué enviar a cloud) \cite{donta2023dccs}.

\subsubsection{Perspectiva 3: Integración \textit{edge caching}}

\begin{figure}[H]
\centering
\includegraphics[width=0.62\linewidth]{Plantilla-Latex-TFG/imagenes/edgecatching.png}
\caption{\textbf{Integración de \textit{edge caching} en un esquema Cloud--Edge.} La figura resume tres funciones típicas en el continuo: (i) \textit{computation offloading} (local/parcial/total), (ii) \textit{edge caching} para reducir latencia y tráfico hacia el núcleo de red y (iii) agregación en el edge como consolidación previa al cloud. Adaptado de Liu et al.~\cite{liu2024iiotframework}.}
\label{fig:edge_caching_liu2024}
\end{figure}

La Figura~\ref{fig:edge_caching_liu2024} muestra cómo el \textbf{edge} puede combinar \textit{caching} y coordinación: almacenar contenido y artefactos cerca del usuario reduce tráfico y latencia, mientras que la consolidación en el edge permite escalar al agrupar múltiples contribuciones antes de enviarlas a la nube. En FL, este patrón se interpreta como (i) \textbf{distribución eficiente del modelo global} (evitando resincronizaciones completas) y (ii) \textbf{reducción del coste de comunicación} al concentrar intercambios en la capa intermedia, manteniendo el cloud como orquestador global~\cite{liu2024iiotframework}.


\subsubsection{Perspectiva 4: Descomposición explícita de servicios por capa}

\begin{figure}[H]
\centering
\includegraphics[width=0.95\linewidth]{Plantilla-Latex-TFG/imagenes/protocols.jpg}
\caption{\textbf{Stack de servicios en la arquitectura continuum.} Capa de Protocolos/Conectividad (MQTT, CoAP, etc.), capa de Datos (ingesta, almacenamiento, sincronización), capa de IA (entrenamiento federado, validación, inferencia). \textit{Nótese: cada capa tiene responsabilidades verticales; la coordinación entre capas es lo crítico.} Adaptado de Liu et al. (2024).}
\label{fig:iioT_liu2024}
\end{figure}

La Figura~\ref{fig:iioT_liu2024} trae el mensaje más importante: \textbf{dibujar capas (Edge-Fog-Cloud) no es suficiente}. \textit{¿Qué hace exactamente cada capa?} Sin responder esto, no hay arquitectura, solo un diagrama bonito.

La figura desglosa:

\begin{itemize}
    \item \textbf{Capa de Protocolos}: MQTT (pub/sub ligero), CoAP (alternativa a HTTP para IoT), sincronización de tiempo, heartbeats.
    
    \item \textbf{Capa de Datos}: buffer local en edge, compresión, sincronización con cloud, manejo de desconexiones.
    
    \item \textbf{Capa de IA/FL}: donde entrenar (edge o cloud), cómo agregar (jerárquico o descentralizado), cuándo comunicar (cada ronda, adaptive, asíncrono).
\end{itemize}

El punto clave: \textbf{la capa de IA no vive en el vacío}. Depende de que la capa de datos entregue actualizaciones sincronizadas y que la capa de protocolos garantice entrega confiable. En dominios biomédicos, un error en la sincronización o pérdida de datos puede invalidar el entrenamiento. Así que estos servicios no son detalles de implementación; son restricciones del diseño \cite{liu2024iiotframework}.

\subsubsection{Síntesis: patrones y compromisos en FL distribuido}

Combinando estas cuatro perspectivas, emergen \textbf{patrones recurrentes}:

\begin{enumerate}
    \item \textbf{Agregación jerárquica}: reduce comunicación, pero añade complejidad en sincronización entre niveles.
    \item \textbf{Caché en el edge}: mejora latencia y energía, pero requiere invalidación coherente cuando el modelo global cambia.
    \item \textbf{Decisiones dinámicas de ubicación}: qué procesar dónde, según capacidad y restricción. Requiere orquestación en tiempo de ejecución.
    \item \textbf{Servicios por capa}: cada capa tiene responsabilidades (protocolo, datos, IA). Las interfaces entre capas son críticas y frecuentemente punto de fallo.
\end{enumerate}

\subsection{Clasificación de los Algoritmos de Aprendizaje Federado}
\label{sec:fl-categories}

Los algoritmos de aprendizaje federado pueden clasificarse en distintas categorías según la relación que existe entre los conjuntos de datos distribuidos en los distintos participantes, atendiendo a dos ejes fundamentales:  
(i) el solapamiento entre los \textbf{individuos o entidades} representados en los datos y  
(ii) el solapamiento entre los \textbf{atributos o características} disponibles en cada nodo.  
Siguiendo esta clasificación, ampliamente aceptada en la literatura \cite{kairouz2021advances,zhang2021survey}, se distinguen tres grandes categorías:

\begin{itemize}
    \item \textbf{Aprendizaje Federado Horizontal (HFL):}  
    Se aplica cuando los distintos participantes disponen de datos correspondientes a \textbf{individuos diferentes}, pero descritos mediante el \textbf{mismo conjunto de atributos}. Es decir, los datos comparten la misma estructura de características, aunque cada nodo posee registros de sujetos distintos.  

    Un ejemplo representativo es el de varios hospitales que registran los mismos atributos clínicos (edad, sexo, constantes vitales, analíticas, etc.) para pacientes diferentes. En este caso, el aprendizaje federado horizontal permite aumentar el tamaño muestral efectivo sin necesidad de compartir los datos sensibles.  

    HFL suele apoyarse en mecanismos de protección de la privacidad como la \textbf{agregación segura}, que impide al servidor observar las contribuciones individuales, y la \textbf{privacidad diferencial}, que introduce ruido controlado en las actualizaciones del modelo para limitar la posible inferencia sobre los datos locales \cite{bonawitz2017practical,bonawitz2019towards}.

    \item \textbf{Aprendizaje Federado Vertical (VFL):}  
    Se emplea cuando los distintos participantes comparten los \textbf{mismos individuos o entidades}, pero cada uno dispone de \textbf{atributos distintos} que describen a dichos individuos. En este escenario, los conjuntos de datos están alineados a nivel de sujetos, pero difieren en el espacio de características.  

    Un caso típico es la colaboración entre un banco y una empresa de comercio electrónico, que poseen información complementaria sobre los mismos clientes (por ejemplo, datos financieros frente a hábitos de consumo).  
    VFL permite combinar estos atributos mediante técnicas de alineamiento seguro de entidades y protocolos criptográficos que evitan la revelación directa de la identidad o de los datos sensibles \cite{cheng2021secureboost,liu2020asymmetrical}.

    \item \textbf{Aprendizaje Federado con Transferencia (FTL):}  
    Este enfoque se aplica cuando existe un solapamiento limitado (o incluso nulo) tanto en los individuos como en los atributos disponibles entre los participantes.  
    En estos casos, se recurre a técnicas de \textit{transfer learning} para aprovechar conocimiento previamente aprendido en dominios relacionados, permitiendo la transferencia parcial de representaciones o parámetros entre modelos \cite{yang2019federated,li2020federated}.
\end{itemize}

\subsection{Tabla de categorización de aprendizaje federado}
\label{sec:fl-table}

La tabla \ref{tab:fl-categorias} resume los principales enfoques del aprendizaje federado según criterios de partición de datos, mecanismos de privacidad, modelos aplicables y estrategias frente a la heterogeneidad, adaptada de \cite{zhang2021survey}.

\begin{table}[H]
\centering
\caption{Clasificación del aprendizaje federado según criterios de partición de datos, privacidad, modelos aplicables y mecanismos frente a la heterogeneidad. Adaptado de \cite{zhang2021survey}.}
\label{tab:fl-categorias}
\begin{tabular}{p{2.5cm} p{3cm} p{4cm} p{4cm}}
\toprule
\textbf{Categoría} & \textbf{Métodos} & \textbf{Ventajas} & \textbf{Aplicaciones} \\
\midrule
\multirow{3}{*}{Partición de datos} 
 & HFL & Incremento del tamaño muestral & Actualización de modelos en móviles \cite{hard2018federated} \\
 & VFL & Aumento de la dimensión de características & Finanzas, salud multi-institucional \cite{cheng2021secureboost} \\
 & FTL & Transferencia entre dominios & Datos pequeños o escasamente solapados \cite{yang2019federated} \\
\midrule
\multirow{3}{*}{Privacidad} 
 & Agregación segura & Evita exponer actualizaciones locales & Deep FL, PATE \cite{bonawitz2017practical,papernot2017semi} \\
 & Encriptación homomórfica & Cálculo sobre datos cifrados & SecureML \cite{mohassel2017secureml} \\
 & Privacidad diferencial & Ruido añadido para limitar filtraciones & Deep learning privado \cite{abadi2016deep} \\
\midrule
\multirow{3}{*}{Modelos aplicables} 
 & Modelos lineales & Simples y eficientes & Regresión lineal, ridge \\
 & Árboles de decisión & Captura relaciones no lineales & Boosting, clasificación \cite{cheng2021secureboost} \\
 & Redes neuronales & Alta capacidad de representación & Reconocimiento de patrones \cite{mcmahan2017communication} \\
\midrule
\multirow{4}{*}{Heterogeneidad} 
 & Comunicación asíncrona & Manejo de retrasos y \textit{stragglers} & Asynchronous FL \cite{xie2019asynchronous} \\
 & Muestreo parcial & Reducción de costes y balanceo & FedAvg con selección \cite{mcmahan2017communication} \\
 & Mecanismos tolerantes a fallos & Prevención de colapso global & Redundancia \\
 & Modelos heterogéneos & Soporte a arquitecturas diversas & FedProx, LG-FedAvg \cite{li2020federated} \\
\bottomrule
\end{tabular}
\end{table}

\subsection{Partición de datos en FL}
\label{sec:fl-partition}
\begin{figure}[H]
  \centering
  \includegraphics[width=1\linewidth]{Plantilla-Latex-TFG/imagenes/Typesfederatedlearning.png}
  \caption{Estrategias de partición de datos en aprendizaje federado: horizontal, vertical y por transferencia. Adaptado de \cite{zhang2021survey}.}
\end{figure}

\subsection{Sistemas de aprendizaje federado en el Computing Continuum}

El despliegue de sistemas de aprendizaje federado en el \emph{computing continuum} (Edge--Fog--Cloud) introduce retos adicionales frente a enfoques centralizados o puramente simulados, especialmente en términos de heterogeneidad de nodos, control del flujo de entrenamiento, comunicaciones distribuidas, privacidad y capacidad de adaptación a arquitecturas jerárquicas. En este contexto, se analizan tres frameworks representativos ---\textit{TensorFlow Federated}, \textit{Flower} y \textit{FLEX}--- cuyas filosofías de diseño y niveles de madurez difieren notablemente. La Tabla~\ref{tab:comparativa-frameworks} recoge una evaluación comparativa cuantitativa basada en los criterios definidos previamente en la Tabla~\ref{tab:criterios-frameworks}.

\subsubsection{TensorFlow Federated (TFF).}
TensorFlow Federated es un framework desarrollado por Google Research con un enfoque fundamentalmente académico, orientado a la investigación en aprendizaje federado y al análisis formal de algoritmos distribuidos~\cite{mcmahan2017communication,bonawitz2019towards}. Su objetivo principal es proporcionar un entorno reproducible para definir cálculos federados y estudiar propiedades como la convergencia, la heterogeneidad de datos no-i.i.d. o la integración de mecanismos de privacidad diferencial.

Desde el punto de vista arquitectónico, TFF se organiza en capas claramente diferenciadas, tal como se muestra en la Figura~\ref{fig:tff_architecture}. En esta arquitectura se distinguen los modelos de aprendizaje y las estrategias federadas, los mecanismos de seguridad y privacidad, y un \emph{runtime} encargado de la ejecución de los cálculos federados. El desarrollador define el comportamiento del sistema de forma declarativa mediante primitivas de alto nivel, mientras que el \emph{runtime} se encarga de orquestar la ejecución sobre clientes simulados.

\begin{figure}[H]
\centering
\includegraphics[width=0.85\linewidth]{Plantilla-Latex-TFG/imagenes/TF.png}
\caption{\textbf{Arquitectura de TensorFlow Federated.} Separación en capas entre modelos y estrategias federadas, mecanismos de seguridad y \emph{runtime} de ejecución. Los clientes son simulados y el flujo de aprendizaje se define de forma declarativa, lo que facilita la experimentación controlada pero limita su aplicación en infraestructuras Edge--Fog--Cloud reales.}
\label{fig:tff_architecture}
\end{figure}

Esta aproximación resulta especialmente adecuada para estudios teóricos y comparativos, pero introduce limitaciones relevantes en términos de viabilidad técnica y flexibilidad arquitectónica cuando se pretende desplegar el sistema sobre nodos físicos heterogéneos, como ocurre en escenarios Edge--Fog--Cloud.

\subsubsection{Flower.}
Flower es un framework de aprendizaje federado \emph{open-source} diseñado explícitamente para su uso en sistemas reales y entornos distribuidos~\cite{beutel2020flower}. A diferencia de TFF, Flower adopta una arquitectura cliente--servidor explícita, donde los distintos componentes del sistema federado se comunican mediante llamadas RPC (\emph{Remote Procedure Calls}), lo que permite desacoplar completamente el entrenamiento local de la coordinación global.

Tal como se ilustra en la Figura~\ref{fig:flower_architecture}, el servidor federado actúa como coordinador del proceso de aprendizaje, gestionando las rondas, seleccionando clientes y aplicando las estrategias de agregación. Los clientes, por su parte, implementan una interfaz bien definida y se comunican con el servidor mediante RPC para recibir el modelo global, entrenar localmente y devolver actualizaciones. Este mecanismo de comunicación resulta especialmente adecuado para entornos distribuidos, ya que permite desplegar clientes en nodos Edge o Fog de forma independiente, incluso en redes heterogéneas.

\begin{figure}[H]
\centering
\includegraphics[width=0.88\linewidth]{Plantilla-Latex-TFG/imagenes/flower-core-framework-architecture.png}
\caption{\textbf{Arquitectura de Flower y bucle de aprendizaje federado.} Comunicación cliente--servidor basada en RPC, donde el servidor coordina las rondas de aprendizaje y los clientes ejecutan el entrenamiento local. Este diseño facilita la integración en infraestructuras distribuidas y la extensión hacia arquitecturas jerárquicas Edge--Fog--Cloud.}
\label{fig:flower_architecture}
\end{figure}

El uso de RPC permite un control explícito del \emph{federated learning loop} y facilita la integración con protocolos y servicios externos, así como la extensión del sistema mediante nodos intermedios o mecanismos de comunicación alternativos. Estas características, junto con su madurez, facilidad de instalación y amplio soporte comunitario, convierten a Flower en una base tecnológica sólida para sistemas federados operativos.

\subsubsection{FLEX.}
FLEX es un framework de aprendizaje federado de carácter experimental, desarrollado en el ámbito académico con el objetivo de maximizar la flexibilidad y la capacidad de experimentación~\cite{flexframework2023}. Su diseño altamente modular permite redefinir prácticamente todos los componentes del flujo federado, incluyendo modelos, estrategias de agregación, protocolos de comunicación y políticas de privacidad.

Esta libertad conceptual lo convierte en una herramienta especialmente adecuada para la investigación y la exploración de nuevas arquitecturas federadas, incluyendo esquemas jerárquicos y asíncronos. Sin embargo, su menor grado de madurez, la escasez de documentación y el reducido soporte comunitario limitan su aplicabilidad como framework base en sistemas Edge--Fog--Cloud complejos y de largo recorrido.

\subsubsection{Síntesis comparativa.}
El análisis realizado pone de manifiesto tres enfoques claramente diferenciados: TensorFlow Federated prioriza el rigor teórico y la simulación controlada; Flower ofrece un equilibrio sólido entre madurez, extensibilidad y viabilidad práctica; y FLEX maximiza la capacidad de experimentación a costa de estabilidad y soporte. 

La comparativa cuantitativa detallada entre los frameworks analizados se presenta en la Tabla~\ref{tab:comparativa-frameworks}, fundamentada en los criterios definidos en la Tabla~\ref{tab:criterios-frameworks}. Dicho análisis sirve de base para la selección tecnológica adoptada en este trabajo y para el diseño de la arquitectura federada propuesta.

\subsection{Síntesis del estado del arte y limitaciones actuales}
\label{sec:fl-synthesis-limitations}
La literatura muestra que el aprendizaje federado ha evolucionado desde el esquema clásico \textit{Edge--Cloud} basado en \textit{FedAvg} \cite{mcmahan2017communication} hacia un marco más amplio orientado a su viabilidad en escenarios reales. Los principales avances se concentran en eficiencia de comunicación, gestión de heterogeneidad, preservación de privacidad y robustez frente a fallos y ataques \cite{kairouz2021advances,li2020federated}.

No obstante, persisten limitaciones estructurales que dificultan su adopción en contextos IoT y biomédicos:

\begin{itemize}
    \item \textbf{Heterogeneidad:} en la práctica, los clientes no tienen datos ni recursos homogéneos. Esto provoca problemas de convergencia, retrasos causados por nodos lentos (\textit{stragglers}) y sesgos, ya que los nodos más rápidos tienden a influir más en el modelo global. El equilibrio entre eficiencia y uso de actualizaciones desactualizadas (\textit{staleness}) sigue siendo un reto abierto.

    \item \textbf{Privacidad y regulación:} aunque los datos no se centralizan, las actualizaciones del modelo pueden filtrar información sensible. Por ello, suelen requerirse técnicas adicionales de protección, que aumentan el coste y pueden degradar el rendimiento, especialmente en datos biomédicos.

    \item \textbf{Privacidad y cumplimiento normativo no garantizados por defecto:} FL evita centralizar datos, pero las actualizaciones pueden filtrar información sensible mediante ataques de reconstrucción o inferencia. Por ello, los mecanismos de privacidad (p.ej., agregación segura y privacidad diferencial) suelen ser necesarios, aunque introducen un coste adicional y pueden afectar al rendimiento del modelo, especialmente con señales biomédicas sensibles.
    
\item \textbf{Seguridad y robustez:} el sistema debe resistir clientes maliciosos (datos o modelos manipulados) y fallos en nodos intermedios. Detectar y mitigar contribuciones anómalas es clave, sin que exista aún una solución general eficiente.

\item \textbf{Coste y escalabilidad:} el uso de dispositivos con recursos limitados obliga a controlar consumo energético y comunicaciones. Ajustar dinámicamente rondas y participantes es necesario, pero depende del contexto y carece de validación estandarizada.

\end{itemize}
Un reto aún poco explorado es el \textbf{diseño arquitectónico completo} de sistemas federados jerárquicos. Aunque los enfoques \textit{Edge--Fog--Cloud} muestran ventajas en latencia y escalabilidad \cite{shi2020communication,liu2022hierarchical,prigent2024flcontinuum,moussa2023decentralised}, la mayoría de trabajos se limitan a formulaciones conceptuales o validaciones parciales. Persisten retos abiertos en orquestación multi-nivel, modelos de confianza en Fog y evaluación reproducible extremo-a-extremo, especialmente en dominios biomédicos.
Por ello, más allá del algoritmo, es clave diseñar \textbf{arquitecturas federadas integrales} que garanticen coordinación, seguridad y fiabilidad, un aspecto crítico en salud conectada.


\section{Detección de Estrés mediante Dispositivos Wearables y Datasets Públicos}
\label{sec:datasets-stress}

El estudio del estrés con datos fisiológicos ha evolucionado desde entornos controlados de laboratorio hacia escenarios reales, gracias a sensores \textit{wearables} y técnicas de aprendizaje automático que permiten una estimación continua y no invasiva en tiempo real \cite{Pinge2024StressWearables,Chacon2025StressRecognition,Rehman2025IoTHealth}.

\subsection{Principales detectores e indicadores de estrés}

La detección automática del estrés mediante dispositivos \textit{wearables} se apoya en la monitorización de distintas señales fisiológicas que reflejan la activación del sistema nervioso autónomo y la respuesta del organismo ante demandas físicas, cognitivas o emocionales. Dado que el estrés no es una magnitud directamente observable, su identificación se realiza de forma indirecta a partir de \textbf{indicadores fisiológicos} que han demostrado una correlación significativa con estados de estrés en numerosos estudios experimentales~\cite{Giakoumis2012,Kim2018}.

Como punto de partida, se analizaron los dispositivos \textit{wearables} y los sensores biométricos más empleados en la literatura para la monitorización del estrés, con el objetivo de comprender en profundidad los mecanismos fisiológicos implicados y evaluar su viabilidad de uso en el contexto de este TFM. A partir de esta revisión se identificaron como señales fisiológicas más relevantes el electrocardiograma (ECG), la fotopletismografía (PPG), la frecuencia cardíaca y su variabilidad (HR/HRV), la actividad electrodérmica (EDA), la acelerometría y la temperatura periférica~\cite{Giakoumis2012,abd2024wearable}.

\paragraph{Electrocardiograma (ECG).}
El electrocardiograma (ECG) es una señal biomédica que registra la actividad eléctrica del corazón mediante electrodos colocados sobre la superficie de la piel, generalmente en el tórax o las extremidades. Cada latido cardíaco genera una secuencia característica de ondas eléctricas asociadas a la contracción y relajación del músculo cardíaco.

A partir del ECG se identifican los denominados complejos QRS, que corresponden a la despolarización ventricular y marcan el instante exacto de cada latido. Esto permite calcular con alta precisión los intervalos de tiempo entre latidos consecutivos, conocidos como intervalos RR.  
Debido a esta precisión temporal, el ECG se considera el \textbf{estándar de referencia clínico} para el análisis de la actividad cardíaca y, en particular, para el estudio de la variabilidad de la frecuencia cardíaca~\cite{Shaffer2017}.

Desde el punto de vista fisiológico, el estrés activa el sistema nervioso simpático, lo que provoca un aumento de la frecuencia cardíaca y una reducción de la variabilidad entre latidos. Por este motivo, los parámetros derivados del ECG han demostrado ser indicadores robustos y fiables del estrés en estudios experimentales y clínicos~\cite{Kim2018}.

\paragraph{Fotopletismografía (PPG).}
La fotopletismografía (PPG) es una técnica óptica no invasiva que mide los cambios en el volumen sanguíneo en los vasos periféricos. Para ello, emplea un emisor de luz (habitualmente LEDs verdes o infrarrojos) y un sensor receptor que detecta la cantidad de luz reflejada o transmitida por el tejido.

Cada latido cardíaco produce una variación en el flujo sanguíneo que se refleja en la señal PPG, permitiendo estimar el ritmo cardíaco y, con técnicas adicionales de procesamiento, la variabilidad de la frecuencia cardíaca.  
La PPG es ampliamente utilizada en \textit{wearables} de muñeca debido a su bajo coste, comodidad y facilidad de integración.

No obstante, la PPG es más sensible al movimiento, al ruido y a factores externos (posición del sensor, temperatura, perfusión sanguínea), lo que introduce incertidumbre en la estimación de parámetros fisiológicos y requiere etapas de filtrado y validación adicionales~\cite{Schafer2013}.

\paragraph{Frecuencia cardíaca (HR) y variabilidad de la frecuencia cardíaca (HRV).}
La frecuencia cardíaca (Heart Rate, HR) representa el número de latidos del corazón por minuto y constituye una de las medidas fisiológicas más básicas del estado cardiovascular.  
Sin embargo, más allá del valor medio de la HR, resulta especialmente relevante analizar cómo varía el tiempo entre latidos consecutivos, lo que se conoce como variabilidad de la frecuencia cardíaca (Heart Rate Variability, HRV).

La HRV refleja el equilibrio dinámico entre las dos ramas del sistema nervioso autónomo: el sistema simpático (asociado a activación y respuesta al estrés) y el parasimpático (asociado a relajación y recuperación).  
En situaciones de estrés, este equilibrio se desplaza hacia la dominancia simpática, lo que se traduce en una \textbf{disminución de la HRV}, indicando una menor capacidad de adaptación fisiológica~\cite{Kim2018}.

Experimentalmente, la HRV se cuantifica mediante diferentes métricas:  
(i) en el dominio temporal (por ejemplo, RMSSD o SDNN),  
(ii) en el dominio frecuencial (componentes LF, HF y su relación LF/HF),  
y (iii) mediante medidas no lineales~\cite{Shaffer2017}.  
Numerosos estudios identifican la HRV como uno de los indicadores más determinantes y sensibles para la detección de estrés~\cite{Kim2018,abd2024wearable}.

\paragraph{Actividad electrodérmica (EDA).}
También conocida como conductancia de la piel, mide los cambios en la capacidad de la piel para conducir corriente eléctrica.  
Estos cambios están directamente relacionados con la actividad de las glándulas sudoríparas, que son controladas exclusivamente por el sistema nervioso simpático.

Desde un punto de vista fisiológico, la EDA constituye un indicador muy directo de la activación emocional y del estrés psicológico, ya que responde rápidamente a estímulos cognitivos o emocionales sin intervención consciente del individuo~\cite{Boucsein2012}.  
Por este motivo, es uno de los marcadores más utilizados en estudios de estrés, ansiedad y carga mental.

En el análisis experimental, la EDA se descompone en un componente tónico (nivel basal de conductancia) y un componente fásico (respuestas rápidas ante estímulos específicos), siendo este último especialmente informativo en protocolos de inducción de estrés~\cite{Boucsein2012}.

\paragraph{Acelerometría.}
La acelerometría mide la aceleración del movimiento corporal en uno o varios ejes espaciales mediante sensores inerciales.  
Aunque no es un indicador fisiológico directo del estrés, su inclusión resulta fundamental para contextualizar el resto de señales biométricas.

En particular, la acelerometría permite distinguir incrementos de la frecuencia cardíaca o de la activación fisiológica debidos a la actividad física de aquellos asociados al estrés psicológico.  
Por ello, suele emplearse como variable de control, para segmentar periodos de reposo, actividad o movimiento excesivo, y para mejorar la robustez de los modelos de detección de estrés~\cite{Giakoumis2012}.

\paragraph{Temperatura periférica.}
La temperatura cutánea periférica refleja el estado de la circulación sanguínea en las extremidades.  
En situaciones de estrés, la activación del sistema nervioso simpático provoca vasoconstricción periférica, lo que puede dar lugar a una disminución de la temperatura de la piel~\cite{Herborn2015}.

Aunque su respuesta es más lenta y menos específica que la de otros indicadores como la EDA o la HRV, la temperatura periférica aporta información complementaria sobre el estado fisiológico general del individuo, especialmente cuando se analiza en combinación con otras señales~\cite{Herborn2015}.

\paragraph{Medición experimental del estrés.}
Dado que el estrés no se mide de forma directa, su evaluación combina señales fisiológicas objetivas con medidas de referencia basadas en la experiencia subjetiva. En laboratorio suele inducirse mediante protocolos estandarizados como el \textit{Trier Social Stress Test}~\cite{Kirschbaum1993}, mientras que en escenarios reales se recurre principalmente a autoinformes y cuestionarios psicológicos validados.

Entre los más utilizados en la literatura destacan cuestionarios como \textbf{PANAS}~\cite{Watson1988}, \textbf{STAI}~\cite{Spielberger1983}, \textbf{SAM}~\cite{Bradley1994} y \textbf{SSSQ}~\cite{Helton2004}, que permiten evaluar emociones, ansiedad y distintos componentes del estado de estrés.  
Estos cuestionarios proporcionan una referencia subjetiva ampliamente aceptada para la anotación de estados de estrés y emociones en estudios experimentales.

Las etiquetas de estrés se construyen integrando señales fisiológicas (medidas objetivas) con autoinformes (medidas subjetivas). Aunque esto introduce variabilidad interindividual, refleja mejor la complejidad del estrés humano y motiva enfoques personalizados y distribuidos como el aprendizaje federado, que se adaptan a cada individuo sin centralizar datos sensibles.

Un recurso fundamental en esta tarea fue la \textbf{Stress in Action Wearables Database (SiA-WD)}, una base de datos abierta que recopila especificaciones técnicas, señales disponibles y evidencia sobre fiabilidad, validez y usabilidad de más de 50 \textit{wearables} de grado comercial y de investigación~\cite{SiA2025}.  

Finalmente, esta revisión se complementó con el meta-análisis de Abd-Alrazaq et al.~\cite{abd2024wearable}, que sintetiza resultados de 19 estudios y 37 estimaciones de rendimiento, mostrando que los indicadores cardíacos y electrodérmicos se encuentran entre los más determinantes, pero también evidenciando una elevada variabilidad interindividual y una tendencia al sobreajuste en datasets pequeños.


\subsection{Principales datasets disponibles}

Tras identificar los principales indicadores fisiológicos del estrés, se revisaron los \textbf{datasets más relevantes de la literatura} para su estudio experimental. Esta revisión permitió evaluar la calidad de las señales, los protocolos de etiquetado y su idoneidad para escenarios distribuidos.

Se priorizaron conjuntos que incluyeran señales fisiológicas relevantes, etiquetas validadas, condiciones realistas o controladas, y potencial de uso en aprendizaje federado por la presencia de múltiples sujetos y variabilidad interindividual.

A continuación, se describen los datasets que han servido de referencia para el diseño experimental y la selección de datos en este TFM.

\textbf{SWEET Study (2025) Smets et al. \cite{Smets2018DigitalPhenotypes}.}  
El estudio SWEET representa uno de los datasets más amplios de monitorización de estrés en condiciones reales. Reunió datos de 1\,002 trabajadores de oficina durante cinco días consecutivos, combinando señales fisiológicas (ECG, conductancia de la piel, temperatura, acelerometría) y datos contextuales de teléfono y encuestas.  
Los resultados mostraron correlaciones significativas entre eventos diarios (como consumo de café o alcohol) y las respuestas fisiológicas, evidenciando la existencia de \textit{fenotipos digitales} individuales de estrés.  
Este trabajo demostró la viabilidad del uso de sensores portátiles para identificar patrones de estrés en entornos naturales, aunque remarcó la necesidad de modelos personalizados. En la actualidad, su acceso no es completamente público y se necesita autorización para su acceso y uso por parte del equipo de investigación. En nuestro caso, como veremos más adelante, se solicitó y se consiguió autorización para su uso.

\textbf{WESAD (Wearable Stress and Affect Detection) Schmidt et al. \cite{schmidt2018wesad}.}  
El dataset WESAD es uno de los conjuntos de datos públicos más utilizados para la detección de estrés debido a su carácter público y disponibilidad.  
Incluye señales multimodales obtenidas de 15 participantes que usaron simultáneamente un sensor de pecho (RespiBAN) y una pulsera (Empatica E4) obteniendo así los siguientes datos: \textit{baseline}, estrés inducido (Trier Social Stress Test) y emoción positiva (videos de humor).  
Contiene registros continuos de ECG, respiración, EMG, EDA, temperatura, acelerómetro y PPG (BVP), junto con cuestionarios psicológicos validados (PANAS, STAI, SAM, SSSQ).  
Su estructura multimodal y su accesibilidad lo han convertido en un estándar de referencia para el desarrollo y comparación de algoritmos de reconocimiento de estrés y emociones.

\textbf{SWELL-KW HRV Koldijk et al. (Radboud University).}\cite{Koldijk2016SWELL}
El dataset SWELL recopila métricas de variabilidad de la frecuencia cardíaca (HRV) de 25 trabajadores de oficina bajo tres condiciones: trabajo normal, presión de tiempo e interrupciones por correo electrónico.  
Cada sesión incluye ventanas deslizantes de 5 minutos con 36 características de HRV calculadas en el dominio temporal, frecuencial y no lineal.  
Los estudios derivados han mostrado una clara diferenciación de estados de estrés basados en HRV, aunque el reducido número de participantes limita su generalización.  
Este dataset es ampliamente utilizado para explorar la relación entre el estrés mental y los patrones cardíacos en contextos de oficina.

\begin{table}[H]
\centering
\caption{Comparativa de los principales datasets para detección de estrés mediante sensores portátiles.}
\label{tab:stress-datasets}
\renewcommand{\arraystretch}{1.25}
\begin{tabular}{p{2.cm} p{1.cm} p{3cm} p{1.5cm} p{2.cm} p{2.cm}}
\toprule
\textbf{Dataset} & \textbf{Sujetos} & \textbf{Datos} & \textbf{Contexto } & \textbf{Etiquetas} & \textbf{Accesibilidad / Ref.} \\
\midrule
\textbf{SWEET (2019)} & 1\,002 &
Fisiológicos (ECG, EDA, Temp., ACC), contextuales (smartphone, EMA, diarios) &
Vida real (5 días continuos) &
Estrés (1–3), placer, arousal, actividad, consumo &
Privado (bajo acuerdo) \cite{Smets2018DigitalPhenotypes} \\

\textbf{WESAD (2018)} & 15 &
Multimodal (ECG, RESP, EMG, EDA, BVP, Temp., ACC) &
Laboratorio controlado (3 fases: baseline, estrés, diversión) &
0 = Baseline, 1 = Estrés, 2 = Alegría &
Público (ACM ICMI) \cite{schmidt2018wesad} \\

\textbf{SWELL-KW HRV (2014)} & 25 &
Derivados de ECG (36 índices HRV: tiempo, frecuencia, no lineal) &
Oficina (tareas reales, 3 condiciones) &
No estrés / Presión temporal / Interrupción &
Público (Radboud Univ.)\cite{Koldijk2016SWELL} \\


\bottomrule
\end{tabular}
\end{table}

\subsection{Discusión y limitaciones}

La comparación entre estos conjuntos de datos evidencia una evolución significativa en la investigación sobre detección de estrés mediante señales fisiológicas.  
Mientras que datasets como \textbf{WESAD} o \textbf{SWELL-KW} fueron diseñados en entornos controlados de laboratorio, el estudio \textbf{SWEET} marca una transición hacia escenarios reales y a gran escala, donde se registran de manera continua las variaciones fisiológicas y contextuales de la vida cotidiana.  

No obstante, la mayoría de los conjuntos disponibles presentan una fuerte \textit{variabilidad interindividual}, reflejando diferencias sustanciales en la respuesta fisiológica al estrés entre sujetos. Esta heterogeneidad limita la capacidad de los modelos generalistas y pone de manifiesto la necesidad de estrategias de \textit{personalización}.  
Diversos trabajos basados en WESAD y SWELL-KW han demostrado que los modelos mejoran considerablemente cuando se aplican técnicas de \textit{fine-tuning} adaptadas a cada individuo, ajustando los parámetros del modelo global a los patrones fisiológicos específicos de cada usuario \cite{abd2024wearable,dahal2023global}.  

En este contexto, las técnicas de \textbf{aprendizaje federado} surgen como una solución prometedora, ya que permiten entrenar modelos personalizados de manera colaborativa sin necesidad de compartir los datos sensibles de los usuarios.  
De este modo, cada dispositivo o nodo local puede adaptar el modelo global a su propia respuesta fisiológica, preservando la privacidad, lo que constituye un paso esencial hacia la detección de estrés verdaderamente contextual, individualizada y ética en entornos reales.





\section{Estudio de las principales tecnologías y herramientas}

%\section{Estudio de Tecnologías}
% Comparación de tecnologías y justificación de las elegidas.
%En este capítulo se presentarán las tecnologías disponibles; en las correspondientes subsecciones de cada Sprint se explicará la elección de las tecnologías que se han elegido y se han tenido en cuenta, y por qué. 
Como se ha definido en los objetivos del proyecto, uno de los propósitos fundamentales de este Trabajo de Fin de Máster es \textbf{diseñar e implementar un prototipo funcional de aprendizaje federado} sobre una arquitectura distribuida basada en el paradigma del \textit{computing continuum}. Para alcanzar este objetivo, resulta imprescindible analizar y comprender las \textbf{tecnologías que permiten materializar dicho sistema}, tanto desde el punto de vista del desarrollo software como del despliegue, la reproducibilidad y la coordinación de componentes distribuidos.

En consecuencia, este capítulo no pretende ofrecer una revisión exhaustiva de todas las tecnologías existentes, sino realizar una \textbf{selección razonada de aquellas herramientas y entornos} que resultan más adecuadas para:
\begin{itemize}
    \item el desarrollo de modelos de \textit{Machine Learning} y aprendizaje federado;
    \item la implementación de arquitecturas distribuidas \textit{Edge--Fog--Cloud};
    \item el despliegue reproducible de servicios en entornos heterogéneos;
    \item y la gestión del ciclo de vida del software en un proyecto experimental de investigación.
\end{itemize}

El análisis de estas tecnologías permite justificar las decisiones adoptadas a lo largo de los distintos \textit{sprints} de desarrollo y proporciona el contexto necesario para comprender la arquitectura final del sistema y su implementación práctica.

\subsection{Lenguajes de desarrollo}
\label{sec:lenguajesdes}

La elección del lenguaje de programación es un factor clave en sistemas de \textit{Machine Learning}, y resulta especialmente crítica en enfoques distribuidos como el aprendizaje federado, donde influyen la disponibilidad de frameworks, la reproducibilidad experimental y la integración con infraestructuras distribuidas.

Informes recientes y encuestas del sector \cite{datacamp2025,stackoverflow2024,tiobe2025} sitúan a lenguajes como Python, R, Java, C++ o Scala entre los más utilizados en ciencia de datos y aprendizaje automático. En este contexto, Python destaca de forma clara como el ecosistema dominante para investigación y prototipado, gracias a su amplio soporte en bibliotecas especializadas y herramientas de desarrollo \cite{geeksforgeeks2025}.

A continuación, se analizan los lenguajes considerados en este TFM y su adecuación para el desarrollo de un sistema de aprendizaje federado.

\begin{itemize}
    \item \textbf{Python:}  
    Python es el lenguaje predominante en el ámbito del \textit{machine learning} y la inteligencia artificial. Su principal fortaleza reside en la amplitud y madurez de su ecosistema de librerías, que incluye frameworks de entrenamiento (\texttt{PyTorch}, \texttt{TensorFlow}), herramientas de modelado clásico (\texttt{Scikit-learn}) y, de forma especialmente relevante para este trabajo, frameworks específicos de aprendizaje federado como \texttt{Flower} y \texttt{TensorFlow Federated}.  
    Además, Python se integra de forma natural con tecnologías de contenedorización, sistemas de mensajería y entornos distribuidos, lo que facilita la implementación rápida de prototipos complejos y reproducibles. Estas características lo convierten en la opción más adecuada para un TFM centrado en arquitectura, experimentación y validación de hipótesis.

    \item \textbf{R:}  
R es un lenguaje ampliamente utilizado en análisis estadístico y validación de modelos, con paquetes consolidados como \texttt{caret}, \texttt{randomForest} o \texttt{tidyverse}. Su fortaleza principal reside en la exploración de datos, el análisis inferencial y la visualización avanzada, siendo especialmente valorado en entornos académicos y clínicos.  
Sin embargo, el soporte nativo para aprendizaje federado, entrenamiento distribuido y despliegues escalables es limitado. En este contexto, R se considera más adecuado como herramienta complementaria para el análisis y la validación posterior de resultados que como lenguaje principal para la implementación de un sistema federado completo.

\item \textbf{C++:}  
C++ destaca por su alto rendimiento, control preciso de memoria y eficiencia computacional. Estas características lo hacen especialmente atractivo para implementar componentes críticos del sistema, como algoritmos de agregación altamente optimizados o módulos ejecutados en dispositivos con recursos muy limitados.  
No obstante, la complejidad del desarrollo, la menor disponibilidad de frameworks de aprendizaje federado de alto nivel y el mayor coste de prototipado y mantenimiento reducen su idoneidad como lenguaje principal en un trabajo de carácter experimental y exploratorio como este TFM.


\item \textbf{Java:}  
Java es ampliamente utilizado en sistemas distribuidos y entornos de producción gracias a la portabilidad que ofrece la \textbf{Java Virtual Machine (JVM)}, que permite ejecutar el mismo código en distintas plataformas e incorpora mecanismos robustos de gestión de memoria, concurrencia y seguridad.  
En aprendizaje federado existen frameworks como \texttt{FATE}, orientados a entornos productivos con altos requisitos de control. No obstante, frente a Python, Java presenta menor flexibilidad y un ciclo de experimentación más costoso, lo que limita su idoneidad para un TFM centrado en prototipado y validación arquitectónica.


\item \textbf{Scala:}  
Scala es un lenguaje que combina programación funcional y orientada a objetos y que se ejecuta directamente sobre la \textbf{JVM}. Esto le permite heredar muchas de las ventajas estructurales de Java, como la portabilidad, la robustez y el acceso al amplio ecosistema de librerías de la plataforma Java.  
Su principal valor añadido reside en su integración nativa con Apache Spark, uno de los motores más utilizados para el procesamiento distribuido de grandes volúmenes de datos. A través de \texttt{MLlib}, la biblioteca de aprendizaje automático de Apache Spark, Scala permite desarrollar \textit{pipelines} de aprendizaje distribuido que ejecutan algoritmos de \textit{machine learning} de forma paralela sobre grandes volúmenes de datos, integrando preprocesamiento, entrenamiento y evaluación dentro de un mismo entorno distribuido y tolerante a fallos. 
Sin embargo, su adopción en aprendizaje federado es todavía limitada, y su curva de aprendizaje es considerablemente mayor que la de Python, lo que dificulta su uso en proyectos donde prima la rapidez de desarrollo y la iteración experimental.


\end{itemize}

Con el fin de objetivar la elección del lenguaje de desarrollo, la Tabla~\ref{tab:comparativa-lenguajes} resume una comparación de los lenguajes analizados atendiendo a criterios técnicos relevantes para este trabajo.

\begin{table}[H]
\centering
\caption{Comparativa de lenguajes de programación para sistemas de aprendizaje federado.}
\label{tab:comparativa-lenguajes}
\renewcommand{\arraystretch}{1.25}
\begin{tabular}{p{2.8cm} p{3cm} p{2.5cm} p{2cm} p{3cm}}
\toprule
\textbf{Lenguaje} & \textbf{Frameworks ML / FL} & \textbf{Soporte FL} & \textbf{Rendimiento} & \textbf{Facilidad de prototipado} \\
\midrule
Python & PyTorch, TensorFlow, Flower, TFF & Muy alto & Medio & Muy alta \\
R & caret, randomForest & Bajo & Medio & Alta (análisis) \\
C++ & Implementaciones propias & Bajo & Muy alto & Baja \\
Java & FATE, DL4J & Medio & Alto & Media \\
Scala & Spark MLlib & Medio & Alto & Baja \\
\bottomrule
\end{tabular}
\end{table}


\begin{figure}[H]
    \centering
    % --- Subfigura: Python ---
    \begin{subfigure}[b]{0.16\linewidth}
        \centering
        \includegraphics[width=\linewidth]{Plantilla-Latex-TFG/imagenes/Python-Logo.png}
        \caption{Python.}
        \label{fig:python}
    \end{subfigure}
    \hfill
    % --- Subfigura: R ---
    \begin{subfigure}[b]{0.16\linewidth}
        \centering
        \includegraphics[width=\linewidth]{Plantilla-Latex-TFG/imagenes/r.jpg}
        \caption{R.}
        \label{fig:r}
    \end{subfigure}
    \hfill
    % --- Subfigura: C++ ---
    \begin{subfigure}[b]{0.16\linewidth}
        \centering
        \includegraphics[width=\linewidth]{Plantilla-Latex-TFG/imagenes/C++.png}
        \caption{C++.}
        \label{fig:cpp}
    \end{subfigure}
    \hfill
    % --- Subfigura: Java ---
    \begin{subfigure}[b]{0.16\linewidth}
        \centering
        \includegraphics[width=\linewidth]{Plantilla-Latex-TFG/imagenes/java.png}
        \caption{Java.}
        \label{fig:java}
    \end{subfigure}
    \hfill
    % --- Subfigura: Scala ---
    \begin{subfigure}[b]{0.16\linewidth}
        \centering
        \includegraphics[width=\linewidth]{Plantilla-Latex-TFG/imagenes/SCALA.png}
        \caption{Scala.}
        \label{fig:scala}
    \end{subfigure}

    \caption{Lenguajes de programación más utilizados en ciencia de datos y aprendizaje automático.}
    \label{fig:lenguajes-programacion}
\end{figure}



\subsection{Entorno de desarrollo, programación y despliegue}

El entorno de desarrollo es clave para garantizar reproducibilidad, portabilidad y control del sistema. Entre las distintas opciones, las tecnologías de virtualización permiten aislar y replicar aplicaciones de forma fiable \cite{turnbull2014docker}.

Pueden distinguirse tres enfoques principales:
\begin{itemize}
    \item \textbf{Máquinas virtuales (VMs):} ofrecen aislamiento completo al simular un sistema operativo íntegro, con alta flexibilidad a costa de mayor consumo de recursos.
    \item \textbf{Contenedores:} proporcionan virtualización ligera compartiendo el kernel del sistema anfitrión, combinando aislamiento, eficiencia y facilidad de despliegue reproducible, especialmente en sistemas distribuidos.
    
    \item \textbf{Ejecución nativa:} simplifica el despliegue, pero reduce la portabilidad y el aislamiento, dificultando la reproducibilidad.
\end{itemize}


En la práctica, destacan dos herramientas principales: \textbf{VirtualBox}, orientada a la virtualización de sistemas operativos completos, y \textbf{Docker}, un motor de contenedores ampliamente adoptado que facilita la creación, despliegue y gestión de aplicaciones empaquetadas. Docker Desktop añade una interfaz integrada que ha favorecido su uso tanto en entornos académicos como industriales.

Como alternativa a Docker, \textbf{Podman} ofrece compatibilidad con su ecosistema (imágenes OCI, \texttt{Dockerfile} y registries estándar), pero introduce diferencias arquitectónicas relevantes. Mientras Docker se apoya en un demonio central persistente (\textit{dockerd}) que opera con privilegios elevados, Podman adopta un enfoque \textit{daemonless}, ejecutando los contenedores como procesos gestionados directamente por el usuario, con implicaciones directas en seguridad y gestión del sistema.

Desde un punto de vista funcional, el uso de un demonio central no es imprescindible, ya que el aislamiento de los contenedores lo proporciona directamente el kernel (namespaces, cgroups y capacidades). Al eliminar este demonio, Podman reduce los componentes privilegiados en ejecución, disminuyendo la superficie de ataque.

Además, Podman permite la ejecución \textit{rootless}, de modo que los contenedores heredan los permisos del usuario que los lanza. Esto limita el impacto de fallos de aislamiento y mejora la trazabilidad y auditoría del sistema al asociar cada contenedor a usuarios y procesos concretos \cite{podman2023}.

Pese a estas consideraciones, en el contexto de este TFM se optó por el uso de \textbf{Docker} debido a su condición de \emph{estándar de facto} en la gestión de contenedores, así como a su amplia adopción, madurez del ecosistema y facilidad de instalación y configuración.
\begin{figure}[H]
    \centering

    \begin{subfigure}[b]{0.28\linewidth}
        \centering
        \includegraphics[width=\linewidth]{Plantilla-Latex-TFG/imagenes/Vbox.jpg}
        \caption{VirtualBox (máquinas virtuales).}
        \label{fig:vbox}
    \end{subfigure}
    \hfill
    \begin{subfigure}[b]{0.28\linewidth}
        \centering
        \includegraphics[width=\linewidth]{Plantilla-Latex-TFG/imagenes/docker.png}
        \caption{Docker (contenedores).}
        \label{fig:docker}
    \end{subfigure}
    \hfill
    \begin{subfigure}[b]{0.28\linewidth}
        \centering
        \includegraphics[width=\linewidth]{Plantilla-Latex-TFG/imagenes/podman.png}
        \caption{Podman (contenedores sin daemon).}
        \label{fig:podman}
    \end{subfigure}

    \caption{Tecnologías de virtualización y contenedores.}
    \label{fig:virtualizacion-general}
\end{figure}

\subsection{Control de versiones y uso de Integración Continua y Despliegue Continuo CI/CD}
\label{sec:Controlversiones}
El uso de sistemas de control de versiones y plataformas de alojamiento en la nube es esencial para el desarrollo de software, ya que permiten gestionar cambios, coordinar el trabajo colaborativo y garantizar la trazabilidad. Integradas con procesos de CI/CD, facilitan la automatización de pruebas y el despliegue ágil en entornos distribuidos \cite{Loeliger2012}.


\begin{itemize}
    \item \textbf{Git:} Sistema de control de versiones distribuido que gestiona de forma precisa el historial de cambios en el código fuente. Su arquitectura descentralizada permite el trabajo con ramas, fusiones y reversión de cambios, asegurando la integridad y trazabilidad de los datos \cite{Chacon2014}.

    \item \textbf{GitHub:} Plataforma de alojamiento en la nube basada en Git que añade herramientas avanzadas de colaboración y gestión de proyectos. Destaca por sus \textit{pull requests}, el seguimiento de incidencias y la integración con \textbf{GitHub Actions}, que permite la automatización de flujos de CI/CD \cite{githubdocs2025}. 
\end{itemize}

No obstante, existen alternativas ampliamente utilizadas en el sector:  

\begin{itemize}
    \item \textbf{GitLab:} Plataforma también basada en Git que ofrece un ecosistema completo de DevOps, integrando de forma nativa el control de versiones, la gestión de proyectos y un potente sistema de integración contínua (CI/CD) \cite{gitlabdocs2025}. Su principal ventaja frente a GitHub es que puede desplegarse en entornos propios, lo que incrementa el control sobre la seguridad y la privacidad de los datos.  

    \item \textbf{Bitbucket:} Desarrollado por Atlassian, combina repositorios Git con integración directa con Jira y Trello, siendo una opción atractiva para equipos que ya utilizan este ecosistema.  

    \item \textbf{Jenkins:} Herramienta de automatización de código abierto muy extendida en la industria. Aunque no es un servicio de alojamiento de código como GitHub o GitLab, se utiliza como sistema de CI/CD para automatizar compilaciones, pruebas y despliegues en prácticamente cualquier entorno de desarrollo \cite{jenkinsdocs2025}.
\end{itemize}


\begin{figure}[H]
    \centering
    \includegraphics[width=0.5\linewidth]{Plantilla-Latex-TFG/imagenes/git.png}
    \caption{Control de versiones con Git.}
    \label{fig:git}
\end{figure}

\begin{figure}[H]
    \centering
    % Subfigura: GitHub
    \begin{subfigure}[b]{0.32\linewidth}
        \centering
        \includegraphics[width=\linewidth]{Plantilla-Latex-TFG/imagenes/github.png}
        \caption{GitHub.}
        \label{fig:github}
    \end{subfigure}
    \hfill
    % Subfigura: GitLab
    \begin{subfigure}[b]{0.32\linewidth}
        \centering
        \includegraphics[width=\linewidth]{Plantilla-Latex-TFG/imagenes/gitlab.png}
        \caption{GitLab.}
        \label{fig:gitlab}
    \end{subfigure}
    \hfill
    % Subfigura: Bitbucket
    \begin{subfigure}[b]{0.32\linewidth}
        \centering
        \includegraphics[width=\linewidth]{Plantilla-Latex-TFG/imagenes/bitbucket.png}
        \caption{Bitbucket.}
        \label{fig:bitbucket}
    \end{subfigure}
    %
    \caption{Plataformas de alojamiento en la nube más utilizadas para control de versiones y colaboración en desarrollo software.  }
    \label{fig:plataformas-nube}
\end{figure}


\subsection{Herramientas de gestión de proyectos}

En el desarrollo de software moderno, las herramientas de gestión de proyectos desempeñan un papel esencial al facilitar la planificación, la comunicación y la coordinación entre los miembros del equipo.  
Su uso promueve un enfoque ágil y colaborativo, mejora la organización de tareas y proporciona una visión clara y actualizada del estado del proyecto, lo que repercute directamente en la calidad del producto final y en la eficiencia del proceso de desarrollo.  

Entre las soluciones más utilizadas se encuentran \textbf{Jira}, \textbf{Trello}, \textbf{GitHub Projects}, \textbf{OpenProject} y \textbf{Redmine}, todas ellas ampliamente empleadas en entornos profesionales y académicos.  
Aunque cada una presenta enfoques diferentes, todas permiten integrar ideas, necesidades y tareas en un entorno visual y colaborativo, adaptándose a distintos estilos de gestión (ágil, híbrido o tradicional) y a equipos de diversos tamaños.

\begin{figure}[H]
    \centering
    % --- Subfigura: Jira ---
    \begin{subfigure}[b]{0.19\linewidth}
        \centering
        \includegraphics[width=\linewidth]{Plantilla-Latex-TFG/imagenes/jira.jpg}
        \caption{Jira.}
        \label{fig:jira}
    \end{subfigure}
    \hfill
    % --- Subfigura: Trello ---
    \begin{subfigure}[b]{0.19\linewidth}
        \centering
        \includegraphics[width=\linewidth]{Plantilla-Latex-TFG/imagenes/trello.png}
        \caption{Trello.}
        \label{fig:trello}
    \end{subfigure}
    \hfill
    % --- Subfigura: GitHub Projects ---
    \begin{subfigure}[b]{0.19\linewidth}
        \centering
        \includegraphics[width=\linewidth]{Plantilla-Latex-TFG/imagenes/githubprojects.png}
        \caption{GitHub Projects.}
        \label{fig:github-projects}
    \end{subfigure}
    \hfill
    % --- Subfigura: OpenProject ---
    \begin{subfigure}[b]{0.19\linewidth}
        \centering
        \includegraphics[width=\linewidth]{Plantilla-Latex-TFG/imagenes/openproj.jpg}
        \caption{OpenProject.}
        \label{fig:openproject}
    \end{subfigure}
    \hfill
    % --- Subfigura: Redmine ---
    \begin{subfigure}[b]{0.19\linewidth}
        \centering
        \includegraphics[width=\linewidth]{Plantilla-Latex-TFG/imagenes/redmine.jpg}
        \caption{Redmine.}
        \label{fig:redmine}
    \end{subfigure}
    
    \caption{Principales plataformas de gestión de proyectos utilizadas en el desarrollo ágil y colaborativo de software.}
    \label{fig:plataformas-gestion-proyectos}
\end{figure}
\begin{itemize}
    \item \textbf{Jira:}  
    Herramienta de referencia para la gestión de proyectos ágiles en equipos medianos y grandes. Soporta \textit{Scrum} y \textit{Kanban}, permite gestionar backlogs, sprints e incidencias, y se integra de forma nativa con herramientas de CI/CD como GitHub, Bitbucket y Jenkins.

    \item \textbf{Trello:}  
    Plataforma ligera basada en \textit{Kanban}, orientada a equipos pequeños y proyectos académicos. Facilita el seguimiento visual de tareas mediante tableros y tarjetas, siendo adecuada para flujos de trabajo simples.

    \item \textbf{GitHub Projects:}  
    Sistema de gestión integrado en GitHub que combina control de versiones y planificación. Permite organizar tareas mediante \textit{issues}, hitos y tableros, manteniendo una trazabilidad directa con el código y los flujos de automatización.

    \item \textbf{OpenProject:}  
    Herramienta de código abierto para equipos medianos o grandes, compatible con metodologías ágiles y tradicionales. Ofrece planificación avanzada mediante diagramas de Gantt y seguimiento estructurado de tareas, habitual en entornos institucionales y de investigación.

    \item \textbf{Redmine:}  
    Plataforma de código abierto flexible y extensible, orientada a la gestión de múltiples proyectos. Proporciona seguimiento de incidencias, control de versiones y planificación temporal, siendo común en entornos académicos y organizaciones con estructuras jerárquicas.
\end{itemize}


\subsection{Protocolos de comunicación en sistemas distribuidos}
\label{sec:protocolos}

En arquitecturas distribuidas, y especialmente en sistemas de aprendizaje federado, la elección del protocolo de comunicación resulta determinante para el rendimiento, la escalabilidad y la robustez del sistema. Estos protocolos gobiernan la forma en que los nodos intercambian parámetros de modelos, métricas y señales de control, afectando directamente a la latencia y al consumo de red \cite{tanenbaum2017distributed}.

Los enfoques clásicos basados en \textbf{HTTP/REST} ofrecen simplicidad y amplia compatibilidad, pero presentan limitaciones en escenarios con comunicaciones frecuentes y bidireccionales, como ocurre en el aprendizaje federado iterativo. En este contexto, alternativas como \textbf{gRPC} permiten implementar llamadas a procedimiento remoto (RPC) de forma eficiente mediante serialización binaria y conexiones persistentes, reduciendo significativamente la latencia y el overhead de comunicación \cite{grpc2023}.

Por otro lado, protocolos orientados a eventos como \textbf{MQTT} adoptan un modelo de publicación/suscripción especialmente adecuado para entornos \textit{Edge--Fog}, donde los nodos pueden ser intermitentes y con recursos limitados. MQTT destaca por su bajo consumo de ancho de banda, soporte de distintos niveles de calidad de servicio (QoS) y elevada escalabilidad, lo que lo convierte en una opción idónea para la coordinación de nodos distribuidos \cite{hunkeler2008mqtt}.

La Tabla~\ref{tab:comparativa-protocolos} sintetiza una comparación de los principales protocolos considerados, atendiendo a su modelo de comunicación, latencia y adecuación al sistema desarrollado en este TFM.


\begin{table}[H]
\centering
\footnotesize
\caption{Comparativa de protocolos de comunicación en sistemas distribuidos y aprendizaje federado.}
\label{tab:comparativa-protocolos}
\renewcommand{\arraystretch}{1.2}
\setlength{\tabcolsep}{3pt}

\begin{tabularx}{\linewidth}{p{2.2cm} X p{1.6cm} p{1.8cm} p{2.2cm}}
\toprule
\textbf{Protocolo} & \textbf{Modelo} & \textbf{Latencia} & \textbf{Escalabilidad} & \textbf{Adecuación} \\
\midrule
HTTP/REST & Cliente--servidor & Media & Media & Media \\
gRPC & RPC binario & Baja & Alta & Muy alta \\
MQTT & Publicación/suscripción & Muy baja & Muy alta & Muy alta \\
WebSockets & Bidireccional persistente & Baja & Media & Baja \\
AMQP & Mensajería robusta & Media & Alta & Media \\
\bottomrule
\end{tabularx}
\end{table}


\subsection{Tecnologías de mensajería y coordinación distribuida}
\label{sec:mensajeria}

Las tecnologías de mensajería constituyen un elemento clave en la coordinación de sistemas distribuidos desacoplados, permitiendo intercambios asíncronos de información entre componentes heterogéneos. En el contexto del aprendizaje federado, estas tecnologías facilitan la comunicación entre nodos Edge, Fog y servidores centrales, incluso en presencia de conectividad intermitente o topologías dinámicas \cite{kleppmann2017ddia}.

Entre las soluciones más ligeras destaca \textbf{MQTT}, diseñado específicamente para entornos con recursos limitados y redes poco fiables. Su bajo overhead y su modelo de publicación/suscripción lo convierten en una opción especialmente adecuada para escenarios IoT y Edge Computing \cite{banks2019mqtt}. 

Otras tecnologías como \textbf{RabbitMQ} o \textbf{Apache Kafka} ofrecen mayores garantías de persistencia y tolerancia a fallos, siendo habituales en arquitecturas empresariales orientadas a procesamiento de eventos a gran escala. No obstante, estas soluciones introducen un mayor coste de despliegue y mantenimiento, que puede resultar excesivo para prototipos experimentales o entornos académicos.

La Tabla~\ref{tab:comparativa-mensajeria} presenta una comparación de las principales tecnologías de mensajería evaluadas, destacando sus diferencias en términos de overhead, calidad de servicio y adecuación a los objetivos de este TFM.


\begin{table}[H]
\centering
\footnotesize
\caption{Comparativa de tecnologías de mensajería para coordinación distribuida.}
\label{tab:comparativa-mensajeria}
\renewcommand{\arraystretch}{1.2}
\setlength{\tabcolsep}{3pt}

\begin{tabularx}{\linewidth}{p{2.4cm} X p{1.8cm} p{1.4cm} p{2.4cm}}
\toprule
\textbf{Tecnología} & \textbf{Paradigma} & \textbf{Overhead} & \textbf{QoS} & \textbf{Adecuación} \\
\midrule
MQTT & Publicación/suscripción ligero & Muy bajo & Sí & Muy alta \\
RabbitMQ & AMQP & Medio & Sí & Media \\
Kafka & Streaming distribuido & Alto & Sí & Baja \\
ZeroMQ & Mensajería directa & Bajo & No & Media \\
\bottomrule
\end{tabularx}
\end{table}

\subsection{Observabilidad y monitorización de sistemas distribuidos}
\label{sec:observabilidad}

La observabilidad es un requisito fundamental en sistemas distribuidos complejos, ya que permite comprender el comportamiento interno del sistema a partir de sus salidas externas, facilitando la detección de errores, el análisis de rendimiento y la validación experimental \cite{sigelman2010dapper}.

Las soluciones tradicionales basadas únicamente en \textbf{logs} resultan insuficientes para arquitecturas distribuidas, donde los flujos de ejecución se propagan a través de múltiples nodos. En este contexto, herramientas modernas como \textbf{Prometheus} y \textbf{OpenTelemetry} permiten recolectar métricas, trazas y eventos de forma estandarizada, proporcionando una visión global y coherente del sistema \cite{prometheus2018}.

Prometheus destaca por su escalabilidad y su integración nativa con arquitecturas basadas en microservicios, mientras que OpenTelemetry actúa como un estándar unificador para la instrumentación de aplicaciones distribuidas, facilitando la interoperabilidad entre distintas herramientas de observabilidad \cite{opentelemetry2023}. Complementariamente, plataformas como \textbf{Grafana} permiten la visualización y análisis de métricas de forma intuitiva, apoyando la toma de decisiones durante el desarrollo y la experimentación.

La Tabla~\ref{tab:comparativa-observabilidad} resume las principales tecnologías consideradas para la observabilidad del sistema, evaluando su escalabilidad, facilidad de integración y adecuación al prototipo desarrollado.


\begin{table}[H]
\centering
\footnotesize
\caption{Comparativa de tecnologías de observabilidad en sistemas distribuidos.}
\label{tab:comparativa-observabilidad}
\renewcommand{\arraystretch}{1.2}
\setlength{\tabcolsep}{3pt}

\begin{tabularx}{\linewidth}{p{2.4cm} p{2.4cm} p{2cm} p{2cm} X}
\toprule
\textbf{Tecnología} & \textbf{Tipo} & \textbf{Escalabilidad} & \textbf{Integración} & \textbf{Adecuación} \\
\midrule
Logging plano & Logs & Baja & Alta & Media \\
Prometheus & Métricas & Muy alta & Alta & Muy alta \\
OpenTelemetry & Trazas y métricas & Muy alta & Alta & Muy alta \\
ELK Stack & Logs centralizados & Alta & Media & Media \\
Grafana & Visualización & Muy alta & Alta & Muy alta \\
\bottomrule
\end{tabularx}
\end{table}


% Detalles del análisis previo al desarrollo, incluidos los requisitos y posiblemente mockups de la aplicación.
 

%%%%%%%%%%
%%%%%%%%%%

\chapter{Análisis de Stakeholders}
\label{chap:stakeholders}

Los \textit{stakeholders} o partes interesadas son todos aquellos agentes que, de forma directa o indirecta, influyen en el desarrollo del proyecto o se ven afectados por sus resultados.  
En este TFM, el producto no consiste en una aplicación final dirigida al usuario, sino en una \textbf{arquitectura de \textit{back-end} basada en aprendizaje federado y \textit{computing continuum}}, concebida como infraestructura de soporte para futuras aplicaciones en distintos dominios.  
El caso de uso inicial que motiva este trabajo se centra en la \textbf{detección y monitorización del estrés en personas mayores}, aunque la arquitectura se ha diseñado con la suficiente flexibilidad para ser reutilizada en otros ámbitos donde la privacidad, la distribución y la escalabilidad de los datos resulten críticas.

\section*{Stakeholders internos}

\begin{description}
  \item[Equipo de desarrollo:]  
  Representado por el autor del TFM, es el responsable del diseño, implementación y validación técnica de la arquitectura.  
  Su interés radica en lograr un sistema funcional, escalable y científicamente sólido que sirva como demostración práctica de los conceptos teóricos estudiados (aprendizaje federado, fog computing, privacidad diferencial, etc.).  
  Para este stakeholder, el proyecto constituye además una oportunidad de aprendizaje avanzado, adquisición de competencias tecnológicas y contribución tangible al progreso científico del grupo de investigación.

\item[Equipo de investigación:]  
Actúa como cliente interno y supervisor técnico-científico del desarrollo.  
Está conformado por las tutoras del TFM junto con otros investigadores.  
Su interés radica en disponer de una \textbf{plataforma experimental robusta} capaz de integrar datos biomédicos entre múltiples instituciones o dispositivos, sin comprometer la privacidad de los participantes.  

Sus objetivos de investigación se articulan en dos líneas principales:
\begin{enumerate}
    \item \textbf{Detectar y monitorizar el estrés en personas mayores mediante tecnologías de \textit{computing continuum}.}  
    Esta línea busca validar el uso de arquitecturas distribuidas que permitan analizar datos fisiológicos en tiempo real, reduciendo la latencia, preservando la privacidad y favoreciendo la toma de decisiones clínicas o preventivas.
    
    \item \textbf{Experimentar con distintas arquitecturas y sistemas de continiuum computing.}  
    El propósito es identificar qué configuraciones, arquitecturas,  frameworks de nodos, esquemas de agregación y modelos de entrenamiento resultan más adecuados para escenarios sensibles, como la salud o el bienestar digital.  
    Este enfoque comparativo permitirá sentar las bases para futuras aplicaciones que aprovechen el potencial del \textit{federated learning} en entornos distribuidos y heterogéneos.
\end{enumerate}

El equipo de investigación tiene, por tanto, un doble rol: definir los requisitos científicos y éticos del sistema, y evaluar su idoneidad como infraestructura reusable en proyectos de salud digital.  
De la validación de este trabajo dependerá la continuidad de la línea de investigación, la publicación de resultados en entornos académicos y la transferencia potencial del conocimiento hacia aplicaciones clínicas o sociales.

\end{description}

\section*{Stakeholders externos}

\begin{description}
  \item[Organismos financiadores:]  
  En este caso, el Gobierno de España, a través del proyecto nacional \textbf{PID2023-149185OB-I00}.  
  Su implicación es estratégica: buscan resultados científicos contrastables que fortalezcan la investigación en inteligencia artificial aplicada a la salud, favorezcan la innovación tecnológica nacional y promuevan el impacto social positivo.  
  Un proyecto exitoso justifica la inversión pública y contribuye a consolidar redes de investigación interdisciplinar, en este caso entre ingeniería, psicología, medicina y ciencia de datos.



 \item[Comunidad científica y académica:]  
  Engloba a investigadores de áreas como informática biomédica, inteligencia artificial, ingeniería del software y ciencias de la salud.  
  Su interés radica en \textbf{validar, reutilizar y extender} la arquitectura propuesta.  
  Esto incluye tanto contextos sanitarios como otros dominios sensibles al dato (telemedicina, bienestar digital, ciberseguridad industrial).

  Dentro de este grupo destacamos al equipo responsable del dataset \textbf{SWEET}.  
  Durante el desarrollo de este TFM se ha establecido contacto formal con ellos y se ha alcanzado un \textbf{acuerdo para poder utilizar dicho dataset en los experimentos del proyecto}.  
  Para este equipo, cualquier avance metodológico o científico basado en su dataset aumenta la visibilidad, validación y relevancia de su trabajo, promoviendo su adopción por la comunidad y fomentando nuevas colaboraciones.
  %Asimismo, grupos como el equipo responsable del dataset \textbf{SWEET} —con el que se han mantenido contactos— 

  \item[Autoridades reguladoras y comités éticos:]  
  Tienen un papel fundamental debido a la naturaleza sensible de los datos tratados (señales fisiológicas y biométricas).  
  Su función no se limita a supervisar la protección de datos personales —cumpliendo la \textbf{LOPD} y el \textbf{GDPR}—, sino que abarca un conjunto más amplio de responsabilidades, tal como recoge el Comité de Ética en Investigación de la Universidad de Granada \cite{ugrcomite}:  
  \begin{itemize}
      \item Evaluar el bienestar, seguridad y dignidad de los participantes;  
      \item Garantizar la proporcionalidad entre riesgos y beneficios;  
      \item Supervisar el consentimiento informado y los procedimientos de anonimización;  
      \item Asegurar la trazabilidad, transparencia y gobernanza ética del proyecto.  
  \end{itemize}
  Su papel es imprescindible para que cualquier despliegue futuro en entornos clínicos o de investigación con sujetos humanos sea legítimo, seguro y ético.
 % Su objetivo es garantizar el cumplimiento de marcos legales como la \textbf{Ley Orgánica de Protección de Datos (LOPD)} y el \textbf{Reglamento General de Protección de Datos (GDPR)}, asegurando que el sistema respete los principios de privacidad por diseño, consentimiento informado y trazabilidad. 
  \item[Usuarios finales:]  
  Aunque el sistema desarrollado es de \textit{back-end}, los beneficiarios últimos de su implementación son varios perfiles:

  \begin{itemize}
      \item \textbf{Personas mayores:} beneficiarias del caso de uso de salud digital. El sistema busca mejorar su bienestar mediante la detección temprana de estados de estrés, facilitando intervenciones preventivas y personalizadas.
      \item \textbf{Cuidadores y personal sanitario:} potenciales usuarios de aplicaciones basadas en esta arquitectura. Su interés radica en recibir información agregada y anonimizada que les permita tomar decisiones clínicas o de acompañamiento más informadas, sin comprometer la privacidad de los pacientes.
      %\item \textbf{Analistas de datos e investigadores:} interesados en explotar la arquitectura como \textit{framework} para entrenar modelos colaborativos entre centros o dispositivos, reduciendo los riesgos asociados al intercambio de datos en crudo.
       \item \textbf{Analistas de datos e investigadores externos:}  
      No forman parte del equipo de investigación del TFM, pero pueden beneficiarse en el futuro tanto de la \textbf{arquitectura y \textit{framework} desarrollados} como de los \textbf{datos experimentales} generados dentro del proyecto financiado.  
      Este grupo puede emplear la infraestructura como base para investigar nuevos algoritmos de aprendizaje federado, evaluar modelos en escenarios heterogéneos o diseñar aplicaciones clínicas basadas en la información obtenida.  
  \end{itemize}
\end{description}

\section*{Conclusiones del análisis de stakeholders}%\section*{Síntesis del análisis}
En ausencia de un cliente comercial directo, los \textit{stakeholders} que influyen de forma más determinante en la definición de los requisitos funcionales y científicos del sistema son el \textbf{equipo de investigación} y el \textbf{equipo de desarrollo}.

El proyecto ha sido concebido con un enfoque \textbf{escalable, reutilizable y transferible}, orientado a su posible adopción futura por instituciones sanitarias, organismos públicos y comunidades científicas internacionales. Esto introduce un conjunto de \textbf{stakeholders potenciales secundarios}, que podrían beneficiarse de la solución en caso de que el sistema alcance el grado de madurez y robustez deseado.

En conjunto, cada grupo de interés aporta una perspectiva específica y complementaria, configurando un ecosistema de colaboración que fortalece la coherencia, la relevancia y la aplicabilidad del proyecto.

\chapter{Metodología de Desarrollo y Planificación del Proyecto}
\label{chap:metodologia}
 %En este capítulo se presentan brevemente las principales metodologías candidatas, tanto tradicionales como ágiles, para finalmente justificar la elección de la metodología adoptada en este TFM, así como la planificación del proyecto.
La metodología de trabajo adoptada para este TFM debe responder a las necesidades reales del proyecto: requisitos que evolucionan conforme avanza la investigación, necesidad de una comunicación frecuente con las tutoras, planificación flexible, y una visualización clara y continua del estado de las tareas. Por ello, en lugar de emplear metodologías tradicionales basadas en fases rígidas, se opta por un enfoque ágil apoyado en dos elementos principales: \textbf{Scrum} y tablero \textbf{Kanban}.

\section{Metodología seleccionada}

\subsection{Justificación del enfoque ágil}

A lo largo del desarrollo de este Trabajo Fin de Máster surgen de manera recurrente nuevos requisitos, se reformulan objetivos parciales y es necesario reajustar prioridades en función de los resultados experimentales obtenidos. Este contexto, caracterizado por la incertidumbre y la evolución progresiva del problema, hace poco adecuado un enfoque de planificación rígido y cerrado desde el inicio.

Por el contrario, este tipo de proyectos encaja de forma natural con los principios establecidos en el \textit{Agile Manifesto} \cite{agileprinciples}, que promueven la adaptación continua al cambio, la entrega incremental de resultados funcionales y la comunicación constante con los \textit{stakeholders}. En este caso, dichas partes interesadas están representadas fundamentalmente por las tutoras del trabajo.

La adopción de un enfoque ágil permite, en particular:

\begin{itemize}
    \item Incorporar cambios en los requisitos sin comprometer de forma rígida la planificación global del proyecto.  
    \item Trabajar mediante iteraciones cortas que producen resultados verificables al final de cada ciclo.  
    \item Mantener una comunicación periódica y estructurada con las tutoras para validar avances y redefinir prioridades.  
    \item Disponer en todo momento de una visión clara del estado del proyecto, reduciendo el riesgo de desviaciones prolongadas.  
\end{itemize}

Estas necesidades se alinean especialmente bien con los principios de \textbf{Scrum}, descritos por Schwaber y Sutherland \cite{schwaber2001agile}, y con la filosofía visual de \textbf{Kanban}, analizada en profundidad en el trabajo de Ahmad et al. \cite{ahmad2013kanban}. Aunque ninguna de estas propuestas fue concebida originalmente para desarrollos individuales, ambas pueden adaptarse de forma efectiva al contexto de un TFM, proporcionando estructura, control y flexibilidad.

\subsection{Scrum como marco de trabajo iterativo}

Scrum es un marco de trabajo orientado a la entrega incremental de valor mediante iteraciones breves y ciclos de retroalimentación continua \cite{schwaber2001agile}. En este proyecto se adopta un modelo de desarrollo ágil \textbf{inspirado en Scrum}, debido a las siguientes razones:

\begin{itemize}
    \item El desarrollo del TFM requiere una \textbf{revisión periódica} con las tutoras, análoga a las revisiones de \textit{sprint}.  
    \item Los requisitos no están completamente definidos desde el inicio y deben \textbf{refinarse de manera continua} conforme avanza la investigación.  
    \item La planificación por iteraciones permite establecer objetivos realistas, verificar resultados parciales y corregir desviaciones de forma temprana.  
    \item El trabajo puede dividirse en unidades manejables que integran análisis, diseño, implementación, validación y documentación.  
\end{itemize}

Dado que se trata de un desarrollo individual, se prescinde de la definición formal de roles como \textit{Scrum Master} o \textit{Product Owner}, manteniéndose únicamente la estructura funcional del ciclo de Scrum: planificación del \textit{sprint}, desarrollo, revisión y reajuste. La duración de cada \textit{sprint} se fija en aproximadamente dos semanas, un periodo adecuado para generar avances medibles y contrastables en cada iteración.

La planificación global del proyecto se articula en \textbf{hitos}, cada uno de los cuales agrupa la consecución total o parcial de distintas historias de usuario (HU) y requisitos no funcionales (RNF). Estas unidades constituyen el \textit{backlog} del proyecto y permiten estructurar el trabajo de forma incremental y trazable.

\subsection{Historias de usuario y estimación del esfuerzo}

Para facilitar la gestión y priorización del trabajo, cada hito se descompone en \textbf{historias de usuario}, entendidas como unidades funcionales que describen una necesidad concreta desde la perspectiva del usuario final o de un actor relevante del sistema.

\begin{itemize}
    \item \textbf{Historias de Usuario (HU):} describen funcionalidades o necesidades que aportan valor al sistema, expresadas de forma clara y orientada a objetivos.
\end{itemize}

En la práctica, cada HU seleccionada para un \textit{sprint} puede descomponerse en una o varias historias de menor tamaño, donde se detalla el trabajo necesario de análisis, diseño, implementación, pruebas, documentación y despliegue. Esta descomposición permite ajustar el alcance de cada iteración y mejorar la precisión en la planificación.

Cada historia o requisito se documenta siguiendo una estructura homogénea que garantiza la trazabilidad y la correcta priorización:

\begin{itemize} 
  \item \textbf{ID:} Identificador único formado por un prefijo (\emph{RNF}, \emph{HU}, \emph{HD}) y un número correlativo.  
  \item \textbf{Nombre:} Denominación breve y representativa.  
  \item \textbf{Descripción:} Explicación concisa de la funcionalidad o tarea.  
  \item \textbf{Prioridad:} Nivel de importancia (\textit{Alta}, \textit{Media} o \textit{Baja}).  
  \item \textbf{Necesidad:} Clasificación en \textit{Esencial}, \textit{Opcional} o \textit{Deseada}.  
  \item \textbf{Alcance:} Componentes o módulos del sistema afectados.  
  \item \textbf{Dependencias:} Historias que deben completarse previamente.  
  \item \textbf{Estado:} Situación actual (\textit{Pendiente}, \textit{En proceso} o \textit{Completado}).  
\end{itemize}

Para la estimación del esfuerzo se emplea un enfoque basado en \textbf{puntos de historia} (\emph{Story Points}), habitual en Scrum. Dado el carácter individual del proyecto y con el objetivo de facilitar la planificación temporal, se adopta una aproximación pragmática en la que \textbf{1 punto de historia (HP) se considera equivalente a 2 horas de trabajo efectivo}. Esta equivalencia permite relacionar directamente el tamaño de las historias con la capacidad real de trabajo en cada \textit{sprint}.

En función de esta estimación, se definen las siguientes categorías de tamaño:

\begin{table}[H]
\centering
\begin{tabular}{|c|c|}
\hline
\textbf{Tamaño} & \textbf{Esfuerzo (HP)} \\ \hline
XS & 2 \\ \hline
S  & 4 \\ \hline
M  & 8 \\ \hline
L  & 12 \\ \hline
XL & 16 \\ \hline
\end{tabular}
\caption{Categorías de esfuerzo en función del tamaño de la historia o tarea.}
\end{table}

Esta clasificación facilita una selección equilibrada de historias por \textit{sprint} y un seguimiento realista del progreso. Al cierre de cada iteración se realiza una \emph{sprint review} con las tutoras y una \emph{retrospectiva} para revisar el \textit{backlog} e identificar mejoras para el siguiente ciclo.

\subsection{Kanban como herramienta de visualización y control de flujo}

Kanban no constituye una metodología de desarrollo en sí misma, sino un \textbf{mecanismo visual de gestión del flujo de trabajo}. Tal y como señalan Ahmad et al. \cite{ahmad2013kanban}, su uso mejora la transparencia del proceso, limita el trabajo en curso y favorece un avance más continuo y controlado.

En este TFM, Kanban se emplea como herramienta complementaria a Scrum con los siguientes objetivos:

\begin{itemize}
    \item Mantener una \textbf{visualización clara e inmediata} del estado de cada historia de usuario.  
    \item Controlar el trabajo simultáneo y reducir la dispersión de esfuerzos.  
    \item Facilitar el seguimiento de los acuerdos alcanzados con las tutoras y el progreso real del proyecto.  
\end{itemize}

El tablero Kanban utilizado se estructura en las siguientes columnas:

\[
\textit{Backlog} \rightarrow \textit{Ready} \rightarrow \textit{In Progress} \rightarrow \textit{In Review} \rightarrow \textit{Done}.
\]

Este sistema proporciona una visión actualizada del flujo de trabajo, facilita la configuración de las reuniones periódicas y refuerza la coherencia entre la planificación de \textit{sprints} y la ejecución efectiva de las tareas.



\section{Planificación y Fases}
\label{sec:planificacion}
Una vez definida la Metodología de Desarrollo, se realizó la siguiente planificación inicial . El proyecto se planificó de tal forma que se pudiera entregar en la primera convocatoria de presentación de TFM's para nuestro máster, en Enero de 2026. 
\subsection{Planificación basada en entregables y criterios de completitud}

Siguiendo la metodología ágil adoptada, los entregables no se completan en fases cerradas, sino que se desarrollan y refinan de forma a lo largo de los sprints.  
Aun así, para disponer de una referencia temporal y facilitar el control del avance, se definen \textbf{fechas objetivo} para cada entregable, entendidas como el momento en el que debe estar lo suficientemente maduro para habilitar el siguiente bloque de trabajo.

A continuación se resumen los entregables principales, su fecha orientativa y los criterios de completitud:

\begin{itemize}

    \item \textbf{E1: Project Kick-off \& Initial Specification}  
    \textit{Fecha objetivo de completitud:} 31 de julio de 2025.  
    \textit{Criterios de completitud:}
    \begin{itemize}
        \item Objetivos generales y específicos formalizados y validados con las tutoras.
        \item Primer conjunto de requisitos no funcionales definido.
        \item Backlog inicial con historias de usuario de alto nivel para especificar requisitos funcionales.
    \end{itemize}

    \item \textbf{E2: Research \& Feasibility Review}  
    \textit{Fecha objetivo de completitud:} 15 de agosto de 2025.  
    \textit{Criterios de completitud:}
    \begin{itemize}
        \item Estado del arte en aprendizaje federado, sensores biométricos y \textit{computing continuum} documentado.
        \item Comparación técnica de frameworks completada (incluyendo Flower).
        \item Dataset seleccionado y validado para uso en escenarios federados.
        \item Primera propuesta de arquitectura viable.
    \end{itemize}

    \item \textbf{E3: MVP — Minimal Viable Prototype}  
    \textit{Fecha objetivo de completitud:} 30 de septiembre de 2025.  
    \textit{Criterios de completitud:}
    \begin{itemize}
        \item Flujo \textit{end–to–end} básico ejecutándose: cliente → broker → servidor.
        \item Preprocesamiento local mínimo implementado.
        \item Entrenamiento federado simple funcionando en entorno controlado.
        \item Métricas iniciales recogidas para validar la viabilidad del enfoque.
    \end{itemize}

    \item \textbf{E4: System Design Consolidation}  
    \textit{Fecha objetivo de completitud:} 30 de octubre de 2025.  
    \textit{Criterios de completitud:}
    \begin{itemize}
        \item Diagramas de arquitectura y flujos de datos refinados.
        \item Definición clara de módulos, responsabilidades y puntos de interacción.
        \item Diseño del pipeline de CI/CD funcional.
        \item Identificación completa de historias de usuario necesarias para la implementación del sistema final.
    \end{itemize}

    \item \textbf{E5: Functional Increments Completed}  
    \textit{Fecha objetivo de completitud:} 21 de noviembre de 2025.  
    \textit{Criterios de completitud:}
    \begin{itemize}
        \item Implementación integrada de todos los módulos principales (FL-agent, broker, agregación Fog, servidor Cloud).
        \item Ejecución federada estable bajo distintos escenarios (heterogeneidad, nodos caídos, cargas variables).
        \item Sistema monitorizable y reproducible.
        \item configuración del entorno para pruebas y validación final.
    \end{itemize}

    \item \textbf{E6: Final Documentation \& Delivery}  
    \textit{Fecha objetivo de completitud:} 22 de diciembre de 2025.  
    \textit{Criterios de completitud:}
    \begin{itemize}
        \item Documentación técnica completada (APIs, despliegue, experimentos, arquitectura).
        \item Comparativa y análisis de resultados finales.
        \item Memoria del TFM completamente redactada.
        \item configuración de la demo, anexos y materiales finales de defensa.
    \end{itemize}

\end{itemize}

\subsection{Dedicación real diagrama de Gantt y correspondencia con la planificación inicial}
En términos generales, la planificación se ha cumplido en los \textbf{hitos principales} definidos al inicio del proyecto, completándose conforme a lo previsto \textbf{todos salvo los dos últimos}. En concreto, el hito \textbf{E6} se cerró con aproximadamente \textbf{una semana de retraso}, y el cierre global del proyecto se extendió \textbf{casi un mes y medio} respecto a la fecha objetivo.

Este desfase no respondió a un bloqueo técnico puntual, sino a una \textbf{decisión deliberada de priorización}: se optó por \textbf{no comprometer} los entregables ya planificados y alcanzados en los sprints anteriores, concentrando el esfuerzo adicional en la fase final de cierre. En ese tramo, la carga de trabajo resultó \textbf{significativamente superior} a la estimada, principalmente por:
(i) la \textbf{incorporación de una cantidad elevada de contenido} a la memoria,
(ii) la \textbf{construcción y consolidación de materiales finales} (resultados, comparativas, anexos y soporte para demo), y
(iii) la \textbf{redacción sistemática y homogénea de los sprints}, asegurando \textbf{coherencia global}, \textbf{trazabilidad} entre historias e hitos, y \textbf{calidad editorial} del documento.

Como consecuencia, se mantuvo el cumplimiento de los objetivos funcionales y técnicos del sistema, asumiendo el ajuste temporal en las últimas fases para garantizar una \textbf{entrega final consistente, completa y verificable}.

\begin{figure}[t]
    \centering
    \includegraphics[width=1\linewidth]{Plantilla-Latex-TFG//imagenes/gantt.png}
    \caption{Dedicación real y correspondencia con la planificación inicial (diagrama de Gantt).}
    \label{fig:gantt_real_vs_plan}
\end{figure}


\chapter{Desarrollo del Proyecto}
\label{chap:desarrollo}
Este capítulo describe cómo se ha llevado a cabo el desarrollo del TFM desde una perspectiva práctica. 
En primer lugar, se presenta el \textit{Sprint 0}, entendido como fase de configuración del proyecto, donde se definen las bases conceptuales, tecnológicas y organizativas. 
Dentro del \textit{Sprint 0} se realiza el análisis de requisitos, en el que se concreta qué se espera del sistema en términos funcionales y de calidad. 
Sobre esta base, se construye el \textit{backlog} mediante historias de usuario y tareas dentro de cada \textit{sprint}. 
Finalmente, en las secciones posteriores del capítulo se detalla el desarrollo iterativo de los distintos \textit{sprints}.

\section{Propuesta de solución}
A partir de los objetivos definidos en los capítulos anteriores, este Trabajo de Fin de Máster plantea el diseño e implementación de un \textbf{prototipo de sistema de aprendizaje federado jerárquico} orientado a la detección de estrés a partir de señales fisiológicas. El trabajo se centra en el diseño de la arquitectura, los mecanismos de coordinación y el ciclo de entrenamiento federado, dejando fuera del alcance la instrumentación física de sensores, pero manteniendo el sistema preparado para su integración futura en un entorno real.

La solución propuesta se basa en una arquitectura distribuida \textbf{Edge--Fog--Cloud}, concebida dentro del paradigma del \textit{computing continuum}. Esta arquitectura permite entrenar modelos de \textit{machine learning} de forma colaborativa sin centralizar los datos biométricos y, al mismo tiempo, ofrece un marco adecuado para estudiar el papel de los nodos Fog como capa intermedia en sistemas federados.

\begin{figure}[H]
  \centering
  \includegraphics[
    width=0.7\textwidth,
    height=0.7\textheight
  ]{Plantilla-Latex-TFG/imagenes/image.png}
  \caption{Arquitectura general del sistema a nivel de nodo dentro del \textit{computing continuum} (sensores en la parte inferior y Cloud en la superior).}
  \label{fig:GArquitecture}
\end{figure}

La Figura~\ref{fig:GArquitecture} resume la arquitectura considerada. Las fuentes de datos se sitúan en la \textbf{parte inferior} y corresponden a sensores y dispositivos \textit{wearables}. En el contexto de este TFM, estos dispositivos se modelan como fuentes abstractas de datos, utilizándose \textit{datasets} de estrés previamente validados para simular el comportamiento de los datos de los clientes. De este modo, el ciclo de aprendizaje federado se inicia directamente a partir de los datos disponibles en cada nodo Edge.

El procesamiento de las señales y el entrenamiento local se realizan en los nodos \textbf{Edge}, garantizando que los datos biométricos no abandonan el ámbito local. Las actualizaciones generadas se envían \textbf{hacia arriba} a los nodos \textbf{Fog}, que actúan como agregadores intermedios y reducen la carga de comunicación hacia el nivel superior. El nivel \textbf{Cloud} coordina el proceso global, gestiona las rondas de entrenamiento y mantiene el modelo federado resultante.

El flujo representado en la figura refleja el carácter iterativo del sistema: las actualizaciones locales \textbf{ascienden} por las capas del continuum, donde son agregadas, y el modelo resultante se \textbf{redistribuye descendiendo} nuevamente hacia los nodos Edge. Este ciclo constituye el núcleo funcional del sistema y define el alcance principal del trabajo, más allá de una simple simulación puntual de clientes.

En esta arquitectura, los nodos \textbf{Edge} representan dispositivos cercanos al usuario responsables del entrenamiento local; los nodos \textbf{Fog} introducen una capa intermedia de agregación y coordinación, pudiendo asumir tareas adicionales como las ya comentadas en el estado del arte; y el nivel \textbf{Cloud} actúa como orquestador global del proceso federado y punto de consolidación del modelo.

Sobre esta base se construye un ciclo de aprendizaje federado que permite entrenar modelos de detección de estrés a partir de \textit{datasets} públicos, emulando escenarios realistas de distribución no homogénea entre clientes, habituales en entornos IoT y de salud conectada.

Uno de los objetivos centrales del trabajo es evaluar de forma crítica el uso de nodos Fog en sistemas de aprendizaje federado. Aunque la bibliografía sugiere que su incorporación puede mejorar la escalabilidad y reducir la latencia, estas ventajas conllevan un aumento de la complejidad del sistema. Por ello, este TFM analiza de manera experimental el impacto de esta capa intermedia sobre la convergencia, la eficiencia de comunicación y la robustez del sistema, abordando además el reto técnico de diseñar e integrar una arquitectura federada con este nivel de complejidad.


\section{Sprint 0}

\label{sect:Sp0}
Enmarcado como una fase previa al desarrollo propiamente dicho, el \textit{Sprint 0} tiene por objetivo preparar el terreno del proyecto ágil, por lo que no es un Sprint al uso. 
Según la literatura técnica, “Sprint 0 es la fase inicial que prepara al equipo, el entorno técnico y el backlog de producto para asegurar un comienzo eficiente” \cite{Sprint0_Definition_PMProjectManagement}. 
Este enfoque responde a la necesidad de acometer acciones clave (como la definición del backlog inicial, la configuración de entornos de desarrollo e infraestructuras y el alineamiento con los \textit{stakeholders}) antes de iniciar los sprints de desarrollo formal según el marco Scrum \cite{Sprint0_Definition_MikeCohn}.  

En el presente TFM, el \textbf{Sprint 0 dio comienzo en la primera semana de abril de 2025}, realizándose las primeras reuniones con las tutoras. 
Durante dichas reuniones, se identificaron los posibles temas de interés para el proyecto y, tras un proceso de discusión inicial, se definieron las primeras tareas fundamentales que se enumeran a continuación:  

\subsection{Tareas realizadas en el Sprint 0}

\begin{itemize}
    \item T-001: Definición de los objetivos generales y específicos del TFM.
    \item T-002: Planificación inicial del proyecto y definición del roadmap de sprints.
    \item T-003: Revisión del estado del arte en aprendizaje federado.
    \item T-004: Revisión del estado del arte en datasets para detección de estrés.
    \item T-005: Revisión del estado del arte en dispositivos biométricos.
    \item T-006: Estudio del paradigma de \textit{Computing Continuum}.
    \item T-009: Estudio y preselección de tecnologías y herramientas.
    \item T-018: Recopilación de requisitos funcionales del sistema.
    \item T-019: Recopilación de requisitos no funcionales del sistema.
    \item T-020: Definición de la arquitectura preliminar del sistema federado.
    \item T-021: Solicitud de acceso al dataset SWEET.
    \item T-022: Diseño y refinamiento de la arquitectura federada (v1 y v2).
    \item T-023: Elección del lenguaje de desarrollo principal.
\end{itemize}
\subsection{Desarrollo del Sprint 0}

Como se ha indicado anteriormente, el Sprint 0 posee un carácter preparatorio y no produce incrementos funcionales del software. Su objetivo es establecer las bases conceptuales, investigadoras, organizativas y técnicas del TFM.

Durante los meses iniciales (abril–junio), coincidiendo con la finalización del curso académico, se llevaron a cabo múltiples reuniones quincenales con las tutoras. En estas sesiones se definió el alcance del proyecto, se revisó el estado del arte, se evaluaron tecnologías y se consolidaron los componentes fundamentales del backlog de producto. Esta fase permitió asentar una visión clara del sistema de aprendizaje federado a desarrollar y determinar sus requisitos funcionales y no funcionales.

A continuación se detallan las tareas completadas durante el Sprint 0:
%añadir las referencias en donde se reflejan las diferentes cosas
\begin{itemize}
  \item \textbf{T-001 – Definición de objetivos del TFM:}  
  Se establecieron los objetivos generales y específicos del proyecto, centrados en el diseño de un sistema de aprendizaje federado jerárquico bajo arquitectura Edge–Fog–Cloud para la detección de estrés mediante señales biométricas.  
  Estos objetivos se definieron tras varias sesiones con las tutoras, y se encuentran detallados en el apartado correspondiente de esta memoria.

  \item \textbf{T-002 – Planificación inicial del proyecto:}  
  Se elaboró el roadmap inicial del TFM, el cronograma preliminar de los sprints y la estructura del backlog.  
  La planificación incluyó dependencias, riesgos iniciales, hitos y priorización de tareas, y se presenta en detalle en el apartado de estimación temporal.

  \item \textbf{T-003 – Revisión del estado del arte en aprendizaje federado:}  
  Se realizó una revisión bibliográfica del aprendizaje federado, cubriendo su motivación inicial, el algoritmo FedAvg y la evolución hacia arquitecturas jerárquicas y escenarios del \textit{computing continuum}.  
Los resultados de esta revisión se recogen en el Capítulo~\ref{chap:estado-del-arte}, Sección~\ref{sec:fl}, incluyendo la descripción del funcionamiento general del FL, la agregación jerárquica y las tendencias actuales, así como una clasificación de algoritmos (Sección~\ref{sec:fl-categories} y Tabla~\ref{tab:fl-categorias}).  


  \item \textbf{T-004 – Revisión del estado del arte en datasets para estrés:}  
  Se analizaron \textit{datasets} de referencia para detección de estrés mediante señales fisiológicas, considerando señales disponibles, protocolos experimentales, calidad de etiquetado y adecuación a escenarios no-i.i.d. La discusión comparativa de \textit{datasets} y sus implicaciones para personalización y aprendizaje federado se presenta en la Sección~\ref{sec:datasets-stress}.

  

\item \textbf{T-005 – Revisión del estado del arte en dispositivos biométricos:}  
Se revisaron sensores y dispositivos \textit{wearables} relevantes para la detección de estrés, evaluando señales fisiológicas clave y consideraciones de usabilidad y validez.  
Esta revisión se integra en la Sección~\ref{sec:datasets-stress}, conectando disponibilidad de señales y limitaciones prácticas con los requisitos de preprocesamiento local y particionado no-i.i.d.


    
  \item \textbf{T-006 – Revisión del estado del arte en Continuum Computing:}  
Se estudió el paradigma del \textit{computing continuum} (Cloud--Fog--Edge) y su relación con sistemas distribuidos y aprendizaje federado jerárquico.  
Los conceptos fundamentales, objetivos y beneficios para FL se presentan en la Sección~\ref{sec:continuum}.

  \item \textbf{T-018 – Recopilación de requisitos funcionales:}  
  Se definieron los requisitos funcionales mediante las HU-001 a la HU-015 basados en la literatura y en las sesiones con las tutoras \ref{Analisisderequisitos}.

  \item \textbf{T-019 – Recopilación de requisitos no funcionales:}  
  Se formalizaron los RNF relativos a seguridad, privacidad, rendimiento, mantenimiento y calidad del código \ref{RequisitosNofuncionales}.

  \item \textbf{T-020 – Definición de la arquitectura preliminar:}  
  Se elaboró el esquema inicial del sistema federado jerárquico y sus flujos de datos.

  \item \textbf{T-009 – Estudio de tecnologías y herramientas:}  
  Se analizó qué candidatos como \textit{frameworks} se podían utilizar para el desarrollo e implementación de este sistema para poder hacer el correspondiente estudio comparativo entre ellos, permitiendo un sistema federado (Flower, PyTorch), sistemas de CI/CD, herramientas de formateo, control de versiones y entornos de despliegue.

\item \textbf{T-023 – Elección de Python como lenguaje de desarrollo:}  

Tras evaluar alternativas (p.\,ej., C++ o TypeScript), se seleccionó \textbf{Python} como lenguaje principal por ser el que mejor equilibra los criterios definidos en la Sección~\ref{sec:lenguajesdes} y la Tabla~\ref{tab:comparativa-lenguajes}: ecosistema maduro de \textit{machine learning}, disponibilidad de frameworks de aprendizaje federado (p.\,ej., \texttt{Flower}, \texttt{TensorFlow Federated}), prototipado ágil e integración sencilla en entornos distribuidos. Esta elección guió la definición de historias de usuario orientadas a validar prototipos federados en los siguientes \textit{sprints}.


\item \textbf{T-021 – Solicitud de acceso al dataset SWEET:}  

Dado el interés en emplear el dataset \textit{SWEET} como uno de los conjuntos de datos principales debido a su atractivo por la cantidad de sujetos que contiene, se llevó a cabo el proceso de solicitud formal de acceso a dicho recurso.  
Al tratarse de un dataset de carácter privado y no disponible de forma abierta, fue necesario contactar directamente con uno de los autores del estudio original, en concreto con \textbf{Jan Cornelis}, coautor del experimento y miembro de la empresa \textit{imec} (Kapeldreef 75, 3001 Leuven, Bélgica), con el fin de solicitar autorización para su uso con fines académicos (autorización que como veremos en el futuro  fué aprobada).

\item \textbf{T-022 – Diseño de la arquitectura del sistema federado:}
En esta tarea se definió la base arquitectónica del sistema de aprendizaje federado con enfoque jerárquico \textit{Edge--Fog--Cloud}. Partiendo de una concepción inicial de alto nivel se evolucionó hacia una \textbf{arquitectura v2} más operativa y adecuada para guiar la implementación incremental en sprints posteriores (Figura~\ref{fig:Arquitecturav2}).

La \textbf{v2} explicita de forma clara:
\begin{itemize}\setlength\itemsep{0.15em}
    \item los \textbf{flujos de información} entre capas (datos biométricos, actualizaciones locales y modelos agregados);
    \item el \textbf{ciclo iterativo} del aprendizaje federado, con redistribución continua del modelo global hacia los nodos Edge;
    \item la \textbf{separación de responsabilidades} por nivel (Edge: entrenamiento local; Fog: agregación intermedia; Cloud: orquestación y agregación global);
    \item los mecanismos de \textbf{comunicación y orquestación} necesarios para soportar un despliegue realista (p.\,ej., broker y servicios de agregación), sin fijar todavía tecnologías concretas.
\end{itemize}

Además, los \textit{wearables} se modelan como fuentes externas de datos, lo que puede entenderse como una cuarta capa de datos. De este modo, queda delimitado el alcance del TFM al plano software del sistema federado (entrenamiento local, comunicación y agregación), pero se deja explícitamente preparada la integración futura de dichos dispositivos sin reestructurar la arquitectura.

La propuesta se mantiene deliberadamente \textbf{agnóstica a herramientas, proveedores y frameworks}, evitando acoplamientos prematuros (p. ej., a un broker concreto, un stack cloud específico o una librería de FL determinada). Esto permite preservar modularidad e intercambiabilidad de componentes y habilita una selección tecnológica informada en sprints posteriores, basada en requisitos ya validados (rendimiento, latencia, consumo, seguridad, coste y mantenibilidad).
\end{itemize}

\begin{figure}[H]
  \centering
  \includegraphics[
    width=1\textwidth
  ]{Plantilla-Latex-TFG/imagenes/Arquitecturav2.png}
  \caption{Arquitectura del sistema v2 a nivel de nodo}
  \label{fig:Arquitecturav2}
\end{figure}



\subsection{Backlog  y Requisitos no Funcionales}
\label{Analisisderequisitos}

El análisis de requisitos constituye una fase esencial para garantizar que el sistema desarrollado satisfaga las necesidades funcionales y los criterios de calidad del proyecto. En este TFM se adopta un enfoque inspirado en Scrum.

Los \textbf{requisitos funcionales} no se enumeran como RF clásicos, sino que se representan mediante un catálogo estructurado de \textbf{historias de usuario (HU)}, al ser el mecanismo estándar en entornos ágiles para modelar funcionalidades. Mantener RF y HU en paralelo introduciría duplicidades e inconsistencias; por ello, en este proyecto las HU \textbf{asumen completamente el papel de requisitos funcionales}.

En cambio, los \textbf{requisitos no funcionales (RNF)} se mantienen como categoría diferenciada, ya que describen propiedades de calidad (seguridad, rendimiento, mantenibilidad, fiabilidad, interoperabilidad y calidad del código) que no se modelan adecuadamente mediante HU.

Tanto las HU como los RNF se documentan siguiendo IEEE 830~\cite{ieee830} para asegurar requisitos claros, verificables y trazables. En la práctica, esto se plasma en un catálogo de HU con campos de seguimiento (p.\,ej., prioridad, alcance y dependencias) y en RNF definidos mediante criterios y métricas comprobables, evitando redundancias y facilitando el control por \textit{sprints}.


\medskip


\subsubsection{Requisitos No Funcionales}
\label{RequisitosNofuncionales}
Los requisitos no funcionales describen las cualidades, restricciones y criterios de calidad que el sistema debe cumplir para garantizar su correcto funcionamiento, mantenimiento y adopción. A continuación se enumeran los requisitos no funcionales definidos para este proyecto:
\begin{itemize}

  \item \textbf{[RNF-001] Seguridad de comunicaciones y secretos.}  
  Todas las comunicaciones entre componentes (Edge--Fog--Servidor) deben ir cifradas y autenticadas, evitando exposición de credenciales.
  \textit{Criterios de verificación:}
  \begin{itemize}
    \item Todo tráfico de control y APIs externas debe usar \textbf{TLS 1.3} (o superior).
    \item Las credenciales (\textit{tokens}, \textit{API keys}, contraseñas) deben inyectarse por \textbf{variables de entorno} o \textbf{secretos}, y \textbf{no} deben aparecer en logs.
    \item El sistema debe soportar rotación de credenciales sin recompilar (reinicio controlado del servicio).
  \end{itemize}

  \item \textbf{[RNF-002] Privacidad y minimización de datos (GDPR/LOPD).}  
  El sistema debe preservar la privacidad evitando la salida de datos crudos y minimizando metadatos identificables.
  \textit{Criterios de verificación:}
  \begin{itemize}
    \item Los datos biomédicos crudos \textbf{no salen del nodo Edge} (validación por inspección de logs y tráfico).
    \item Los sujetos se referencian mediante \textbf{identificadores pseudonimizados} (p.\,ej., \textit{subject\_id} aleatorio), sin nombre/DNI/email en manifiestos ni resultados.
    \item Los artefactos exportados (métricas, manifiestos, resultados) deben excluir campos identificables y documentar su esquema.
  \end{itemize}

  \item \textbf{[RNF-003] Calidad del modelo (rendimiento y no-regresión).}  
  El entrenamiento federado debe producir un modelo global útil, sin degradaciones severas frente a referencias razonables.
  \textit{Criterios de verificación:}
  \begin{itemize}
    \item Para una misma configuración y partición de datos, el rendimiento del modelo global (métrica principal definida por el experimento) debe quedar \textbf{a menos de 15 puntos porcentuales} del mejor modelo local.
    \item En ejecuciones repetidas con la misma \textit{seed} y configuración, la variación del rendimiento final no debe exceder \textbf{$\pm 6$ puntos porcentuales} (cuando aplique).
  \end{itemize}

\item \textbf{[RNF-004] Reproducibilidad experimental (configuración declarativa).}  
Cada \textbf{arquitectura federada} debe poder reconstruirse de forma \textbf{determinista} a partir de su fichero de configuración y del manifiesto generado, garantizando que un mismo escenario se despliega siempre con la misma topología, roles y asignaciones.
\textit{Criterios de verificación:}
\begin{itemize}
  \item Toda ejecución debe generar un \textbf{manifiesto de experimento} que incluya, como mínimo: versión del código/imagen, configuración efectiva (resuelta), \textit{seed} empleada, topología (Edge/Fog/Servidor), asignación de sujetos a nodos y parámetros de entrenamiento.
  \item A partir del mismo manifiesto, debe ser posible \textbf{relanzar} la ejecución obteniendo \textbf{la misma arquitectura y plan de asignación} (sujetos$\rightarrow$clientes$\rightarrow$Fog, umbrales y parámetros), sin intervención manual.
\end{itemize}

\textit{Nota:} debido a la naturaleza distribuida del sistema y a la existencia de \textbf{condiciones de carrera} asociadas al umbral $K$ (número de clientes requeridos antes de activar la agregación en un nodo Fog), el \textbf{orden temporal exacto} de llegada de mensajes y, por tanto, el trazado fino de ciertos eventos de ronda puede variar entre ejecuciones. Por ello, este requisito se orienta a asegurar la \textbf{reproducibilidad estructural} del expe

  \item \textbf{[RNF-005] Observabilidad nativa (métricas, logs y trazas).}  
  El sistema debe exponer telemetría para diagnóstico de rendimiento y fallos (sin depender de depuración manual).
  \textit{Criterios de verificación:}
  \begin{itemize}
    \item Publicación de métricas clave (p.\,ej., latencias, throughput, estado de ronda, contadores de error) en Prometheus y paneles básicos en Grafana.
    \item Trazas distribuidas o correlación de eventos (p.\,ej., \textit{round\_id}) para seguir el flujo Edge--Fog--Servidor.
    \item Logs estructurados con niveles (\textit{INFO/WARN/ERROR}) y sin datos sensibles.
  \end{itemize}

  \item \textbf{[RNF-006] Fiabilidad y tolerancia a desconexiones.}  
  El sistema debe mantener estabilidad ante desconexiones de clientes/nodos Fog y fallos transitorios.
  \textit{Criterios de verificación:}
  \begin{itemize}
    \item La detección de nodo no disponible (por \textit{heartbeat}/timeout) debe producirse en menos tiempo \textbf{t}.
    \item El sistema debe continuar la ronda si la pérdida de nodos Fog no supera el \textbf{25\%}; si supera ese umbral, debe abortar de forma controlada y registrar el motivo.
    \item Tras un fallo recuperable, el sistema debe dejar el estado en condición consistente (sin rondas ``a medias'').
  \end{itemize}

  \item \textbf{[RNF-007] Rendimiento (latencia y coste operativo).}  
  El sistema debe completar rondas con tiempos razonables en entornos de laboratorio (máquinas personales/VM).
  \textit{Criterios de verificación:}
  \begin{itemize}
    \item La sobrecarga de instrumentación (métricas/trazas) no debe aumentar el tiempo medio de ronda en más de \textbf{15\%} frente a ejecución sin telemetría.
    \item La serialización y envío de parámetros debe mantenerse por debajo de un tamaño máximo configurado (p.\,ej., \textbf{50--200 MB} por ronda, según arquitectura/modelo).
  \end{itemize}

  \item \textbf{[RNF-008] Mantenibilidad y calidad del código}  
  El sistema debe permitir la incorporación de nuevas funcionalidades sin afectar al funcionamiento global.  
  \textit{Criterios de verificación:}
  \begin{itemize}
    \item Se debe alcanzar una cobertura de pruebas unitarias superior al 70\,\%.
  \item \textbf{Linting (PEP~8):} el análisis estático del código debe reportar \textbf{0 violaciones severas} del estándar de estilo, garantizando consistencia y legibilidad \cite{pep8}.
  \item \textbf{Principios SOLID:} la implementación debe respetar SOLID para mejorar \textbf{mantenibilidad}, \textbf{modularidad} y facilidad de extensión/refactorización del sistema \cite{martin2008clean}.
  \end{itemize}
  
  \item \textbf{[RNF-009] Portabilidad y despliegue reproducible.}  
  El sistema debe desplegarse de forma consistente mediante contenedores, reduciendo dependencias del entorno.
  \textit{Criterios de verificación:}
  \begin{itemize}
    \item Despliegue soportado con contenedores (p.\,ej., Docker/Podman) y \textit{compose} equivalente.
    \item Un despliegue ``mínimo'' (1 servidor + 1 Fog + 1 Edge) debe completarse en \textbf{$\leq 15$ minutos} siguiendo documentación.
  \end{itemize}

  \item \textbf{[RNF-010] Interoperabilidad (APIs y mensajería).}  
  El sistema debe integrarse con herramientas externas y facilitar exportación de resultados.
  \textit{Criterios de verificación:}
  \begin{itemize}
    \item APIs de control/consulta con contratos claros (JSON) y códigos HTTP correctos cuando aplique.
    \item La mensajería interna debe documentar \textbf{tópicos/canales}, esquema de mensajes y versiones.
    \item Exportación de configuración y resultados en \textbf{JSON/YAML}.
  \end{itemize}

  \item \textbf{[RNF-011] Usabilidad en CLI y operabilidad.}  
  El sistema debe ser usable sin interfaz gráfica, priorizando DX (developer experience) y errores accionables.
  \textit{Criterios de verificación:}
  \begin{itemize}
    \item Comandos con \texttt{--help}, ejemplos mínimos y mensajes de error accionables (incluyendo \textit{exit codes}).
    \item Parámetros críticos configurables sin tocar código (ficheros de config + variables de entorno).
  \end{itemize}

  \item \textbf{[RNF-012] Idioma y documentación.}  
  El sistema y la documentación técnica asociada deben estar en inglés.
  \textit{Criterios de verificación:}
  \begin{itemize}
    \item README y guías de ejecución en inglés; nombres de métricas y logs en inglés.
  \end{itemize}

\end{itemize}




\subsubsection{Historias de usuario}
\label{HistoriasdeUsuario}

Cada historia de usuario se documenta con los siguientes campos:

\begin{itemize}\setlength\itemsep{0.15em}
    \item \textbf{ID}: identificador único.
    \item \textbf{Nombre}: título breve de la funcionalidad.
    \item \textbf{Descripción}: formulación en lenguaje natural desde la perspectiva del actor.
    \item \textbf{Prioridad/Necesidad}: peso relativo (Alta/Media/Baja; Esencial/Importante/Deseada).
    \item \textbf{Tamaño}: esfuerzo estimado (XS--XL) en puntos de historia (HP).
    \item \textbf{Alcance}: componentes y dimensiones implicadas (Edge/Fog/Cloud, privacidad, monitorización, etc.).
    \item \textbf{Dependencias}: historias previas necesarias.
    \item \textbf{Estado}: completado / en proceso / pendiente /bloqueada.
\end{itemize}

A continuación se listan las historias del \textit{backlog} y su estado final al cierre del proyecto:

\begin{table}[H]
\centering
\caption{Descripción de la historia de usuario HU-001}
\label{tab:hu001}
\begin{tabular}{|l|p{12cm}|}
\hline
\textbf{ID} & HU-001 \\ \hline
\textbf{Tamaño} & M \\ \hline
\textbf{Nombre} & Preprocesamiento local automático en el Edge \\ \hline
\textbf{Descripción} & 
Como usuario del sistema, quiero que cada nodo Edge pueda preprocesar mis datos automáticamente (limpieza, normalización y anonimización) para garantizar que la información enviada al servidor federado cumple los requisitos de calidad y privacidad. \\ \hline
\textbf{Prioridad} & Alta \\ \hline
\textbf{Necesidad} & Esencial \\ \hline
\textbf{Alcance} & Módulo de preprocesamiento en nodos Edge; calidad de datos; seguridad y privacidad de la información biomédica. \\ \hline
\textbf{Dependencias} & - \\ \hline
\textbf{Estado} & Completado \\ \hline
\end{tabular}
\end{table}



\begin{table}[H]
\centering
\caption{Descripción de la historia de usuario HU-002}
\label{tab:hu002}
\begin{tabular}{|l|p{12cm}|}
\hline
\textbf{ID} & HU-002 \\ \hline
\textbf{Tamaño} & L \\ \hline
\textbf{Nombre} & Incorporación segura de nodos al sistema \\ \hline
\textbf{Descripción} & 
Como administrador del sistema, quiero que los nodos Edge y Fog se registren y autentiquen automáticamente mediante certificados o tokens, para poder garantizar que solo dispositivos de confianza participen en el entrenamiento federado. \\ \hline
\textbf{Prioridad} & Alta \\ \hline
\textbf{Necesidad} & Esencial \\ \hline
\textbf{Alcance} & Gestión de registro y autenticación de nodos; seguridad de acceso; fiabilidad operativa frente a nodos maliciosos o no autorizados. \\ \hline
\textbf{Dependencias} & HU-001 \\ \hline
\textbf{Estado} & Pendiente \\ \hline
\end{tabular}
\end{table}



\begin{table}[H]
\centering
\caption{Descripción de la historia de usuario HU-003}
\label{tab:hu003}
\begin{tabular}{|l|p{12cm}|}
\hline
\textbf{ID} & HU-003 \\ \hline
\textbf{Tamaño} & M \\ \hline
\textbf{Nombre} & Configuración de capacidades Fog \\ \hline
\textbf{Descripción} & 
Como operador del sistema, quiero poder activar diferentes capacidades en los nodos Fog (agregación parcial, filtrado, cacheo de modelos, balanceo, adición de ruido) mediante archivos JSON/YAML para adaptar la red federada a diferentes escenarios de carga. \\ \hline
\textbf{Prioridad} & Alta \\ \hline
\textbf{Necesidad} & Importante \\ \hline
\textbf{Alcance} & Nodos Fog como capa intermedia; configuración modular de agregación y filtrado; mantenibilidad y extensibilidad de la arquitectura. \\ \hline
\textbf{Dependencias} & HU-002 \\ \hline
\textbf{Estado} & Completado \\ \hline
\end{tabular}
\end{table}



\begin{table}[H]
\centering
\caption{Descripción de la historia de usuario HU-004}
\label{tab:hu004}
\begin{tabular}{|l|p{12cm}|}
\hline
\textbf{ID} & HU-004 \\ \hline
\textbf{Tamaño} & M \\ \hline
\textbf{Nombre} & Distribución del modelo global \\ \hline
\textbf{Descripción} & 
Como cliente federado, quiero recibir automáticamente el modelo global actualizado tras cada ronda para poder continuar el entrenamiento con la versión más reciente. \\ \hline
\textbf{Prioridad} & Alta \\ \hline
\textbf{Necesidad} & Esencial \\ \hline
\textbf{Alcance} & Servidor central de orquestación; mecanismo de distribución de modelos; actualización coordinada de modelos locales en clientes y nodos Fog. \\ \hline
\textbf{Dependencias} & HU-001, HU-002 \\ \hline
\textbf{Estado} & Completado \\ \hline
\end{tabular}
\end{table}



\begin{table}[H]
\centering
\caption{Descripción de la historia de usuario HU-005}
\label{tab:hu005}
\begin{tabular}{|l|p{12cm}|}
\hline
\textbf{ID} & HU-005 \\ \hline
\textbf{Tamaño} & M \\ \hline
\textbf{Nombre} & Selección inteligente de clientes \\ \hline
\textbf{Descripción} & 
Como servidor federado, quiero seleccionar automáticamente los clientes más adecuados para cada ronda en función de disponibilidad, reputación, ancho de banda o heterogeneidad para mejorar la calidad del entrenamiento. \\ \hline
\textbf{Prioridad} & Alta \\ \hline
\textbf{Necesidad} & Importante \\ \hline
\textbf{Alcance} & Algoritmo de selección de clientes; políticas de disponibilidad y reputación; eficiencia del entrenamiento y estabilidad de la convergencia. \\ \hline
\textbf{Dependencias} & HU-003, HU-004 \\ \hline
\textbf{Estado} & Pendiente \\ \hline
\end{tabular}
\end{table}



\begin{table}[H]
\centering
\caption{Descripción de la historia de usuario HU-006}
\label{tab:hu006}
\begin{tabular}{|l|p{12cm}|}
\hline
\textbf{ID} & HU-006 \\ \hline
\textbf{Tamaño} & M \\ \hline
\textbf{Nombre} & Monitorización de métricas \\ \hline
\textbf{Descripción} & 
Como investigador, quiero ver métricas de entrenamiento en tiempo real (pérdida, latencias, \textit{throughput}, convergencia) para evaluar la estabilidad y rendimiento del sistema. \\ \hline
\textbf{Prioridad} & Media \\ \hline
\textbf{Necesidad} & Importante \\ \hline
\textbf{Alcance} & Subsistema de monitorización y métricas; panel de observabilidad; análisis de rendimiento del entrenamiento federado. \\ \hline
\textbf{Dependencias} & HU-004 \\ \hline
\textbf{Estado} & Completado \\ \hline
\end{tabular}
\end{table}



\begin{table}[H]
\centering
\caption{Descripción de la historia de usuario HU-007}
\label{tab:hu007}
\begin{tabular}{|l|p{12cm}|}
\hline
\textbf{ID} & HU-007 \\ \hline
\textbf{Tamaño} & M \\ \hline
\textbf{Nombre} & Modo síncrono/asíncrono configurable \\ \hline
\textbf{Descripción} & 
Como administrador del sistema, quiero cambiar entre entrenamiento síncrono y asíncrono sin necesidad de detener los nodos para permitir distintas estrategias de convergencia. \\ \hline
\textbf{Prioridad} & Media \\ \hline
\textbf{Necesidad} & Importante \\ \hline
\textbf{Alcance} & Módulo de orquestación de rondas; configuración de modos de sincronización; flexibilidad en la gestión del entrenamiento federado. \\ \hline
\textbf{Dependencias} & HU-004, HU-005 \\ \hline
\textbf{Estado} & Deprecada \\ \hline
\end{tabular}
\end{table}



\begin{table}[H]
\centering
\caption{Descripción de la historia de usuario HU-008}
\label{tab:hu008}
\begin{tabular}{|l|p{12cm}|}
\hline
\textbf{ID} & HU-008 \\ \hline
\textbf{Tamaño} & XL \\ \hline
\textbf{Nombre} & Particionado avanzado de \textit{datasets} \\ \hline
\textbf{Descripción} & 
Como investigador, quiero generar particiones no-i.i.d. configurables para simular escenarios realistas de distribución heterogénea en FL. \\ \hline
\textbf{Prioridad} & Alta \\ \hline
\textbf{Necesidad} & Esencial \\ \hline
\textbf{Alcance} & Módulo de particionado de datos; gestión de sujetos y ventanas temporales; soporte explícito para distribuciones no-i.i.d. en entornos federados. \\ \hline
\textbf{Dependencias} & HU-001 \\ \hline
\textbf{Estado} & Completado \\ \hline
\end{tabular}
\end{table}



\begin{table}[H]
\centering
\caption{Descripción de la historia de usuario HU-009}
\label{tab:hu009}
\begin{tabular}{|l|p{12cm}|}
\hline
\textbf{ID} & HU-009 \\ \hline
\textbf{Tamaño} & M \\ \hline
\textbf{Nombre} & Aprendizaje incremental local \\ \hline
\textbf{Descripción} & 
Como usuario, quiero que mi dispositivo siga aprendiendo de nuevos datos aunque no esté participando activamente en rondas federadas, para mantener actualizado el modelo local. \\ \hline
\textbf{Prioridad} & Media \\ \hline
\textbf{Necesidad} & Importante \\ \hline
\textbf{Alcance} & Módulo de entrenamiento local incremental en el cliente; gestión de datos en tiempo real; coherencia entre modelo local y modelo global federado. \\ \hline
\textbf{Dependencias} & HU-004 \\ \hline
\textbf{Estado} & Completado \\ \hline
\end{tabular}
\end{table}



\begin{table}[H]
\centering
\caption{Descripción de la historia de usuario HU-010}
\label{tab:hu010}
\begin{tabular}{|l|p{12cm}|}
\hline
\textbf{ID} & HU-010 \\ \hline
\textbf{Tamaño} & L \\ \hline
\textbf{Nombre} & Explicación del modelo federado \\ \hline
\textbf{Descripción} & 
Como investigador, quiero obtener reportes automáticos de interpretabilidad del sistema federado (Grad-CAM, SHAP, LIME) para evaluar sesgos y calidad del modelo final. \\ \hline
\textbf{Prioridad} & Media \\ \hline
\textbf{Necesidad} & Importante \\ \hline
\textbf{Alcance} & Módulo de explicabilidad y auditoría de modelos; generación de informes automáticos; soporte a la validación científica de los resultados. \\ \hline
\textbf{Dependencias} & HU-004, HU-006 \\ \hline
\textbf{Estado} & Pendiente \\ \hline
\end{tabular}
\end{table}



\begin{table}[H]
\centering
\caption{Descripción de la historia de usuario HU-011}
\label{tab:hu011}
\begin{tabular}{|l|p{12cm}|}
\hline
\textbf{ID} & HU-011 \\ \hline
\textbf{Tamaño} & XL \\ \hline
\textbf{Nombre} & Seguridad y privacidad reforzada \\ \hline
\textbf{Descripción} & 
Como usuario final, quiero que mis datos estén protegidos mediante cifrado y mecanismos de privacidad diferencial para evitar filtraciones o reconstrucción de información sensible. \\ \hline
\textbf{Prioridad} & Alta \\ \hline
\textbf{Necesidad} & Esencial \\ \hline
\textbf{Alcance} & Módulos de cifrado extremo a extremo y preservación de privacidad; cumplimiento de requisitos normativos; protección frente a ataques de reconstrucción. \\ \hline
\textbf{Dependencias} & HU-001, HU-002 \\ \hline
\textbf{Estado} & Pendiente \\ \hline
\end{tabular}
\end{table}



\begin{table}[H]
\centering
\caption{Descripción de la historia de usuario HU-012}
\label{tab:hu012}
\begin{tabular}{|l|p{12cm}|}
\hline
\textbf{ID} & HU-012 \\ \hline
\textbf{Tamaño} & L \\ \hline
\textbf{Nombre} & Tolerancia a fallos \\ \hline
\textbf{Descripción} & 
Como administrador, quiero que el sistema detecte fallos de nodos y recupere automáticamente la ronda sin pérdida significativa de clientes, para mantener la estabilidad del entrenamiento federado. \\ \hline
\textbf{Prioridad} & Alta \\ \hline
\textbf{Necesidad} & Esencial \\ \hline
\textbf{Alcance} & Gestión de fallos y reconexión de nodos; mecanismos de reequilibrado de carga; robustez operativa del sistema federado. \\ \hline
\textbf{Dependencias} & HU-005, HU-006 \\ \hline
\textbf{Estado} & Pendiente \\ \hline
\end{tabular}
\end{table}



\begin{table}[H]
\centering
\caption{Descripción de la historia de usuario HU-013}
\label{tab:hu013}
\begin{tabular}{|l|p{12cm}|}
\hline
\textbf{ID} & HU-013 \\ \hline
\textbf{Tamaño} & L \\ \hline
\textbf{Nombre} & Entrenamiento híbrido \\ \hline
\textbf{Descripción} & 
Como investigador, quiero combinar aprendizaje federado horizontal y técnicas de \textit{transfer learning} distribuidas para mejorar la calidad del modelo final en escenarios complejos. \\ \hline
\textbf{Prioridad} & Media \\ \hline
\textbf{Necesidad} & Importante \\ \hline
\textbf{Alcance} & Configuración de topologías híbridas y flujos avanzados de entrenamiento; soporte a escenarios complejos de datos y dominios. \\ \hline
\textbf{Dependencias} & HU-004 \\ \hline
\textbf{Estado} & Pendiente \\ \hline
\end{tabular}
\end{table}



\begin{table}[H]
\centering
\caption{Descripción de la historia de usuario HU-014}
\label{tab:hu014}
\begin{tabular}{|l|p{12cm}|}
\hline
\textbf{ID} & HU-014 \\ \hline
\textbf{Tamaño} & L \\ \hline
\textbf{Nombre} & Ranking de contribución \\ \hline
\textbf{Descripción} & 
Como investigador, quiero conocer la contribución de cada cliente (por valores de Shapley u otros) para identificar qué nodos aportan mayor valor durante el entrenamiento. \\ \hline
\textbf{Prioridad} & Media \\ \hline
\textbf{Necesidad} & Importante \\ \hline
\textbf{Alcance} & Módulo de análisis de contribución de clientes; cálculo de métricas de influencia; generación de rankings persistentes para auditoría y análisis. \\ \hline
\textbf{Dependencias} & HU-004, HU-006 \\ \hline
\textbf{Estado} & Parcialmente \\ \hline
\end{tabular}
\end{table}



\begin{table}[H]
\centering
\caption{Descripción de la historia de usuario HU-015}
\label{tab:hu015}
\begin{tabular}{|l|p{12cm}|}
\hline
\textbf{ID} & HU-015 \\ \hline
\textbf{Tamaño} & M \\ \hline
\textbf{Nombre} & Gestión de hiperparámetros \\ \hline
\textbf{Descripción} & 
Como investigador, quiero poder modificar, versionar y comparar conjuntos de hiperparámetros para estudiar cómo influyen en la convergencia y estabilidad del sistema federado. \\ \hline
\textbf{Prioridad} & Media \\ \hline
\textbf{Necesidad} & Importante \\ \hline
\textbf{Alcance} & Gestor centralizado de configuraciones de entrenamiento; versionado de hiperparámetros; soporte a experimentación reproducible. \\ \hline
\textbf{Dependencias} & HU-006 \\ \hline
\textbf{Estado} & Pendiente \\ \hline
\end{tabular}
\end{table}



\begin{table}[H]
\centering
\caption{Descripción de la historia de usuario HU-016}
\label{tab:hu016}
\begin{tabular}{|l|p{12cm}|}
\hline
\textbf{ID} & HU-016 \\ \hline
\textbf{Tamaño} & S \\ \hline
\textbf{Nombre} & Consolidación de requisitos funcionales \\ \hline
\textbf{Descripción} & 
Como equipo de desarrollo, quiero consolidar y formalizar los requisitos funcionales del sistema en forma de historias de usuario trazables, para disponer de un \textit{product backlog} coherente que guíe la planificación de los sprints. \\ \hline
\textbf{Prioridad} & Alta \\ \hline
\textbf{Necesidad} & Esencial \\ \hline
\textbf{Alcance} & Definición y organización del \textit{product backlog}; alineamiento con los objetivos del TFM; trazabilidad entre requisitos, arquitectura e iteraciones de desarrollo. \\ \hline
\textbf{Dependencias} & HU-003, HU-008 (revisión de capacidades y flujo de datos definidos en la arquitectura preliminar). \\ \hline
\textbf{Estado} & Completado \\ \hline
\end{tabular}
\end{table}



\begin{table}[H]
\centering
\caption{Descripción de la historia de usuario HU-017}
\label{tab:hu017}
\begin{tabular}{|l|p{12cm}|}
\hline
\textbf{ID} & HU-017 \\ \hline
\textbf{Tamaño} & S \\ \hline
\textbf{Nombre} & Definición de requisitos no funcionales \\ \hline
\textbf{Descripción} & 
Como equipo de desarrollo, quiero identificar y formalizar los requisitos no funcionales (seguridad, privacidad, rendimiento, fiabilidad, mantenibilidad, interoperabilidad, calidad de código), para disponer de criterios de calidad objetivos que se apliquen a todas las historias de usuario. \\ \hline
\textbf{Prioridad} & Alta \\ \hline
\textbf{Necesidad} & Esencial \\ \hline
\textbf{Alcance} & Catálogo de requisitos no funcionales; definición de métricas y criterios de verificación; asociación de RNF a las historias de usuario mediante el campo \emph{Alcance}. \\ \hline
\textbf{Dependencias} & HU-016 (backlog funcional consolidado). \\ \hline
\textbf{Estado} & Completado \\ \hline
\end{tabular}
\end{table}



\begin{table}[H]
\centering
\caption{Descripción de la historia de usuario HU-018}
\label{tab:hu018}
\begin{tabular}{|l|p{12cm}|}
\hline
\textbf{ID} & HU-018 \\ \hline
\textbf{Tamaño} & M \\ \hline
\textbf{Nombre} & Definición de la arquitectura preliminar Edge--Fog--Cloud \\ \hline
\textbf{Descripción} & 
Como equipo de desarrollo, quiero definir una arquitectura preliminar basada en niveles Edge--Fog--Cloud, especificando el flujo de datos y las responsabilidades principales de cada módulo, para asegurar que las historias de usuario se implementan sobre una base técnica consistente y extensible. \\ \hline
\textbf{Prioridad} & Alta \\ \hline
\textbf{Necesidad} & Esencial \\ \hline
\textbf{Alcance} & Modelo arquitectónico del sistema; definición de nodos Edge, Fog y Cloud; flujos de datos y puntos de integración; soporte a futuras decisiones de despliegue en un \textit{computing continuum}. \\ \hline
\textbf{Dependencias} & HU-016, HU-017 (alineamiento con backlog funcional y requisitos de calidad). \\ \hline
\textbf{Estado} & Completado \\ \hline
\end{tabular}
\end{table}



\begin{table}[H]
\centering
\caption{Descripción de la historia de usuario HU-019}
\label{tab:hu019}
\begin{tabular}{|l|p{12cm}|}
\hline
\textbf{ID} & HU-019 \\ \hline
\textbf{Tamaño} & S \\ \hline
\textbf{Nombre} & Estandarización automática del código \\ \hline
\textbf{Descripción} &
Como equipo de desarrollo, quiero que el código del proyecto se formatee y valide automáticamente siguiendo estándares reconocidos (PEP~8), para garantizar consistencia, legibilidad y calidad desde las primeras fases del desarrollo. \\ \hline
\textbf{Prioridad} & Media \\ \hline
\textbf{Necesidad} & Importante \\ \hline
\textbf{Alcance} &
Calidad de código; mantenibilidad; configuración de herramientas de formateo y análisis estático (\texttt{black}, \texttt{flake8}). \\ \hline
\textbf{Dependencias} & -- \\ \hline
\textbf{Estado} & Completado \\ \hline
\end{tabular}
\end{table}



\begin{table}[H]
\centering
\caption{Descripción de la historia de usuario HU-020}
\label{tab:hu020}
\begin{tabular}{|l|p{12cm}|}
\hline
\textbf{ID} & HU-020 \\ \hline
\textbf{Tamaño} & S \\ \hline
\textbf{Nombre} & Configuración de integración continua (CI) \\ \hline
\textbf{Descripción} &
Como equipo de desarrollo, quiero disponer de una canalización de integración continua que ejecute automáticamente pruebas básicas y verificaciones de calidad en cada cambio del repositorio, para detectar errores de forma temprana y asegurar la estabilidad del proyecto. \\ \hline
\textbf{Prioridad} & Alta \\ \hline
\textbf{Necesidad} & Esencial \\ \hline
\textbf{Alcance} &
Infraestructura CI/CD; ejecución automática de tests; validación de estilo; repositorio GitHub y GitHub Actions. \\ \hline
\textbf{Dependencias} & HU-019 \\ \hline
\textbf{Estado} & Completado \\ \hline
\end{tabular}
\end{table}



\begin{table}[H]
\centering
\caption{Descripción de la historia de usuario HU-021}
\label{tab:hu021}
\begin{tabular}{|l|p{12cm}|}
\hline
\textbf{ID} & HU-021 \\ \hline
\textbf{Tamaño} & S \\ \hline
\textbf{Nombre} & Configuración del entorno de desarrollo en Python \\ \hline
\textbf{Descripción} &
Como equipo de desarrollo, quiero configurar un entorno de desarrollo en Python reproducible (gestión de dependencias, entornos virtuales y scripts de inicialización) para asegurar que todo el equipo trabaja con la misma configuración. \\ \hline
\textbf{Prioridad} & Alta \\ \hline
\textbf{Necesidad} & Esencial \\ \hline
\textbf{Alcance} &
Definición del entorno Python; gestión de dependencias (\texttt{requirements.txt}/\texttt{pyproject.toml}); scripts de setup; documentación de instalación y uso. \\ \hline
\textbf{Dependencias} & HU-019 \\ \hline
\textbf{Estado} & Pendiente \\ \hline
\end{tabular}
\end{table}



\begin{table}[H]
\centering
\caption{Descripción de la historia de usuario HU-022}
\label{tab:hu022}
\begin{tabular}{|l|p{12cm}|}
\hline
\textbf{ID} & HU-022 \\ \hline
\textbf{Tamaño} & M \\ \hline
\textbf{Nombre} & Estudio de frameworks de Aprendizaje Federado \\ \hline
\textbf{Descripción} &
Como equipo de desarrollo, quiero analizar distintos frameworks de aprendizaje federado disponibles en Python para evaluar su viabilidad técnica y seleccionar la solución más adecuada como base del sistema propuesto. \\ \hline
\textbf{Prioridad} & Alta \\ \hline
\textbf{Necesidad} & Esencial \\ \hline
\textbf{Alcance} &
Evaluación tecnológica de frameworks de FL; definición de criterios de comparación (arquitectura, facilidad de uso, documentación, extensibilidad, comunidad); soporte a la toma de decisiones técnicas del proyecto. \\ \hline
\textbf{Dependencias} & -- \\ \hline
\textbf{Estado} & Completado \\ \hline
\end{tabular}
\end{table}



\begin{table}[H]
\centering
\caption{Descripción de la historia de usuario HU-023}
\label{tab:hu023}
\begin{tabular}{|l|p{12cm}|}
\hline
\textbf{ID} & HU-023 \\ \hline
\textbf{Tamaño} & S \\ \hline
\textbf{Nombre} & Estudio de TensorFlow Federated \\ \hline
\textbf{Descripción} &
Como equipo de desarrollo, quiero analizar y prototipar un ejemplo básico con la librería TensorFlow Federated para evaluar su idoneidad como framework de aprendizaje federado dentro del proyecto. \\ \hline
\textbf{Prioridad} & Media \\ \hline
\textbf{Necesidad} & Importante \\ \hline
\textbf{Alcance} &
Prototipado mínimo con TensorFlow Federated; análisis de compatibilidad con el entorno Python; evaluación de complejidad, documentación y posibilidades de extensión; comparación con otros frameworks FL. \\ \hline
\textbf{Dependencias} & HU-022 \\ \hline
\textbf{Estado} & Bloqueada \\ \hline
\end{tabular}
\end{table}



\begin{table}[H]
\centering
\caption{Descripción de la historia de usuario HU-024}
\label{tab:hu024}
\begin{tabular}{|l|p{12cm}|}
\hline
\textbf{ID} & HU-024 \\ \hline
\textbf{Tamaño} & M \\ \hline
\textbf{Nombre} & Revisión bibliográfica de frameworks de Aprendizaje Federado \\ \hline
\textbf{Descripción} &
Como equipo de desarrollo, quiero realizar una búsqueda complementaria de literatura reciente y herramientas emergentes de aprendizaje federado para detectar alternativas no contempladas inicialmente y extraer criterios de comparación relevantes que guíen la selección de tecnologías. \\ \hline
\textbf{Prioridad} & Alta \\ \hline
\textbf{Necesidad} & Esencial \\ \hline
\textbf{Alcance} &
Revisión de literatura FL; identificación de frameworks emergentes; extracción de criterios de comparación; soporte a la toma de decisiones tecnológicas. \\ \hline
\textbf{Dependencias} & HU-022 \\ \hline
\textbf{Estado} & Completado \\ \hline
\end{tabular}
\end{table}



\begin{table}[H]
\centering
\caption{Descripción de la historia de usuario HU-025}
\label{tab:hu025}
\begin{tabular}{|l|p{12cm}|}
\hline
\textbf{ID} & HU-025 \\ \hline
\textbf{Tamaño} & M \\ \hline
\textbf{Nombre} & Análisis del framework FLEX (UGR) \\ \hline
\textbf{Descripción} &
Como equipo de desarrollo, quiero analizar el framework FLEX propuesto por un grupo de investigación de la UGR, evaluando su orientación hacia flexibilidad experimental y seguridad/privacidad, para determinar su viabilidad como alternativa complementaria a otros frameworks estudiados. \\ \hline
\textbf{Prioridad} & Alta \\ \hline
\textbf{Necesidad} & Importante \\ \hline
\textbf{Alcance} &
Análisis técnico del framework FLEX; evaluación de flexibilidad experimental; capacidades de seguridad y privacidad; comparativa preliminar con otros frameworks FL. \\ \hline
\textbf{Dependencias} & HU-024 \\ \hline
\textbf{Estado} & Completado \\ \hline
\end{tabular}
\end{table}



\begin{table}[H]
\centering
\caption{Descripción de la historia de usuario HU-026}
\label{tab:hu026}
\begin{tabular}{|l|p{12cm}|}
\hline
\textbf{ID} & HU-026 \\ \hline
\textbf{Tamaño} & S \\ \hline
\textbf{Nombre} & Consolidación de criterios técnicos para evaluación de frameworks \\ \hline
\textbf{Descripción} &
Como equipo de desarrollo, quiero consolidar un conjunto inicial de criterios técnicos para evaluar frameworks de FL (curva de aprendizaje, modularidad, extensibilidad, soporte a privacidad/seguridad, documentación y madurez del ecosistema) para estructurar la comparativa de soluciones. \\ \hline
\textbf{Prioridad} & Alta \\ \hline
\textbf{Necesidad} & Esencial \\ \hline
\textbf{Alcance} &
Definición de criterios de evaluación; definición de métricas de comparación; matriz de decisión para frameworks; soporte a la evaluación técnica y trazabilidad. \\ \hline
\textbf{Dependencias} & HU-024, HU-025 \\ \hline
\textbf{Estado} & Completado \\ \hline
\end{tabular}
\end{table}



\begin{table}[H]
\centering
\caption{Descripción de la historia de usuario HU-027}
\label{tab:hu027}
\begin{tabular}{|l|p{12cm}|}
\hline
\textbf{ID} & HU-027 \\ \hline
\textbf{Tamaño} & XS \\ \hline
\textbf{Nombre} & Contacto y coordinación de reunión técnica con equipo FLEX \\ \hline
\textbf{Descripción} &
Como equipo de desarrollo, quiero establecer contacto con los investigadores del framework FLEX (Mario García-Márquez y equipo UGR) y programar una reunión técnica para profundizar en el estado de madurez del framework y su posible aplicabilidad al proyecto TFM. \\ \hline
\textbf{Prioridad} & Media \\ \hline
\textbf{Necesidad} & Importante \\ \hline
\textbf{Alcance} &
Contacto con equipo investigador; coordinación de reunión técnica; recopilación de información técnica; evaluación de viabilidad en el contexto del TFM. \\ \hline
\textbf{Dependencias} & HU-025, HU-026 \\ \hline
\textbf{Estado} & Completado \\ \hline
\end{tabular}
\end{table}



\begin{table}[H]
\centering
\caption{Descripción de la historia de usuario HU-028}
\label{tab:hu028}
\begin{tabular}{|l|p{12cm}|}
\hline
\textbf{ID} & HU-028 \\ \hline
\textbf{Tamaño} & S \\ \hline
\textbf{Nombre} & Documentación de decisiones técnicas y evidencias (Sprint 1) \\ \hline
\textbf{Descripción} &
Como equipo de desarrollo, quiero registrar conclusiones y evidencias de la fase exploratoria de frameworks (notas de revisión, enlaces a recursos, resumen de hallazgos) para mantener trazabilidad de las decisiones técnicas en sprints posteriores. \\ \hline
\textbf{Prioridad} & Media \\ \hline
\textbf{Necesidad} & Importante \\ \hline
\textbf{Alcance} &
Documentación de decisiones; registro de evidencias de evaluación; trazabilidad de criterios de decisión; repositorio centralizado de hallazgos técnicos. \\ \hline
\textbf{Dependencias} & HU-024, HU-025, HU-026, HU-027 \\ \hline
\textbf{Estado} & Completado \\ \hline
\end{tabular}
\end{table}



\begin{table}[H]
\centering
\caption{Descripción de la historia de usuario HU-029}
\label{tab:hu029}
\begin{tabular}{|l|p{12cm}|}
\hline
\textbf{ID} & HU-029 \\ \hline
\textbf{Tamaño} & M \\ \hline
\textbf{Nombre} & Estudio del framework FLOWER \\ \hline
\textbf{Descripción} &
Como equipo de desarrollo, quiero analizar y prototipar un ejemplo básico con el framework FLOWER para evaluar su viabilidad como base del sistema de aprendizaje federado propuesto. \\ \hline
\textbf{Prioridad} & Alta \\ \hline
\textbf{Necesidad} & Esencial \\ \hline
\textbf{Alcance} &
Ejecución de ejemplos básicos con FLOWER; evaluación de facilidad de integración, flexibilidad arquitectónica y soporte comunitario; comparación con TensorFlow Federated para la toma de decisiones técnicas. \\ \hline
\textbf{Dependencias} & HU-022 \\ \hline
\textbf{Estado} & Completado \\ \hline
\end{tabular}
\end{table}



\begin{table}[H]
\centering
\caption{Descripción de la historia de usuario HU-030}
\label{tab:hu030}
\begin{tabular}{|l|p{12cm}|}
\hline
\textbf{ID} & HU-030 \\ \hline
\textbf{Tamaño} & M \\ \hline
\textbf{Nombre} & Implementación de modelos baseline sobre ECG5000 \\ \hline
\textbf{Descripción} &
Como equipo de desarrollo, quiero implementar y evaluar modelos baseline centralizados sobre el dataset ECG5000 para disponer de una referencia cuantitativa con la que comparar el rendimiento de las configuraciones federadas. \\ \hline
\textbf{Prioridad} & Alta \\ \hline
\textbf{Necesidad} & Importante \\ \hline
\textbf{Alcance} &
Implementación de modelos (MLP/CNN) en entrenamiento centralizado; definición de particiones Train/Validation/Test; generación de métricas base (accuracy, loss) para comparación posterior con FedAvg. \\ \hline
\textbf{Dependencias} & HU-020 \\ \hline
\textbf{Estado} & Completado \\ \hline
\end{tabular}
\end{table}



\begin{table}[H]
\centering
\caption{Descripción de la historia de usuario HU-031}
\label{tab:hu031}
\begin{tabular}{|l|p{12cm}|}
\hline
\textbf{ID} & HU-031 \\ \hline
\textbf{Tamaño} & M \\ \hline
\textbf{Nombre} & Evaluación técnica y decisión sobre viabilidad del framework FLEX \\ \hline
\textbf{Descripción} &
Como equipo de desarrollo, quiero celebrar una reunión técnica con Mario García-Márquez para discutir en detalle el estado actual del framework FLEX, su aplicabilidad al TFM y su viabilidad técnica como solución para el entrenamiento federado. \\ \hline
\textbf{Prioridad} & Alta \\ \hline
\textbf{Necesidad} & Esencial \\ \hline
\textbf{Alcance} &
Reunión técnica con investigadores FLEX; evaluación de madurez del framework; análisis de compatibilidad; toma de decisión sobre su incorporación o descarte en el proyecto. \\ \hline
\textbf{Dependencias} & HU-027 \\ \hline
\textbf{Estado} & Completado \\ \hline
\end{tabular}
\end{table}



\begin{table}[H]
\centering
\caption{Descripción de la historia de usuario HU-032}
\label{tab:hu032}
\begin{tabular}{|l|p{12cm}|}
\hline
\textbf{ID} & HU-032 \\ \hline
\textbf{Tamaño} & S \\ \hline
\textbf{Nombre} & Documentación de decisiones técnicas y conclusiones (Sprint 2) \\ \hline
\textbf{Descripción} &
Como equipo de desarrollo, quiero registrar las conclusiones técnicas, recomendaciones y evidencias obtenidas durante el Sprint 2 (incluyendo la reunión con el equipo FLEX) para garantizar la trazabilidad de las decisiones adoptadas y servir de referencia para futuros sprints. \\ \hline
\textbf{Prioridad} & Media \\ \hline
\textbf{Necesidad} & Importante \\ \hline
\textbf{Alcance} &
Documentación de conclusiones Sprint 2; registro de recomendaciones técnicas; trazabilidad de decisiones; repositorio centralizado de aprendizajes y hallazgos. \\ \hline
\textbf{Dependencias} & HU-031 \\ \hline
\textbf{Estado} & Completado \\ \hline
\end{tabular}
\end{table}



\begin{table}[H]
\centering
\caption{Descripción de la historia de usuario HU-033}
\label{tab:hu033}
\begin{tabular}{|l|p{12cm}|}
\hline
\textbf{ID} & HU-033 \\ \hline
\textbf{Tamaño} & M \\ \hline
\textbf{Nombre} & Análisis de asincronía en el framework Flower \\ \hline
\textbf{Descripción} &
Como equipo de desarrollo, quiero analizar si Flower permite comportamientos asíncronos o tolerantes a la no disponibilidad de clientes para determinar si es viable implementar la arquitectura Edge--Fog--Cloud sin mecanismos externos. \\ \hline
\textbf{Prioridad} & Alta \\ \hline
\textbf{Necesidad} & Esencial \\ \hline
\textbf{Alcance} &
Evaluación técnica del modelo de ejecución de Flower; sincronía de rondas; requisitos de disponibilidad (RNF-005); impacto sobre diseño jerárquico Edge--Fog--Cloud. \\ \hline
\textbf{Dependencias} & HU-029 \\ \hline
\textbf{Estado} & Completado \\ \hline
\end{tabular}
\end{table}



\begin{table}[H]
\centering
\caption{Descripción de la historia de usuario HU-034}
\label{tab:hu034}
\begin{tabular}{|l|p{12cm}|}
\hline
\textbf{ID} & HU-034 \\ \hline
\textbf{Tamaño} & M \\ \hline
\textbf{Nombre} & Diseño arquitectónico definiendo tecnologías en cada capa \\ \hline
\textbf{Descripción} &
Como equipo de desarrollo, quiero definir una arquitectura Edge--Fog--Cloud asignando tecnologías concretas a cada capa para que el diseño sea implementable y trazable respecto a los requisitos del sistema. \\ \hline
\textbf{Prioridad} & Alta \\ \hline
\textbf{Necesidad} & Esencial \\ \hline
\textbf{Alcance} &
Arquitectura tecnológica del continuum; asignación explícita de comunicación (MQTT) y coordinación FL (Flower); definición de responsabilidades por capa; mantenibilidad y extensibilidad. \\ \hline
\textbf{Dependencias} & HU-018, HU-033 \\ \hline
\textbf{Estado} & Completado \\ \hline
\end{tabular}
\end{table}



\begin{table}[H]
\centering
\caption{Descripción de la historia de usuario HU-035}
\label{tab:hu035}
\begin{tabular}{|l|p{12cm}|}
\hline
\textbf{ID} & HU-035 \\ \hline
\textbf{Tamaño} & S \\ \hline
\textbf{Nombre} & Estudio y selección de broker MQTT \\ \hline
\textbf{Descripción} &
Como equipo de desarrollo, quiero evaluar distintos brokers MQTT para seleccionar uno que permita comunicación asíncrona estable entre Edge y Fog y que sea fácil de desplegar en un entorno contenerizado. \\ \hline
\textbf{Prioridad} & Alta \\ \hline
\textbf{Necesidad} & Esencial \\ \hline
\textbf{Alcance} &
Selección tecnológica del broker MQTT; criterios de evaluación (madurez, facilidad de despliegue, soporte, adecuación Edge--Fog); soporte a la arquitectura híbrida. \\ \hline
\textbf{Dependencias} & HU-034 \\ \hline
\textbf{Estado} & Completado \\ \hline
\end{tabular}
\end{table}



\begin{table}[H]
\centering
\caption{Descripción de la historia de usuario HU-036}
\label{tab:hu036}
\begin{tabular}{|l|p{12cm}|}
\hline
\textbf{ID} & HU-036 \\ \hline
\textbf{Tamaño} & S \\ \hline
\textbf{Nombre} & Implementación y puesta a punto de broker MQTT contenerizado \\ \hline
\textbf{Descripción} &
Como equipo de desarrollo, quiero desplegar un broker MQTT contenerizado y verificar su funcionamiento para habilitar el intercambio asíncrono de actualizaciones entre Edge y Fog dentro del prototipo. \\ \hline
\textbf{Prioridad} & Alta \\ \hline
\textbf{Necesidad} & Esencial \\ \hline
\textbf{Alcance} &
Despliegue con Docker/Docker Compose; exposición de puertos; configuración mínima inicial; pruebas de publicación/suscripción; base para integración Edge--Fog. \\ \hline
\textbf{Dependencias} & HU-035 \\ \hline
\textbf{Estado} & Completado \\ \hline
\end{tabular}
\end{table}



\begin{table}[H]
\centering
\caption{Descripción de la historia de usuario HU-037}
\label{tab:hu037}
\begin{tabular}{|l|p{12cm}|}
\hline
\textbf{ID} & HU-037 \\ \hline
\textbf{Tamaño} & M \\ \hline
\textbf{Nombre} & Diseño a bajo nivel del prototipo (clases y responsabilidades) \\ \hline
\textbf{Descripción} &
Como equipo de desarrollo, quiero definir el diseño a bajo nivel del sistema (principales clases, responsabilidades e interacciones) para preparar la implementación del primer prototipo funcional de la arquitectura híbrida. \\ \hline
\textbf{Prioridad} & Alta \\ \hline
\textbf{Necesidad} & Esencial \\ \hline
\textbf{Alcance} &
Diseño software del prototipo; identificación de módulos (EdgeClient, FogAggregator, MQTTAdapter, FlowerClient/Server, ModelRepository); mantenibilidad; trazabilidad con HU-034. \\ \hline
\textbf{Dependencias} & HU-034, HU-036 \\ \hline
\textbf{Estado} & Pendiente \\ \hline
\end{tabular}
\end{table}



\begin{table}[H]
\centering
\caption{Descripción de la historia de usuario HU-038}
\label{tab:hu038}
\begin{tabular}{|l|p{12cm}|}
\hline
\textbf{ID} & HU-038 \\ \hline
\textbf{Tamaño} & M \\ \hline
\textbf{Nombre} & Análisis exploratorio del dataset WESAD \\ \hline
\textbf{Descripción} &
Como investigador, quiero analizar la estructura, variables y formato del dataset WESAD para evaluar su idoneidad como base experimental del sistema y estimar la viabilidad de reproducir resultados reportados en la literatura. \\ \hline
\textbf{Prioridad} & Alta \\ \hline
\textbf{Necesidad} & Importante \\ \hline
\textbf{Alcance} &
Análisis de dataset; variables fisiológicas; segmentación temporal; configuración de particionado no-i.i.d.; definición de baseline para reproducibilidad. \\ \hline
\textbf{Dependencias} & HU-001, HU-008 \\ \hline
\textbf{Estado} & Pendiente \\ \hline
\end{tabular}
\end{table}



\begin{table}[H]
\centering
\caption{Descripción de la historia de usuario HU-039}
\label{tab:hu039}
\begin{tabular}{|l|p{12cm}|}
\hline
\textbf{ID} & HU-039 \\ \hline
\textbf{Tamaño} & M \\ \hline
\textbf{Nombre} & Análisis exploratorio del dataset SWELL \\ \hline
\textbf{Descripción} &
Como investigador, quiero analizar la estructura, variables y formato del dataset SWELL para evaluar su adecuación al sistema propuesto y comparar su potencial experimental frente a WESAD. \\ \hline
\textbf{Prioridad} & Alta \\ \hline
\textbf{Necesidad} & Importante \\ \hline
\textbf{Alcance} &
Análisis de dataset; variables fisiológicas y contexto; segmentación y etiquetas; estimación de complejidad; criterios de selección del dataset principal del proyecto. \\ \hline
\textbf{Dependencias} & HU-001, HU-008 \\ \hline
\textbf{Estado} & Pendiente \\ \hline
\end{tabular}
\end{table}



\begin{table}[H]
\centering
\caption{Descripción de la historia de usuario HU-040}
\label{tab:hu040}
\begin{tabular}{|l|p{12cm}|}
\hline
\textbf{ID} & HU-040 \\ \hline
\textbf{Tamaño} & M \\ \hline
\textbf{Nombre} & Diseño de baseline en WESAD/SWELL \\ \hline
\textbf{Descripción} &
Como investigador, quiero implementar un baseline reproducible sobre el dataset seleccionado (WESAD o SWELL) para comprobar si se alcanzan resultados comparables a los reportados y disponer de una referencia sólida para el experimento federado. \\ \hline
\textbf{Prioridad} & Alta \\ \hline
\textbf{Necesidad} & Esencial \\ \hline
\textbf{Alcance} &
Implementación de pipeline de entrenamiento centralizado; métricas y validación; reproducibilidad experimental; base comparativa para posteriores rondas federadas. \\ \hline
\textbf{Dependencias} & HU-038, HU-039 \\ \hline
\textbf{Estado} & Pendiente \\ \hline
\end{tabular}
\end{table}



\begin{table}[H]
\centering
\caption{Descripción de la historia de usuario HU-041}
\label{tab:hu041}
\begin{tabular}{|l|p{12cm}|}
\hline
\textbf{ID} & HU-041 \\ \hline
\textbf{Tamaño} & S \\ \hline
\textbf{Nombre} & Planificación del plan de reproducibilidad experimental (baseline) \\ \hline
\textbf{Descripción} &
Como investigador, quiero definir un plan de evaluación para reproducir resultados centralizados como condición previa fundamental para la validación del aprendizaje federado posterior, estableciendo las líneas base contra las que comparar el desempeño de la arquitectura distribuida. \\ \hline
\textbf{Prioridad} & Alta \\ \hline
\textbf{Necesidad} & Esencial \\ \hline
\textbf{Alcance} &
Plan de reproducibilidad experimental; definición de datasets base; establecimiento de métricas de referencia centralizadas; protocolo de evaluación reproducible; documentación de metodología. \\ \hline
\textbf{Dependencias} & HU-030, HU-040 \\ \hline
\textbf{Estado} & En proceso \\ \hline
\end{tabular}
\end{table}



\begin{table}[H]
\centering
\caption{Descripción de la historia de usuario HU-042}
\label{tab:hu042}
\begin{tabular}{|l|p{12cm}|}
\hline
\textbf{ID} & HU-042 \\ \hline
\textbf{Tamaño} & L \\ \hline
\textbf{Nombre} & Desarrollo de nodos cliente (Edge) \\ \hline
\textbf{Descripción} &
Como equipo de desarrollo, quiero implementar una primera versión funcional del nodo cliente en la capa Edge, que permita entrenar un modelo localmente y publicar las actualizaciones del modelo mediante MQTT, con el fin de validar el flujo Edge--Fog dentro de la arquitectura propuesta. \\ \hline
\textbf{Prioridad} & Alta \\ \hline
\textbf{Necesidad} & Esencial \\ \hline
\textbf{Alcance} &
Implementación inicial del nodo Edge (\texttt{FLClientMQTT}); entrenamiento local básico; serialización de pesos del modelo; publicación de actualizaciones vía MQTT; verificación de conectividad con el broker y el nodo Fog. \\ \hline
\textbf{Dependencias} & HU-034, HU-036, HU-037 \\ \hline
\textbf{Estado} & Completado \\ \hline
\end{tabular}
\end{table}



\begin{table}[H]
\centering
\caption{Descripción de la historia de usuario HU-043}
\label{tab:hu043}
\begin{tabular}{|l|p{12cm}|}
\hline
\textbf{ID} & HU-043 \\ \hline
\textbf{Tamaño} & M \\ \hline
\textbf{Nombre} & Desarrollo del nodo cliente para SWELL \\ \hline
\textbf{Descripción} &
Como investigador, quiero adaptar el nodo cliente Edge al dataset SWELL, incorporando su carga de datos, segmentación temporal y configuración de etiquetas, para permitir entrenamientos locales coherentes con el escenario experimental definido. \\ \hline
\textbf{Prioridad} & Alta \\ \hline
\textbf{Necesidad} & Importante \\ \hline
\textbf{Alcance} &
Integración del pipeline de datos SWELL en el nodo Edge; segmentación en ventanas temporales; configuración de etiquetas; construcción de estructuras de datos compatibles con el entrenamiento local; alineamiento con el esquema de particionado definido. \\ \hline
\textbf{Dependencias} & HU-042, HU-039 \\ \hline
\textbf{Estado} & Completado \\ \hline
\end{tabular}
\end{table}



\begin{table}[H]
\centering
\caption{Descripción de la historia de usuario HU-044}
\label{tab:hu044}
\begin{tabular}{|l|p{12cm}|}
\hline
\textbf{ID} & HU-044 \\ \hline
\textbf{Tamaño} & L \\ \hline
\textbf{Nombre} & Desarrollo de nodos Fog (recepción y agregación parcial) \\ \hline
\textbf{Descripción} &
Como equipo de desarrollo, quiero implementar un nodo Fog capaz de recibir actualizaciones de múltiples clientes Edge a través de MQTT y realizar una agregación parcial básica, con el objetivo de validar el papel del Fog como capa intermedia en el flujo Edge--Fog--Cloud. \\ \hline
\textbf{Prioridad} & Alta \\ \hline
\textbf{Necesidad} & Esencial \\ \hline
\textbf{Alcance} &
Implementación inicial del nodo Fog (\texttt{BrokerFog}/\texttt{FogClient}); suscripción a tópicos MQTT; almacenamiento temporal de actualizaciones; agregación parcial simple de modelos; publicación del modelo agregado hacia la capa superior. \\ \hline
\textbf{Dependencias} & HU-034, HU-036, HU-037 \\ \hline
\textbf{Estado} & Completado \\ \hline
\end{tabular}
\end{table}



\begin{table}[H]
\centering
\caption{Descripción de la historia de usuario HU-045}
\label{tab:hu045}
\begin{tabular}{|l|p{12cm}|}
\hline
\textbf{ID} & HU-045 \\ \hline
\textbf{Tamaño} & M \\ \hline
\textbf{Nombre} & Desarrollo de nodos Fog para SWELL \\ \hline
\textbf{Descripción} &
Como investigador, quiero adaptar el nodo Fog al flujo experimental basado en el dataset SWELL, garantizando que las actualizaciones generadas por clientes entrenados con SWELL puedan ser recibidas y agregadas de forma coherente. \\ \hline
\textbf{Prioridad} & Alta \\ \hline
\textbf{Necesidad} & Importante \\ \hline
\textbf{Alcance} &
Adaptación del nodo Fog al esquema de datos de SWELL; soporte a particiones por sujeto; cálculo de pesos de agregación basados en tamaño muestral; validación funcional del flujo Edge--Fog con datos reales. \\ \hline
\textbf{Dependencias} & HU-044, HU-043, HU-037 \\ \hline
\textbf{Estado} & Completado \\ \hline
\end{tabular}
\end{table}



\begin{table}[H]
\centering
\caption{Descripción de la historia de usuario HU-046}
\label{tab:hu046}
\begin{tabular}{|l|p{12cm}|}
\hline
\textbf{ID} & HU-046 \\ \hline
\textbf{Tamaño} & L \\ \hline
\textbf{Nombre} & Particionado avanzado de \textit{datasets} \\ \hline
\textbf{Descripción} &
Como investigador, quiero implementar el particionado avanzado de \textit{datasets} (no-i.i.d., por sujeto y por ventanas temporales) para ejecutar experimentos federados realistas y reproducibles. \\ \hline
\textbf{Prioridad} & Alta \\ \hline
\textbf{Necesidad} & Esencial \\ \hline
\textbf{Alcance} &
Implementación de particionadores configurables; validación de distribuciones; exportación de splits; integración en el pipeline experimental. \\ \hline
\textbf{Dependencias} & HU-008, HU-038, HU-039, HU-040 \\ \hline
\textbf{Estado} & Completado \\ \hline
\end{tabular}
\end{table}



\begin{table}[H]
\centering
\caption{Descripción de la historia de usuario HU-047}
\label{tab:hu047}
\begin{tabular}{|l|p{12cm}|}
\hline
\textbf{ID} & HU-047 \\ \hline
\textbf{Tamaño} & L \\ \hline
\textbf{Nombre} & Desarrollo de modelos baseline centralizados \\ \hline
\textbf{Descripción} &
Como investigador, quiero implementar y evaluar modelos baseline centralizados clásicos (regresión logística, SVM, Random Forest, XGBoost y MLP) para disponer de referencias cuantitativas con las que comparar posteriormente el rendimiento del aprendizaje federado. \\ \hline
\textbf{Prioridad} & Alta \\ \hline
\textbf{Necesidad} & Esencial \\ \hline
\textbf{Alcance} &
Pipeline de entrenamiento centralizado; definición de métricas (Accuracy y Macro F1); validación cruzada; evaluación en WESAD y SWELL; generación de tablas comparativas para experimentos federados posteriores. \\ \hline
\textbf{Dependencias} & HU-037, HU-040 \\ \hline
\textbf{Estado} & Completado \\ \hline
\end{tabular}
\end{table}



\begin{table}[H]
\centering
\caption{Descripción de la historia de usuario HU-048}
\label{tab:hu048}
\begin{tabular}{|l|p{12cm}|}
\hline
\textbf{ID} & HU-048 \\ \hline
\textbf{Tamaño} & L \\ \hline
\textbf{Nombre} & Servidor Cloud Flower \\ \hline
\textbf{Descripción} &
Como equipo de desarrollo, quiero implementar el servidor central en la capa Cloud usando \textit{Flower} para orquestar rondas federadas y realizar la agregación global del modelo. \\ \hline
\textbf{Prioridad} & Alta \\ \hline
\textbf{Necesidad} & Esencial \\ \hline
\textbf{Alcance} &
Implementación del servidor Flower; configuración de la estrategia de agregación global (FedAvg); gestión de rondas y distribución del modelo global. \\ \hline
\textbf{Dependencias} & HU-042, HU-043, HU-044, HU-045, HU-037, HU-047 \\ \hline
\textbf{Estado} & Completado \\ \hline
\end{tabular}
\end{table}



\begin{table}[H]
\centering
\caption{Descripción de la historia de usuario HU-049}
\label{tab:hu049}
\begin{tabular}{|l|p{12cm}|}
\hline
\textbf{ID} & HU-049 \\ \hline
\textbf{Tamaño} & L \\ \hline
\textbf{Nombre} & Fog como cliente Flower \\ \hline
\textbf{Descripción} &
Como equipo de desarrollo, quiero adaptar los nodos Fog para actuar como clientes federados de \textit{Flower}, permitiendo su participación estable en las rondas coordinadas por el servidor Cloud. \\ \hline
\textbf{Prioridad} & Alta \\ \hline
\textbf{Necesidad} & Esencial \\ \hline
\textbf{Alcance} &
Adaptación de los nodos Fog a la interfaz cliente de Flower; gestión de rondas; envío y recepción de parámetros hacia/desde la capa Cloud. \\ \hline
\textbf{Dependencias} & HU-044, HU-045, HU-048 \\ \hline
\textbf{Estado} & Completado \\ \hline
\end{tabular}
\end{table}



\begin{table}[H]
\centering
\caption{Descripción de la historia de usuario HU-050}
\label{tab:hu050}
\begin{tabular}{|l|p{12cm}|}
\hline
\textbf{ID} & HU-050 \\ \hline
\textbf{Tamaño} & M \\ \hline
\textbf{Nombre} & Estandarización de pesos/serialización Fog--Cloud \\ \hline
\textbf{Descripción} &
Como equipo de desarrollo, quiero estandarizar el formato de los pesos y la serialización de actualizaciones entre Fog y Cloud para asegurar un intercambio de parámetros consistente y reproducible. \\ \hline
\textbf{Prioridad} & Alta \\ \hline
\textbf{Necesidad} & Importante \\ \hline
\textbf{Alcance} &
Normalización de estructuras de pesos; definición de formato de intercambio; validación de compatibilidad entre Fog y Cloud. \\ \hline
\textbf{Dependencias} & HU-048, HU-049 \\ \hline
\textbf{Estado} & Completado \\ \hline
\end{tabular}
\end{table}



\begin{table}[H]
\centering
\caption{Descripción de la historia de usuario HU-051}
\label{tab:hu051}
\begin{tabular}{|l|p{12cm}|}
\hline
\textbf{ID} & HU-051 \\ \hline
\textbf{Tamaño} & M \\ \hline
\textbf{Nombre} & Validación end-to-end reproducible en SWELL \\ \hline
\textbf{Descripción} &
Como investigador, quiero ejecutar y validar un flujo \textit{end-to-end} reproducible sobre el dataset SWELL para confirmar la estabilidad del ciclo Edge--Fog--Cloud. \\ \hline
\textbf{Prioridad} & Alta \\ \hline
\textbf{Necesidad} & Esencial \\ \hline
\textbf{Alcance} &
Configuración del experimento reproducible; ejecución de rondas federadas completas; verificación de métricas y consistencia de resultados. \\ \hline
\textbf{Dependencias} & HU-042, HU-043, HU-049, HU-050 \\ \hline
\textbf{Estado} & Completado \\ \hline
\end{tabular}
\end{table}



\begin{table}[H]
\centering
\caption{Descripción de la historia de usuario HU-052}
\label{tab:hu052}
\begin{tabular}{|l|p{12cm}|}
\hline
\textbf{ID} & HU-052 \\ \hline
\textbf{Tamaño} & S \\ \hline
\textbf{Nombre} & Test del sistema y de clases \\ \hline
\textbf{Descripción} &
Como investigador, quiero probar y validar de forma automática que las clases y componentes del sistema funcionan conforme a los criterios definidos, con el fin de garantizar la corrección, estabilidad y mantenibilidad del sistema. \\ \hline
\textbf{Prioridad} & Media \\ \hline
\textbf{Necesidad} & Importante \\ \hline
\textbf{Alcance} &
Diseño e implementación de pruebas unitarias y de integración para las clases principales del sistema; validación de casos de uso críticos; ejecución automática de tests y verificación de resultados esperados. \\ \hline
\textbf{Dependencias} & HU-048, HU-051 \\ \hline
\textbf{Estado} & Pendiente \\ \hline
\end{tabular}
\end{table}




\begin{table}[H]
\centering
\caption{Descripción de la historia de usuario HU-053}
\label{tab:hu053}
\begin{tabular}{|l|p{12cm}|}
\hline
\textbf{ID} & HU-053 \\ \hline
\textbf{Tamaño} & M \\ \hline
\textbf{Nombre} & Análisis de viabilidad de Federated Transfer Learning \\ \hline
\textbf{Descripción} &
Como investigador, quiero analizar la viabilidad de incorporar Federated Transfer Learning sobre la arquitectura actual e identificar los cambios necesarios para su adopción. \\ \hline
\textbf{Prioridad} & Alta \\ \hline
\textbf{Necesidad} & Esencial \\ \hline
\textbf{Alcance} &
Análisis conceptual del FTL; evaluación de compatibilidad con el flujo Edge--Fog--Cloud; implicaciones a nivel de modelos, datasets y coordinación de rondas y clientes. \\ \hline
\textbf{Dependencias} &
HU-042, HU-043, HU-044, HU-045, HU-037, HU-047 \\ \hline
\textbf{Estado} & Completado \\ \hline
\end{tabular}
\end{table}




\begin{table}[H]
\centering
\caption{Descripción de la historia de usuario HU-054}
\label{tab:hu054}
\begin{tabular}{|l|p{12cm}|}
\hline
\textbf{ID} & HU-054 \\ \hline
\textbf{Tamaño} & M \\ \hline
\textbf{Nombre} & Análisis de compatibilidad entre datasets para transferencia \\ \hline
\textbf{Descripción} &
Como investigador, quiero analizar la compatibilidad entre distintos datasets fisiológicos mediante modelos baseline centralizados, para evaluar hasta qué punto es viable la transferencia de conocimiento entre dominios.\\ \hline
\textbf{Prioridad} & Alta \\ \hline
\textbf{Necesidad} & Esencial \\ \hline
\textbf{Alcance} &
Evaluación sobre baselines con los datasets WESAD y SWELL; análisis de similitud de etiquetas y parámetros; identificación de limitaciones de transferencia.\\ \hline
\textbf{Dependencias} &
HU-047, HU-053 \\ \hline
\textbf{Estado} & Completado \\ \hline
\end{tabular}
\end{table}


\begin{table}[H]
\centering
\caption{Descripción de la historia de usuario HU-055}
\label{tab:hu055}
\begin{tabular}{|l|p{12cm}|}
\hline
\textbf{ID} & HU-055 \\ \hline
\textbf{Tamaño} & S \\ \hline
\textbf{Nombre} & Validación metodológica del particionado por sujeto \\ \hline
\textbf{Descripción} &
Como investigador, quiero revisar y validar el esquema de particionamiento de los datos (train, validation y test) para identificar y valorar problemas de solapamiento o \textit{data leakage}. \\ \hline
\textbf{Prioridad} & Alta \\ \hline
\textbf{Necesidad} & Esencial \\ \hline
\textbf{Alcance} &
Análisis del particionamiento actual; detección de solapamientos; identificación de \textit{data leakage} y necesidades de corrección. \\ \hline
\textbf{Dependencias} &
HU-054 \\ \hline
\textbf{Estado} & Completada \\ \hline
\end{tabular}
\end{table}


\begin{table}[H]
\centering
\caption{Descripción de la historia de usuario HU-056}
\label{tab:hu056}
\begin{tabular}{|l|p{12cm}|}
\hline
\textbf{ID} & HU-056 \\ \hline
\textbf{Tamaño} & S \\ \hline
\textbf{Nombre} & Redefinición del particionado experimental por sujeto \\ \hline
\textbf{Descripción} &
Como investigador, quiero redefinir el particionado experimental por sujeto para corregir problemas detectados (solapamientos o \textit{data leakage}) y asegurar la validez de los resultados. \\ \hline
\textbf{Prioridad} & Alta \\ \hline
\textbf{Necesidad} & Esencial \\ \hline
\textbf{Alcance} &
Revisión y ajuste de los splits por sujeto; reglas de separación; documentación de cambios y criterios de particionado. \\ \hline
\textbf{Dependencias} & HU-055 \\ \hline
\textbf{Estado} & Completada \\ \hline
\end{tabular}
\end{table}

\begin{table}[H]
\centering
\caption{Descripción de la historia de usuario HU-057}
\label{tab:hu057}
\begin{tabular}{|l|p{12cm}|}
\hline
\textbf{ID} & HU-057 \\ \hline
\textbf{Tamaño} & M \\ \hline
\textbf{Nombre} & Distribuidor de dataset configurable por sujeto \\ \hline
\textbf{Descripción} &
Como investigador, quiero disponer de un distribuidor de dataset configurable mediante ficheros externos para asignar sujetos a nodos federados de forma flexible, reproducible y desacoplada del código. \\ \hline
\textbf{Prioridad} & Alta \\ \hline
\textbf{Necesidad} & Esencial \\ \hline
\textbf{Alcance} &
Diseño de un esquema de configuración declarativo; implementación del módulo de lectura y validación de configuraciones; materialización del dataset particionado conforme a la asignación definida. \\ \hline
\textbf{Dependencias} & HU-056 \\ \hline
\textbf{Estado} & Completada \\ \hline
\end{tabular}
\end{table}

\begin{table}[H]
\centering
\caption{Descripción de la historia de usuario HU-058}
\label{tab:hu058}
\begin{tabular}{|l|p{12cm}|}
\hline
\textbf{ID} & HU-058 \\ \hline
\textbf{Tamaño} & M \\ \hline
\textbf{Nombre} & Adaptación de los componentes al nuevo esquema de distribución \\ \hline
\textbf{Descripción} &
Como investigador, quiero adaptar los componentes del sistema federado al nuevo esquema de distribución del dataset para asegurar la correcta ejecución del entrenamiento y evaluación en todos los niveles de la arquitectura. \\ \hline
\textbf{Prioridad} & Alta \\ \hline
\textbf{Necesidad} & Esencial \\ \hline
\textbf{Alcance} &
Integración del distribuidor de dataset en los nodos Edge, Fog y Server; ajuste de los flujos de entrenamiento y evaluación al nuevo particionado; verificación del funcionamiento extremo a extremo del sistema federado. \\ \hline
\textbf{Dependencias} &
HU-057 \\ \hline
\textbf{Estado} & Pendiente \\ \hline
\end{tabular}
\end{table}

\begin{table}[H]
\centering
\caption{Descripción de la historia de usuario HU-059}
\label{tab:hu059}
\begin{tabular}{|l|p{12cm}|}
\hline
\textbf{ID} & HU-059 \\ \hline
\textbf{Tamaño} & S \\ \hline
\textbf{Nombre} & Testing del sistema de particionado y arquitectura \\ \hline
\textbf{Descripción} &
Como investigador, quiero validar mediante pruebas automáticas el sistema de particionado y la arquitectura federada para garantizar que las configuraciones se materialicen correctamente y el plan de ejecución sea consistente. \\ \hline
\textbf{Prioridad} & Media \\ \hline
\textbf{Necesidad} & Importante \\ \hline
\textbf{Alcance} &
Desarrollo de tests unitarios para la validación de configuraciones; verificación de la correcta carga y materialización de arquitecturas federadas; comprobación del correcto funcionamiento del plan de ejecución. \\ \hline
\textbf{Dependencias} &
HU-057, HU-058 \\ \hline
\textbf{Estado} & Pendiente \\ \hline
\end{tabular}
\end{table}


\begin{table}[H]
\centering
\caption{Descripción de la historia de usuario HU-060}
\label{tab:hu060}
\begin{tabular}{|l|p{12cm}|}
\hline
\textbf{ID} & HU-060 \\ \hline
\textbf{Tamaño} & M \\ \hline
\textbf{Nombre} & Implementación de TLS \\ \hline
\textbf{Descripción} &
Como administrador del sistema, quiero asegurar las comunicaciones mediante TLS para proteger la confidencialidad e integridad de los datos y cumplir los requisitos de seguridad. \\ \hline
\textbf{Prioridad} & Alta \\ \hline
\textbf{Necesidad} & Esencial \\ \hline
\textbf{Alcance} &
Seguridad en la capa de comunicaciones; cifrado extremo a extremo; gestión de certificados y canales seguros. \\ \hline
\textbf{Dependencias} & HU-002 \\ \hline
\textbf{Estado} & Pendiente \\ \hline
\end{tabular}
\end{table}



\begin{table}[H]
\centering
\caption{Descripción de la historia de usuario HU-061}
\label{tab:hu061}
\begin{tabular}{|l|p{12cm}|}
\hline
\textbf{ID} & HU-061 \\ \hline
\textbf{Tamaño} & M \\ \hline
\textbf{Nombre} & Configuración declarativa de la arquitectura federada \\ \hline
\textbf{Descripción} &
Como investigador, quiero definir la arquitectura federada y sus parámetros de ejecución mediante ficheros de configuración declarativos para poder instanciar y comparar distintos escenarios experimentales de forma reproducible. \\ \hline
\textbf{Prioridad} & Media \\ \hline
\textbf{Necesidad} & Importante \\ \hline
\textbf{Alcance} &
Definición de una configuración global del sistema; descripción declarativa de nodos Fog y clientes Edge; centralización de parámetros de modelo y entrenamiento; validaciones de consistencia. \\ \hline
\textbf{Dependencias} & HU-057, HU-058 \\ \hline
\textbf{Estado} & Pendiente \\ \hline
\end{tabular}
\end{table}



\begin{table}[H]
\centering
\caption{Descripción de la historia de usuario HU-062}
\label{tab:hu062}
\begin{tabular}{|l|p{12cm}|}
\hline
\textbf{ID} & HU-062 \\ \hline
\textbf{Tamaño} & L \\ \hline
\textbf{Nombre} & Estudio del estándar OpenTelemetry \\ \hline
\textbf{Descripción} &
Como investigador, quiero analizar el estándar OpenTelemetry para definir cómo instrumentar el sistema federado y habilitar trazabilidad extremo a extremo. \\ \hline
\textbf{Prioridad} & Media \\ \hline
\textbf{Necesidad} & Importante \\ \hline
\textbf{Alcance} &
Revisión del estándar; identificación de puntos de instrumentación; criterios de integración en Edge--Fog--Server. \\ \hline
\textbf{Dependencias} & - \\ \hline
\textbf{Estado} & Completada \\ \hline
\end{tabular}
\end{table}


\begin{table}[H]
\centering
\caption{Descripción de la historia de usuario HU-063}
\label{tab:hu063}
\begin{tabular}{|l|p{12cm}|}
\hline
\textbf{ID} & HU-063 \\ \hline
\textbf{Tamaño} & L \\ \hline
\textbf{Nombre} & Exploración de herramientas de trazado distribuido \\ \hline
\textbf{Descripción} &
Como investigador, quiero evaluar herramientas de trazado distribuido compatibles con OpenTelemetry para seleccionar una solución de observabilidad adecuada. \\ \hline
\textbf{Prioridad} & Media \\ \hline
\textbf{Necesidad} & Importante \\ \hline
\textbf{Alcance} &
Análisis de herramientas (Jaeger, Zipkin); validación de compatibilidad OTEL; definición de un esquema de integración preliminar. \\ \hline
\textbf{Dependencias} & HU-061 \\ \hline
\textbf{Estado} & Completada \\ \hline
\end{tabular}
\end{table}


\begin{table}[H]
\centering
\caption{Descripción de la historia de usuario HU-064}
\label{tab:hu064}
\begin{tabular}{|l|p{12cm}|}
\hline
\textbf{ID} & HU-064 \\ \hline
\textbf{Tamaño} & XL \\ \hline
\textbf{Nombre} & Análisis preliminar del dataset SWEET \\ \hline
\textbf{Descripción} &
Como investigador, quiero analizar el dataset SWEET para caracterizar su estructura, tamaño y viabilidad experimental en el sistema federado. \\ \hline
\textbf{Prioridad} & Media \\ \hline
\textbf{Necesidad} & Importante \\ \hline
\textbf{Alcance} &
Exploración de estructura por sujetos y sesiones; estimación de requisitos computacionales; evaluación de viabilidad. \\ \hline
\textbf{Dependencias} & - \\ \hline
\textbf{Estado} & Pendiente \\ \hline
\end{tabular}
\end{table}


\begin{table}[H]
\centering
\caption{Descripción de la historia de usuario HU-065}
\label{tab:hu065}
\begin{tabular}{|l|p{12cm}|}
\hline
\textbf{ID} & HU-065 \\ \hline
\textbf{Tamaño} & M \\ \hline
\textbf{Nombre} & Entrenamiento centralizado en servidor de un modelo avanzado sobre SWEET \\ \hline
\textbf{Descripción} &
Como investigador, quiero entrenar de forma centralizada un modelo avanzado sobre el dataset SWEET para establecer un baseline comparativo frente al aprendizaje federado. \\ \hline
\textbf{Prioridad} & Media \\ \hline
\textbf{Necesidad} & Importante \\ \hline
\textbf{Alcance} &
Entrenamiento centralizado; evaluación base de rendimiento; comparación con resultados federados. \\ \hline
\textbf{Dependencias} & HU-064 \\ \hline
\textbf{Estado} & Pendiente \\ \hline
\end{tabular}
\end{table}


\begin{table}[H]
\centering
\caption{Descripción de la historia de usuario HU-066}
\label{tab:hu066}
\begin{tabular}{|l|p{12cm}|}
\hline
\textbf{ID} & HU-066 \\ \hline
\textbf{Tamaño} & L \\ \hline
\textbf{Nombre} & Finalización del sistema federado SWELL \\ \hline
\textbf{Descripción} &
Como investigador, quiero completar la integración del sistema federado SWELL con telemetría y métricas para disponer de un entorno final operativo y observable. \\ \hline
\textbf{Prioridad} & Alta \\ \hline
\textbf{Necesidad} & Importante \\ \hline
\textbf{Alcance} &
Integración de componentes Edge--Fog--Server; instrumentación de telemetría; validación del flujo extremo a extremo. \\ \hline
\textbf{Dependencias} & HU-004, HU-006 \\ \hline
\textbf{Estado} & Completada \\ \hline
\end{tabular}
\end{table}


\begin{table}[H]
\centering
\caption{Descripción de la historia de usuario HU-067}
\label{tab:hu067}
\begin{tabular}{|l|p{12cm}|}
\hline
\textbf{ID} & HU-067 \\ \hline
\textbf{Tamaño} & M \\ \hline
\textbf{Nombre} & Testing del sistema final \\ \hline
\textbf{Descripción} &
Como investigador, quiero realizar pruebas integrales del sistema federado para validar su funcionamiento extremo a extremo. \\ \hline
\textbf{Prioridad} & Media \\ \hline
\textbf{Necesidad} & Importante \\ \hline
\textbf{Alcance} &
Tests de integración y sistema; validación de componentes MQTT y de agregación; verificación de consistencia. \\ \hline
\textbf{Dependencias} & HU-064 \\ \hline
\textbf{Estado} & Completada \\ \hline
\end{tabular}
\end{table}


\begin{table}[H]
\centering
\caption{Descripción de la historia de usuario HU-068}
\label{tab:hu068}
\begin{tabular}{|l|p{12cm}|}
\hline
\textbf{ID} & HU-068 \\ \hline
\textbf{Tamaño} & L \\ \hline
\textbf{Nombre} & Modelos neuronales avanzados para estrés fisiológico \\ \hline
\textbf{Descripción} &
Como investigador, quiero implementar modelos neuronales avanzados para mejorar la clasificación de estrés fisiológico en el sistema federado. \\ \hline
\textbf{Prioridad} & Media \\ \hline
\textbf{Necesidad} & Importante \\ \hline
\textbf{Alcance} &
Diseño de arquitecturas profundas; ajuste de regularización; integración con el pipeline de entrenamiento. \\ \hline
\textbf{Dependencias} & - \\ \hline
\textbf{Estado} & Pendiente \\ \hline
\end{tabular}
\end{table}


\begin{table}[H]
\centering
\caption{Descripción de la historia de usuario HU-069}
\label{tab:hu069}
\begin{tabular}{|l|p{12cm}|}
\hline
\textbf{ID} & HU-069\\ \hline
\textbf{Tamaño} & M \\ \hline
\textbf{Nombre} & Implementación de métricas específicas para aprendizaje federado \\ \hline
\textbf{Descripción} &
Como investigador, quiero definir e implementar métricas específicas de aprendizaje federado para monitorizar rondas, rendimiento y participación de clientes. \\ \hline
\textbf{Prioridad} & Media \\ \hline
\textbf{Necesidad} & Importante \\ \hline
\textbf{Alcance} &
Definición de métricas FL; integración con Prometheus; dashboards para visualización. \\ \hline
\textbf{Dependencias} & HU-006 \\ \hline
\textbf{Estado} & Pendiente \\ \hline
\end{tabular}
\end{table}


\begin{table}[H]
\centering
\caption{Descripción de la historia de usuario HU-070}
\label{tab:hu070}
\begin{tabular}{|l|p{12cm}|}
\hline
\textbf{ID} & HU-070 \\ \hline
\textbf{Tamaño} & L \\ \hline
\textbf{Nombre} & Automatización del despliegue de arquitecturas federadas \\ \hline
\textbf{Descripción} &
Como investigador, quiero desplegar arquitecturas federadas de forma automatizada a partir de configuraciones para ejecutar experimentos reproducibles sin intervención manual. \\ \hline
\textbf{Prioridad} & Media \\ \hline
\textbf{Necesidad} & Importante \\ \hline
\textbf{Alcance} &
Launcher config-driven; coordinación de procesos; aplicación de manifiestos de despliegue. \\ \hline
\textbf{Dependencias} & HU-061 \\ \hline
\textbf{Estado} & Pendiente \\ \hline
\end{tabular}
\end{table}


\begin{table}[H]
\centering
\caption{Descripción de la historia de usuario HU-071}
\label{tab:hu071}
\begin{tabular}{|l|p{12cm}|}
\hline
\textbf{ID} & HU-071 \\ \hline
\textbf{Tamaño} & L \\ \hline
\textbf{Nombre} & Implementación de stack completo de observabilidad containerizado \\ \hline
\textbf{Descripción} &
Como investigador, quiero desplegar un stack completo de observabilidad containerizado para monitorizar métricas y trazas del sistema federado. \\ \hline
\textbf{Prioridad} & Media \\ \hline
\textbf{Necesidad} & Importante \\ \hline
\textbf{Alcance} &
Jaeger, Prometheus, Grafana y OTEL Collector; integración de telemetría distribuida. \\ \hline
\textbf{Dependencias} & HU-062, HU-064, HU-069\\ \hline
\textbf{Estado} & Pendiente \\ \hline
\end{tabular}
\end{table}

\newpage
\section{Sprint 1}
\label{sect:Sp1}
El Sprint 1 marca el inicio del trabajo de desarrollo propiamente dicho. Su objetivo principal no fue todavía implementar la funcionalidad completa del sistema, sino reducir la incertidumbre técnica existente tras el Sprint 0 y validar qué herramientas y frameworks resultaban viables para materializar la arquitectura propuesta.
De acuerdo con la planificación realizada, el Sprint 1 se estructuró en torno a las historias de usuario recogidas en la Tabla~\ref{tab:sprint1}.
\begin{table}[H]
\centering
\footnotesize
\caption{Sprint 1: Historias de usuario y tareas asociadas}
\label{tab:sprint1}
\renewcommand{\arraystretch}{1.25}
\setlength{\tabcolsep}{4pt}

\begin{tabular}{|l|p{4.2cm}|c|p{6.2cm}|}
\hline
\textbf{ID} & \textbf{Historia de usuario} & \textbf{Tamaño} & \textbf{Tareas principales (resumen)} \\ \hline

HU-019 &
Estandarización automática del código &
S &
\begin{itemize}\setlength\itemsep{0.2em}
  \item Configurar formateo automático conforme a \textit{PEP~8}.
  \item Integrar \texttt{black} y \texttt{flake8} para validación de estilo.
\end{itemize}
\\ \hline

HU-020 &
Configuración de integración continua (CI) &
S &
\begin{itemize}\setlength\itemsep{0.2em}
  \item Crear repositorio en GitHub.
  \item Definir estrategia de ramas (\texttt{main}/\texttt{develop}).
  \item Configurar GitHub Actions para validaciones automáticas.
\end{itemize}
\\ \hline

HU-021 &
Configuración del entorno de desarrollo en Python &
S &
\begin{itemize}\setlength\itemsep{0.2em}
  \item Seleccionar versión de Python compatible con los frameworks evaluados.
  \item Configurar el entorno con \textbf{Python 3.10}.
  \item Definir dependencias iniciales del proyecto.
\end{itemize}
\\ \hline

HU-022 &
Estudio de frameworks de Aprendizaje Federado &
M &
\begin{itemize}\setlength\itemsep{0.2em}
  \item Exploración inicial de frameworks prometedores.
  \item Primer prototipado simplificado para validar integración y ejecución.
\end{itemize}
\\ \hline

HU-023 &
Estudio de TensorFlow Federated &
S &
\begin{itemize}\setlength\itemsep{0.2em}
  \item Instalar y configurar \textit{TensorFlow Federated}.
  \item Desarrollar prototipo mínimo cliente--servidor.
  \item Ejecutar pruebas preliminares con un dataset sencillo.
\end{itemize}
\\ \hline

HU-024 &
Revisión bibliográfica de frameworks de Aprendizaje Federado &
M &
\begin{itemize}\setlength\itemsep{0.2em}
  \item Búsqueda complementaria de literatura y herramientas emergentes.
  \item Extracción de criterios de comparación relevantes.
\end{itemize}
\\ \hline

HU-025 &
Análisis del framework FLEX (UGR) &
M &
\begin{itemize}\setlength\itemsep{0.2em}
  \item Identificar y analizar FLEX como alternativa.
  \item Evaluar alineación con el proyecto (seguridad, privacidad, madurez).
\end{itemize}
\\ \hline

HU-026 &
Consolidación de criterios técnicos para evaluación de frameworks &
S &
\begin{itemize}\setlength\itemsep{0.2em}
  \item Definir criterios técnicos (madurez, soporte, extensibilidad, etc.).
  \item Preparar rúbrica de evaluación (puntuación 1--5).
\end{itemize}
\\ \hline

HU-027 &
Contacto y coordinación de reunión técnica con equipo FLEX &
XS &
\begin{itemize}\setlength\itemsep{0.2em}
  \item Contactar con el investigador responsable.
  \item Planificar reunión técnica para el siguiente sprint.
\end{itemize}
\\ \hline

HU-028 &
Documentación de decisiones técnicas y evidencias (Sprint 1) &
S &
\begin{itemize}\setlength\itemsep{0.2em}
  \item Documentar estructura del repo y decisiones técnicas.
  \item Mantener trazabilidad en \texttt{README} y notas del sprint.
\end{itemize}
\\ \hline

\end{tabular}
\end{table}
\subsection{Desarrollo del Sprint 1}

\paragraph{HU-020: Configuración de la infraestructura de control de versiones}
Durante el Sprint 1 se inició la configuración de la infraestructura básica del proyecto. En el marco de esta historia de usuario se seleccionó \textbf{GitHub} como plataforma de control de versiones y se configuró el repositorio inicial, definiendo una estructura de ramas que separa el código estable del desarrollo en curso. Esta estrategia, basada en el uso de \textit{Git} como sistema de control de versiones distribuido, sigue las recomendaciones recogidas en el material docente \emph{Aprende Git} \cite{aprendeGitUGR}, proporcionando un flujo de trabajo sencillo, robusto y adecuado para proyectos académicos y de investigación.

\paragraph{HU-019: Estandarización automática del código}
De forma complementaria, se incorporaron herramientas de formateo automático y análisis estático del código con el objetivo de garantizar la calidad, coherencia y mantenibilidad del software desde las primeras iteraciones del proyecto. En concreto, se estableció el cumplimiento del estándar \textit{PEP~8} \cite{pep8}, que define una guía oficial de estilo para el código Python y proporciona recomendaciones ampliamente aceptadas en cuanto a nomenclatura, estructura y legibilidad. La configuración de \texttt{black} se definió mediante el archivo \texttt{pyproject.toml}, estableciendo una longitud máxima de línea de 88 caracteres y fijando como versiones objetivo de Python las comprendidas entre 3.8 y 3.11. De forma coherente con esta configuración, \texttt{flake8} se configuró mediante el archivo \texttt{.flake8}, adoptando la misma longitud máxima de línea y deshabilitando advertencias que entran en conflicto con el estilo impuesto por \texttt{black} o que no resultan relevantes en el contexto del proyecto.

\paragraph{HU-021: Configuración del entorno de desarrollo en Python}
A continuación, se abordó la configuración del entorno de desarrollo del proyecto. Se fijó \textbf{Python 3.10} como versión base, atendiendo a criterios de compatibilidad con los frameworks de aprendizaje federado que serían evaluados en fases posteriores, así como a la estabilidad y madurez del ecosistema de librerías científicas y de aprendizaje automático asociadas a dicha versión. Con el objetivo de aislar las dependencias del proyecto, evitar interferencias con otras instalaciones del sistema y facilitar una futura replicación del entorno mediante contenedores Docker, se optó por el uso de un entorno virtual local gestionado mediante el módulo estándar \texttt{venv}. Las dependencias iniciales del proyecto se documentaron en un fichero \texttt{requirements.txt}, facilitando la reproducibilidad del entorno \cite{pythonVenvDocs}.

\paragraph{HU-023: Estudio de TensorFlow Federated}
Como parte del trabajo planificado para el Sprint 1, se abordó el estudio específico del framework \textit{TensorFlow Federated}. En el marco de esta historia de usuario se procedió a la instalación y configuración inicial de la librería, así como al desarrollo de un prototipo mínimo basado en una arquitectura clásica cliente--servidor. El objetivo de este prototipo no fue obtener un sistema funcional completo, sino validar aspectos básicos de integración, ejecución y comunicación entre los nodos participantes.

Para estas pruebas preliminares se empleó el conjunto de datos \textit{ECG5000}, perteneciente al archivo UCR de clasificación de series temporales \cite{Dau2019UCR}. La elección de este dataset respondió a su tamaño moderado y a su estructura sencilla, lo que permitió centrar el análisis en el comportamiento del framework y no en la complejidad del modelo o de los datos. Este enfoque facilitó la identificación temprana de limitaciones técnic


\paragraph{HU-024: Revisión bibliográfica de frameworks de Aprendizaje Federado}
En paralelo a las tareas de infraestructura, se llevó a cabo una revisión bibliográfica adicional centrada en frameworks de aprendizaje federado, con el objetivo de identificar alternativas relevantes y comprender el estado actual del arte. Como resultado de esta revisión se identificó, entre otros trabajos, un framework denominado \emph{FLEX}, desarrollado por un grupo de investigación de la Universidad de Granada.

\paragraph{HU-025: Análisis del framework FLEX (UGR)}
El framework \emph{FLEX} fue analizado de forma más específica, evaluando su orientación a la experimentación con mecanismos avanzados de seguridad y privacidad, así como su posible encaje dentro de los objetivos del proyecto. Dado el interés suscitado por esta propuesta, se estableció contacto con uno de los investigadores responsables del framework y se planificó una reunión técnica para el siguiente sprint (HU-027), con el fin de profundizar en su diseño, grado de madurez y aplicabilidad práctica.

\paragraph{HU-027: Contacto y coordinación de reunión técnica con el equipo FLEX}
Como parte de esta historia de usuario, se estableció contacto con el equipo desarrollador del framework FLEX y se coordinó una reunión técnica para el Sprint 2, orientada a contrastar la viabilidad real del framework y obtener una visión directa sobre su estado de desarrollo.

\paragraph{HU-022: Estudio exploratorio de frameworks de Aprendizaje Federado}
Finalmente, se inició un estudio general de distintos frameworks de aprendizaje federado desde un enfoque exploratorio. Este estudio no tuvo como objetivo una implementación completa, sino la identificación temprana de limitaciones técnicas, barreras de entrada y posibles incompatibilidades con la arquitectura propuesta, permitiendo descartar de forma eficiente aquellas soluciones que no resultaban adecuadas para el contexto del proyecto.

\paragraph{HU-026: Consolidación de criterios técnicos de evaluación}
Como complemento al análisis exploratorio del estado del arte, se definió y consolidó un conjunto de \textbf{criterios técnicos de evaluación} orientados a comparar frameworks de Aprendizaje Federado en condiciones homogéneas. Esta formalización responde a una necesidad práctica: los frameworks disponibles presentan diferencias sustanciales en \textit{arquitectura}, \textit{mecanismos soportados} y \textit{nivel de madurez}, por lo que una elección basada únicamente en impresiones, ejemplos puntuales o documentación promocional introduciría sesgos y dificultaría la reproducibilidad del proceso de selección.

La Tabla~\ref{tab:criterios-frameworks} recoge estos criterios como un \textbf{constructo propio} adaptado al contexto del proyecto, construido mediante una \textbf{síntesis} de dimensiones recurrentes identificadas en la literatura que compara herramientas y entornos de FL. En particular, los trabajos revisados suelen evaluar (i) aspectos operativos de adopción (\textit{instalación, estabilidad, documentación y soporte}), (ii) capacidades del framework para experimentar con estrategias, topologías y configuraciones (\textit{flexibilidad y extensibilidad}), y (iii) propiedades críticas del entrenamiento federado en despliegues reales, como la heterogeneidad \textit{non-i.i.d.} y las consideraciones de privacidad/seguridad \cite{Riedel2024Frameworks,Boekhoff2025Comparison,Elshair2024Simulators,FedEvalSurvey2023,NonIIDReview2024}. 

Adicionalmente, para evitar que los criterios quedasen desconectados de prácticas consolidadas en ingeniería del software, se alinearon con dimensiones generales de calidad (p.\ ej., compatibilidad, mantenibilidad, fiabilidad y seguridad) recogidas en \textbf{ISO/IEC 25010}, reinterpretándolas para el caso específico de bibliotecas y plataformas FL \cite{ISO25010}. Bajo este enfoque, \emph{viabilidad técnica} e \emph{integración con el ecosistema Python} capturan la factibilidad de adopción e interoperabilidad; \emph{madurez} y \emph{soporte comunitario/documentación} aproximan estabilidad y mantenibilidad; \emph{flexibilidad arquitectónica} refleja la capacidad de operar en topologías centralizadas o jerárquicas (alineadas con los escenarios Edge--Fog--Cloud tratados en el estado del arte); mientras que \emph{soporte no-i.i.d.} y \emph{privacidad y seguridad} incorporan retos considerados estructurales en FL, especialmente relevantes en dominios sensibles y entornos IoT \cite{kairouz2021advances,li2020federated}. Finalmente, todos los criterios se plantean para ser evaluados en una \textbf{escala ordinal 1--5}, con el objetivo de facilitar una comparación consistente entre alternativas sin asumir una precisión artificial en métricas que dependen del contexto de uso.


\paragraph{HU-028: Documentación de decisiones técnicas y evidencias del Sprint 1}
Todas las decisiones técnicas, conclusiones y evidencias generadas durante el Sprint 1 fueron recopiladas y documentadas de forma sistemática, garantizando la trazabilidad del proceso de toma de decisiones y facilitando la continuidad del proyecto en iteraciones posteriores.

\subsubsection{Resultados del \textit{Sprint} 1}

Los resultados obtenidos durante el \textit{Sprint} 1 se presentan organizados por historia de usuario, de acuerdo con la planificación recogida en la Tabla~\ref{tab:sprint1}. En esta sección se resumen los resultados alcanzados y el estado final de cada historia, evitando repetir los detalles técnicos ya descritos en la fase de desarrollo.

\paragraph{HU-019: Estandarización automática del código}
La HU-019 quedó completada, estableciéndose un sistema de validación automática del estilo del código integrado en el flujo de trabajo del proyecto. Como resultado, el código del proyecto mantiene una estructura homogénea y consistente desde las primeras iteraciones, sentando una base sólida para su mantenibilidad a medio y largo plazo.

\paragraph{HU-020: Configuración de integración continua (CI)}
Como resultado de la HU-020 se definió una infraestructura funcional de control de versiones e integración continua. El repositorio quedó estructurado con un flujo de ramas estable y se implementó una canalización inicial de CI que valida automáticamente el código antes de su integración en la rama principal. Esta infraestructura constituye el punto de partida para la incorporación progresiva de pruebas más avanzadas en sprints posteriores.  
La configuración inicial de dicha canalización se muestra en la Figura~\ref{fig:cicd-pipeline}.

\begin{figure}[H]
\centering
\includegraphics[width=\linewidth]{Plantilla-Latex-TFG/imagenes/CICDv1.png}
\caption{Configuración inicial de la canalización de Integración Continua (CI) en GitHub Actions}
\label{fig:cicd-pipeline}
\end{figure}

\paragraph{HU-021: Configuración del entorno de desarrollo en Python}
La HU-021 se completó con la definición de un entorno de desarrollo estable y reproducible basado en Python~3.10. El entorno quedó documentado junto con sus dependencias, permitiendo su recreación consistente en otros sistemas y facilitando la continuidad del desarrollo en los sprints posteriores.

\paragraph{HU-022: Estudio de frameworks de Aprendizaje Federado}
La HU-022 permitió realizar una primera exploración de frameworks de aprendizaje federado con un enfoque preliminar. Como resultado, se identificaron herramientas potencialmente viables y se detectaron limitaciones técnicas tempranas, quedando esta historia parcialmente completada y planificada para su continuación en iteraciones posteriores.

\paragraph{HU-023: Estudio de TensorFlow Federated}
Como resultado de la HU-023 se validó la viabilidad inicial del framework \textit{TensorFlow Federated} mediante un prototipo mínimo basado en una arquitectura clásica cliente--servidor, representada en la Figura~\ref{fig:Arqbasic}. Las pruebas preliminares se realizaron empleando el conjunto de datos \textit{ECG5000}, perteneciente al archivo UCR de clasificación de series temporales.

No obstante, durante la experimentación se detectaron problemas de estabilidad en la comunicación entre los nodos cliente y el servidor federado, lo que impidió completar una validación funcional del framework en este sprint. Como consecuencia, la historia quedó parcialmente completada y su resolución se trasladó al Sprint~2.

\begin{figure}[H]
\centering
\includegraphics[width=\linewidth]{Plantilla-Latex-TFG/imagenes/ArquitecturaFlpruebabasica.png}
\caption{Arquitectura clásica cliente--servidor empleada en las pruebas preliminares de aprendizaje federado}
\label{fig:Arqbasic}
\end{figure}

\paragraph{HU-024: Revisión bibliográfica de frameworks de Aprendizaje Federado}
Como resultado de la HU-024 se llevó a cabo una revisión bibliográfica complementaria que permitió identificar criterios relevantes de comparación y detectar alternativas de interés dentro del ecosistema de aprendizaje federado.

\paragraph{HU-025 y HU-027: Análisis de FLEX y coordinación de reunión técnica}
En el marco de la HU-025 se analizó el framework \emph{FLEX}, concluyéndose que su enfoque resulta conceptualmente alineado con los objetivos del proyecto, aunque su grado de madurez limita su adopción inmediata. Como resultado de la HU-027 se estableció contacto con el equipo desarrollador y se planificó una reunión técnica para profundizar en su evaluación durante el siguiente sprint.

\paragraph{HU-026: Consolidación de criterios técnicos para evaluación de frameworks}
Como resultado de la HU-026 se definió una rúbrica estructurada de criterios técnicos para la evaluación comparativa de frameworks de aprendizaje federado. Estos criterios proporcionan una base objetiva y reproducible que será empleada sistemáticamente en sprints posteriores. La Tabla~\ref{tab:criterios-frameworks} recoge el conjunto de criterios definidos.

\begin{table}[H]
\centering
\footnotesize
\caption{Criterios definidos para la evaluación comparativa de frameworks de Aprendizaje Federado (la evaluación posterior emplea una escala 1--5)}
\label{tab:criterios-frameworks}

\renewcommand{\arraystretch}{1.15}
\setlength{\tabcolsep}{4pt}

\begin{tabularx}{\linewidth}{@{}p{3.6cm} >{\RaggedRight\arraybackslash}X @{}}
\toprule
\textbf{Criterio} & \textbf{Descripción} \\
\midrule
Viabilidad técnica
& Capacidad del framework para ser instalado, configurado y ejecutado de forma estable en el entorno del proyecto. \\

Madurez del framework
& Grado de desarrollo del framework, estabilidad del código, número de versiones disponibles y continuidad del proyecto. \\

Soporte comunitario y documentación
& Calidad y disponibilidad de documentación, ejemplos, tutoriales y nivel de actividad de la comunidad de usuarios y desarrolladores. \\

Flexibilidad arquitectónica
& Capacidad del framework para adaptarse a distintas topologías de aprendizaje federado, incluyendo arquitecturas centralizadas, jerárquicas y Edge--Fog--Cloud. \\

Soporte para múltiples algoritmos de ML
& Posibilidad de emplear distintos modelos y algoritmos de aprendizaje automático, como \textit{redes neuronales} o árboles de decisión. \\

Soporte para escenarios no-i.i.d.
& Herramientas o facilidades para gestionar distribuciones de datos heterogéneas y no independientes entre los nodos participantes. \\

Privacidad y seguridad
& Disponibilidad de mecanismos orientados a la preservación de la privacidad y la seguridad, como privacidad diferencial, cifrado o defensas frente a ataques. \\

Capacidad de experimentación
& Facilidad para modificar parámetros, estrategias de agregación y flujos de entrenamiento con fines de investigación experimental. \\

Extensibilidad
& Posibilidad de extender o personalizar el framework mediante la incorporación de módulos o componentes propios. \\

Integración con el ecosistema Python
& Compatibilidad e integración con librerías habituales del ecosistema Python, como NumPy, PyTorch o TensorFlow. \\
\bottomrule
\end{tabularx}
\end{table}

\paragraph{HU-028: Documentación de decisiones técnicas y evidencias}
Finalmente, como resultado de la HU-028 se documentaron de forma sistemática las decisiones técnicas adoptadas, las configuraciones realizadas y las evidencias generadas durante el Sprint~1, garantizando la trazabilidad del proceso de desarrollo.



\subsection{Sprint Review (Sprint 1)}
El Sprint 1 sirvió para aterrizar el trabajo de evaluación de frameworks de aprendizaje federado y evidenció un aprendizaje clave de planificación: las \textbf{estimaciones iniciales} no capturaban bien la \textbf{incertidumbre técnica} asociada a validar stacks FL en escenarios reales. En particular, HU-021 resultó demasiado amplia y se refinó en historias más acotadas, mientras que HU-022 quedó bloqueada por problemas operativos y se trasladó al Sprint 2.

\vspace{0.15em}
\paragraph{Qué se logró.}
Se estableció una base técnica y metodológica para evaluar frameworks FL, clarificando criterios y necesidades de validación. Además, se produjo un refinamiento del backlog: HU-021 se descompuso en historias más pequeñas (HU-022, HU-023 y una posible HU para FLEX), incrementando control y trazabilidad del avance.

\vspace{0.15em}
\paragraph{Qué fue bien.}
Se planificó correctamente trabajo \textbf{estructural y habilitador} propio de una primera iteración (criterios, preparación técnica y delimitación del problema), evitando forzar implementación prematura sin una base mínima.

\vspace{0.15em}
\paragraph{Qué se detectó / qué fue mal.}
Se subestimó el \textbf{tamaño y alcance} de HU-021 al plantearla como una evaluación “única”, sin granularidad suficiente para controlar riesgos. También se infravaloró el \textbf{riesgo técnico} asociado a HU-022 (TensorFlow Federated): no se reservó margen explícito para bloqueos de ejecución y restricciones cliente--servidor típicas en validaciones de frameworks, lo que impidió completar su validación en el sprint.

\vspace{0.15em}
\paragraph{Acciones y siguientes pasos.}
Mantener historias de evaluación con \textbf{finalidad acotada} y criterio de salida verificable. Continuar con HU-023 y reabordar HU-022 con validación incremental (prueba mínima primero; escalado después), para evitar consumo de tiempo sin retorno.




\section{Sprint 2}
\label{sect:Sp2}
El Sprint 2 se centró íntegramente en la \textbf{evaluación práctica} de frameworks y herramientas de aprendizaje federado, con un objetivo decisivo: \textbf{seleccionar el stack tecnológico} sobre el que se desarrollará el sistema en los siguientes sprints. A diferencia del Sprint 1, orientado a reducir incertidumbre inicial y preparar el entorno, esta iteración tuvo un carácter \textbf{analítico y resolutivo}, ya que sus resultados determinaban qué tecnologías podían sostener el desarrollo posterior.

De acuerdo con la planificación realizada, el Sprint 2 se estructuró en torno a las historias de usuario recogidas en la Tabla~\ref{tab:sprint2}. Estas historias se estimaron en función de su tamaño relativo y se seleccionaron por su impacto directo en la toma de decisiones tecnológicas.

\begin{table}[H]
\centering
\footnotesize
\caption{Sprint 2: Historias de usuario y tareas asociadas}
\label{tab:sprint2}
\renewcommand{\arraystretch}{1.25}
\setlength{\tabcolsep}{4pt}

\begin{tabular}{|l|p{4.6cm}|c|p{5.8cm}|}
\hline
\textbf{ID} & \textbf{Historia de usuario} & \textbf{Tamaño} & \textbf{Tareas principales (resumen)} \\ \hline

HU-029 &
Estudio del framework Flower &
M &
\begin{itemize}\setlength\itemsep{0.2em}
  \item Prototipo federado mínimo (cliente--servidor).
  \item Ejecución con clientes múltiples y estrategia FedAvg.
  \item Validación de estabilidad y resultados.
\end{itemize}
\\ \hline

HU-030 &
Implementación de modelos baseline sobre ECG5000 &
M &
\begin{itemize}\setlength\itemsep{0.2em}
  \item Definir baselines centralizados comparables.
  \item Entrenar y reportar resultados con 5-fold CV.
\end{itemize}
\\ \hline

HU-031 &
Evaluación técnica y decisión sobre viabilidad de FLEX &
M &
\begin{itemize}\setlength\itemsep{0.2em}
  \item Reunión técnica con el equipo desarrollador.
  \item Contrastar soporte real para Edge--Fog--Cloud.
  \item Decisión de adopción/descartado.
\end{itemize}
\\ \hline

HU-032 &
Documentación de decisiones técnicas y conclusiones (Sprint 2) &
S &
\begin{itemize}\setlength\itemsep{0.2em}
  \item Registrar evidencias y conclusiones.
  \item Consolidar tablas comparativas y resultados.
\end{itemize}
\\ \hline
\end{tabular}
\end{table}


\subsection{Desarrollo del Sprint 2}

\paragraph{HU-031: Evaluación técnica y decisión sobre viabilidad del framework FLEX}
Se celebró una reunión técnica con Mario García-Márquez (equipo desarrollador de \emph{FLEX}) para contrastar su viabilidad en el contexto del TFM. Se presentó la arquitectura Edge--Fog--Cloud propuesta y se discutió su encaje como base tecnológica. Como conclusión de la reunión, se determinó que \emph{FLEX} \textbf{no resulta viable} para una arquitectura jerárquica con nodos Fog en su estado actual, debido a que el framework no está completamente funcional y mantiene características clave sin implementar. Aun así, se destacó su enfoque conceptual orientado a la experimentación (p.\ ej., estrategias de agregación) como línea interesante para trabajos futuros.

\paragraph{HU-029: Estudio del framework Flower}
Se desarrolló un prototipo federado mínimo con \textit{Flower} para validar integración, ejecución y entrenamiento distribuido. El trabajo se orientó a confirmar la viabilidad práctica del framework en un escenario funcional de aprendizaje federado con múltiples clientes y agregación tipo \textit{FedAvg}, priorizando reproducibilidad y estabilidad de ejecución.

\paragraph{HU-030: Implementación de modelos baseline sobre ECG5000}
Se implementaron modelos baseline centralizados sobre ECG5000 para disponer de una referencia comparativa directa frente a los resultados federados. El objetivo fue aislar el impacto del entrenamiento distribuido respecto al rendimiento del modelo centralizado equivalente, manteniendo arquitecturas comparables.

\paragraph{HU-032: Documentación de decisiones técnicas y conclusiones (Sprint 2)}
Se documentaron las evidencias y conclusiones del sprint para garantizar trazabilidad: decisiones de descarte/adopción de frameworks, tablas comparativas y resultados experimentales obtenidos en ECG5000.


\subsection{Resultados del Sprint 2}

Los resultados se presentan organizados por historia de usuario, resumiendo el estado final alcanzado sin repetir los detalles técnicos del desarrollo.

\paragraph{HU-031: Evaluación de FLEX}
Se concluyó que \emph{FLEX} no puede adoptarse como tecnología base dentro del alcance temporal del TFM por falta de soporte operativo para una arquitectura jerárquica Edge--Fog--Cloud en su estado actual. Queda como alternativa potencial para trabajos futuros cuando alcance mayor madurez.

\paragraph{HU-029: Estudio de Flower}
\textit{Flower} permitió implementar un escenario federado funcional sobre ECG5000 con tres clientes y estrategia \textit{FedAvg}, confirmando su viabilidad como base del desarrollo. Se identificó como limitación relevante que el flujo de agregación requiere disponibilidad completa de clientes bajo la configuración empleada, lo que entra en conflicto con RNF-005 (Fiabilidad) y obliga a plantear mecanismos de tolerancia a fallos o adaptar la orquestación del entrenamiento.

\paragraph{HU-030: Baselines centralizados ECG5000}
Se obtuvieron resultados baseline centralizados mediante 5-fold cross-validation, que se emplearon como referencia para cuantificar la degradación de rendimiento asociada al entrenamiento federado.

\paragraph{HU-032: Documentación}
Se consolidaron evidencias y conclusiones del sprint, incluyendo tablas comparativas de frameworks y resultados de rendimiento centralizado vs federado.


\begin{table}[H]
\centering
\footnotesize
\caption{Tabla comparativa de frameworks de Aprendizaje Federado (valoración 1--5)}
\label{tab:comparativa-frameworks}

\renewcommand{\arraystretch}{1.25}
\setlength{\tabcolsep}{4pt}

\begin{tabular}{@{} l c c c @{}}  
\toprule
\textbf{Criterio} 
& \textbf{TensorFlow Federated} 
& \textbf{Flower} 
& \textbf{FLEX} \\
\midrule
Viabilidad técnica & 2 & 5 & - \\
Madurez del framework & 3 & 5 & 2 \\
Soporte comunitario y documentación & 3 & 5 & 1 \\
Facilidad de instalación y configuración & 3 & 5 & 2 \\
Flexibilidad arquitectónica & 2 & 4 & 5 \\
Soporte para múltiples algoritmos de ML & 3 & 3 & 3 \\
Soporte para escenarios no-i.i.d. & 3 & 4 & - \\
Privacidad y seguridad & 4 & 4 & 5 \\
Capacidad de experimentación & 2 & 4 & 5 \\
Extensibilidad & 3 & 5 & 4 \\
Integración con el ecosistema Python & 3 & 5 & 3 \\
Soporte para tolerancia a fallos & 2 & 2 & - \\
Adecuación a arquitectura Edge--Fog--Cloud & 2 & 3 & - \\
\bottomrule
\end{tabular}
\end{table}


\begin{table}[H]
\centering
\footnotesize
\caption{Comparativa de rendimiento en ECG5000: modelo centralizado vs aprendizaje federado (FedAvg) sobre Flower. Media de exactitud mediante validación cruzada de 5 particiones (5-fold cross-validation).}
\label{tab:ECG5000-baseline-vs-fedavg}

\renewcommand{\arraystretch}{1.25}
\setlength{\tabcolsep}{4pt}

\begin{tabular}{@{} l c c c @{}}  
\toprule
\textbf{Arquitectura} 
& \textbf{Baseline Centralizado} 
& \textbf{Federated (FedAvg)} 
& \textbf{Diferencia} \\
\midrule
MLP simple (2 capas densas) & 98.1\% & 96.6\% & $-1.5\%$ \\
CNN ligera (Conv--ReLU--Pool) & 99.0\% & 97.8\% & $-1.2\%$ \\
CNN profunda (2 Conv + Dense) & 99.3\% & 98.0\% & $-1.3\%$ \\
\bottomrule
\end{tabular}
\end{table}


\subsection{Sprint Review (Sprint 2)}
El Sprint 2 permitió \textbf{cerrar la incertidumbre tecnológica} y tomar una decisión fundamentada sobre el stack del proyecto: \textbf{Flower} se confirmó como opción principal, mientras que \textit{TensorFlow Federated} y \emph{FLEX} se descartaron como base por inviabilidad práctica (TFF) y falta de madurez/soporte operativo para Edge--Fog--Cloud (FLEX).

\vspace{0.15em}
\paragraph{Qué se logró.}
Se validó Flower con un prototipo funcional y resultados comparables (baseline vs federado), reduciendo el riesgo del proyecto al fijar una tecnología viable para el desarrollo posterior. En paralelo, se identificó con claridad que TFF no era ejecutable de forma estable en el contexto del proyecto y que FLEX no ofrecía el soporte operativo necesario para la arquitectura objetivo.

\vspace{0.15em}
\paragraph{Qué fue bien.}
Se dimensionó correctamente el esfuerzo para \textbf{validar Flower con un prototipo operativo}, alcanzando el objetivo de cerrar la decisión tecnológica con evidencias y no solo con revisión teórica.

\vspace{0.15em}
\paragraph{Qué se detectó / qué fue mal.}
Se subestimó el coste de \textbf{bloqueos operativos} asociados a frameworks no maduros o difíciles de depurar (especialmente TFF), lo que consumió tiempo sin retorno experimental. Esto confirmó que, en evaluaciones de frameworks, el cuello de botella no es solo conceptual, sino de \textbf{ejecutabilidad}.

\vspace{0.15em}
\paragraph{Acciones y siguientes pasos.}
Formalizar un \textbf{criterio de salida} para evaluaciones: si un framework no ejecuta un experimento mínimo en un tiempo acotado, se descarta. En paralelo, comenzar a diseñar una estrategia práctica en Flower para cumplir RNF-005, incorporando tolerancia a indisponibilidad (clientes opcionales, agregación parcial, orquestación más robusta o mecanismos propios).

\section{Sprint 3}
\label{sect:Sp3}
El Sprint 3 se centró en resolver el principal bloqueo identificado en el Sprint 2: la \textbf{viabilidad arquitectónica} de una solución Edge--Fog--Cloud cuando el framework seleccionado (\textit{Flower}) impone restricciones de ejecución. Con \textit{Flower} ya consolidado como tecnología base, el objetivo del sprint fue definir una forma \textbf{implementable} de materializar la arquitectura sin comprometer requisitos clave, prestando especial atención a asincronía, tolerancia a fallos y escalabilidad.

De acuerdo con la planificación realizada, el Sprint 3 se estructuró en torno a las historias de usuario recogidas en la Tabla~\ref{tab:sprint3}. Estas historias se seleccionaron por su carácter decisivo para traducir el diseño conceptual del sistema a decisiones tecnológicas concretas.

\begin{table}[H]
\centering
\footnotesize
\caption{Sprint 3: Historias de usuario y tareas asociadas}
\label{tab:sprint3}
\renewcommand{\arraystretch}{1.25}
\setlength{\tabcolsep}{4pt}

\begin{tabular}{|l|p{4.6cm}|c|p{5.8cm}|}
\hline
\textbf{ID} & \textbf{Historia de usuario} & \textbf{Tamaño} & \textbf{Tareas principales (resumen)} \\ \hline

HU-033 &
Análisis de asincronía en el framework Flower &
M &
\begin{itemize}\setlength\itemsep{0.2em}
  \item Analizar modelo de ejecución y sincronía por rondas.
  \item Evaluar impacto en tolerancia a fallos y Edge intermitente.
  \item Concluir viabilidad de Edge--Fog--Cloud ``nativo''.
\end{itemize}
\\ \hline

HU-034 &
Diseño arquitectónico definiendo tecnologías en cada capa &
M &
\begin{itemize}\setlength\itemsep{0.2em}
  \item Asignar responsabilidades por capa (Edge/Fog/Cloud).
  \item Definir desacoplos necesarios respecto al framework.
  \item Formalizar arquitectura híbrida implementable.
\end{itemize}
\\ \hline

HU-035 &
Estudio y selección de broker MQTT &
S &
\begin{itemize}\setlength\itemsep{0.2em}
  \item Comparar alternativas (madurez, ligereza, soporte, despliegue).
  \item Seleccionar broker adecuado para Edge--Fog.
\end{itemize}
\\ \hline

HU-036 &
Implementación y puesta a punto de broker MQTT contenerizado &
S &
\begin{itemize}\setlength\itemsep{0.2em}
  \item Desplegar broker con Docker.
  \item Validar operación y preparar integración Edge--Fog.
\end{itemize}
\\ \hline
\end{tabular}
\end{table}


\subsection{Desarrollo del Sprint 3}

\paragraph{HU-033: Análisis de asincronía en el framework Flower}
Se iteró sobre el prototipo experimental basado en ECG5000 con el objetivo de evaluar si \textit{Flower} podía sostener, por sí mismo, la arquitectura Edge--Fog--Cloud. El análisis permitió comprobar que el framework impone un flujo de entrenamiento por rondas con una dinámica predominantemente síncrona, lo que limita escenarios realistas con nodos Edge intermitentes y dificulta introducir de forma directa una capa Fog intermedia sin bloquear el proceso global.

\paragraph{HU-034: Diseño arquitectónico definiendo tecnologías en cada capa}
A partir de las restricciones observadas, se descartó una implementación ``todo en Flower'' para la arquitectura completa y se exploró un \textbf{desacoplamiento tecnológico entre capas}. El objetivo fue mantener la coherencia del diseño Edge--Fog--Cloud, evitando forzar al framework a asumir responsabilidades fuera de su modelo de ejecución.

\paragraph{HU-035: Estudio y selección de broker MQTT}
Se evaluó la introducción de un sistema de mensajería asíncrona para gestionar la comunicación Edge--Fog de forma independiente al ciclo de entrenamiento global. Se compararon distintos brokers MQTT considerando criterios de madurez, simplicidad, soporte comunitario y facilidad de despliegue.

\paragraph{HU-036: Implementación y puesta a punto de broker MQTT contenerizado}
Se implementó un broker MQTT contenerizado como base técnica para la comunicación asíncrona entre Edge y Fog. El despliegue se preparó para su integración progresiva en sprints posteriores.


\subsection{Resultados del Sprint 3}

Los resultados se presentan organizados por historia de usuario, resumiendo el estado final alcanzado sin repetir el desarrollo técnico.

\paragraph{HU-033: Análisis de asincronía en Flower}
Se confirmó que el modelo de ejecución del framework restringe la implementación directa de una arquitectura jerárquica con nodos intermitentes. Este resultado acotó el papel real de \textit{Flower} dentro del sistema y justificó la necesidad de introducir un mecanismo asíncrono entre Edge y Fog.

\paragraph{HU-034: Diseño arquitectónico detallado}
Se definió una arquitectura híbrida Edge--Fog--Cloud en la que:
(i) la comunicación y agregación parcial entre nodos Edge y Fog se gestiona de forma asíncrona mediante \textit{MQTT}, y
(ii) la coordinación de rondas y la agregación global del modelo se mantiene bajo \textit{Flower} en la capa Cloud.
Este diseño desacopla la disponibilidad de los nodos Edge del entrenamiento global, trasladando la estabilización a la capa Fog como agregador regional y participante fiable frente al servidor central.

\begin{figure}[H]
\centering
\includegraphics[height=0.9\textheight,width=1.0\linewidth]{Plantilla-Latex-TFG/imagenes/arquitectura2.png}
\caption{Arquitectura tecnológica Edge--Fog--Cloud con comunicación MQTT y Flower}
\label{fig:ArquitecturaTecnologia}
\end{figure}


\paragraph{HU-035: Selección del broker MQTT}
Se seleccionó \textit{Eclipse Mosquitto} por su ligereza, madurez y facilidad de despliegue, resultando adecuado para escenarios Edge--Fog con recursos limitados (Tabla~\ref{tab:comparativa-mqtt}).

\paragraph{HU-036: Broker MQTT contenerizado}
El broker quedó desplegado y operativo mediante contenedores Docker, quedando preparado para su integración en el Sprint 4.
\begin{table}[H]
\centering
\footnotesize
\caption{Comparativa de brokers MQTT y justificación de la elección tecnológica}
\label{tab:comparativa-mqtt}

\renewcommand{\arraystretch}{1.25}
\setlength{\tabcolsep}{3pt}

\begin{tabularx}{\linewidth}{@{} p{3.2cm} p{4.3cm} >{\RaggedRight\arraybackslash}X @{}}
\toprule
\textbf{Broker MQTT} 
& \textbf{Características principales} 
& \textbf{Justificación y adecuación al proyecto} \\
\midrule

\textbf{Eclipse Mosquitto} 
& Broker MQTT ligero y open-source (Eclipse Foundation). Soporte MQTT 3.1/3.1.1/5.0. Bajo consumo de recursos y despliegue sencillo con Docker. 
& Adecuado para Edge--Fog por simplicidad y bajo consumo, facilitando depuración e integración en un entorno experimental \cite{mosquitto_docs, mosquitto_project, mqtt_comparison_hivemq}. \\

\textbf{EMQX} 
& Broker escalable orientado a despliegues de gran escala. Soporte avanzado de clustering y balanceo. 
& Madurez elevada, pero complejidad operativa excesiva para el alcance del TFM \cite{emqx_docs, mqtt_comparison_emqx}. \\

\textbf{HiveMQ (Community)} 
& Broker robusto con amplia documentación y soporte comercial. Limitaciones en la edición comunitaria. 
& Buena documentación, pero restricciones y enfoque comercial reducen su idoneidad en un entorno académico \cite{hivemq_docs}. \\

\textbf{VerneMQ} 
& Broker distribuido basado en Erlang/OTP, orientado a alta disponibilidad y clustering. 
& Arquitectura sólida, pero despliegue/configuración más complejos de lo necesario en esta fase \cite{vernemq_docs}. \\

\textbf{RabbitMQ (plugin MQTT)} 
& Broker generalista con soporte MQTT vía plugin adicional. 
& No es MQTT nativo y añade complejidad/consumo, menos adecuado para Edge/Fog \cite{rabbitmq_mqtt}. \\
\bottomrule
\end{tabularx}
\end{table}




\subsection{Sprint Review (Sprint 3)}
El Sprint 3 cumplió su objetivo: transformar una limitación técnica detectada en el Sprint 2 en una \textbf{decisión arquitectónica implementable}. El sprint fue valioso porque permitió cerrar un supuesto erróneo (que Flower cubriría nativamente todas las responsabilidades) y reemplazarlo por una separación tecnológica coherente por capas (MQTT para mensajería Edge--Fog, Flower para agregación en Cloud).

\vspace{0.15em}
\paragraph{Qué se logró.}
Se definió una arquitectura más realista para el contexto Edge--Fog--Cloud, identificando que la asincronía y la indisponibilidad propias del Edge requieren un mecanismo de mensajería desacoplado. Como resultado, se fijó la combinación MQTT + Flower como base tecnológica, separando comunicación (Edge--Fog) y coordinación FL (Cloud).

\vspace{0.15em}
\paragraph{Qué fue bien.}
Se dimensionó correctamente el esfuerzo para \textbf{desbloquear el diseño} mediante análisis técnico y cierre de una decisión tecnológica concreta. El resultado redujo incertidumbre arquitectónica y habilitó planificación incremental para implementación.

\vspace{0.15em}
\paragraph{Qué se detectó / qué fue mal.}
Se subestimó el coste de \textbf{validar restricciones operativas} del framework al intentar extender Flower más allá de escenarios estándar, lo que obligó a replantear la estrategia durante el sprint. Esto evidenció que parte del trabajo debía reformularse como prueba mínima por capas, no como extensión directa del framework.

\vspace{0.15em}
\paragraph{Acciones y siguientes pasos.}
Definir criterios de salida para validaciones arquitectónicas: prueba mínima que evidencie tolerancia a asincronía e indisponibilidad Edge--Fog. Priorizar integración incremental: (i) comunicación MQTT Edge--Fog, (ii) agregación regional en Fog, (iii) acoplamiento final con Flower en Cloud.





\section{Sprint 4}
\label{sect:Sp4}
En el Sprint 4, el proyecto avanza desde la definición conceptual hacia un \textbf{diseño orientado a la implementación}, con el objetivo de concretar cómo debe operar el sistema federado en la práctica. El foco de esta iteración no es aún la ejecución completa, sino la definición clara de los componentes, sus responsabilidades y el flujo de comunicación necesario para construir un primer prototipo funcional.

El sistema se concibe como una arquitectura de \textbf{aprendizaje federado jerárquico} basada en el paradigma Edge--Fog--Cloud, inspirada en propuestas previas de la literatura. En este esquema, los clientes Edge entrenan modelos de forma local sobre sus datos privados, los nodos Fog introducen una capa intermedia de agregación regional y el servidor central coordina la agregación global del modelo.

El flujo de una ronda federada, ilustrado en la Figura~\ref{fig:flow_diagram}, se estructura de forma secuencial: los nodos Edge generan actualizaciones locales y las envían mediante un mecanismo de comunicación ligera; los nodos Fog recogen y agregan un conjunto de contribuciones antes de reenviarlas al servidor central; finalmente, el servidor realiza la agregación global y redistribuye el modelo actualizado, cerrando el ciclo de entrenamiento.

\begin{figure}[H]
\centering
\begin{tikzpicture}[node distance=2cm, auto]
    \tikzstyle{process} = [rectangle, rounded corners, minimum width=3cm, minimum height=1cm, text centered, draw=black, fill=blue!20]
    \tikzstyle{arrow} = [thick,->,>=stealth]

    \node (start) [process] {Inicio};
    \node (step1) [process, below of=start] {Clientes Edge entrenan y publican pesos vía MQTT};
    \node (step2) [process, below of=step1] {FogBroker agrega \(K\) actualizaciones (promedio ponderado)};
    \node (step3) [process, below of=step2] {FogClientSwell envía modelo parcial vía Flower gRPC};
    \node (step4) [process, below of=step3] {Servidor agrega parciales y publica modelo global};
    \node (step5) [process, below of=step4] {Modelo global se difunde vía MQTT a clientes Edge};
    \node (end)   [process, below of=step5] {Fin de ronda};

    \draw [arrow] (start) -- (step1);
    \draw [arrow] (step1) -- (step2);
    \draw [arrow] (step2) -- (step3);
    \draw [arrow] (step3) -- (step4);
    \draw [arrow] (step4) -- (step5);
    \draw [arrow] (step5) -- (end);
\end{tikzpicture}
\caption{Diagrama de flujo del sistema federado propuesto.}
\label{fig:flow_diagram}
\end{figure}

\begin{table}[H]
\centering
\footnotesize
\caption{Sprint 4: Historias de usuario y tareas asociadas}
\label{tab:sprint4}
\renewcommand{\arraystretch}{1.25}
\setlength{\tabcolsep}{4pt}

\begin{tabular}{|l|p{4.7cm}|c|p{5.7cm}|}
\hline
\textbf{ID} & \textbf{Historia de usuario} & \textbf{Tamaño} & \textbf{Tareas principales (resumen)} \\ \hline

HU-037 &
Diseño a bajo nivel del prototipo (clases y responsabilidades) &
M &
\begin{itemize}\setlength\itemsep{0.2em}
  \item Definir clases principales y responsabilidades por capa.
  \item Relacionar clases con bloques de la arquitectura (Fig.~5.9).
  \item Preparar UML y diagrama de secuencia del flujo híbrido.
\end{itemize}
\\ \hline

HU-038 &
Análisis exploratorio del dataset WESAD &
M &
\begin{itemize}\setlength\itemsep{0.2em}
  \item Revisar estructura, señales y etiquetas.
  \item Evaluar partición por sujeto y viabilidad no-i.i.d.
  \item Generar visualizaciones preliminares.
\end{itemize}
\\ \hline

HU-039 &
Análisis exploratorio del dataset SWELL &
M &
\begin{itemize}\setlength\itemsep{0.2em}
  \item Revisar estructura, variables/sesiones y etiquetas.
  \item Evaluar partición por sujeto/sesión y viabilidad no-i.i.d.
  \item Generar visualizaciones preliminares.
\end{itemize}
\\ \hline

HU-040 &
Diseño de baseline en WESAD/SWELL &
S &
\begin{itemize}\setlength\itemsep{0.2em}
  \item Definir pipeline base (ventanas + features + modelos).
  \item Definir métricas (Accuracy, Macro-F1) y protocolos (hold-out + CV).
\end{itemize}
\\ \hline

HU-041 &
Planificación del plan de reproducibilidad experimental (baseline) &
S &
\begin{itemize}\setlength\itemsep{0.2em}
  \item Definir criterios de reproducibilidad y trazabilidad.
  \item Preparar estructura de experimentos y registro de resultados.
\end{itemize}
\\ \hline
\end{tabular}
\end{table}
\subsection{Desarrollo del Sprint 4}
\paragraph{HU-037: Diseño a bajo nivel del prototipo (clases y responsabilidades)}
Durante el Sprint 4 se llevó la arquitectura híbrida Edge--Fog--Cloud a una \textbf{versión beta simplificada} orientada a un primer prototipo implementable. El objetivo no fue completar todavía una integración de extremo a extremo, sino definir un conjunto mínimo de clases con responsabilidades claras, de forma que el sistema pueda evolucionar de manera incremental sin bloquear la experimentación.

El diseño se apoyó en dos decisiones de base: (i) desacoplar la comunicación Edge--Fog mediante un canal asíncrono (MQTT) para responder las limitaciones del framework Flower, y (ii) mantener la coordinación federada global en Cloud mediante \textit{Flower}, tratando los nodos Fog como participantes estables. A partir de ello, el trabajo se centró en definir \textbf{qué misión tiene cada clase} y \textbf{cómo se relaciona con los bloques operativos} del sistema: \emph{recepción de actualizaciones}, \emph{agregación parcial}, \emph{gestión de disponibilidad} (en Fog) y \emph{comunicación federada global} (en Cloud).

A nivel de responsabilidades, se estableció el siguiente reparto:
(i) en \textbf{Edge}, entrenamiento local y publicación de actualizaciones;
(ii) en \textbf{Fog}, recepción asíncrona, \emph{buffering}, control de disponibilidad y agregación parcial;
y (iii) en \textbf{Cloud}, coordinación de rondas y agregación global del modelo.



\begin{figure}[H]
\centering
\includegraphics[width=0.62\linewidth]{Plantilla-Latex-TFG/imagenes/Fogarq.png}
\caption{\textbf{Arquitectura edge--fog--cloud empleada como referencia.} Adaptado de Liu et al.~\cite{liu2024iiotframework}.}
\label{fig:edge-fog-server_liu2024}
\end{figure}

\paragraph{HU-038: Análisis exploratorio del dataset WESAD}
En paralelo al diseño software, se inició el análisis exploratorio del dataset \textit{WESAD} con el objetivo de caracterizar el dominio de datos disponible y anticipar implicaciones prácticas para el pipeline experimental. El trabajo se centró en revisar la \textbf{estructura del dataset}, las \textbf{señales/variables disponibles}, el \textbf{esquema de etiquetas} y el \textbf{formato de almacenamiento}, con el fin de evaluar su adecuación al escenario del proyecto y detectar posibles limitaciones tempranas (p.\ ej., rangos, valores faltantes, coherencia temporal o inconsistencias). Para ello se desarrollaron scripts en \texttt{Python} basados en \texttt{pandas} y \texttt{NumPy} para inspección y manipulación de datos, y \texttt{matplotlib} para generar visualizaciones preliminares de apoyo al análisis.

\paragraph{HU-039: Análisis exploratorio del dataset SWELL}
De forma análoga, se llevó a cabo el análisis exploratorio del dataset \textit{SWELL} aplicando \textbf{el mismo procedimiento} y con el mismo objetivo: caracterizar su estructura, variables, etiquetas y formato, y valorar su adecuación al contexto del proyecto antes de definir el baseline. Se emplearon igualmente scripts en \texttt{Python} apoyados en \texttt{pandas} y \texttt{NumPy} para la inspección y tratamiento de datos, junto con \texttt{matplotlib} para visualizaciones exploratorias preliminares orientadas a detectar patrones generales, rangos y posibles inconsistencias.


\paragraph{HU-040: Diseño de baseline en WESAD/SWELL}
Se definió el diseño del \textbf{baseline centralizado} como referencia experimental previa al aprendizaje federado. El objetivo fue establecer un punto de comparación \textbf{reproducible} que permitiera, posteriormente, cuantificar el impacto tanto de la descentralización como de la arquitectura híbrida propuesta. Además, este baseline se planteó como un criterio práctico de decisión para seleccionar con qué dataset iniciar el desarrollo: disponer de un pipeline base común y verificable en \textit{WESAD} y \textit{SWELL} permite comparar su viabilidad experimental y elegir el primer conjunto de datos sobre el que construir la implementación federada.


\paragraph{HU-041: Planificación del plan de reproducibilidad experimental (baseline)}
Finalmente, se definió un \textbf{plan de reproducibilidad experimental} asociado al baseline, estableciendo de forma explícita cómo registrar de manera sistemática las \textbf{configuraciones} del experimento, las \textbf{decisiones de preprocesado}, las \textbf{semillas} empleadas y los \textbf{resultados} obtenidos. Este esquema de registro garantiza la \textbf{trazabilidad} entre datasets, transformaciones aplicadas, modelos entrenados y métricas reportadas, de modo que los experimentos puedan replicarse y compararse de forma consistente a lo largo de los siguientes sprints.




\subsection{Resultados del Sprint 4}

%De forma complementaria, se formalizó el flujo operativo completo del sistema mediante un \textbf{diagrama de secuencia} (Figura~\ref{fig:seq-sprint4}), que describe el ciclo federado híbrido extremo a extremo: entrenamiento local en Edge, publicación de actualizaciones vía MQTT, agregación parcial en Fog y agregación global coordinada en Cloud, seguida de la redistribución del modelo global hacia los nodos Edge.

Los resultados obtenidos durante el Sprint 4 se presentan organizados por historia de usuario, de acuerdo con la planificación recogida en la Tabla~\ref{tab:sprint4}.

\paragraph{HU-037: Diseño a bajo nivel del prototipo (clases y responsabilidades)}
La HU-037 quedó completada con la definición del \textbf{diseño beta} del prototipo híbrido y la trazabilidad explícita entre componentes software y responsabilidades por capa. Como evidencias del diseño se elaboraron:
(i) el diagrama UML de clases (Figura~\ref{fig:uml-sprint4}), y
(ii) el diagrama de secuencia del ciclo híbrido (Figura~\ref{fig:seq-sprint4}).

La descomposición en clases cubre los bloques operativos requeridos en la capa Fog:
\textbf{recepción de actualizaciones} (suscripción MQTT y entrada de mensajes),
\textbf{agregación parcial} (promedio ponderado sobre un umbral $K$ de contribuciones),
y \textbf{gestión de disponibilidad} (buffering, control de contribuciones válidas y tolerancia a inactividad parcial).
De forma separada, la \textbf{comunicación federada global} se mantiene en Cloud mediante \textit{Flower}, operando sobre participantes Fog más estables que los clientes Edge individuales.

\begin{figure}[H]
    \centering
    \includegraphics[width=0.5\linewidth]{Plantilla-Latex-TFG/imagenes/diagramauml1.png}
    \caption{Diagrama UML de clases del prototipo híbrido Edge--Fog--Cloud. Se representan los componentes mínimos (\texttt{FLClientMQTT}, \texttt{BrokerFog}, \texttt{FogClient}, \texttt{MQTTFedAvg}) y su interacción mediante MQTT y Flower.}
    \label{fig:uml-sprint4}
\end{figure}

El diagrama de secuencia formaliza el flujo extremo a extremo: entrenamiento local y publicación desde Edge, recepción y agregación parcial en Fog, agregación global coordinada en Cloud, y redistribución del modelo global hacia Edge.

\begin{figure}[H]
    \centering
    \includegraphics[width=\linewidth]{Plantilla-Latex-TFG/imagenes/diagramasecuenciav1.png}
    \caption{Diagrama de secuencia del ciclo federado híbrido. El cliente Edge entrena y publica actualizaciones; Fog recibe y agrega parcialmente; Cloud coordina rondas y agrega globalmente (Flower), redistribuyendo el modelo global a Edge.}
    \label{fig:seq-sprint4}
\end{figure}

\paragraph{Misión de cada clase (trazabilidad con los bloques del sistema)}
Para completar la trazabilidad del diseño, se fijó la responsabilidad de las clases principales del prototipo:

\begin{itemize}
    \item \textbf{\texttt{FLClientMQTT} (Edge):} ejecuta \textbf{entrenamiento local} y publica actualizaciones (pesos/gradientes serializados) al broker MQTT.
    \item \textbf{\texttt{BrokerFog} (Fog):} implementa la \textbf{recepción asíncrona} y el \textbf{buffering} de contribuciones; desencadena la agregación parcial cuando se alcanza un umbral $K$ de actualizaciones válidas.
    \item \textbf{\texttt{MQTTFedAvg} / \texttt{weighted\_average} (Fog):} aplica la \textbf{agregación parcial} (promedio ponderado) sobre las contribuciones recibidas.
    \item \textbf{\texttt{FogClient} (Fog $\rightarrow$ Cloud):} consume la salida agregada del Fog y participa como \textbf{cliente estable} en la federación global.
    \item \textbf{\texttt{Server} (Cloud, Flower):} coordina rondas y realiza la \textbf{agregación global} del modelo a partir de contribuciones de Fog.
    \item \textbf{\texttt{DataModel}/\texttt{ECGModel}:} encapsula la arquitectura del modelo y la lógica de entrenamiento/evaluación, aislándola del mecanismo de comunicación y agregación.
\end{itemize}

\paragraph{HU-038: Análisis exploratorio del dataset WESAD}
La HU-038 quedó en proceso con una primera caracterización del dataset \textit{WESAD} (señales disponibles, formato y etiquetas), identificando los puntos críticos para su posterior uso experimental: definición de ventanas temporales, alineamiento de etiquetas y partición por sujeto para simular escenarios no-i.i.d.

\paragraph{HU-039: Análisis exploratorio del dataset SWELL}
La HU-039 quedó en proceso con una revisión preliminar de la organización del dataset \textit{SWELL} (estructura por sujetos/sesiones y variables disponibles), con especial foco en la viabilidad de particiones por sujeto/sesión que permitan construir distribuciones heterogéneas relevantes para aprendizaje federado.

\paragraph{HU-040: Diseño de baseline en WESAD/SWELL}
La HU-040 quedó definida a nivel de diseño, estableciendo un pipeline homogéneo para ambos datasets con el fin de obtener una referencia centralizada previa a la federación. El baseline se estructuró en: segmentación por ventanas, extracción de características estadísticas, preprocesamiento estándar y entrenamiento/evaluación con modelos clásicos, reportando métricas \textit{Accuracy} y \textit{Macro-F1}.

\paragraph{HU-041: Planificación del plan de reproducibilidad experimental (baseline)}
La HU-041 quedó definida con un plan de reproducibilidad orientado a asegurar trazabilidad entre datos, transformaciones y resultados. Se establecieron como elementos mínimos a registrar: versión de dataset, criterios de partición, semillas, configuración de ventanas, lista de características, hiperparámetros, métricas y protocolo de evaluación, de forma que los resultados reportados puedan ser replicados en iteraciones posteriores.




\paragraph{Visualización exploratoria.}
Como apoyo al análisis, se generaron visualizaciones preliminares del dataset .


\subsection{Sprint Review (Sprint 4)}


\section{Sprint 5}
\label{sect:Sp5}

%Coincide con la entrada del mes de SeptiembreInicio de desarrollo clase cliente y fog parte mqttrealización de las historias de usuario  desarrollo HU-033 desarrollo de nodos clientes: HU 034 desarrollo nodo cliente SWELL hu 35 desarrollo nodos fog  HU36 desarrollo nodos fog SWELL, HU37 desarrollo distribuidor del dataset. HU38 desarrollo modelos baselines basados en el estado del arterESULTADOS Desarrollo baseline models

El Sprint 5 marca el paso desde el \textbf{diseño a bajo nivel} (Sprint 4) hacia la \textbf{implementación ejecutable} del primer prototipo. En el Sprint 4 se definieron las \textbf{clases, responsabilidades e interacciones} del sistema (Figura~\ref{fig:uml-sprint4}) y se formalizó el \textbf{flujo operativo} extremo a extremo (Figura~\ref{fig:seq-sprint4}), pero \textbf{no se implementó todavía} el comportamiento completo asociado a dichas clases.

Por tanto, el objetivo de este sprint fue \textbf{convertir el diseño en código funcional} y empezar a validar el flujo Edge--Fog con comunicación MQTT. En particular, el foco del Sprint 5 se puso en:
(i) implementar el cliente Edge que entrena y publica actualizaciones, 
(ii) implementar el Fog que \textbf{recibe, bufferiza y agrega parcialmente} dichas actualizaciones, y 
(iii) preparar el sistema para operar con un dataset real (SWELL), lo que obliga a adaptar la capa de datos (carga, segmentación y etiquetas) sin cambiar el papel arquitectónico de los componentes.
\begin{table}[H]
\centering
\footnotesize
\caption{Sprint 5: Historias de usuario y tareas asociadas}
\label{tab:sprint5}
\renewcommand{\arraystretch}{1.25}
\setlength{\tabcolsep}{4pt}

\begin{tabular}{|l|p{4.8cm}|c|p{5.6cm}|}
\hline
\textbf{ID} & \textbf{Historia de usuario} & \textbf{Tamaño} & \textbf{Tareas principales (resumen)} \\ \hline

HU-042 &
Desarrollo de nodos cliente (Edge) &
L &
\begin{itemize}\setlength\itemsep{0.2em}
  \item Implementar \texttt{FLClientMQTT} (entrenamiento local + publicación MQTT).
  \item Serializar/deserializar pesos (\texttt{state\_dict} $\leftrightarrow$ JSON).
  \item Suscripción a \texttt{fl/global\_model} e integración con \texttt{load\_state\_dict}.
  \item Gestión básica de reconexión MQTT y logs de ronda.
\end{itemize}
\\ \hline

HU-043 &
Desarrollo del nodo cliente para SWELL &
M &
\begin{itemize}\setlength\itemsep{0.2em}
  \item Implementar \texttt{load\_swell\_dataset()} y adaptación del formato $(X,y)$.
  \item Normalización (p.\ ej.\ Z-score) y generación de mini-batches.
  \item Verificar compatibilidad dimensional con el modelo y el \texttt{DataLoader}.
\end{itemize}
\\ \hline

HU-044 &
Desarrollo de nodos Fog (recepción y agregación parcial) &
L &
\begin{itemize}\setlength\itemsep{0.2em}
  \item Implementar \texttt{BrokerFog} (suscripción a \texttt{fl/updates}).
  \item \emph{Buffering} por región con \texttt{defaultdict(list)}.
  \item Validación mínima de mensajes (esquema JSON) y control de errores.
  \item Agregación parcial ponderada (\texttt{weighted\_average}, $K=3$) y publicación en \texttt{fl/partial}.
\end{itemize}
\\ \hline

HU-045 &
Desarrollo de nodos Fog para SWELL &
M &
\begin{itemize}\setlength\itemsep{0.2em}
  \item Alinear el esquema de updates SWELL con el pipeline de \texttt{BrokerFog}.
  \item Verificar apilado/compatibilidad de tensores en agregación parcial.
  \item Ajustar metadatos mínimos (p.\ ej.\ \texttt{num\_samples}) para ponderación.
\end{itemize}
\\ \hline

HU-046 &
Particionado avanzado de \textit{datasets} &
L &
\begin{itemize}\setlength\itemsep{0.2em}
  \item Implementar particionado reproducible con semilla fija.
  \item Definir protocolo de validación cruzada (p.\ ej.\ \texttt{StratifiedKFold}).
  \item Exportar splits/configuración para trazabilidad (JSON/YAML).
\end{itemize}
\\ \hline

HU-047 &
Desarrollo de modelos baseline centralizados &
L &
\begin{itemize}\setlength\itemsep{0.2em}
  \item Entrenar baselines centralizados (LR, SVM, MLP, XGBoost, RF).
  \item Evaluar con métricas (Accuracy, Macro-F1) y protocolo homogéneo.
  \item Registrar resultados y configuración de ejecución (para comparar con federado).
\end{itemize}
\\ \hline

\end{tabular}
\end{table}
\subsection{Desarrollo del Sprint 5}

\vspace{0.15em}
\paragraph{HU-047: Desarrollo de modelos baseline centralizados.}
Para la implementación de los modelos baseline centralizados se empleó el ecosistema estándar de \textit{machine learning} en Python, siguiendo  el plan de reproducibilidad experimental definido en el Sprint 4. El preprocesado de los datos y la gestión de las particiones de entrenamiento, validación y test se realizaron mediante \textbf{NumPy} y \textbf{Pandas}. Los modelos clásicos considerados (Regresión Logística, SVM, MLP y Random Forest) se implementaron utilizando \textbf{scikit-learn} (librería de python), además del uso de sus utilidades para validación cruzada y el cálculo homogéneo de métricas de evaluación (Accuracy y Macro-F1). El modelo \textbf{XGBoost} se entrenó con la librería homónima, y se usaron las semillas e hiperparámetros expuestos en el plan de reproducibilidad. 


\vspace{0.15em}
\paragraph{HU-046: Particionado de \textit{datasets}.}
Se implementó un esquema de particionado \textbf{reproducible} para soportar comparaciones consistentes entre enfoques (baseline vs federado) y entre ejecuciones. En concreto, se empleó validación cruzada estratificada mediante \texttt{StratifiedKFold} con \texttt{random\_state=65}, registrando la configuración de particionado para poder reconstruir exactamente los \emph{splits} utilizados. Esta pieza es crítica porque evita que diferencias de rendimiento se deban a particiones distintas y no al efecto real de la descentralización o la comunicación.

\vspace{0.15em}
\paragraph{Decisión: priorizar SWELL para el primer prototipo Edge--Fog--Server.}
Con (i) un baseline cuantitativo centralizado (HU-047) y (ii) un particionado reproducible (HU-046), se decidió construir la \textbf{primera integración funcional} del prototipo Edge--Fog--Server sobre \textit{SWELL}. La elección se fundamenta en dos criterios prácticos: 
\textbf{(a)} SWELL permite modelar un escenario con \textbf{más clientes}, lo que resulta más representativo para evaluar el comportamiento del flujo federado; y 
\textbf{(b)} los baselines obtenidos en este sprint ofrecían una referencia \textbf{especialmente nítida} para medir el impacto de la descentralización y la comunicación MQTT. En particular, las métricas elevadas observadas en SWELL eran coherentes con resultados reportados en implementaciones de referencia (p.ej., Mortensen et al.~\cite{Mortensen2023SWELL}), lo que reforzó su idoneidad como punto de partida experimental.



\vspace{0.15em}
\paragraph{HU-042: Desarrollo de nodos cliente (Edge).}
Se implementó la clase \texttt{FLClientMQTT} como cliente Edge operativo, integrando: (i) bucle de entrenamiento local con PyTorch \texttt{DataLoader}, (ii) optimización con \texttt{Adam(lr=1e-3)}, y (iii) criterio de pérdida \texttt{BCEWithLogitsLoss} para clasificación binaria. Tras completar la ronda local, el cliente ejecuta \texttt{state\_dict()} y publica la actualización en el tópico \texttt{fl/updates}. Para cerrar el ciclo, el cliente queda suscrito a \texttt{fl/global\_model} y, al recibir un modelo global, lo integra mediante \texttt{load\_state\_dict()} antes de iniciar la siguiente ronda, alineándose con el flujo definido en el diagrama de secuencia (Figura~\ref{fig:seq-sprint4}).

\vspace{0.15em}
\paragraph{HU-044: Desarrollo de nodos Fog (recepción y agregación parcial).}
Se implementó el componente Fog \texttt{BrokerFog} como punto de recepción asíncrona y agregación parcial. La recepción se realiza mediante suscripción MQTT a \texttt{fl/updates}, donde el callback \texttt{on\_update()} deserializa los mensajes JSON, valida su esquema mínimo (campos y dimensiones esperadas) y los almacena en \texttt{buffers}. Para permitir concurrencia por regiones sin bloquear el flujo, los buffers se gestionan con \texttt{defaultdict(list)}.

Cuando el número de actualizaciones válidas alcanza un umbral $K$ (en esta versión, $K=3$), \texttt{weighted\_average()} ejecuta una agregación parcial ponderada. Para ello, los parámetros se alinean y apilan mediante \texttt{np.stack()} y se aplica promedio ponderado por \texttt{num\_samples}. Esta operación se corresponde con una FedAvg local:
\begin{equation}
\label{eq:fed-avg}
w_{\text{agg}} = \sum_{i=1}^{K} \frac{n_i}{\sum_{j=1}^{K} n_j}\, w_i,
\end{equation}
donde $w_i$ son los pesos del cliente $i$ y $n_i$ su número de muestras locales. El resultado parcial se publica en \texttt{fl/partial} para que el componente \texttt{FogClient} lo consuma y lo pueda reenviar al ciclo global (Fog $\rightarrow$ Cloud). A diferencia del Sprint 4, donde el Fog era un diseño, en el Sprint 5 se implementa \textbf{la lógica de operatividad}: \emph{buffering}, validación de mensajes, umbral $K$ y publicación de agregados.

\vspace{0.15em}
\paragraph{HU-043 y HU-045: Adaptación del cliente y del Fog a SWELL.}
Se desarrollaron versiones específicas para operar sobre SWELL manteniendo intacta la arquitectura y los roles de los componentes. La adaptación consistió en integrar una función \texttt{load\_swell\_dataset()} dentro de la inicialización del cliente, encargada de devolver tensores $(X,y)$ con el formato esperado por el modelo y el \texttt{DataLoader}. Sobre esa base, el cliente aplica normalización (Z-score) y construye mini-batches (tamaño 32) compatibles con el entrenamiento.

En el lado Fog, la adaptación se limitó a asegurar que el esquema de mensajes y las dimensiones de los pesos serializados se mantuviesen consistentes para que la agregación parcial sea directa. Es decir, \textbf{no se cambió el rol del Fog}, sino que se alineó el formato de entrada (updates) generado por la capa de datos de SWELL para que \texttt{BrokerFog} pueda apilar y promediar sin problemas.

\paragraph{Por qué es necesario ``adaptar un componente a un dataset''.}
La adaptación a un dataset no implica redefinir la arquitectura, sino \textbf{inyectar una capa de datos invariante}. Un cliente Edge no puede entrenar ni producir actualizaciones si no se define explícitamente cómo se construyen las muestras $(X,y)$ y cómo \texttt{DataLoader} accede a ellas. Por ello, adaptar a SWELL significa implementar \texttt{load\_swell\_dataset()} para devolver datos en el formato y dimensionalidad esperados por el modelo, manteniendo sin cambios la lógica de \texttt{train\_round()} y el protocolo MQTT.


 %Desarrollo de nodos Fog


\subsubsection{Resultados del Sprint 5}

\paragraph{HU-046: Particionado avanzado de \textit{datasets}.}
La historia quedó \textbf{completada} con un esquema de particionado reproducible (semillas fijas y \texttt{StratifiedKFold}). Este resultado habilita comparaciones consistentes en sprints posteriores (mismo split para baseline y federado) y reduce el riesgo de conclusiones no replicables por variación en particiones.

\paragraph{HU-047: Modelos baseline centralizados.}
La historia quedó \textbf{completada} con la obtención de referencias cuantitativas centralizadas para los datasets \textit{WESAD} y \textit{SWELL}, empleando un pipeline homogéneo de preprocesado, entrenamiento y evaluación. Los resultados obtenidos, resumidos en la Tabla~\ref{tab:baseline-unificado}, se sitúan en términos generales en el rango alto de la literatura para modelos clásicos, pese a tratarse de baselines relativamente sencillos y sin ajuste fino exhaustivo.

En el caso de \textit{WESAD}, los modelos basados en \textit{Random Forest} y \textit{XGBoost} alcanzan valores de \emph{accuracy} y \emph{F1 macro} coherentes con los reportados en trabajos previos con enfoques clásicos. Este comportamiento es consistente con el estado del arte y refuerza la validez del pipeline experimental empleado bajo el esquema de particionado utilizado en este sprint.

Sin embargo, el caso del dataset \textit{SWELL} mostró un comportamiento \textbf{anormalmente favorable}, con métricas extremadamente elevadas incluso para modelos clásicos sin mecanismos explícitos de regularización avanzada. Como se observa en la Tabla~\ref{tab:baseline-unificado}, se alcanzan valores superiores al 90\% de \emph{accuracy} con \textit{Random Forest}, lo que resulta llamativo en comparación con la complejidad real del problema de detección de estrés.

Al contrastar estos resultados con estudios previos ampliamente citados (en particular el trabajo de Mortensen et al.~\cite{Mortensen2023SWELL}, cuyos resultados se ilustran en la Figura~\ref{fig:mortensen-comparativa}) se comprobó que las métricas obtenidas eran sorprendentemente cercanas. En una primera interpretación, esta coincidencia se consideró un indicio de reproducibilidad del baseline y de correcta implementación del pipeline.

No obstante, como se analiza en profundidad en el Sprint 7, esta aparente reproducibilidad terminó revelándose como un \textbf{indicador de un problema metodológico estructural}. En concreto, tanto los resultados obtenidos en este trabajo como los reportados por Mortensen et al.\ comparten la ausencia de un particionado estricto por sujeto entre los conjuntos de entrenamiento y evaluación. Esta configuración introduce \emph{data leakage} en datasets fisiológicos, donde múltiples muestras altamente correlacionadas proceden del mismo individuo, y conduce a métricas artificialmente infladas que no reflejan la capacidad real de generalización del modelo.

\begin{figure}[H]
    \centering
    \includegraphics[width=\linewidth]{Plantilla-Latex-TFG/imagenes/mortensencomparativa.png}
    \caption{Comparativa de resultados reportados por Mortensen et al.\ sobre el dataset SWELL WESAD comparación de los modelos del estado del arte.}
    \label{fig:mortensen-comparativa}
\end{figure}


\begin{table}[H]
\centering
\footnotesize
\caption{Resultados de modelos baseline centralizados sobre WESAD y SWELL (baseline preliminar)}
\label{tab:baseline-unificado}

\renewcommand{\arraystretch}{1.25}
\setlength{\tabcolsep}{6pt}

\begin{tabular}{@{} l l c c @{}}  
\toprule
\textbf{Dataset} 
& \textbf{Modelo} 
& \textbf{Accuracy (\%)} 
& \textbf{F1 Macro} \\
\midrule

\multirow{5}{*}{\textit{WESAD}} 
& Logistic Regression & 83.0 & 0.821 \\
& SVM                 & 84.6 & 0.839 \\
& MLP                 & 85.1 & 0.847 \\
& XGBoost             & 85.8 & 0.854 \\
& Random Forest       & \textbf{86.2} & \textbf{0.859} \\
\midrule

\multirow{5}{*}{\textit{SWELL}} 
& Logistic Regression & 91.3 & 0.905 \\
& SVM                 & 90.1 & 0.895 \\
& MLP                 & 84.4 & 0.842 \\
& XGBoost             & 92.9 & 0.938 \\
& Random Forest       & \textbf{93.2} & \textbf{0.945} \\

\bottomrule
\end{tabular}

\vspace{0.5em}
\begin{minipage}{0.95\linewidth}
\footnotesize
\textit{Nota: Los resultados corresponden a un baseline centralizado preliminar, obtenido bajo un esquema de particionado no estricto por sujeto. Como se analiza en el Sprint 7, este esquema introduce fuga de información en datasets fisiológicos y conduce a métricas artificialmente infladas, especialmente en SWELL. Estos valores se presentan únicamente como referencia inicial para el desarrollo del sistema.}
\end{minipage}
\end{table}


\paragraph{HU-042: Desarrollo de nodos cliente (Edge).}
La historia quedó \textbf{completada} con una implementación funcional del cliente \texttt{FLClientMQTT}. El cliente ejecuta entrenamiento local por rondas, serializa parámetros y publica updates en \texttt{fl/updates}, además de recibir e integrar modelos globales desde \texttt{fl/global\_model}. Esto proporciona un nodo Edge operativo y listo para integrarse en el flujo Edge--Fog.

\paragraph{HU-044: Desarrollo de nodos Fog (recepción y agregación parcial).}
La historia quedó \textbf{completada} con un \texttt{BrokerFog} operativo que recibe updates de múltiples clientes, los bufferiza por región y ejecuta agregación parcial al alcanzar $K$ contribuciones válidas, publicando el agregado en \texttt{fl/partial}. Con ello, la capa Fog deja de ser únicamente un diseño y pasa a ser un componente ejecutable dentro del prototipo.

\paragraph{HU-043 y HU-045: Adaptación a SWELL (cliente y Fog).}
Ambas historias quedaron \textbf{completadas}. Se logró operar con un dataset real (SWELL) manteniendo la arquitectura y responsabilidades intactas. La adaptación se materializa como una capa de carga y preprocesado que hace compatibles los datos con el \texttt{DataLoader} y el formato de updates, permitiendo que el flujo MQTT (Edge--Fog) funcione sin cambios estructurales en los componentes.





\subsection{Sprint Review (Sprint 5)}
El Sprint 5 consolidó el paso de \textbf{diseño a prototipo ejecutable}, dejando preparada una base experimental (baselines + particionado) y validando el flujo Edge--Fog en un escenario realista con \textit{SWELL}. Aun así, el sprint también destapó un riesgo metodológico que condicionaría la interpretación de resultados posteriores.

\vspace{0.15em}
\paragraph{Qué se logró.}
Se establecieron \textbf{baselines centralizados} y un esquema de particionado reproducible, habilitando comparaciones controladas. Además, se implementaron componentes Edge y Fog en una arquitectura Edge--Fog funcional, verificando intercambio de actualizaciones y agregación parcial. La adaptación a \textit{SWELL} se realizó sin cambios estructurales, reforzando que la arquitectura es \textbf{agnóstica al dataset}.

\vspace{0.15em}
\paragraph{Qué fue bien.}
El enfoque de construir primero un \textbf{marco experimental} (baselines + particionado) fue acertado porque reduce ambigüedad en sprints posteriores. También fue positiva la validación de la capa Fog como punto de agregación regional, anticipando la integración jerárquica completa.

\vspace{0.15em}
\paragraph{Qué se detectó / qué fue mal.}
Durante la evaluación aparecieron métricas anormalmente elevadas en \textit{SWELL}, señalando un posible problema metodológico en el particionado (riesgo que no se cerró dentro del sprint). Además, la adaptación a \textit{SWELL} requirió más esfuerzo del estimado, afectando a la planificación temporal.

\vspace{0.15em}
\paragraph{Acciones y siguientes pasos.}
En el siguiente sprint se prioriza: (i) integración completa Fog--Cloud para cerrar el flujo end-to-end y (ii) revisión/corrección del particionado para garantizar validez experimental antes de extraer conclusiones comparativas.






\section{Sprint 6}
\label{sect:Sp6}
 El objetivo del Sprint 6 fue cerrar la \textbf{integración de la capa Cloud} es decir que esta arquitectura permita que el servidor central \textit{Flower} interactúe con los nodos Fog sin conocer los detalles de la agregación regional MQTT subyacente, manteniendo la abstracción de cliente federado estándar. Para ello, el desarrollo se estructuró en torno a la implementación del servidor central federado con comunicación dual (gRPC + MQTT), la creación del componente puente que permite a los nodos Fog actuar como clientes \textit{Flower}, la consolidación del broker de agregación regional y, finalmente, la validación funcional del sistema completo sobre un escenario experimental controlado.

\begin{table}[H]
\centering
\footnotesize
\caption{Sprint 6: Integración del servidor central federado}
\label{tab:sprint6}
\renewcommand{\arraystretch}{1.2}
\setlength{\tabcolsep}{4pt}
\begin{tabular}{|l|p{4.8cm}|c|p{5.6cm}|}
\hline
\textbf{ID} & \textbf{Historia de usuario} & \textbf{Tamaño} & \textbf{Tareas principales (resumen)} \\ \hline

HU-048 &
Servidor Cloud Flower &
L &
\begin{itemize}\setlength\itemsep{0.15em}
  \item Implementar servidor central \textit{Flower} y configuración de rondas.
  \item Definir estrategia de agregación global e inicialización del modelo.
  \item Distribuir el modelo global a clientes federados.
\end{itemize}
\\ \hline

HU-049 &
Fog como cliente Flower &
L &
\begin{itemize}\setlength\itemsep{0.15em}
  \item Adaptar el nodo Fog como cliente federado \textit{Flower}.
  \item Enviar agregados parciales al servidor y recibir el modelo global.
  \item Propagar el modelo actualizado hacia la capa Edge.
\end{itemize}
\\ \hline

HU-050 &
Estandarización de pesos/serialización Fog--Cloud &
M &
\begin{itemize}\setlength\itemsep{0.15em}
  \item Unificar formato de pesos entre PyTorch y \textit{Flower}.
  \item Normalizar serialización/deserialización Fog--Cloud.
  \item Verificar coherencia dimensional de los parámetros.
\end{itemize}
\\ \hline

HU-051 &
Validación E2E reproducible en SWELL &
M &
\begin{itemize}\setlength\itemsep{0.15em}
  \item Configurar escenario E2E con tres nodos Fog sobre \textit{SWELL}.
  \item Definir partición controlada y semillas reproducibles.
  \item Ejecutar un experimento federado mínimo de validación.
\end{itemize}
\\ \hline

HU-052 &
Test del sistema y clases&
S &
\begin{itemize}\setlength\itemsep{0.15em}
  \item Incorporar tests para asegurar el correcto funcionamiento del sistema.
\end{itemize}
\\ \hline

\end{tabular}
\end{table}

\subsection{Desarrollo del Sprint 6}
\vspace{0.15em}
\paragraph{HU-048: Servidor central con estrategia MQTTFedAvg.}

Se implementó un servidor central de aprendizaje federado utilizando el framework \textit{Flower}, incorporando una estrategia personalizada denominada \texttt{MQTTFedAvg}. Esta estrategia extiende el enfoque clásico de agregación federada para integrar, además de la comunicación estándar del servidor, un canal de intercambio basado en MQTT.

El servidor se encarga de coordinar las rondas federadas, agregar los modelos locales y publicar el modelo global resultante tras cada iteración mediante el sistema de mensajería. Esta integración permite desacoplar la coordinación del entrenamiento del mecanismo de distribución del modelo, facilitando su difusión hacia distintos niveles de la arquitectura.

Para la validación inicial del flujo completo, el servidor se configuró con parámetros básicos que permiten iniciar rondas federadas con un número reducido de clientes y un número limitado de iteraciones, facilitando la verificación del funcionamiento global del sistema.

\vspace{0.15em}
\paragraph{HU-049: FogClient como puente MQTT--gRPC.}
Se desarrolló un componente en la capa Fog que actúa como intermediario entre los clientes Edge y el servidor central, conectando la comunicación basada en MQTT con el flujo federado gestionado por el servidor. Este componente recibe actualizaciones locales procedentes de los clientes Edge, las agrupa a nivel regional y las reenvía al servidor como una única contribución federada.

El comportamiento del nodo Fog está controlado por el \textbf{factor \(K\)}, que define el \textbf{número mínimo de clientes Edge cuyas actualizaciones deben recibirse antes de iniciar una comunicación con el servidor central}. En la práctica, el nodo Fog espera a que al menos \(K\) clientes hayan participado en una ronda regional de agregación antes de enviar el resultado al servidor, reduciendo así la frecuencia de comunicación y actuando como un punto de consolidación intermedio.

Una vez alcanzado este umbral, el agregado parcial se envía al servidor central para su integración en el modelo global. Tras recibir el modelo actualizado, el nodo Fog lo redistribuye a los clientes Edge, manteniendo la coherencia del ciclo federado en la arquitectura Edge--Fog--Server.

\vspace{0.15em}
\paragraph{HU-050: Consolidación del BrokerFog con agregación regional.}
Se completó la implementación del componente \texttt{broker\_fog.py}, que actúa como \textbf{agregador regional} recibiendo actualizaciones de clientes Edge vía MQTT y generando agregados parciales. Las funcionalidades clave implementadas incluyen:

\begin{itemize}
    \item \textbf{Acumulación por región}: Uso de \texttt{defaultdict} para mantener buffers independientes por región, acumulando actualizaciones hasta alcanzar $K$ por zona
    \item \textbf{Agregación ponderada}: Implementación de la función \texttt{weighted\_average} que computa el promedio ponderado de los parámetros del modelo considerando el número de muestras de cada cliente Edge
    \item \textbf{Publicación de agregados parciales}: Una vez acumuladas $K$ actualizaciones, el broker ejecuta la agregación y publica el resultado en el topic MQTT \texttt{fl/partial}, donde es consumido por \texttt{FogClient}
    \item \textbf{Serialización JSON}: Conversión de arrays NumPy a listas Python para permitir la transmisión por MQTT, manteniendo compatibilidad con el resto del sistema
\end{itemize}

Este componente cierra el ciclo de doble agregación característico de la arquitectura: agregación regional en Fog (promedio de $K$ clientes Edge) y agregación global en Cloud (FedAvg sobre agregados parciales de todas las regiones).

\vspace{0.15em}
\paragraph{HU-051: Validación \emph{end-to-end} sobre SWELL.}
Con los componentes Edge (\texttt{client.py}), Fog (\texttt{broker\_fog.py} + \texttt{fog\_flower\_client.py}) y Cloud (\texttt{server.py}) integrados, se ejecutó la primera validación funcional del flujo federado completo. La configuración experimental empleada fue:

\begin{itemize}
    \item \textbf{Dataset}: SWELL (detección de estrés en entornos de oficina)
    \item \textbf{Arquitectura}: 3 regiones Fog simuladas, cada una coordinando clientes Edge locales dividiendo el conjunto train val y test en 3.
    \item \textbf{Parámetros federados}: Configuración base con las rondas definidas en el servidor
    \item \textbf{Distribución de datos}: Split controlado entre regiones Fog para verificar coordinación
    \item \textbf{Objetivo}: Validar que el flujo\\
    \textbf{Edge→BrokerFog→FogClient→Servidor→FogClient→Edge }\\se ejecuta correctamente
\end{itemize}



\vspace{0.15em}
\paragraph{HU-052: Testing.}
 Se implementaron 56 tests automatizados usando \texttt{pytest}, organizados en suites que cubren:

\begin{itemize}
    \item \textbf{Tests de integración end-to-end} (\texttt{test\_integration.py}): Validan el flujo completo simulando múltiples rondas federadas
    \item \textbf{Tests unitarios de componentes}: Servidor (\texttt{test\_server.py}), \\clientes MQTT (\texttt{test\_mqtt\_components.py}), broker Fog, estrategia de agregación
    \item \textbf{Tests de modelo y utilidades}: Validación de conversión de parámetros, operaciones sobre \texttt{state\_dict}, carga de datos
    \item \textbf{Mocking de dependencias}: Simulación de brokers MQTT, conexiones de red y modelos para testing aislado
\end{itemize}

\subsection{Resultados del Sprint 6}
El Sprint 6 culminó con la obtención del \textbf{primer prototipo funcional completo del sistema de aprendizaje federado híbrido}, validando por primera vez el flujo \emph{end-to-end} Edge--Fog--Cloud propuesto en la arquitectura del proyecto. Los resultados se presentan organizados por historia de usuario, resumiendo el estado final alcanzado y las evidencias funcionales obtenidas, sin repetir los detalles de implementación descritos en el desarrollo.

\vspace{0.15em}
\paragraph{HU-048: Servidor central con estrategia MQTTFedAvg.}
La historia quedó \textbf{completada} con la implementación de un servidor central federado operativo basado en \textit{Flower}, capaz de coordinar rondas federadas y ejecutar agregación global del modelo. El servidor incorpora una estrategia personalizada (\texttt{MQTTFedAvg}) que permite publicar el modelo global agregado mediante MQTT tras cada ronda, habilitando la comunicación descendente hacia los nodos Fog y Edge.

\vspace{0.15em}
\paragraph{HU-049: FogClient como cliente federado.}
La historia quedó \textbf{completada} con la adaptación de los nodos Fog como clientes federados estables frente al servidor central. Cada nodo Fog participa correctamente en el ciclo federado global, reenviando agregados parciales al servidor Cloud y propagando el modelo global actualizado hacia los clientes Edge tras cada ronda, cerrando así el ciclo completo de actualización.

\vspace{0.15em}
\paragraph{HU-050: Estandarización de pesos y serialización Fog--Cloud.}
La historia quedó \textbf{completada} con la unificación del formato de los parámetros del modelo y del protocolo de serialización/deserialización entre Fog y Cloud. Este resultado garantiza la compatibilidad entre las representaciones de PyTorch y \textit{Flower}, evitando errores de conversión y asegurando la coherencia dimensional de los modelos intercambiados.


\vspace{0.15em}

\paragraph{HU-051: Validación \emph{end-to-end} en SWELL.}
La historia quedó \textbf{completada} tras ejecutar satisfactoriamente el \textbf{primer flujo federado completo} sobre el dataset \textit{SWELL} en un escenario con \textbf{tres nodos Fog} actuando como clientes federados. El objetivo de esta validación no fue maximizar el rendimiento del modelo, sino \textbf{demostrar viabilidad funcional} y coherencia del ciclo completo de entrenamiento jerárquico \textit{Edge--Fog--Cloud}.

Como se detalla en el \textbf{diagrama de secuencia} (Fig.~\ref{fig:diagrama_secuenciass6}), una ronda federada se compone de las siguientes fases, que permiten separar claramente cómputo local, agregación regional y agregación global:

\begin{itemize}\setlength\itemsep{0.25em}
    \item \textbf{Fase 1: entrenamiento local (Edge).} Cada cliente ejecuta \texttt{train\_round()}, entrenando durante $N$ épocas sobre sus datos asignados. El resultado de esta fase es un conjunto de \textbf{pesos actualizados} (o actualización equivalente) que encapsula el aprendizaje local.

    \item \textbf{Fase 2: publicación de actualizaciones (Edge $\rightarrow$ MQTT).} Tras el entrenamiento, el cliente publica sus pesos en el tópico \texttt{fl/updates} mediante \texttt{publish(fl/updates, weights)}. El mensaje incluye metadatos mínimos (\texttt{client\_id}, \texttt{round}) para permitir trazabilidad, y los parámetros del modelo (\texttt{weights}) como carga útil.

    \item \textbf{Fase 3: acumulación y agregación regional (Fog).} El \texttt{BrokerFog} recibe los mensajes (\texttt{on\_message()}) y los acumula en buffers por región (\texttt{buffers[region]}). La agregación regional está gobernada por el umbral $K$ (en este experimento, $K=3$): si el número de actualizaciones recibidas es menor que $K$, el broker espera; cuando se alcanza $K$, ejecuta una agregación \textbf{ponderada} (\texttt{weighted\_average}) para producir un \textbf{agregado parcial} representativo de la región.
    
    \item \textbf{Fase 4: envío del agregado parcial (Fog $\rightarrow$ MQTT).} Una vez construido el agregado regional, el broker lo publica en \texttt{fl/partial} incluyendo (\texttt{region}, \texttt{round}, \texttt{partial\_weights}). Esta fase reduce la frecuencia de comunicación hacia Cloud y materializa la idea de \textbf{dos niveles de agregación}: primero regional, después global.

    \item \textbf{Fase 5: ronda Flower / agregación global (Fog $\rightarrow$ Cloud).} El componente \texttt{FogClient} actúa como \textbf{puente} entre MQTT y el protocolo de Flower: consume \texttt{fl/partial} y participa en la ronda federada estándar llamando a \texttt{fit(parameters, config)} vía gRPC. En el servidor, la estrategia \texttt{MQTTFedAvg} ejecuta \texttt{aggregate\_fit(results)}, combinando los agregados parciales procedentes de múltiples regiones mediante \textbf{FedAvg} para obtener el \textbf{modelo global}.

    \item \textbf{Fase 6: redistribución del modelo global (Cloud $\rightarrow$ MQTT $\rightarrow$ Edge).} Tras la agregación global, el servidor publica el modelo en \texttt{fl/global\_model} con (\texttt{round}, \texttt{global\_weights}, \texttt{metrics}). Los clientes reciben el mensaje (\texttt{on\_message}) y actualizan su estado local (\texttt{load\_state\_dict(global\_weights)}), quedando sincronizados y listos para la siguiente ronda.

    \item \textbf{Fase 7: evaluación opcional (Cloud).} De forma adicional, \texttt{FogClient} puede invocar \texttt{evaluate(parameters, config)} y el servidor ejecuta \texttt{aggregate\_evaluate()} para obtener métricas globales agregadas, separando así el ciclo de optimización del ciclo de evaluación.
\end{itemize}

Bajo esta configuración, se ejecutó un experimento de \textbf{10 rondas federadas} sobre \textit{SWELL}. El resultado observado fue un \textbf{76{,}43\%} de \emph{accuracy} en test, lo que confirma que el sistema \textbf{aprende de forma incremental} incorporando conocimiento distribuido, y, sobre todo, que la arquitectura jerárquica propuesta es \textbf{operativa, reproducible y coherente} en ejecución \emph{end-to-end} (entrenamiento local, agregación Fog, agregación global y redistribución del modelo).



\vspace{0.15em}
\paragraph{HU-052: Testing.}
Tras la implementación de los test, como resultado nos dio una arquitectura funcional y con garantías de que funciona correctamente. Ésta se ha incorporado al flujo de CI/CD de Github para validarse automáticamente cuando se realiza un \textit{merge} sobre la rama \textit{main}.
\begin{table}[H]
\centering
\caption{Número de tests automatizados implementados en la HU-052.}
\label{tab:tests-hu052}
\renewcommand{\arraystretch}{1.2}
\setlength{\tabcolsep}{6pt}

\begin{tabular}{|l|l|r|}
\hline
\textbf{Suite de tests (HU-052)} & \textbf{Archivo} & \textbf{Nº de tests} \\
\hline
Integración end-to-end & \texttt{test\_integration.py} & 10 \\
Componentes unitarios -- Servidor & \texttt{test\_server.py} & 10 \\
Componentes unitarios -- MQTT y Fog & \texttt{test\_mqtt\_components.py} & 12 \\
Modelo & \texttt{test\_model.py} & 12 \\
Utilidades & \texttt{test\_utils.py} & 8 \\
\hline
\textbf{Total HU-052} & & \textbf{52} \\
\hline
\end{tabular}
\end{table}



\vspace{0.15em}
\paragraph{Resultado global del sprint.}
En conjunto, el Sprint 6 cierra la fase de implementación estructural del sistema y valida la viabilidad técnica del flujo federado híbrido Edge--Fog--Cloud. Como resultado del \textit{sprint} el sistema queda con el siguiente diagrama de clases (Figura~\ref{fig:arquitectura-federada}).

\begin{landscape}
\begin{figure}[p]
    \centering
    \includegraphics[width=\linewidth]{Plantilla-Latex-TFG/imagenes/arquitectura_actual.png}
    \caption{Diagrama de clases de la arquitectura federada jerárquica SWELL}
    \label{fig:arquitectura-federada}
\end{figure}
\end{landscape}

\newpage
\clearpage
\begin{figure}[p]
  \centering
  \includegraphics[width=\textwidth,height=\textheight,keepaspectratio]{Plantilla-Latex-TFG/imagenes/diagramasecuenciaUMLs6.png}
  \caption{Diagrama de secuencia Ronda de Entrenamiento Federado}
  \label{fig:diagrama_secuenciass6}
\end{figure}
\clearpage

\newpage
\subsubsection{Descripción de componentes por capa}

\paragraph{Capa Cloud -- Servidor Federado.}
\begin{itemize}
    \item \textbf{SwellServer}: Punto de entrada del servidor central. Inicializa el modelo global, configura la estrategia de agregación y lanza el servidor Flower mediante \texttt{fl.server.start\_server()}. Gestiona la evaluación global del modelo tras cada ronda.
    
    \item \textbf{MQTTFedAvgSwell}: Estrategia de agregación que extiende \texttt{FedAvg} de Flower. Implementa \texttt{aggregate\_fit()} para combinar los pesos parciales recibidos de las regiones Fog y publica el modelo global resultante vía MQTT en el topic \texttt{fl/global\_model}.
\end{itemize}

\paragraph{Capa Fog -- Agregación Regional.}
\begin{itemize}
    \item \textbf{FogBroker}: Agregador regional que actúa como intermediario MQTT. Acumula \texttt{K} actualizaciones de clientes locales en \texttt{buffers}, aplica \texttt{weighted\_average()} para generar un modelo parcial regional, y lo publica en \texttt{fl/partial}. Retransmite el modelo global a los clientes Edge.
    
    \item \textbf{FogClientSwell}: Bridge híbrido que implementa tanto \texttt{BaseMQTTComponent} como \\ \texttt{fl.client.NumPyClient}. Recibe modelos parciales por MQTT y los envía al servidor central mediante Flower gRPC (\texttt{fit()}). Es el componente clave que conecta los dos protocolos de comunicación.
\end{itemize}

\paragraph{Capa Edge -- Clientes Locales.}
\begin{itemize}
    \item \textbf{BaseMQTTComponent}: Clase base abstracta que encapsula la lógica de conexión MQTT. Proporciona métodos comunes como \texttt{on\_message()}, \texttt{publish\_json()} y gestión de suscripciones a topics.
    
    \item \textbf{SwellFLClientMQTT}: Cliente de entrenamiento local que hereda de \texttt{BaseMQTTComponent}. Carga los datos particionados por sujeto, ejecuta \texttt{train\_one\_round()} con PyTorch, y publica los pesos actualizados en \texttt{fl/updates}. Espera el modelo global para sincronizar cada ronda.
\end{itemize}

\subsection{Sprint Review}

El Sprint 6 representa un \textbf{hito técnico fundamental} dentro del proyecto, ya que permite cerrar por primera vez el flujo completo de aprendizaje federado híbrido Edge--Fog--Cloud y validar su viabilidad desde un punto de vista funcional.

\vspace{0.15em}
\paragraph{Qué se logró.}
Se integró con éxito la capa Cloud basada en \textit{Flower} con los nodos Fog, que pasan a actuar como clientes federados estables. Como resultado, el sistema es capaz de ejecutar rondas federadas completas, incluyendo entrenamiento local en Edge, agregación regional en Fog y agregación global en Cloud, con redistribución del modelo actualizado hacia las capas inferiores. Este sprint cierra la fase de implementación estructural del sistema y proporciona el primer prototipo federado \emph{end-to-end} plenamente operativo.

\vspace{0.15em}
\paragraph{Qué fue bien.}
La decisión de centrar el sprint en tareas de integración y validación estructural resultó adecuada, al tratarse de un punto bloqueante para el resto del backlog. La arquitectura híbrida diseñada en sprints anteriores se validó en ejecución real, confirmándose que el uso combinado de MQTT y gRPC permite desacoplar eficazmente la asincronía del Edge del ciclo federado global. Asimismo, la incorporación temprana de mecanismos de testing y trazabilidad ha facilitado la depuración y ha aumentado la confianza en la estabilidad del sistema.

\vspace{0.15em}
\paragraph{Qué se detectó / limitaciones actuales.}
El objetivo del sprint no fue optimizar el rendimiento del modelo ni explorar configuraciones complejas, sino validar el funcionamiento técnico del flujo federado. En consecuencia, la configuración experimental empleada (número reducido de rondas, distribución controlada de datos y ausencia de ajuste fino de hiperparámetros) limita la interpretación cuantitativa de los resultados obtenidos. Estas limitaciones se asumen como parte natural de una fase de integración inicial.

\vspace{0.15em}
\paragraph{Siguientes pasos.}
Con la arquitectura completamente operativa, los siguientes sprints pueden centrarse en la evaluación sistemática del aprendizaje federado probando diferentes arquitecturas, dado que ya hemos implementado la más básica de ellas el aprendizaje horizontal. 

\vspace{0.15em}
\paragraph{Conclusión global del sprint.}
Desde el punto de vista de planificación, el Sprint 6 desbloquea el avance del proyecto hacia fases de análisis y validación experimental. Desde el punto de vista técnico, confirma que el diseño híbrido Edge--Fog--Cloud es viable y ejecutable, estableciendo una base sólida para el resto del trabajo desarrollado en el TFM.




\section{Sprint 7}
\label{sect:Sp7}
En el \textit{Sprint 7} se planteó un cambio de enfoque, pasando del estudio del \textbf{aprendizaje federado horizontal} a la exploración de escenarios más avanzados, en particular \textbf{Federated Transfer Learning (FTL)}, apoyándose en la arquitectura híbrida Edge--Fog--Cloud validada en el Sprint 6.

Hasta ese momento, el sistema se había evaluado exclusivamente bajo un esquema de aprendizaje federado horizontal, en el que todos los clientes comparten un mismo espacio de características y entrenan modelos homogéneos sobre subconjuntos distintos de datos. El objetivo del sprint era analizar cómo extender este enfoque para permitir la participación de clientes con \textbf{datasets heterogéneos}, potencialmente diferentes en distribución, características y tamaño.

Durante este análisis se identificó un \textbf{problema metodológico crítico} en el pipeline experimental, lo que obligó a replantear el enfoque y a priorizar la corrección y validez de la evaluación antes de avanzar hacia escenarios multi-dataset.


\begin{table}[H]
\centering
\footnotesize
\caption{Sprint 7: Análisis de Federated Transfer Learning y validación metodológica}
\label{tab:sprint7}
\renewcommand{\arraystretch}{1.2}
\setlength{\tabcolsep}{4pt}
\begin{tabular}{|l|p{4.8cm}|c|p{5.6cm}|}
\hline
\textbf{ID} & \textbf{Historia de usuario} & \textbf{Tamaño} & \textbf{Tareas principales (resumen)} \\ \hline

HU-053 &
Análisis de viabilidad de Federated Transfer Learning &
M &
\begin{itemize}\setlength\itemsep{0.15em}
  \item Analizar el encaje de FTL en la arquitectura Edge--Fog--Cloud validada.
  \item Identificar componentes reutilizables y adaptaciones necesarias a nivel de modelo y flujo de datos.
  \item Delimitar requisitos mínimos para escenarios multi-dataset.
\end{itemize}
\\ \hline

HU-054 &
Análisis de compatibilidad entre datasets para transferencia &
M &
\begin{itemize}\setlength\itemsep{0.15em}
  \item Ejecutar baselines centralizados sobre WESAD y SWELL con un protocolo homogéneo.
  \item Evaluar similitudes y diferencias relevantes entre datasets.
  \item Detectar posibles incoherencias metodológicas durante la comparación.
\end{itemize}
\\ \hline

HU-055 &
Validación metodológica del particionado por sujeto &
L &
\begin{itemize}\setlength\itemsep{0.15em}
  \item Revisar el esquema de particionado train/validation/test utilizado.
  \item Identificar fugas de información por solapamiento de sujetos.
  \item Definir y probar un particionado provisional estricto por sujeto/cliente.
\end{itemize}
\\ \hline

\end{tabular}
\end{table}

\subsection{Desarrollo del Sprint 7}

\vspace{0.15em}
\paragraph{HU-053: Análisis de viabilidad de Federated Transfer Learning.}
Se analizó la viabilidad del Federated Transfer Learning partiendo de la arquitectura Edge--Fog--Cloud ya validada. A diferencia del aprendizaje federado horizontal, FTL no puede asumir homogeneidad ni en el espacio de características ni en la estructura del modelo entre los distintos clientes participantes.

En este análisis se distinguieron explícitamente dos niveles. Por un lado, los elementos puramente infraestructurales (comunicación, orquestación de rondas y coordinación entre capas Edge, Fog y Cloud) que permanecen independientes del tipo de datos y del modelo. Por otro lado, los componentes dependientes del supuesto de homogeneidad (estructura del modelo, formato de los parámetros y criterios de agregación—) que dejan de ser válidos en escenarios multi-dataset.

El resultado de este análisis fue constatar que la infraestructura federada existente es, en gran medida, reutilizable, pero que la transferencia entre dominios heterogéneos requiere decisiones explícitas sobre qué partes del modelo son compartibles, cómo se define la transferencia entre datasets y bajo qué criterios debe evaluarse su efectividad.

\vspace{0.15em}
\paragraph{HU-054: Análisis de compatibilidad entre datasets para transferencia.}

Como paso previo a escenarios de \textit{Federated Transfer Learning}, se evaluaron modelos baseline centralizados sobre los datasets \textit{WESAD} y \textit{SWELL} siguiendo un protocolo homogéneo. Este análisis permitió comparar su comportamiento en dominios fisiológicos distintos y detectar diferencias relevantes en la estructura y distribución de los datos.

Durante este proceso se observaron rendimientos anómalamente elevados, especialmente en \textit{SWELL}, lo que motivó una revisión crítica del pipeline experimental. El análisis reveló que los esquemas de particionado empleados no garantizaban una separación por sujetos, sino únicamente por muestras, introduciendo solapamientos problemáticos en datos fisiológicos.

La revisión de implementaciones de referencia, como el trabajo de Mortensen et al.~\cite{Mortensen2023SWELL}, accesible en \footnote{\url{https://www.kaggle.com/code/roy2004/mortensen-model-on-swell-dataset-diff-amt-data}}, confirmó que este problema también estaba presente en estudios previos. El acceso a los datos en bruto de \textit{SWELL} permitió identificar un subconjunto de sujetos válidos y sentar las bases para redefinir un particionado estricto por sujeto, condición necesaria para una evaluación metodológicamente correcta.


\vspace{0.15em}
\paragraph{HU-055: Validación metodológica del particionado por sujeto.}
A partir de las conclusiones obtenidas en el análisis de compatibilidad, se identificó formalmente un \textbf{problema metodológico crítico}: el esquema de particionado train/validation/test utilizado hasta ese momento no garantizaba disjunción por \textbf{sujeto}, introduciendo una fuga de información (\emph{data leakage}) en la evaluación.

Este tipo de fuga permite que el modelo aprenda patrones específicos de un individuo durante el entrenamiento y los reutilice indirectamente durante la fase de test, produciendo métricas artificialmente infladas que no reflejan la capacidad real de generalización.

Para contrastar este efecto, se definió una función provisional de particionado estricto por sujeto/cliente, forzando la disjunción completa entre los conjuntos de entrenamiento, validación y test. Al repetir los experimentos bajo este nuevo esquema, los valores de los resultados disminuyeron de forma significativa y mostraron un comportamiento mucho más coherente con el rendimiento esperado en escenarios reales, confirmando empíricamente el impacto del problema detectado.


\subsection{Resultados del Sprint 7}

\paragraph{HU-053: Análisis de viabilidad de Federated Transfer Learning.}
El análisis concluyó que la \textbf{infraestructura federada existente es mayoritariamente reutilizable}, pero que la transferencia entre datasets heterogéneos requiere cambios explícitos en el modelo y en la lógica de agregación. En particular, se identificaron tres requisitos clave para soportar escenarios de \textit{Federated Transfer Learning (FTL)}:

\begin{enumerate}
    \item \textbf{Separación explícita entre encoder y head}:  
    El modelo debe modularizarse para distinguir entre componentes compartibles (encoder) y componentes locales (head), permitiendo transferir únicamente representaciones genéricas y mantener la especialización privada en cada cliente.
    
    \item \textbf{Agregación selectiva de parámetros}:  
    Los mecanismos de agregación deben operar exclusivamente sobre los parámetros transferibles, evitando mezclar pesos específicos de cada cliente y reduciendo el volumen de comunicación.
    
    \item \textbf{Validación de compatibilidad entre clientes}:  
    Es necesario incorporar metadatos descriptivos del modelo en las comunicaciones para verificar la compatibilidad arquitectónica antes de realizar la agregación.
\end{enumerate}

La Tabla~\ref{tab:ftl-reutilizacion} resume qué componentes del sistema pueden reutilizarse directamente y cuáles requieren adaptación para soportar FTL.


\begin{table}[H]
\centering
\small
\begin{tabular}{|l|c|l|}
\hline
\textbf{Componente} & \textbf{Reutilizable} & \textbf{Observaciones} \\
\hline
\multicolumn{3}{|c|}{\textit{Infraestructura de comunicación}} \\
\hline
BaseMQTTComponent & \checkmark & Agnóstico al contenido de mensajes \\
FogBroker (topics, buffering) & \checkmark & Orquestación independiente del modelo \\
SwellServer (gRPC, rondas) & \checkmark & Coordinación reutilizable \\
\hline
\multicolumn{3}{|c|}{\textit{Lógica de agregación}} \\
\hline
FogBroker.weighted\_average() & $\sim$ & Requiere agregación selectiva por capas \\
MQTTFedAvgSwell.aggregate\_fit() & $\sim$ & Necesita estrategia FTL específica \\
FogClientSwell.fit() & $\sim$ & Filtrar parámetros transferibles \\
\hline
\multicolumn{3}{|c|}{\textit{Modelo y serialización}} \\
\hline
SwellMLP & $\times$ & Arquitectura fija, incompatible con FTL \\
publish\_weights() & $\times$ & No separa encoder de classification head \\
get/set\_parameters() & $\times$ & Asume correspondencia 1:1 entre clientes \\
\hline
\end{tabular}
\caption{Análisis de reutilización de componentes para Federated Transfer Learning}
\label{tab:ftl-reutilizacion}
\end{table}

\paragraph{HU-054: Análisis de compatibilidad entre datasets para transferencia.}
La tarea como resultado fue detectar dicho error crítico pese a no haber sido resuelto satisfactoriamente, lo que haría que quede bloqueada por el momento.

\paragraph{HU-055: Validación metodológica del particionado por sujeto.}
A partir de las conclusiones obtenidas en el análisis de compatibilidad, se identificó formalmente un \textbf{problema metodológico crítico}: el esquema de particionado train/validation/test utilizado hasta ese momento no garantizaba disjunción por \textbf{sujeto}, introduciendo una fuga de información (\emph{data leakage}) en la evaluación.

Este tipo de fuga permite que el modelo aprenda patrones específicos de un individuo durante el entrenamiento y los reutilice indirectamente durante la fase de test, produciendo métricas artificialmente infladas que no reflejan la capacidad real de generalización.

Para contrastar este efecto, se definió una función provisional de particionado estricto por sujeto/cliente, forzando la disjunción completa entre los conjuntos de entrenamiento, validación y test. La Tabla~\ref{tab:leakage-comparison} muestra la comparativa de resultados.

\begin{table}[H]
\centering
\begin{tabular}{|l|c|c|c|c|}
\hline
\multirow{2}{*}{\textbf{Modelo}} & \multicolumn{2}{c|}{\textbf{Sin disjunción (leakage)}} & \multicolumn{2}{c|}{\textbf{Con disjunción por sujeto}} \\
\cline{2-5}
 & Accuracy & Macro-F1 & Accuracy & Macro-F1 \\
\hline
Logistic Regression & 91.3 \% & 0.905 & 67.9\% & 0.431 \\
Random Forest & 93.2\% & 0.945 & 67.1\% & 0.571 \\
\hline
\end{tabular}
\caption{Comparativa de métricas en SWELL con y sin particionado estricto por sujeto}
\label{tab:leakage-comparison}
\end{table}

Al repetir los experimentos bajo el nuevo esquema de validación cruzada agrupada por sujeto (5-fold Group CV), los valores de los resultados disminuyeron de forma significativa:

\begin{itemize}
    \item \textbf{Logistic Regression}: de 99.58\% a \textbf{66.9\% $\pm$ 1.4\%} (caída de 32.7 puntos)
    \item \textbf{Random Forest}: de 99.62\% a \textbf{67.0\% $\pm$ 1.4\%} (caída de 32.6 puntos)
\end{itemize}


\subsection{Sprint Review 7}
El Sprint 7 supuso un \textbf{punto de inflexión metodológico} en el proyecto: se inició la exploración de escenarios más avanzados (Federated Transfer Learning, FTL) apoyándose en la arquitectura híbrida Edge--Fog--Cloud ya validada, pero el sprint terminó priorizando la \textbf{corrección experimental} al detectarse un problema crítico de particionado con \emph{data leakage} por sujeto.

\vspace{0.15em}
\paragraph{Qué se logró.}
Se completó el análisis de viabilidad de FTL (HU-053), concluyendo que la infraestructura de mensajería y orquestación es en gran medida reutilizable, pero que la transferencia multi-dataset exige modularización del modelo (encoder/head), agregación selectiva y validaciones de compatibilidad.
Además, el sprint permitió \textbf{identificar formalmente} un defecto de evaluación (HU-055): la partición por muestras (sin disjunción por sujeto) producía métricas infladas y conclusiones no fiables. Se diseñó y aplicó un particionado provisional estricto por sujeto/cliente, verificando empíricamente la magnitud del sesgo.
En consecuencia, se detectó que el análisis de compatibilidad entre datasets (HU-054) quedaba \textbf{bloqueado} hasta corregir el pipeline experimental.

\vspace{0.15em}
\paragraph{Qué fue bien.}
Se actuó correctamente al \textbf{detener el avance hacia FTL} cuando se evidenció un fallo de validez metodológica: se priorizó asegurar un protocolo experimental sólido antes de invertir esfuerzo en escenarios multi-dataset.
Asimismo, el sprint aportó claridad arquitectónica: se distinguieron con nitidez los componentes \textbf{agnósticos al modelo/dataset} (comunicación, rondas, brokers) frente a los dependientes de homogeneidad (serialización 1:1, agregación completa de pesos), dejando un mapa claro de qué habría que adaptar para FTL.

\vspace{0.15em}
\paragraph{Qué se detectó / limitaciones actuales.}
El pipeline previo no garantizaba disjunción por sujeto, introduciendo \emph{data leakage} especialmente crítico en señales fisiológicas. Esto invalida comparativas previas y obliga a reinterpretar resultados anteriores bajo la nueva evidencia.
Además, la línea FTL requiere decisiones de diseño no triviales (qué parámetros se comparten, cómo se evalúa transferencia, compatibilidad entre espacios de características), por lo que no era realista avanzar sin antes estabilizar el protocolo experimental.

\vspace{0.15em}
\paragraph{Siguientes pasos.}
El backlog inmediato debe centrarse en \textbf{hacer el particionado reproducible y configurable por sujeto} e integrarlo de forma transversal (configuración + materialización + tests), para después retomar (i) comparativas multi-dataset con protocolo homogéneo y (ii) la exploración de FTL sobre una base metodológicamente válida.


\section{Sprint 8}
\label{sect:Sp8}
El \textit{Sprint 8}, de duración inferior a la inicialmente planificada, se centró en abordar un riesgo experimental identificado en el Sprint 7 relacionado con el esquema de particionado de los datos en el contexto federado. Como objetivo principal, se diseñó e integró un \textbf{distribuidor de dataset a nivel de sujeto}, garantizando un particionado correcto, reproducible y configurable entre clientes.

Desde el punto de vista del backlog, este trabajo implicó una reformulación de la HU-008 y su descomposición en historias más específicas (HU-056, HU-057 y HU-058). Adicionalmente, esta redefinición permitió avanzar en los objetivos de la HU-015, incorporando una configuración global basada en ficheros \texttt{YAML}/\texttt{JSON} que centraliza tanto el particionado de datos como la definición de parámetros del sistema.

Las historias de usuario abordadas durante este sprint se resumen en la Tabla~\ref{tab:sprint8}.
\begin{table}[H]
\centering
\footnotesize
\caption{Sprint 8: Diseño e integración del distribuidor de dataset por sujeto}
\label{tab:sprint8}
\renewcommand{\arraystretch}{1.2}
\setlength{\tabcolsep}{4pt}
\begin{tabular}{|l|p{4.8cm}|c|p{5.6cm}|}
\hline
\textbf{ID} & \textbf{Historia de usuario} & \textbf{Tamaño} & \textbf{Tareas principales (resumen)} \\ \hline

HU-008 &
Particionado avanzado de \textit{datasets} &
XL &
\begin{itemize}\setlength\itemsep{0.15em}
  \item Revisar la definición original de la historia a la luz del riesgo experimental identificado.
  \item Refinar el alcance para contemplar explícitamente la partición estricta por sujeto.
  \item Descomponer la historia en unidades funcionales más manejables (HU-056, HU-057, HU-058).
\end{itemize}
\\ \hline

HU-056 &
Redefinición del particionado experimental a nivel de sujeto &
S &
\begin{itemize}\setlength\itemsep{0.15em}
  \item Analizar el esquema de particionado previo y sus limitaciones metodológicas.
  \item Definir formalmente la unidad sujeto como criterio mínimo de partición.
  \item Establecer la asignación de sujetos a clientes federados sin solapamiento.
\end{itemize}
\\ \hline

HU-057 &
Distribuidor de dataset configurable mediante \texttt{YAML}/\texttt{JSON} &
M &
\begin{itemize}\setlength\itemsep{0.15em}
  \item Diseñar el formato de configuración externa para la distribución de sujetos.
  \item Implementar el módulo de lectura y validación de configuraciones.
  \item Desarrollar el mecanismo de reparto de datos conforme a la configuración definida.
\end{itemize}
\\ \hline

HU-058 &
Adaptación de los componentes al nuevo esquema de distribución &
M &
\begin{itemize}\setlength\itemsep{0.15em}
  \item Integrar el distribuidor de dataset en los nodos Edge, Fog y Server.
  \item Ajustar los flujos de entrenamiento y evaluación al nuevo particionado.
  \item Verificar el funcionamiento extremo a extremo del sistema federado.
\end{itemize}
\\ \hline

HU-015 &
Gestión de hiperparámetros &
M &
\begin{itemize}\setlength\itemsep{0.15em}
  \item Introducir una configuración global del sistema mediante ficheros \texttt{YAML}/\texttt{JSON}.
  \item Centralizar parámetros relevantes de entrenamiento y arquitectura.
  \item Sentar las bases para el versionado y comparación de configuraciones experimentales.
\end{itemize}
\\ \hline

HU-059 &
Testing del sistema de particionado y arquitectura &
S &
\begin{itemize}\setlength\itemsep{0.15em}
  \item Desarrollar tests unitarios para validación de configuraciones.
  \item Verificar la carga y materialización de arquitecturas federadas.
  \item Asegurar el correcto funcionamiento del plan de ejecución.
\end{itemize}
\\ \hline

\end{tabular}
\end{table}


\subsection{Desarrollo del Sprint 8}

El Sprint 8 se planteó como una respuesta directa al riesgo experimental identificado en el Sprint 7, relacionado con el esquema de particionado de los datos. Se detectó que la asignación de muestras de un mismo sujeto a distintos clientes federados podía introducir fugas de información y sesgos relevantes en la evaluación.

Como punto de partida, se refinó conceptualmente la historia de usuario HU-008, inicialmente formulada como un particionado avanzado de datasets. Este refinamiento permitió concretar el problema desde un punto de vista metodológico y establecer la partición estricta a nivel de sujeto como un requisito esencial para garantizar la validez de los experimentos federados.

\vspace{0.15em}
\paragraph{HU-015/056: Configuración centralizada de hiperparámetros y particionado experimental por sujeto.}

Se consolidó un sistema unificado de configuración que permite gestionar de forma centralizada tanto los \textbf{hiperparámetros del entrenamiento} como el \textbf{esquema de particionado experimental}. Este enfoque garantiza coherencia, reproducibilidad y desacopla la definición del experimento de la lógica de ejecución.

La configuración se estructura mediante clases específicas que encapsulan los parámetros relevantes del experimento. En particular, \texttt{DatasetConfig} agrupa la información asociada al dataset (nombre, directorio, modalidades y proporciones de split), mientras que \texttt{ModelConfig} define los parámetros esenciales del modelo (tipo y dimensión de entrada). La carga y validación global se realiza a través de la función \texttt{load\_architecture\_config}, que verifica entre otras, condiciones de consistencia como la correcta normalización de los splits y la existencia de al menos un cliente por nodo Fog.

Se redefinió formalmente el esquema de particionado estableciendo el \textbf{sujeto como unidad mínima de asignación}, evitando cualquier solapamiento de muestras entre conjuntos de entrenamiento, validación y prueba. Implementando la clase de \texttt{FederatedConfig}, que centraliza los parámetros de particionado: sujetos disponibles, semilla aleatoria, proporciones \textit{train/val/test}, tipo de escalado y estrategia de partición.

El parámetro \texttt{split\_strategy} permite seleccionar entre dos modalidades:
\begin{itemize}
    \item \texttt{"global"}: cada sujeto se asigna íntegramente a un único conjunto (train, val o test).
    \item \texttt{"per\_subject"}: cada sujeto dispone de un split interno train/val/test, manteniendo la disjunción a nivel de cliente federado.
\end{itemize}

La asignación determinista de sujetos a splits se implementa mediante la función \texttt{\_split\_subjects}, controlada por una semilla configurable, lo que asegura la reproducibilidad completa del particionado entre ejecuciones.


\vspace{0.15em}
\paragraph{HU-057: Implementación del distribuidor de dataset configurable.}

Una vez definido el nuevo esquema de particionado, se implementó un \textbf{distribuidor de dataset} que permite materializar la partición de los datos de forma flexible y reproducible. Este mecanismo se apoya en ficheros de configuración externos en formato \texttt{YAML} o \texttt{JSON}, donde se especifica la asignación de sujetos a los distintos nodos Fog. La Figura~\ref{fig:hu057_dataset_distributor} resume el flujo general de este proceso.

El sistema interpreta esta configuración y genera automáticamente la estructura de datos federada correspondiente, aplicando valores por defecto cuando es necesario y garantizando la coherencia del experimento.



\begin{figure}[H]
\centering
\footnotesize

\begin{tikzpicture}[
  node distance=8mm and 10mm,
  box/.style={draw, rounded corners, align=left, inner sep=6pt, text width=6.2cm},
  arrow/.style={-{Stealth[length=2mm]}, thick}
]

% --- Entrada
\node[box] (cfg) {
  \textbf{Entrada: configuración declarativa (YAML/JSON)}\\
  \vspace{0.15em}
  Define \textbf{dataset}, \textbf{split} y \textbf{federation}
  (modo manual o automático).
};

% --- Lectura de configuración
\node[box, below=of cfg] (read) {
  \textbf{Lectura de configuración}\\
  \vspace{0.15em}
  Función \texttt{\_read\_config}:\\
  carga parámetros y aplica \textbf{valores por defecto}.
};

% --- Planificación
\node[box, below=of read] (plan) {
  \textbf{Planificación de la partición}\\
  \vspace{0.15em}
  Función \texttt{plan\_and\_materialize\_swell\_federated}:\\
  determina la asignación \textbf{sujeto $\rightarrow$ Fog}.
};

% --- Estrategias
\node[box, below=of plan] (strategy) {
  \textbf{Selección de estrategia de materialización}\\
  \vspace{0.15em}
  \texttt{per\_subject} o \texttt{global}
  según la configuración.
};

% --- Materialización
\node[box, below=of strategy] (mat) {
  \textbf{Materialización del dataset federado}\\
  \vspace{0.15em}
  Generación de splits agregados\\
  (\texttt{train/val/test}) por nodo Fog.
};

% --- Salida
\node[box, below=of mat] (out) {
  \textbf{Salida: estructura federada + \texttt{manifest.json}}\\
  \vspace{0.15em}
  Directorios por Fog, subdirectorios por sujeto y
  metadatos completos del experimento.
};

% --- Flechas
\draw[arrow] (cfg) -- (read);
\draw[arrow] (read) -- (plan);
\draw[arrow] (plan) -- (strategy);
\draw[arrow] (strategy) -- (mat);
\draw[arrow] (mat) -- (out);

\end{tikzpicture}

\caption{\textbf{Flujo de HU-057: configuración declarativa $\rightarrow$ planificación $\rightarrow$ materialización del dataset federado.}
La partición del dataset se define externamente mediante YAML/JSON, se planifica de forma determinista y se materializa en una estructura federada reproducible consumida por la arquitectura Edge--Fog--Server.}
\label{fig:hu057_dataset_distributor}
\end{figure}
\vspace{0.15em}

\paragraph{HU-058: Adaptación de los componentes del sistema.}
Con el distribuidor de dataset implementado, se adaptaron los distintos componentes del sistema (Edge, Fog y Server) para trabajar con el nuevo esquema de distribución, garantizando que cada nodo federado utilice únicamente los datos que le han sido asignados.

Los clientes Edge se ajustaron para cargar automáticamente los conjuntos de datos particionados y configurar el modelo de acuerdo con la información recibida. 

Por su parte, los nodos Fog incorporaron un mecanismo de agregación regional configurable, permitiendo controlar cuántas actualizaciones locales se combinan antes de ser enviadas al servidor central. Este ajuste facilita la adaptación del sistema a distintas topologías y cargas de trabajo, y permite responder a limitaciones de conectividad en dispositivos del continuo \textit{Edge--Fog--Cloud} (véase el requisito de reproducibilidad/configuración declarativa, \textbf{RNF-004}, Sección~\ref{RequisitosNofuncionales}).



\vspace{0.15em}
\paragraph{HU-059: Testing del sistema de particionado y arquitectura.}

Se implementó un conjunto de tests unitarios e integración para validar el correcto funcionamiento del sistema de configuración y particionado. Estos tests se organizaron en el módulo \texttt{test\_federated\_architecture.py}.

Adicionalmente, el módulo \texttt{test\_datasets.py} incluye tests para las funciones de particionado por sujetos, verificando la validación de parámetros y el manejo de errores cuando se solicitan más clientes que sujetos disponibles.

\subsection{Resultados del Sprint 8}

El Sprint 8 culminó con la resolución del problema de particionado de una manera integrada sobre configuración.

\vspace{0.15em}
\paragraph{HU-056: Redefinición del particionado experimental a nivel de sujeto.}
La historia quedó \textbf{completada}. Se formalizó e implementó un esquema de particionado estrictamente disjunto por sujeto mediante la clase \texttt{FederatedConfig} y la función \texttt{\_split\_subjects}. El sistema soporta dos estrategias (\texttt{per\_subject} y \texttt{global}), permitiendo flexibilidad según los requisitos experimentales.


\vspace{0.15em}
\paragraph{HU-057: Implementación del distribuidor de dataset configurable.}
La historia quedó \textbf{completada}. Se implementó el módulo \texttt{swell\_federated.py} con la función \\ \texttt{plan\_and\_materialize\_swell\_federated} que procesa configuraciones YAML/JSON y genera artefactos materializados (manifest). Un ejemplo de configuración ilustra la estructura declarativa:

\begin{lstlisting}[language=yaml, basicstyle=\scriptsize\ttfamily]
dataset:
  data_dir: data/SWELL
  modalities: [computer, facial, posture, physiology]
split:
  train: 0.7
  val: 0.15
  test: 0.15
  strategy: per_subject
federation:
  mode: manual
  manual_assignments:
    fog_0: [1, 2, 3, 4, 5, 6]
    fog_1: [7, 8, 9, 10, 11, 12]
\end{lstlisting}

La estructura de salida generada organiza los datos por nodo Fog y sujeto:
\begin{verbatim}
federated_runs/swell/run_42/
├── manifest.json
├── fog_0/
│   ├── train.npz, val.npz, test.npz
│   └── subject_1/, subject_2/, ...
\end{verbatim}


\vspace{0.15em}
\paragraph{HU-058: Adaptación de los componentes del sistema.}
La historia quedó \textbf{completada}. Los componentes Edge (\texttt{SwellFLClientMQTT}), Fog (\texttt{BrokerFog}) y la capa de orquestación (\texttt{federated\_architecture.py}) fueron adaptados para trabajar con el nuevo sistema de configuración.

\vspace{0.15em}
\paragraph{HU-015: Gestión centralizada de hiperparámetros.}
La historia quedó \textbf{completada}. Se definieron las dataclasses \texttt{DatasetConfig},  \texttt{ModelConfig} y \texttt{FederatedArchitecture} que centralizan todos los parámetros de configuración. La función \texttt{load\_architecture\_config} proporciona validación automática de consistencia.

\vspace{0.15em}
\paragraph{HU-059: Testing del sistema de particionado }
La historia quedó \textbf{completada}. Se implementaron tests unitarios en \texttt{test\_federated\_architecture.py} \texttt{test\_federated\_architecture.py} & Unitario & 2 & Configuraciones YAML/JSON que validan la carga de configuraciones JSON/YAML, y la distribución de los correcta de los dataset.



\subsection{Sprint Review (Sprint 8)}
El Sprint 8 supuso un punto de inflexión metodológico al corregir de forma estructural el problema de particionado detectado previamente. El sprint consolidó el enfoque de \textbf{arquitectura y experimento definidos por configuración}, eliminando fugas por sujeto y reforzando reproducibilidad.

\vspace{0.15em}
\paragraph{Qué se logró.}
Se redefinió el particionado para imponer \textbf{disjunción por sujeto} y se integró un distribuidor configurable (YAML/JSON) capaz de planificar y materializar splits federados reproducibles. Además, se adaptaron Edge/Fog/Server al nuevo esquema y se incorporaron tests específicos para validar configuración y particionado.

\vspace{0.15em}
\paragraph{Qué fue bien.}
Se atacó correctamente el problema en su raíz: no “parchear métricas”, sino rediseñar el pipeline para que la evaluación sea metodológicamente válida. También fue un acierto centralizar la configuración (dataset/split/federation/hiperparámetros), reforzando el enfoque \textbf{architecture-as-configuration}.

\vspace{0.15em}
\paragraph{Qué se detectó / qué fue mal.}
Hasta este sprint se asumía la corrección del particionado sin verificación explícita por sujeto, lo que explica métricas infladas y conclusiones potencialmente erróneas en sprints anteriores. Esto evidenció la necesidad de introducir validaciones sistemáticas en la ruta crítica experimental.

\vspace{0.15em}
\paragraph{Acciones y siguientes pasos.}
A partir de aquí, los experimentos federados pueden considerarse comparables bajo condiciones controladas. El siguiente paso natural es explotar la configuración declarativa para \textbf{comparar arquitecturas} (nº Fog, distribución, $K$) y empezar a instrumentar observabilidad real.


\section{Sprint 9}
\label{sect:Sp9}
Sobre esta base metodológica consolidada, el \textit{Sprint 9} se planteó como una extensión natural del trabajo previo, con el objetivo de reforzar tres aspectos clave del sistema federado: la \textbf{observabilidad y trazabilidad en ejecución}, la \textbf{configuración declarativa de arquitecturas federadas} y la \textbf{evaluación de nuevos datasets} con vistas a escenarios experimentales futuros.

Paralelamente, la recepción del dataset \textit{SWEET} abrió una línea de trabajo enfocada en la caracterización de nuevas fuentes de datos fisiológicos, con el fin de evaluar su viabilidad para experimentos federados a mayor escala.

Las historias de usuario planificadas y trabajadas durante este sprint se recogen en la Tabla~\ref{tab:sprint9}.
\begin{table}[H]
\centering
\footnotesize
\caption{Sprint 9: Observabilidad del sistema federado y análisis preliminar de SWEET}
\label{tab:sprint9}
\renewcommand{\arraystretch}{1.2}
\setlength{\tabcolsep}{4pt}
\begin{tabular}{|l|p{4.8cm}|c|p{5.6cm}|}
\hline
\textbf{ID} & \textbf{Historia de usuario} & \textbf{Tamaño} & \textbf{Tareas principales (resumen)} \\ \hline

HU-061 &
Configuración declarativa de la arquitectura federada &
XL &
\begin{itemize}\setlength\itemsep{0.15em}
  \item Definir esquema de configuración YAML/JSON para arquitectura federada.
  \item Implementar generador de manifiesto con asignaciones de nodos y clientes.
  \item Desarrollar validaciones de consistencia de configuración.
  \item Integrar configuración con scripts de preparación de datos.
\end{itemize}
\\ \hline

HU-062 &
Estudio del estándar \textit{OpenTelemetry} &
L &
\begin{itemize}\setlength\itemsep{0.15em}
  \item Analizar los principios y componentes del estándar OpenTelemetry.
  \item Estudiar su encaje en arquitecturas Edge--Fog--Server.
  \item Identificar puntos de instrumentación relevantes en el flujo federado.
\end{itemize}
\\ \hline

HU-063 &
Exploración de herramientas de trazado distribuido y diseño de esquema de integración preliminar  &
L &
\begin{itemize}\setlength\itemsep{0.15em}
  \item Evaluar herramientas compatibles con OpenTelemetry (Jaeger, Zipkin).
  \item Analizar Jaeger como sistema de visualización de trazas.
  \item Valorar su integración futura mediante contenedores Docker.
  \item Diseño de esquema de integración preliminar.
\end{itemize}
\\ \hline

HU-064 &
Análisis preliminar del dataset SWEET &
XL &
\begin{itemize}\setlength\itemsep{0.15em}
  \item Recepción y exploración estructural del dataset SWEET.
  \item Análisis de su organización por sujetos y sesiones.
  \item Evaluación de viabilidad experimental y requisitos computacionales.
\end{itemize}
\\ \hline

\end{tabular}
\end{table}


\subsection{Desarrollo del Sprint 9}
\vspace{0.15em}
\paragraph{Motivación para la observabilidad.}

La motivación para abordar esta línea de trabajo surgió tras una charla técnica en la asignatura \textit{PGPI} impartida por Salvador Moreno (\textit{AI Technical Lead} en Loomis). En dicha sesión se presentó un sistema basado en \textbf{agentes distribuidos}, en el que la \textbf{observabilidad} y el \textbf{trazado de ejecución} resultan esenciales para comprender el comportamiento del sistema en producción.

Durante la charla se expuso el uso de \textbf{estándares de trazado distribuido} para \textbf{monitorizar flujos complejos}, \textbf{identificar cuellos de botella} y \textbf{depurar interacciones} entre componentes. Este planteamiento es directamente análogo a la arquitectura \textbf{Edge--Fog--Server} propuesta en este TFM, donde el \textbf{entrenamiento federado} se reparte entre múltiples nodos y procesos que intercambian actualizaciones de modelo de forma \textbf{asíncrona}.

A partir de esta referencia, se evaluó su aplicabilidad al prototipo desarrollado, identificando \textbf{OpenTelemetry} y su ecosistema de herramientas como una base adecuada para dotar al sistema de \textbf{capacidades de observabilidad}.

\vspace{0.15em}
\paragraph{HU-061: Configuración declarativa de la arquitectura federada.}

Se implementó un sistema de \textbf{configuración declarativa} para definir la \textbf{arquitectura federada completa} mediante ficheros \textbf{YAML/JSON}, orientado a garantizar \textbf{reproducibilidad} y \textbf{automatización} de escenarios experimentales complejos.

El flujo se apoya en tres piezas clave:
\begin{itemize}
    \item \textbf{Esquema de configuración}: describe \textbf{parámetros globales}, \textbf{dataset}, \textbf{split} y \textbf{asignaciones Fog}.
    \item \textbf{Validaciones de consistencia}: verifican \textbf{coherencia} antes de materializar datos (claves, rangos de split, sujetos y asignaciones).
    \item \textbf{Generador de manifiesto}: procesa la configuración y genera un \textbf{\texttt{manifest.json}} con asignaciones completas y metadatos.
\end{itemize}

La configuración se organiza en:
\begin{itemize}
    \item \textbf{Dataset}: directorio, modalidades, sujetos y etiquetado.
    \item \textbf{Split}: proporciones \textit{train/val/test}, \textit{seed}, escalado y estrategia (global/per-subject).
    \item \textbf{Federation}: modo (manual/auto), nº de nodos Fog y reglas de asignación de sujetos.
\end{itemize}


% ============================================================

\begin{figure}[H]
\centering
\footnotesize

\begin{tikzpicture}[
  node distance=8mm and 10mm,
  box/.style={draw, rounded corners, align=left, inner sep=6pt, text width=6.2cm},
  smallbox/.style={draw, rounded corners, align=left, inner sep=6pt, text width=6.2cm},
  arrow/.style={-{Stealth[length=2mm]}, thick}
]

% --- Entrada
\node[box] (cfg) {
  \textbf{Entrada: configuración declarativa (YAML/JSON)}\\
  \vspace{0.15em}
  \textbf{Objetivo:} \textbf{reproducibilidad} y \textbf{automatización} de escenarios federados.
};

% --- Bloques de secciones
\node[smallbox, below=of cfg] (sections) {
  \textbf{Secciones principales del fichero:}
  \begin{itemize}\setlength\itemsep{0.15em}
    \item \textbf{Dataset}: \textit{data\_dir}, modalidades, sujetos, etiquetado.
    \item \textbf{Split}: \textit{train/val/test}, \textit{seed}, \textit{scaler}, estrategia (global/per-subject).
    \item \textbf{Federation}: modo (manual/auto), nº Fog, asignaciones, porcentajes auto.
  \end{itemize}
};

% --- Validaciones
\node[box, below=of sections] (checks) {
  \textbf{Validaciones de consistencia (pre-ejecución)}\\
  \vspace{0.15em}
  Comprobaciones automáticas de \textbf{coherencia}:
  \begin{itemize}\setlength\itemsep{0.15em}
    \item claves obligatorias y tipos (YAML/JSON),
    \item rangos y sumas de splits,
    \item sujetos disponibles y no solapados,
    \item asignaciones Fog/Edge válidas.
  \end{itemize}
};

% --- Generador
\node[box, below=of checks] (gen) {
  \textbf{Generador de manifiesto}\\
  \vspace{0.15em}
  Script: \texttt{prepare\_swell\_federated.py --config <yaml>}\\
  Procesa la configuración y \textbf{materializa} el experimento.
};

% --- Salida manifiesto
\node[box, below=of gen] (manifest) {
  \textbf{Salida: \texttt{manifest.json} + directorio de particiones}\\
  \vspace{0.15em}
  Contiene:
  \begin{itemize}\setlength\itemsep{0.15em}
    \item configuración completa del experimento,
    \item splits globales (train/val/test),
    \item mapeo \textbf{sujetos $\rightarrow$ Fog},
    \item mapeo \textbf{clientes Edge por Fog},
    \item metadatos (nº sujetos, features, muestras).
  \end{itemize}
};

% --- Consumo en ejecución
\node[box, below=of manifest] (run) {
  \textbf{Consumo en el ciclo federado (Edge--Fog--Server)}\\
  \vspace{0.15em}
  El entrenamiento utiliza el manifiesto para \textbf{instanciar determinísticamente}
  la arquitectura y las particiones de datos en cada ejecución.
};

% --- Flechas
\draw[arrow] (cfg) -- (sections);
\draw[arrow] (sections) -- (checks);
\draw[arrow] (checks) -- node[right, align=left]{\textbf{OK}} (gen);
\draw[arrow] (gen) -- (manifest);
\draw[arrow] (manifest) -- (run);

% --- Rama de error (opcional)
\node[smallbox, right=of checks, text width=5.8cm] (error) {
  \textbf{Si falla una validación:}\\
  se aborta la generación y se reportan
  \textbf{errores accionables} (clave faltante, split inválido, sujeto no existente, etc.).
};
\draw[arrow] (checks) -- node[above]{\textbf{ERROR}} (error);

\end{tikzpicture}

\caption{\textbf{Flujo de HU-061: configuración declarativa $\rightarrow$ validación $\rightarrow$ manifiesto reproducible.}
La arquitectura federada se define en YAML/JSON, se valida antes de materializar datos y se transforma en un
\texttt{manifest.json} consumido por el ciclo Edge--Fog--Server.}
\label{fig:hu059_config_flow}
\end{figure}


\vspace{0.15em}
\paragraph{HU-062: Estudio del estándar OpenTelemetry.}

Se llevó a cabo un análisis del estándar \textit{OpenTelemetry} \cite{opentelemetry2023} con el objetivo de determinar cómo instrumentar el sistema federado sin alterar su lógica funcional. OpenTelemetry es un framework de observabilidad \textit{vendor-neutral} que proporciona APIs, SDKs y herramientas para la recolección de telemetría (trazas, métricas y logs) en sistemas distribuidos.

El análisis se centró en los tres pilares fundamentales del estándar:
\begin{itemize}
    \item \textbf{Trazas (\textit{traces})}: Permiten seguir el flujo de una operación a través de múltiples servicios, capturando la latencia y las dependencias entre componentes.
    \item \textbf{Métricas (\textit{metrics})}: Proporcionan mediciones numéricas agregadas sobre el comportamiento del sistema (throughput, latencias, tasas de error).
    \item \textbf{Logs}: Registros de eventos puntuales que complementan la información de trazas y métricas.
\end{itemize}

En el contexto del sistema federado, se identificaron los siguientes puntos de instrumentación como prioritarios:
\begin{itemize}
    \item \textbf{Capa Edge}: Inicio y finalización del entrenamiento local, serialización y envío de actualizaciones de modelo.
    \item \textbf{Capa Fog}: Recepción de actualizaciones, proceso de agregación regional, publicación de agregados parciales.
    \item \textbf{Capa Server}: Ciclos de federación global, recepción de agregados Fog, actualización del modelo global.
    \item \textbf{Comunicación MQTT}: Tiempos de publicación y suscripción, latencias de mensajes entre capas.
\end{itemize}

\vspace{0.15em}
\paragraph{HU-063: Exploración de herramientas de trazado distribuido.}

De forma complementaria al estudio del estándar, se evaluó la herramienta \textit{Jaeger} \cite{jaeger2023} como sistema de visualización de trazas compatible con OpenTelemetry. Jaeger es un proyecto de código abierto originado en Uber Technologies y actualmente registrado en la CNCF (\textit{Cloud Native Computing Foundation}).

El análisis incluyó los siguientes aspectos:
\begin{itemize}
    \item \textbf{Arquitectura}: Jaeger soporta múltiples modos de despliegue, desde configuraciones \textit{all-in-one} para desarrollo, hasta arquitecturas distribuidas con almacenamiento persistente (Elasticsearch, Cassandra).
    \item \textbf{Integración}: La instrumentación puede realizarse mediante el SDK de OpenTelemetry para Python (\texttt{opentelemetry-sdk}), exportando trazas al colector de Jaeger.
    \item \textbf{Despliegue}: Se verificó la disponibilidad de una imagen Docker oficial (\texttt{jaegertracing/all-in-one}) que simplifica significativamente la puesta en marcha del sistema.
\end{itemize}

Con base en esta información, se diseñó  un esquema de integración preliminar donde cada componente del sistema federado (clientes Edge, brokers Fog, servidor central) exportaría spans a un colector centralizado, permitiendo visualizar el flujo completo de un ciclo de federación. 
\vspace{0.15em}
\paragraph{HU-064: Análisis preliminar del dataset SWEET.}

Tras recibir acceso al dataset \textit{SWEET} (\textit{Stress at Work: Exploring Emotional Tracking}), se realizó una inspección de su estructura de almacenamiento y contenido. Este dataset constituye una de las colecciones más extensas disponibles para el estudio del estrés laboral mediante señales fisiológicas.

La caracterización reveló los siguientes aspectos estructurales:
\begin{itemize}
    \item \textbf{Volumen de sujetos}: El dataset contiene datos de \textbf{1\,002 sujetos} en total, distribuidos en dos carpetas correspondientes a sesiones experimentales distintas: la primera con \textbf{468 sujetos} y la segunda con \textbf{544 sujetos}.
    \item \textbf{Tamaño por sujeto}: El volumen medio de información por sujeto se sitúa en torno a los \textbf{560\,MB}, lo que implica un volumen total del orden de \textbf{500--600\,GB}.
    \item \textbf{Modalidades}: El dataset incluye múltiples señales fisiológicas (electrocardiograma, actividad electrodérmica, temperatura) junto con datos contextuales y etiquetas de estrés.
    \item \textbf{Organización temporal}: Los datos están segmentados por sesiones de trabajo, con marcas temporales que permiten su alineación.
\end{itemize}

Este tamaño supone un incremento de \textbf{dos órdenes de magnitud} respecto a los datasets empleados previamente (SWELL-KW con 25 sujetos, WESAD con 15 sujetos) y conlleva un coste elevado en términos de almacenamiento, tiempos de carga y ejecución de experimentos federados.

\subsection{Resultados del Sprint 9}

% quedaron deliberadamente en estado de \textbf{bloquedas} para su implementación en fases posteriores.
El Sprint 9 concluyó con avances significativos en la configuración declarativa de arquitecturas federadas, el análisis de observabilidad y la caracterización de nuevas fuentes de datos. No obstante, varias historias se mantuvieron deliberadamente en estado de \textbf{bloqueadas}, al depender de prerrequisitos técnicos o requerir un esfuerzo adicional que podía comprometer el cierre funcional del sistema en el Sprint 10.

En particular, se mantuvo bloqueada la implementación de un \textbf{modelo SWEET de mayor capacidad} (HU-065), condicionada por la necesidad de completar primero un análisis y preparación exhaustiva del dataset (estructura por sujetos/sesiones, viabilidad experimental y requisitos computacionales), así como por la priorización del caso SWELL como escenario principal de validación.

De igual forma, quedaron bloqueadas las historias asociadas al refuerzo de \textbf{seguridad y privacidad} (HU-002, HU-011, HU-060). Estas tareas implicaban introducir mecanismos de \textbf{registro y autenticación segura de nodos} (HU-002), el \textbf{cifrado extremo-a-extremo mediante TLS} (HU-060) y técnicas adicionales de \textbf{preservación de privacidad} (HU-011). Su integración requería modificar transversalmente la capa de comunicaciones y la gestión de credenciales/certificados, por lo que se pospusieron para una fase posterior al cierre del prototipo.

Finalmente, se dejó planificada para etapas posteriores la incorporación de \textbf{tolerancia a fallos} y \textbf{reequilibrado/balanceo de carga en Fog}, al tratarse de capacidades que exigen definir políticas de \textit{timeouts}, reintentos y recuperación consistente de estado en un entorno asíncrono, y que dependen además de mecanismos de monitorización y trazabilidad ya consolidados en la fase final del proyecto (Sprint 10).

\vspace{0.15em}
\paragraph{HU-061: Configuración declarativa de la arquitectura federada.}
La historia quedó \textbf{completada a nivel de implementación}. Se desarrolló un sistema completo de configuración declarativa que permite definir \textbf{no solo parámetros de ejecución, sino la arquitectura federada completa} mediante ficheros YAML. Este sistema incluye:
\begin{itemize}
    \item Un esquema de configuración estructurado que soporta tanto asignaciones manuales como automáticas de sujetos a nodos Fog.
    \item Un generador de manifiesto que procesa las configuraciones y crea estructuras de datos preparadas para el aprendizaje federado.
    \item Validaciones de consistencia que verifican la coherencia de las configuraciones antes de la ejecución.
\end{itemize}
Este sistema se integró completamente con los scripts de preparación de datos, permitiendo la reproducción determinista de experimentos federados con diferentes arquitecturas y asignaciones de sujetos.

La implementación efectiva de la configuración de arquitectura, más allá de simples parámetros, dota al framework de una \textbf{potencialidad y versatilidad excepcional}, permitiendo:
\begin{itemize}
    \item \textbf{Experimentación reproducible}: Instanciar escenarios idénticos con diferentes topologías de red sin modificar código.
    \item \textbf{Escalabilidad flexible}: Adaptar dinámicamente el número de nodos Fog y su distribución de sujetos según necesidades experimentales.
    \item \textbf{Comparación sistemática}: Evaluar el impacto de diferentes arquitecturas federadas bajo las mismas condiciones de datos y modelo.
    \item \textbf{Validación robusta}: Verificar hipótesis sobre la influencia de la distribución de datos en el rendimiento del aprendizaje federado.
\end{itemize}

Esta capacidad de configuración arquitectónica declarativa representa un avance significativo que posiciona el framework como una herramienta versátil para investigación en aprendizaje federado aplicado a escenarios biomédicos.

\vspace{0.15em}
\paragraph{HU-062: Estudio del estándar OpenTelemetry.}
La historia quedó \textbf{completada a nivel de análisis}. Se documentaron los principios del estándar OpenTelemetry y se identificaron los puntos de instrumentación relevantes en la arquitectura Edge--Fog--Server. El análisis estableció que:
\begin{itemize}
    \item La instrumentación mediante spans permitiría capturar el flujo completo de un ciclo de federación.
    \item Los atributos de \textit{span} podrían incluir metadatos como identificador de cliente, ronda de federación, tamaño de actualización y métricas de entrenamiento local.
    \item La propagación de contexto entre componentes requeriría adaptar los mensajes MQTT para incluir \textit{headers} de trazado.
\end{itemize}
La implementación efectiva de la instrumentación se pospuso a sprints posteriores.
\vspace{0.15em}
\paragraph{HU-063: Exploración de herramientas de trazado distribuido.}
La historia quedó \textbf{completada a nivel de evaluación}. Se validó Jaeger como herramienta adecuada para la visualización de trazas en el sistema federado:
\begin{itemize}
    \item La imagen Docker \texttt{jaegertracing/all-in-one:latest} permite un despliegue inmediato sin dependencias adicionales.
    \item La interfaz web de Jaeger (puerto 16686 por defecto) proporciona visualización de trazas, análisis de latencias y grafos de dependencias.
    \item La integración con OpenTelemetry se realiza mediante el exportador OTLP (\textit{OpenTelemetry Protocol}).
\end{itemize}
Se dejó preparada la configuración base para su activación en experimentos futuros.

\vspace{0.15em}
\paragraph{HU-064: Análisis preliminar del dataset SWEET.}
Lo único que se hizo fue caracterizar exhaustivamente la estructura del dataset, identificando tanto su potencial experimental como las restricciones asociadas a su volumen:

\begin{table}[H]
\centering
\footnotesize
\caption{Comparativa de datasets disponibles para experimentación federada}
\label{tab:dataset_comparison}
\renewcommand{\arraystretch}{1.2}
\begin{tabular}{|l|c|c|c|}
\hline
\textbf{Dataset} & \textbf{Sujetos} & \textbf{Tamaño aprox.} & \textbf{Viabilidad actual} \\ \hline
SWELL-KW & 25 & $\sim$6\,GB & Alta \\ \hline
WESAD & 15 & $\sim$17\,GB & Alta \\ \hline
SWEET & 1\,002 & $\sim$560\,GB & Baja (limitaciones HW) \\ \hline
\end{tabular}
\end{table}

Dadas las \textbf{limitaciones hardware del entorno de desarrollo} (16\,GB RAM, almacenamiento SSD limitado) y el hecho de que el proyecto se encontraba en la fase \textbf{E5: Functional Increments Completed}, se determinó que la integración completa de SWEET no era viable sin comprometer los plazos establecidos. Por estos motivos, la historia quedó \textbf{bloqueada}
No obstante, quedaría relegada al backlog por si en futuras iteraciones pudiera ser retomada.

\subsection{Sprint Review (Sprint 9)}
El Sprint 9 consolidó el salto de \textbf{“prototipo que funciona”} a \textbf{plataforma de experimentación}, reforzando la configuración declarativa de arquitecturas (HU-061) y dejando preparada la hoja de ruta de observabilidad (HU-062/HU-063). En paralelo, se acotó la viabilidad del dataset SWEET (HU-064) sin comprometer el cierre funcional del Sprint 10.

\vspace{0.15em}
\paragraph{Qué se logró.}
Se implementó la configuración declarativa para describir \textbf{la arquitectura completa} (topología, nodos Fog, clientes y asignaciones), habilitando reproducción y comparación sistemática sin cambios de código. Además, se analizó OpenTelemetry y se evaluó Jaeger como backend de trazas, identificando puntos de instrumentación relevantes. Finalmente, se caracterizó SWEET, cuantificando su coste (IO/almacenamiento) y su baja viabilidad en el entorno actual.

\vspace{0.15em}
\paragraph{Qué fue bien.}
La priorización fue correcta: se invirtió en \textbf{infraestructura experimental} (configuración arquitectónica) con impacto directo en reproducibilidad y comparabilidad. También fue un acierto frenar la integración de SWEET a tiempo, evitando una “deuda experimental” que hubiera puesto en riesgo el cierre del prototipo.

\vspace{0.15em}
\paragraph{Qué se detectó / qué fue mal.}
La observabilidad quedó en fase de análisis/diseño (sin instrumentación real), por lo que su utilidad inmediata para diagnóstico aún era limitada. Además, SWEET evidenció que la planificación de datasets debe condicionarse por un \textbf{presupuesto hardware/IO} explícito.

\vspace{0.15em}
\paragraph{Acciones y siguientes pasos.}
Convertir el diseño en \textbf{instrumentación mínima end-to-end} (spans + métricas base) para validar el flujo real. Formalizar un checklist de viabilidad por dataset antes de implementación, y mantener SWEET en backlog con estrategia alternativa (submuestreo, streaming o preparación offline).




\section{Sprint 10}
\label{sect:Sp10}

El \textit{Sprint 10} tuvo una duración aproximada de \textbf{un sprint y medio} y constituyó el \textbf{último sprint de desarrollo del proyecto}, cerrando el ciclo de implementación funcional del sistema de aprendizaje federado. El trabajo se articuló en torno a cinco objetivos interrelacionados: (i) la \textbf{implantación del sistema de monitorización de métricas} mediante OpenTelemetry, Prometheus y Grafana, (ii) la \textbf{finalización e integración end-to-end del sistema federado SWELL}, (iii) la \textbf{validación del sistema mediante ejecuciones reproducibles} sobre distintas arquitecturas federadas, (iv) el \textbf{refinamiento de la configuración de capacidades Fog} para alinearla con los requisitos de observabilidad y despliegue, y (v) el \textbf{testing del sistema completo} como garantía de estabilidad y calidad.

Durante la planificación se priorizó, de forma deliberada, el \textbf{cierre funcional y verificable} del sistema por encima de extensiones avanzadas que, aun siendo relevantes, requerían madurez adicional de infraestructura y decisiones de diseño transversales. Por este motivo, se mantuvieron fuera del alcance del sprint (o en estado \textbf{bloqueado}) las historias relativas a la \textbf{implementación de un modelo SWELL de mayor capacidad} (HU-065) y al refuerzo de \textbf{seguridad y privacidad} (HU-002, HU-011, HU-060), incluyendo la \textbf{incorporación segura de nodos} y el \textbf{endurecimiento de las medidas de seguridad del sistema} frente a accesos no autorizados. Asimismo, se dejaron planificadas para fases posteriores la incorporación de \textbf{tolerancia a fallos} y \textbf{mecanismos de balanceo de carga entre nodos Fog}, con el objetivo de aumentar la resiliencia del despliegue y optimizar el reparto dinámico de trabajo en escenarios con conectividad y carga heterogéneas.

En paralelo, el sprint recuperó y completó historias previamente iniciadas que eran críticas para consolidar la plataforma: desde el punto de vista del \textit{backlog}, se abordó la \textbf{HU-006} (Monitorización de métricas), pendiente de implementación completa en sprints anteriores. Adicionalmente, se refinó la \textbf{HU-003} (Configuración de capacidades Fog) para integrar telemetría en los nodos intermedios y se completó la \textbf{HU-004} (Distribución del modelo global) incorporando trazabilidad de extremo a extremo en el flujo de publicación y propagación del modelo.

Las historias de usuario abordadas durante el \textit{sprint} 10 se recogen en las Tablas~\ref{tab:sprint10a,tab:sprint10b}.



\begin{table}[H]
\centering
\footnotesize
\caption{Sprint 10: Observabilidad, finalización del sistema SWELL y validación (I)}
\label{tab:sprint10a}
\renewcommand{\arraystretch}{1.2}
\setlength{\tabcolsep}{4pt}
\begin{tabular}{|l|p{4.8cm}|c|p{5.6cm}|}
\hline
\textbf{ID} & \textbf{Historia de usuario} & \textbf{Tamaño} & \textbf{Tareas principales (resumen)} \\ \hline

HU-006 &
Monitorización de métricas &
M &
\begin{itemize}\setlength\itemsep{0.15em}
  \item Implementar módulo \texttt{telemetry.py} con soporte OpenTelemetry.
  \item Desarrollar módulo \texttt{prometheus\_metrics.py} para métricas FL.
  \item Desplegar stack de observabilidad (Jaeger, Prometheus, Grafana).
  \item Diseñar dashboard con métricas de rounds, accuracy, latencias y throughput.
\end{itemize}
\\ \hline

HU-003 &
Configuración de capacidades Fog &
M &
\begin{itemize}\setlength\itemsep{0.15em}
  \item Extender configuración YAML para incluir parámetros de telemetría.
  \item Instrumentar broker Fog con spans y métricas de agregación.
  \item Configurar exportación de métricas por región.
\end{itemize}
\\ \hline

HU-004 &
Distribución del modelo global &
M &
\begin{itemize}\setlength\itemsep{0.15em}
  \item Añadir trazabilidad al flujo de publicación del modelo global.
  \item Implementar propagación de contexto entre servidor y clientes.
  \item Verificar recepción y aplicación del modelo en cada round.
\end{itemize}
\\ \hline

HU-066 &
Finalización del sistema federado SWELL &
L &
\begin{itemize}\setlength\itemsep{0.15em}
  \item Completar cliente \texttt{swell.py} con telemetría integrada.
  \item Finalizar servidor \texttt{swell.py} con métricas Prometheus.
  \item Integrar todos los componentes Edge--Fog--Server.
\end{itemize}
\\ \hline

HU-062 &
Validación con diferentes arquitecturas &
M &
\begin{itemize}\setlength\itemsep{0.15em}
  \item Diseñar configuraciones con diferentes topologías (2-3 nodos Fog).
  \item Ejecutar experimentos federados end-to-end.
  \item Verificar funcionamiento de telemetría y métricas en tiempo real.
\end{itemize}
\\ \hline

\end{tabular}
\end{table}
\begin{table}[H]
\centering
\footnotesize
\caption{Sprint 10: Observabilidad, finalización del sistema SWELL y validación (II)}
\label{tab:sprint10b}
\renewcommand{\arraystretch}{1.2}
\setlength{\tabcolsep}{4pt}
\begin{tabular}{|l|p{4.8cm}|c|p{5.6cm}|}
\hline
\textbf{ID} & \textbf{Historia de usuario} & \textbf{Tamaño} & \textbf{Tareas principales (resumen)} \\ \hline

HU-067 &
Testing del sistema final &
M &
\begin{itemize}\setlength\itemsep{0.15em}
  \item Desarrollar tests de integración end-to-end.
  \item Implementar tests de componentes MQTT (cliente, broker, servidor).
  \item Verificar consistencia de parámetros y matemáticas de agregación.
\end{itemize}
\\ \hline

HU-068 &
Implementación de modelos de red neuronal avanzados para clasificación de estrés fisiológico &
L &
\begin{itemize}\setlength\itemsep{0.15em}
  \item Desarrollar arquitecturas CNN profundas inspiradas en literatura reciente.
  \item Implementar modelos con atención y mecanismos de regularización avanzada.
  \item Optimizar hiperparámetros para datasets biomédicos.
  \item Integrar modelos con pipeline federado existente.
\end{itemize}
\\ \hline

HU-069&
Implementación de métricas específicas para aprendizaje federado &
M &
\begin{itemize}\setlength\itemsep{0.15em}
  \item Definir métricas FL personalizadas (rounds, accuracy global, loss, active clients).
  \item Implementar métricas por región y cliente.
  \item Integrar con Prometheus para scraping automático.
  \item Diseñar dashboards específicos para monitoreo federado.
\end{itemize}
\\ \hline

HU-070 &
Desarrollo de launcher config-driven para despliegue automatizado de arquitecturas federadas &
L &
\begin{itemize}\setlength\itemsep{0.15em}
  \item Implementar script \texttt{run\_architecture\_from\_config.py} para carga de configuraciones YAML/JSON.
  \item Desarrollar sistema de planificación de ejecución y lanzamiento coordinado de procesos.
  \item Aplicar manifiestos de despliegue para servidores, brokers y clientes.
  \item Gestionar entornos y terminación ordenada de procesos.
\end{itemize}
\\ \hline

HU-071 &
Implementación de stack completo de observabilidad containerizado &
L &
\begin{itemize}\setlength\itemsep{0.15em}
  \item Configurar stack Docker con Jaeger, Prometheus, Grafana y OTEL Collector.
  \item Implementar métricas específicas para aprendizaje federado.
  \item Desarrollar dashboards preconfigurados para monitoreo en tiempo real.
  \item Integrar telemetría distribuida con propagación de contexto.
\end{itemize}
\\ \hline

\end{tabular}
\end{table}

\subsection{Desarrollo del Sprint 10}

El desarrollo del Sprint 10 se estructuró en líneas de trabajo paralelas que convergieron en un sistema federado completamente funcional, observable y validado experimentalmente.

\vspace{0.15em}
\paragraph{Observabilidad y métricas extremo-a-extremo (HU-006, HU-069, HU-071, HU-003, HU-004).}
Se integró un sistema completo de \textbf{observabilidad extremo-a} extremo para el ciclo Edge–Fog–Server, combinando trazado distribuido y métricas en tiempo real. Este sistema permite analizar no solo los resultados finales del aprendizaje federado, sino también \textbf{cómo y cuándo interactúan los distintos componentes} durante cada ronda.

Como se ilustra en la Figura~\ref{fig:monitoring_stack}, se añadió capacidad de monitoreo al framework principal y se definieron métricas clave para el aprendizaje federado (rondas de entrenamiento, precisión global, número de clientes activos, tiempos de respuesta y rendimiento) que pueden ser exportadas y visualizadas de forma centralizada.  
Los nodos Fog se configuraron para incorporar \textbf{monitoreo por región}, permitiendo medir la recepción, almacenamiento temporal y agregación de actualizaciones locales. Asimismo, se implementó un seguimiento completo de la \textbf{distribución del modelo global}, verificando que este se publique y aplique correctamente en cada ronda.

El sistema de observabilidad se apoya en cuatro componentes principales:

\textbf{Telemetría distribuida:} se desarrolló un módulo simplificado de instrumentación que permite integrar el seguimiento de operaciones en el sistema federado, degradando de forma segura cuando alguna herramienta no está disponible.

\textbf{Seguimiento con Jaeger:} se configuró una herramienta de trazado distribuido para rastrear operaciones entre Edge, Fog y Server, incluyendo la recepción de actualizaciones y la generación automática de métricas de rendimiento.

\textbf{Métricas con Prometheus:} se implementaron métricas específicas para aprendizaje federado, recolectadas desde el servidor central, los nodos Fog y los clientes, incluyendo contadores de rondas, precisión y tiempos de entrenamiento.

\textbf{Visualización con Grafana:} se diseñó un panel de control con vistas del estado global del aprendizaje federado, nodos Fog, distribución de datos, rendimiento del entrenamiento y métricas del broker, con actualización automática en tiempo real.

El stack completo se despliega como un conjunto de \textbf{seis servicios interconectados}, incluyendo broker de mensajes, recolector de telemetría, almacenamiento de métricas y herramientas de visualización.  
Los nodos Fog se ampliaron con configuración específica de monitoreo, incorporando seguimiento de recepción, agregación y publicación de resultados parciales, así como métricas de capacidad y actividad por región. Finalmente, el envío del modelo global se instrumentó con trazado continuo, conectando servidor y clientes mediante mensajes marcados para su seguimiento.


\begin{figure}[H]
    \centering
    \includegraphics[width=.78\linewidth]{monitoring_stack.png}
    \caption{Entorno de monitoreo y observabilidad del sistema federado Edge--Fog--Server.}
    \label{fig:monitoring_stack}
\end{figure}
\vspace{0.15em}
\paragraph{HU-069: Implementación de métricas específicas para aprendizaje federado.}
Se definió un catálogo de métricas federadas en \texttt{src/flower\_basic/prometheus\_metrics.py} y se instrumentaron los componentes críticos para registrar métricas en tiempo real. En el \textbf{servidor} se registran rondas completadas, \emph{accuracy} y pérdida globales, clientes activos, duración de ronda y tiempos de agregación, además de los tamaños globales de train/val/test y la matriz de confusión final cuando hay evaluación centralizada (\texttt{src/flower\_basic/servers/swell.py}). En el \textbf{broker Fog} se contabilizan actualizaciones recibidas, tamaño del \emph{buffer}, agregaciones y parciales publicados, clientes únicos por región y contribución por cliente, junto con estadísticas del modelo agregado (norma/media/desviación) (\texttt{src/flower\_basic/brokers/fog.py}). En cada \textbf{cliente} se exponen tamaños de dataset, duración del entrenamiento local, pérdida y precisión locales por ronda (\texttt{src/flower\_basic/clients/swell.py}). Las métricas se exponen mediante endpoint HTTP compatible con Prometheus (puertos configurables con \texttt{METRICS\_PORT\_*}) y se realiza \emph{push} a Pushgateway al cierre para preservar las series.


\vspace{0.15em}
\paragraph{HU-066: Finalización del sistema federado SWELL.}

Se completaron todos los componentes del sistema SWELL con telemetría y métricas integradas:
\begin{itemize}
    \item \texttt{src/flower\_basic/clients/swell.py} (515 líneas): Cliente federado MQTT con spans para entrenamiento local.
    \item \texttt{src/flower\_basic/servers/swell.py} (439 líneas): Servidor central con estrategia \texttt{MQTTFedAvgSwell}.
    \item \texttt{src/flower\_basic/clients/fog\_bridge\_swell.py}: Bridge que conecta la capa Fog con el servidor Flower.
    \item \texttt{src/flower\_basic/brokers/fog.py}: Broker regional con agregación parametrizable.
\end{itemize}

\vspace{0.15em}
\paragraph{HU-062: Validación con diferentes arquitecturas.}
Se validó el sistema completo mediante la ejecución de experimentos end-to-end con distintas configuraciones arquitectónicas, variando el número de nodos Fog, clientes por región y el parámetro de agregación $K$.

Estas ejecuciones permitieron verificar no solo el correcto funcionamiento del aprendizaje federado, sino también la consistencia del sistema de observabilidad, la trazabilidad de las rondas y la robustez del despliegue automatizado. Entre las configuraciones evaluadas se incluyen:

\textbf{Configuración simple} (\texttt{federated\_architecture\_simple.yaml}):
\begin{itemize}
    \item 3 nodos Fog con 1 cliente por nodo, $K = 1$, 5 rondas.
\end{itemize}

\textbf{Configuración avanzada} (\texttt{swell\_federated\_10runs.yaml}):
\begin{itemize}
    \item 3 nodos Fog con 8, 8 y 7 clientes por nodo respectivamente, $K = 3$, 5 rondas.
\end{itemize}

La Figura~\ref{fig:arch_swell_10runs} muestra un ejemplo de la arquitectura correspondiente a la configuración avanzada.

\begin{figure}[H]
    \centering
    \includegraphics[width=0.82\linewidth]{Plantilla-Latex-TFG/imagenes/arquiotecturabonita.png}
    \caption{Ejemplo de arquitectura correspondiente a la configuración \texttt{swell\_federated\_10runs.yaml}.}
    \label{fig:arch_swell_10runs}
\end{figure}

\vspace{0.15em}
\paragraph{HU-068 Modelos neuronales más avanzados para estrés fisiológico.}
Se diseñó e integró una arquitectura de red neuronal específica para la clasificación de estrés fisiológico, inspirada en trabajos recientes sobre procesamiento de señales biomédicas. La arquitectura se orienta al tratamiento de datos fisiológicos tabulares y combina transformaciones lineales, activaciones no lineales y técnicas de regularización con el objetivo de mejorar la capacidad de generalización del modelo y reducir el riesgo de sobreajuste.

Como arquitectura base se empleó un modelo inspirado en Mortensen et al.\ \cite{Mortensen2023SWELL}, consistente en un perceptrón multicapa sencillo pero efectivo, compuesto por dos capas lineales ocultas (input$\rightarrow$128 y 128$\rightarrow$64) con funciones de activación ReLU y regularización mediante \textit{Dropout} (0.2), seguido de una capa de salida para clasificación binaria (estrés vs.\ no estrés). La estructura implementada se representa en la Figura~\ref{fig:swellmlp_architecture}.

\begin{figure}[H]
    \centering
    \includegraphics[width=0.6\linewidth]{Arquitectuared.png}
    \caption{Arquitectura de la red neuronal SwellMLP empleada para la clasificación de estrés fisiológico.}
    \label{fig:swellmlp_architecture}
\end{figure}


\paragraph{HU-070 Automatización del despliegue de arquitecturas federadas.}
Se desarrolló un \textbf{launcher basado en configuración} (\texttt{run\_architecture\_from\_config.py}) que permite desplegar de forma automatizada arquitecturas federadas completas a partir de ficheros YAML/JSON. El script interpreta el manifiesto experimental, prepara particiones de datos cuando es necesario y coordina el lanzamiento ordenado de servidores, nodos Fog y clientes Edge.

El launcher incorpora soporte para la carga automática de manifiestos de clientes, la distribución de configuraciones mediante MQTT y la preparación controlada de los conjuntos de datos, además de realizar validaciones previas sobre la disponibilidad de datos y el ajuste dinámico de parámetros en función del número de clientes activos. Este componente abstrae la complejidad del despliegue manual y permite ejecutar experimentos reproducibles con distintas configuraciones sin necesidad de modificar el código, facilitando la escalabilidad y el proceso de validación del sistema federado.

La Figura~\ref{fig:diagramaflujo} muestra el diagrama de flujo que ejemplifica el proceso completo de gestión del experimento, desde la generación del conjunto de datos y del manifiesto de clientes hasta la coordinación del lanzamiento y la ejecución de los distintos componentes. Este gestor constituye una primera aproximación a un sistema de orquestación experimental, que en futuras versiones podría ser sustituido por soluciones basadas en herramientas de infraestructura como \textit{Terraform} o \textit{Kubernetes}.
\clearpage
\begin{figure}[p]
    \centering
    \includegraphics[width=1.12\linewidth]{Plantilla-Latex-TFG/imagenes/flujo.png}
    \caption{Diagrama de flujo del proceso de gestión y despliegue automatizado de arquitecturas federadas.}
    \label{fig:diagramaflujo}
\end{figure}
\clearpage
\paragraph{HU-067: Testing del sistema final.}

Se diseñó y amplió una \textbf{suite de testing integral} para validar el sistema federado tanto de extremo a extremo como por componentes. En la versión actual se han implementado \textbf{147 tests} distribuidos en \textbf{más de 20 módulos}, combinando pruebas unitarias e integración. La suite se concentra en la ruta crítica SWELL (cliente, servidor y particionado federado) y en los bloques esenciales de comunicación y orquestación.

El desarrollo se centró en los componentes principales:
\begin{itemize}
    \item \textbf{Flujo end-to-end}: validación de todo el ciclo federado, desde la publicación de updates hasta la agregación global y el avance por rondas. Se comprueba que los formatos de parámetros son consistentes y que el intercambio conserva dimensiones y tipos.
    \item \textbf{Capa MQTT/Broker Fog}: verificación de la acumulación regional, el disparo de agregación al alcanzar el umbral $K$, la gestión de buffers y la tolerancia a mensajes malformados o incompletos.
    \item \textbf{Servidor central}: pruebas de la agregación global, la conversión entre formatos (state\_dict/ndarray), el comportamiento con y sin MQTT y la ejecución con múltiples rondas.
    \item \textbf{Cliente SWELL}: validación del entrenamiento local, sincronización con el modelo global, aplicación segura de estados recibidos y manejo de timeouts en ausencia de actualizaciones.
    \item \textbf{Datasets y particionados SWELL}: comprobación de la carga, normalización y particionado por sujeto, garantizando que las divisiones respetan la lógica experimental y evitan fugas de datos.
\end{itemize}

En conjunto, la suite garantiza \textbf{corrección funcional}, \textbf{consistencia matemática} del algoritmo FedAvg y \textbf{robustez operativa} ante escenarios degradados, cubriendo la configuración objetivo del sistema federado.



\subsection{Resultados del Sprint 10}


El Sprint 10 culminó con un sistema federado basado en \textbf{SWELL} operativo y preparado para experimentación, incorporando \textbf{observabilidad extremo-a-extremo}, \textbf{automatización del despliegue} y una \textbf{validación por testing} a nivel de sistema.

\paragraph{Observabilidad y métricas extremo-a-extremo (HU-006, HU-069, HU-071, HU-003, HU-004).}
\textbf{Resultado: observabilidad extremo-a-extremo operativa y verificable.} Se consolidó un sistema de monitorización que combina \textbf{métricas Prometheus} y \textbf{trazas Jaeger} para visualizar en tiempo real el ciclo Edge--Fog--Server. Las métricas se instrumentaron de forma homogénea en servidor, broker Fog y clientes, registrando rondas completadas, \emph{accuracy} y pérdida globales, clientes activos, duración de ronda y tiempos de agregación; además de métricas por región (buffers, actualizaciones, agregaciones, clientes únicos y estadísticas del modelo agregado) y por cliente (tamaños de dataset, duración de entrenamiento local, pérdida y precisión). Todas estas series se exponen vía endpoint HTTP para scraping automático.

\begin{itemize}\setlength\itemsep{0.15em}
    \item \textbf{Dashboard en Grafana (Prometheus):} paneles globales con rondas, accuracy/loss, clientes activos y duración de ronda; vistas por región con updates, tamaño de buffer, agregaciones y volumen de parciales; y vistas por cliente con tiempos de entrenamiento y métricas locales. Esto permite identificar cuellos de botella, desbalances entre regiones y evolución del rendimiento.
    \item \textbf{Trazas distribuidas con Jaeger:} seguimiento de eventos críticos (recepción de actualizaciones, agregación Fog/global y distribución del modelo) para comprobar causalidad y latencias reales entre componentes.
    \item \textbf{Instrumentación homogénea}: cada ronda deja evidencia en métricas y trazas, facilitando diagnóstico operativo y validación experimental.
\end{itemize}

En conjunto, el sistema proporciona \textbf{visibilidad integrada y en tiempo real} del aprendizaje federado, combinando indicadores operativos (Grafana) y trazabilidad de ejecución (Jaeger).

\begin{landscape}
\begin{figure}[H]
    \centering
    \includegraphics[width=\linewidth,height=0.92\textheight,keepaspectratio]{Plantilla-Latex-TFG/imagenes/Grafana.png}
    \caption{\textbf{Dashboard de Grafana del sistema federado.} Panel de monitorización en tiempo real alimentado por Prometheus, mostrando métricas globales (rondas, clientes activos, latencias) y vistas por región (Fog).}
    \label{fig:grafana_dashboard_sprint10}
\end{figure}
\end{landscape}


\begin{figure}[H]
    \centering
    \includegraphics[width=\linewidth]{Plantilla-Latex-TFG/imagenes/jaeger.png}
    \caption{\textbf{Trazas distribuidas en Jaeger para una ronda de aprendizaje federado.} 
    Visualización de los \textit{spans} entre Edge, Fog y servidor central. 
    Se observa en tiempo real la recepción de actualizaciones en la función \texttt{function\_fog} 
    y el número de servicios desplegados (tres nodos Fog activos en la arquitectura).}
    \label{fig:jaeger_traces_sprint10}
\end{figure}

\vspace{0.15em}
\paragraph{HU-066: Sistema federado SWELL final.}
\textbf{Estado final: completado, orquestado y trazable.}
Se consolidó el sistema federado SWELL como un \textbf{pipeline end-to-end ejecutable} que integra, de forma coherente, \textbf{orquestación por configuración}, \textbf{ciclo federado jerárquico} (Edge--Fog--Cloud) y \textbf{observabilidad transversal}.
Tal y como refleja el diagrama UML final (Fig.~\ref{fig:diagramumlfinal}), el cierre de la historia no se limita a “componentes funcionando”, sino a una \textbf{arquitectura completa} en la que cada capa dispone de responsabilidades bien definidas y conectadas por contratos de comunicación explícitos cumpliendo así los principios SOLID\cite{\cite{martin2008clean}} RNF-008 :

\begin{itemize}\setlength\itemsep{0.15em}
    \item \textbf{Control / Orquestación (config-driven).}
    Se dispone de un punto de entrada (\texttt{RunArchitectureFromConfig}) que carga la arquitectura desde YAML/JSON, valida consistencia, materializa particiones y construye un \textbf{plan de ejecución reproducible} (\texttt{FederatedArchitectureUtils} + \texttt{FederatedArchitecture}). Este bloque permite instanciar topologías heterogéneas (nº de Fog, clientes por región y umbrales $K$) sin modificar el código.

    \item \textbf{Capa Edge (entrenamiento local sobre MQTT).}
    Los clientes \texttt{SwellFLClientMQTT} (sobre \texttt{BaseMQTTComponent}) ejecutan entrenamiento local por ronda y publican sus actualizaciones en \texttt{fl/updates}. Además, consumen \texttt{fl/global\_model} para sincronizar su estado con el modelo global tras cada ronda, cerrando el ciclo de aprendizaje distribuido.

    \item \textbf{Capa Fog (agregación regional y puente MQTT--Flower).}
    El \texttt{FogBroker} acumula actualizaciones en buffers y aplica agregación regional controlada por umbrales (\texttt{K} / \texttt{K\_MAP}), publicando agregados parciales en \texttt{fl/partial}. A su vez, \texttt{FogClientSwell} actúa como \textbf{bridge híbrido}: consume parciales por MQTT y participa en la ronda federada global mediante \texttt{fl.client.NumPyClient} (\texttt{fit/evaluate}) vía gRPC contra el servidor central.

    \item \textbf{Capa Cloud (agregación global y distribución descendente).}
    El servidor \texttt{SwellServer} ejecuta el proceso federado global con una estrategia personalizada \texttt{MQTTFedAvgSwell} (extensión de \texttt{FedAvg}) que agrega resultados (\texttt{aggregate\_fit}) y publica el \textbf{modelo global} en \texttt{fl/global\_model} para su redistribución hacia Fog y Edge. Esta publicación materializa el canal descendente MQTT y garantiza sincronización global entre rondas.

    \item \textbf{Datos y modelo (SWELL + SwellMLP).}
    El modelo \texttt{SwellMLP} (PyTorch) queda integrado de forma homogénea en Edge/Fog/Cloud, mientras que el soporte de datos se articula mediante \texttt{SwellFederatedDataset} para planificar/materializar particiones y \texttt{SwellDataset} para utilidades de carga y particionado por sujeto. Esto asegura consistencia entre preparación de datos, entrenamiento local y evaluación.

    \item \textbf{Observabilidad (cross-cutting).}
    La telemetría se integra como preocupación transversal mediante el módulo \texttt{Telemetry}, consumido por los componentes críticos (\texttt{SwellFLClientMQTT}, \texttt{FogBroker}, \texttt{FogClientSwell}, \texttt{MQTTFedAvgSwell}). De este modo, el sistema no solo “funciona”, sino que deja \textbf{evidencia trazable} del flujo de mensajes y de la ejecución por rondas.
\end{itemize}

\begin{landscape}
\begin{figure}[H]
\centering
\includegraphics[width=\linewidth,height=0.92\textheight,keepaspectratio]{Plantilla-Latex-TFG/imagenes/diagramaUmlSP10.png}
\caption{\textbf{Diagrama UML final del sistema federado SWELL.}}
\label{fig:diagramumlfinal}
\end{figure}
\end{landscape}

\vspace{0.15em}
\paragraph{HU-062 Validación experimental con arquitecturas múltiples.}
\textbf{Estado final: validado mediante ejecuciones end-to-end.} Se verificó el sistema con configuraciones reproducibles que variaron número de nodos Fog, clientes por región y el umbral de agregación $K$, confirmando:
\begin{itemize}\setlength\itemsep{0.15em}
    \item Consistencia funcional del aprendizaje federado y de los intercambios MQTT.
    \item Correcta trazabilidad de rondas y visibilidad del comportamiento por región.
    \item Robustez del despliegue automatizado en escenarios simple y avanzado.
\end{itemize}
La Figura~\ref{fig:arch_swell_10runs} aporta la evidencia visual de la arquitectura avanzada ejecutada. Y en la siguiente \ref{fig:logs_swell_10runs} se presenta, a nivel de logs, el resultado obtenido en ella, que, como podemos ver, corresponde de igual manera a lo visualizable en Grafana.
\begin{figure}[H]
    \centering
    \includegraphics[width=0.76\linewidth]{image.png}
    \caption{Evidencia por logs de la ejecución de la arquitectura avanzada (HU-062): Accuracy sobre server final.}
    \label{fig:logs_swell_10runs}
\end{figure}

\vspace{0.15em}
\paragraph{HU-068: Modelo neuronal para estrés fisiológico.}
\textbf{Estado final: implementado e integrado.} Se incorporó el modelo \texttt{SwellMLP} para clasificación binaria de estrés, quedando disponible como arquitectura de entrenamiento federado/centralizado. La Figura~\ref{fig:swellmlp_architecture} documenta la estructura final del modelo. La cuál permitió obtener resultados entre 85.2\% y 92.3\% según semillas y condiciones de carrera.

\vspace{0.15em}
\paragraph{HU-070: Automatización del despliegue de arquitecturas federadas.}
\textbf{Estado final: operativo y reutilizable.} Se dejó disponible un launcher basado en configuración capaz de:
\begin{itemize}\setlength\itemsep{0.15em}
    \item Interpretar ficheros YAML/JSON y manifiestos experimentales.
    \item Coordinar el arranque ordenado de servidor, Fog y Edge.
    \item Preparar particiones/datasets cuando procede y validar prerequisitos de ejecución.
    \item Distribuir configuraciones y parámetros mediante MQTT para ejecutar escenarios sin cambios de código.
\end{itemize}
La Figura~\ref{fig:diagramaflujo} queda como evidencia del flujo completo de orquestación experimental.

\vspace{0.15em}
\paragraph{HU-067: Testing del sistema final.}
La historia quedó \textbf{completada}. Se consolidó una suite de testing orientada a validar tanto el \textbf{correcto funcionamiento interno de los componentes críticos} como el \textbf{comportamiento global end-to-end} del sistema federado. En esta versión se ejecutan \textbf{160 tests} (con \textbf{156 pasados} y \textbf{4 omitidos} al depender de datasets reales), y la cobertura global del núcleo del sistema alcanza el \textbf{74\%}. Este valor es representativo porque se calcula sobre los módulos esenciales del framework (excluyendo scripts y utilidades auxiliares). La única excepción destacable es \textbf{telemetría}, que queda por detrás debido a la complejidad de su mockeo y la dependencia de servicios externos; aun así, la \textbf{media de cobertura en los módulos principales es alta} y refleja una validación sólida del sistema. El alcance completo de la cobertura funcional se detalla en la Tabla~\ref{tab:test_coverage_full}, y el resumen de cobertura por módulo principal se presenta en la Tabla~\ref{tab:coverage_summary}.

En términos de validación, la suite verifica de manera sistemática:

\begin{itemize}
    \item \textbf{Consistencia de formatos y contratos de mensajería}: se comprueba que los mensajes intercambiados entre cliente, broker Fog y servidor mantienen estructura, campos obligatorios y compatibilidad de serialización/deserialización, evitando fallos por incompatibilidad entre componentes.
    \item \textbf{Preservación de dimensiones y estabilidad de parámetros}: se valida que los tensores/arrays conservan dimensionalidad, orden y tipado a lo largo del ciclo federado (extracción, empaquetado, transporte, reconstrucción y actualización), garantizando que las actualizaciones son aplicables sin corrupción de estado.
    \item \textbf{Correctitud matemática de la agregación (FedAvg)}: se ejecutan pruebas con pesos y escenarios controlados para asegurar que la agregación corresponde a la combinación ponderada esperada, validando la implementación de la estrategia.
    \item \textbf{Ejecución en múltiples rondas y manejo de errores}: se comprueba el flujo completo en rondas sucesivas e incluye escenarios de error (mensajes malformados, timeouts y respuestas vacías), verificando que el sistema mantiene coherencia y no rompe el protocolo.
    \item \textbf{Comparación con el baseline centralizado}: se contrasta entrenamiento federado frente a un entrenamiento centralizado en condiciones comparables, validando que el pipeline federado produce resultados coherentes y reproducibles.
\end{itemize}

\begin{table}[H]
\centering
\footnotesize
\caption{Cobertura funcional de tests del sistema federado (tabla completa)}
\label{tab:test_coverage_full}
\renewcommand{\arraystretch}{1.2}
\begin{tabular}{llcl}
\hline
\textbf{Módulo de test} & \textbf{Tipo} & \textbf{Tests} & \textbf{Componentes validados} \\
\hline
\texttt{scripts/test\_sweet\_federated.py} & Unitario & 3 & Pipeline SWEET (tests básicos) \\
\texttt{test\_baseclient.py} & Unitario & 1 & Cliente base MQTT \\
\texttt{test\_baseline\_comparison.py} & Unitario/Integración & 17 & Comparativa centralizado vs federado \\
\texttt{test\_datasets.py} & Unitario & 14 & Carga y particionado de datasets \\
\texttt{test\_entrypoints\_modular.py} & Unitario & 1 & Puntos de entrada del sistema \\
\texttt{test\_federated\_architecture.py} & Unitario & 2 & Configuraciones YAML/JSON \\
\texttt{test\_federated\_architecture\_swell.py} & Unitario & 4 & Arquitectura SWELL y plan de ejecución \\
\texttt{test\_fog\_bridge\_swell.py} & Unitario & 4 & Bridge MQTT--Flower SWELL \\
\texttt{test\_fog\_flower\_client.py} & Unitario & 7 & Bridge genérico MQTT--Flower \\
\texttt{test\_integration.py} & Integración & 10 & Flujo end-to-end, rondas federadas \\
\texttt{test\_logging\_and\_telemetry.py} & Unitario & 10 & Logging y telemetría (OTEL) \\
\texttt{test\_model.py} & Unitario & 12 & ECGModel, SwellMLP, parámetros \\
\texttt{test\_mqtt\_components.py} & Integración & 12 & Broker Fog, agregación regional \\
\texttt{test\_prometheus\_metrics.py} & Unitario & 4 & Métricas Prometheus \\
\texttt{test\_real\_datasets.py} & Integración & 3 & SWELL con datos reales (omitidos) \\
\texttt{test\_run\_architecture\_from\_config.py} & Unitario & 2 & Lanzador por configuración \\
\texttt{test\_server.py} & Unitario/Integración & 10 & Servidor central, MQTTFedAvg \\
\texttt{test\_swell\_client.py} & Unitario/Integración & 11 & Cliente SWELL, sincronización global \\
\texttt{test\_swell\_dataset\_helpers.py} & Unitario & 4 & Normalización y carga SWELL \\
\texttt{test\_swell\_federated.py} & Unitario & 7 & Particionado federado SWELL \\
\texttt{test\_swell\_server.py} & Unitario & 6 & Servidor SWELL y evaluación final \\
\texttt{test\_utils.py} & Unitario & 16 & Utilidades, serialización, métricas \\
\hline
\textbf{Total} & & \textbf{160} & \textbf{Sistema federado completo} \\
\hline
\end{tabular}
\end{table}

\begin{table}[H]
\centering
\footnotesize
\caption{Resumen de cobertura por módulo principal}
\label{tab:coverage_summary}
\renewcommand{\arraystretch}{1.2}
\begin{tabular}{lcc}
\hline
\textbf{Componente} & \textbf{Archivo} & \textbf{Cobertura} \\
\hline
Cliente SWELL & \texttt{clients/swell.py} & 75\% \\
Servidor SWELL & \texttt{servers/swell.py} & 74\% \\
Fog bridge genérico & \texttt{fog\_flower\_client.py} & 83\% \\
Telemetría & \texttt{telemetry.py} & 59\% \\
Datasets SWELL & \texttt{datasets/swell.py} & 72\% \\
Particionado SWELL & \texttt{datasets/swell\_federated.py} & 72\% \\
Arquitectura federada & \texttt{federated\_architecture.py} & 81\% \\
Logging & \texttt{logging\_utils.py} & 100\% \\
\hline
\textbf{Cobertura global (core)} & & \textbf{74\%} \\
\hline
\end{tabular}
\end{table}

\vspace{0.5em}


\subsection{Sprint Review (Sprint 10)}
El Sprint 10 cerró el ciclo de desarrollo funcional entregando un sistema federado SWELL \textbf{operativo, observable, reproducible y validado por testing}. La decisión clave fue priorizar un \textbf{cierre verificable} (end-to-end + métricas + tests) frente a extensiones transversales de alto riesgo (seguridad avanzada, SWEET a gran escala, tolerancia a fallos), garantizando un resultado final defendible y medible.

\vspace{0.15em}
\paragraph{Qué se logró.}
Se consolidó una \textbf{plataforma completa de experimentación} sobre SWELL con ejecución end-to-end Edge--Fog--Server, incorporando:
(i) un stack de \textbf{observabilidad real} (Prometheus/Grafana + Jaeger) integrado en el flujo federado;
(ii) \textbf{automatización config-driven} para desplegar arquitecturas reproducibles sin modificar código;
(iii) \textbf{validación experimental} ejecutando configuraciones con múltiples topologías y umbral de agregación $K$ parametrizable; y
(iv) una \textbf{validación por testing} extensiva (160 tests; 156 pasados; 4 omitidos por datasets reales) con \textbf{74\% de cobertura en el core}, verificando contratos MQTT, serialización, preservación dimensional y correctitud matemática de FedAvg.

\vspace{0.15em}
\paragraph{Qué fue bien.}
El sprint ejecutó una convergencia muy sólida de líneas críticas:
\begin{itemize}\setlength\itemsep{0.25em}
  \item \textbf{Observabilidad extremo-a-extremo}: las rondas dejan evidencia en métricas y trazas, permitiendo diagnóstico y verificación de causalidad/latencias entre componentes.
  \item \textbf{Reproducibilidad operativa}: el despliegue basado en configuración reduce fricción, facilita repetir experimentos y comparar topologías de forma sistemática.
  \item \textbf{Validación multi-arquitectura}: ejecuciones end-to-end con distintos despliegues confirmaron estabilidad del protocolo, del flujo de agregación y de la instrumentación.
  \item \textbf{Calidad y robustez}: la suite de tests refuerza confianza en el núcleo del sistema, cubriendo tanto unitarios como integración en la ruta crítica.
\end{itemize}

\vspace{0.15em}
\paragraph{Qué se detectó / qué fue mal.}
Al durar $\sim 1.5$ sprints, quedó patente que la \textbf{integración transversal} (telemetría + despliegue + tests) tiene un coste no lineal y tiende a absorber capacidad adicional por interdependencias. Además, la \textbf{telemetría} es el área con mayor dificultad de \textbf{test/mocking} (dependencia de servicios externos y propagación de contexto), lo que limita cobertura y aumenta el esfuerzo de mantenimiento.
Finalmente, quedó confirmado que \textbf{seguridad/privacidad} y \textbf{resiliencia} requieren decisiones arquitectónicas globales y una integración de la cual no se dispone de tiempo si se quiere cumplir con la planificación de entregables \cite{sec:planificacion} .

\vspace{0.15em}
\paragraph{Acciones y siguientes pasos.}
Las mejoras y líneas de evolución posteriores al prototipo (endurecimiento de seguridad, privacidad, tolerancia a fallos y escalado de datasets como SWEET mediante estrategias alternativas) se recogen en la sección de \textit{Trabajos futuros} (Sección~\ref{sec:trabajo_futuro}).




\chapter{Estimación de Costes}
\label{chap:costes}

Para dimensionar la dedicación y los costes del proyecto, se ha tomado como referencia el estándar académico, según el cual cada crédito ECTS equivale aproximadamente a entre 25 y 30 horas de trabajo \cite{UGR_ECTS}. Por consiguiente, se estima que la realización del TFM requerirá una dedicación total de entre 600 y 700 horas, lo que permite planificar adecuadamente la carga de trabajo y los recursos necesarios para alcanzar los objetivos propuestos.
   


\subsection{Equipamiento}

El desarrollo del sistema propuesto requiere una infraestructura mínima que permita reproducir de forma realista los tres niveles de la arquitectura Edge--Fog--Cloud.  
El objetivo no es disponer de infraestructura dedicada exclusivamente al proyecto, sino emplear equipamiento representativo y suficiente para el prototipado, las pruebas funcionales y la evaluación experimental.

Inicialmente se contempló el uso explícito de dispositivos móviles y nodos Fog basados en Raspberry Pi. Sin embargo, a medida que el alcance técnico del proyecto quedó definido, dichos dispositivos no fueron necesarios en la fase experimental final.  
Por este motivo, \textbf{no se imputan económicamente al coste final del proyecto}.

A continuación se detalla el equipamiento efectivamente utilizado y su justificación.

\paragraph{Ordenador portátil / PC (nivel Fog--Cloud local).}
El equipo principal actúa como entorno de desarrollo, nodo Fog de pruebas y servidor local para simulaciones del nivel Cloud.  
Debe disponer de capacidad de virtualización, memoria suficiente y almacenamiento adecuado para la ejecución de contenedores, entrenamiento de modelos y análisis de resultados.

\paragraph{Pantalla secundaria.}
La pantalla adicional no participa directamente en el procesamiento, pero mejora significativamente la productividad durante el desarrollo, permitiendo la visualización simultánea de código, logs, métricas y paneles de observabilidad.

\paragraph{Servidor Cloud on--premise (cedido por MYDASS).}
El nivel Cloud del sistema federado se ejecuta sobre un servidor físico \textbf{on--premise} cedido por el grupo de investigación MYDASS.  
Este servidor soporta la agregación global de modelos, la coordinación de rondas federadas y la ejecución de experimentos distribuidos.

El servidor dispone de las siguientes características principales:
\begin{itemize}
    \item CPU: Intel Xeon Silver 4210 (10 núcleos, 20 hilos)
    \item Memoria RAM: 16 GB
    \item Sistema operativo: Ubuntu 20.04 LTS
\end{itemize}

Para estimar su valor económico, se ha tomado como referencia el precio de mercado de servidores con especificaciones equivalentes.  
Por ejemplo, un servidor Dell PowerEdge R540 con procesador Intel Xeon Silver 4210 y 16\,GB de RAM presenta un precio aproximado de \textbf{2.300 €} en el mercado actual\footnote{\href{https://www.pccomponentes.com/dell-poweredge-r540-xeon-silver-4210-16gb}{PCComponentes – Dell PowerEdge R540 Xeon Silver 4210}}.

Dado que el servidor no se adquiere específicamente para este TFM y forma parte de la infraestructura del grupo de investigación, su coste se prorratea en función del periodo efectivo de uso.

\paragraph{Prorrateo del coste del equipamiento.}
El equipamiento utilizado no se adquiere exclusivamente para el TFM, por lo que su coste se prorratea considerando:
\begin{itemize}
    \item una vida útil estimada de 4–5 años para el servidor y el portátil,
    \item un periodo efectivo de uso en el TFM de aproximadamente 6 meses,
    \item el valor actual de mercado.
\end{itemize}

Aplicando un modelo de amortización lineal, se imputa aproximadamente un \textbf{10–15\%} del valor del equipamiento al proyecto.



\begin{table}[H]
\centering
\footnotesize
\caption{Equipamiento utilizado en el proyecto y coste prorrateado.}
\label{tab:equipamiento_prorrateado}

\renewcommand{\arraystretch}{1.15}
\setlength{\tabcolsep}{4pt}

\begin{tabular}{|p{3.2cm}|p{4.8cm}|p{2.4cm}|p{2.6cm}|}
\hline
\textbf{Componente} &
\textbf{Modelo / referencia} &
\textbf{Coste total} &
\textbf{Coste imputado} \\
\hline

Ordenador portátil / PC &
\href{https://www.pccomponentes.com/lenovo-v15-g2-itl-intel-core-i3-1115g4-8gb-256gb-ssd-156}{Lenovo V15 G2} &
255\,€ &
31{,}9\,€ \\
\hline

Dispositivos móviles (x3) &
\href{https://eurofertas.eu/products/smartphone-pocophone-xiaomi-poco-m4-pro-6-6-octa-core-4-gb-ram-64-gb-azul}{Xiaomi Poco M4 Pro} &
617{,}43\,€ &
0\,€ \\
\hline

Raspberry Pi (x2) &
\href{https://store.rokland.com/products/raspberry-pi-4-model-b-2gb-no-quantity-restrictions}{Raspberry Pi 4 Model B (2GB)} &
90\,€ &
0\,€ \\
\hline

Pantalla secundaria &
\href{https://www.monitors.com/products/hp-z24u-1c4z6aa}{HP Z24u G3 24''} &
280\,€ &
20\,€ \\
\hline

Servidor Cloud on--premise &
\href{https://www.pccomponentes.com/dell-poweredge-r540-xeon-silver-4210-16gb}{Dell PowerEdge R540 (equivalente)} &
2.300\,€ &
300\,€ \\
\hline

\multicolumn{3}{|r|}{\textbf{Coste total imputado al proyecto}} &
\textbf{351{,}9\,€} \\
\hline

\end{tabular}
\end{table}


Dado que la dedicación estimada para el TFM se encuentra entre 600 y 700 horas \cite{UGR_ECTS} debido a los 24 créditos de este TFM, y considerando el coste horario calculado anteriormente, el coste total aproximado del personal asociado al proyecto se presenta en el Cuadro \ref{tab:coste-proyecto}.  

\begin{table}[H]
\centering
\begin{tabular}{|p{4cm}|p{3cm}|p{3cm}|}
\hline
\textbf{Concepto} & \textbf{Mínimo} & \textbf{Máximo} \\
\hline
Horas estimadas de dedicación & 600 & 700 \\
\hline
Coste por hora (empresa) & 16,6 €/h & 20,2 €/h \\
\hline
Coste total aproximado del personal para el proyecto & 9.960 € & 14.140 € \\
\hline
\end{tabular}
\caption{Estimación del coste total de personal asociado al proyecto considerando la dedicación estimada.}
\label{tab:coste-proyecto}
\end{table}

\subsection{Software}
Todo el software empleado en el desarrollo del proyecto es de carácter libre o gratuito, con la única excepción de la licencia de \textit{Windows}, la cual se encuentra ya incluida en el precio del equipo presupuestado en el apartado anterior. 
Por tanto, el coste adicional asociado al uso de software o a la adquisición de datos para este proyecto es de \textbf{0 euros}.

\subsection{Coste de Electricidad }

Con objeto de reflejar un escenario realista de consumo energético que incluya climatización (aire acondicionado y calefacción) y otros usos domésticos, se adopta como referencia el gasto medio mensual de electricidad de un hogar en España. 
De acuerdo con estimaciones de la OCU para el PVPC, el gasto medio en \textbf{agosto de 2025} fue de \textbf{67,74\,€/mes} \cite{OCU_Precio_Luz_2025}. 
Adicionalmente, el precio medio mensual del PVPC en \textbf{agosto de 2025} se situó en \textbf{0,1331\,€/kWh} \cite{CosteEnergia_PVPC_2025}, consistente con un rango de facturación mensual típico en 2025.

Dado que la planificación del proyecto abarca aproximadamente \textbf{6 meses} (julio–diciembre de 2025), y con el fin de mantener una hipótesis conservadora y homogénea, se toma el valor de agosto como \emph{proxy} representativo mensual. 
Así, el coste eléctrico total imputable al periodo del proyecto se estima como:
\[
\text{Coste total electricidad (hogar)} \approx 6 \times 67{,}74\,\text{€} \;=\; \textbf{406{,}44\,€}.
\]
\subsection{Presupuesto Total Estimado}

Una vez considerados los distintos apartados de gastos; equipamiento, personal, software y electricidad. Es posible calcular una estimación global del presupuesto necesario para la realización del proyecto. 
En la Tabla \ref{tab:coste-total} se detallan los importes aproximados por cada concepto y el sumatorio final.

\begin{table}[H]
\centering
\begin{tabular}{|p{5cm}|p{3cm}|p{3cm}|}
\hline
\textbf{Concepto} & \textbf{Coste mínimo} & \textbf{Coste máximo} \\
\hline
Equipamiento & 351,9 € & 351,9 € \\
\hline
Coste de personal & 9.960 € & 14.140 € \\
\hline
Software & 0 € & 0 € \\
\hline
Electricidad & 406,44 € & 406,44 € \\
\hline
\textbf{Total estimado} & \textbf{10.718,34 €} & \textbf{14.898,34 €} \\
\hline
\end{tabular}
\caption{Resumen del presupuesto total estimado del proyecto.}
\label{tab:coste-total}
\end{table}


En conclusión, el presupuesto global del proyecto se sitúa en un rango aproximado de entre \textbf{10.700 € y 14.900 €}, siendo el coste de personal el factor dominante, mientras que el equipamiento y la electricidad representan una fracción reducida del total.

%Respecto a las conclusiones, deben tener tres partes: Frase o párrafo que resuma el trabajo realizado. Revisión de todos los objetivos específicos. Se pone uno por uno y se va diciendo si se ha cumplido totalmente o no, se resume cómo y se indica el lugar de la memoria donde se aborda ese objetivo.Valoración personal del trabajo realizado y lo aprendido, y lo que has utilizado de lo visto en el máster.En cuanto a los trabajos futuros, se enumeran tareas que no se hayan podido hacer o que te gustaría hacer más adelante. Para cada una se puede explicar muy brévemente cómo se abordaría en el futuro.

\chapter{Conclusiones}
\label{chap:conclusiones}

Este TFM ha materializado un \textbf{framework experimental de Aprendizaje Federado} alineado con el paradigma del \textit{Computing Continuum}, incorporando una capa \textit{Fog} como nivel de agregación intermedia para superar el patrón cliente--servidor clásico. El resultado es una plataforma modular capaz de ejecutar experimentos reproducibles en topologías \textit{Edge--Fog--Cloud}, con \textbf{despliegue por configuración} (YAML/JSON), \textbf{observabilidad extremo-a-extremo} (métricas y trazas) y \textbf{validación por testing}, aplicada a un caso biomédico realista de detección de estrés con el dataset SWELL.

Desde una perspectiva incremental, los sprints previos (\textit{Sprint 1~\ref{sect:Sp1}--Sprint 9~\ref{sect:Sp9}}) fueron consolidando los pilares de arquitectura, comunicación (MQTT), modularidad y reproducibilidad, mientras que el \textbf{Sprint 10 (sección ~\ref{sect:Sp10})} cerró el ciclo de desarrollo funcional integrando la observabilidad, la automatización de ejecución y la validación integral del sistema (véanse Fig.~\ref{fig:monitoring_stack}, Fig.~\ref{fig:grafana_dashboard_sprint10} y Fig.~\ref{fig:jaeger_traces_sprint10}).

\vspace{0.35em}
\section{Revisión de objetivos}
\label{sec:revision_objetivos}
\vspace{0.25em}
\subsection*{Objetivo general}
\textbf{Cumplimiento: completo.} Se diseñó, desarrolló e implementó un sistema federado \textit{Edge--Fog--Cloud} con capa \textit{Fog} intermedia, utilizándose como plataforma experimental para ejecutar escenarios reproducibles, instrumentados y trazables.  
\textit{Evidencia:} Capítulos de Arquitectura e Implementación; cierre de desarrollo en el Sprint 10 (HU-066, HU-070, HU-006/HU-069/HU-071).

\vspace{0.35em}
\subsection*{Objetivos específicos}

\begin{enumerate}\setlength\itemsep{0.40em}

\item \textbf{Revisar y analizar el estado del arte.}  
\textbf{Cumplimiento: completo.} Se revisaron los fundamentos de FL, el \textit{computing continuum} y los patrones \textit{Edge/Fog/Cloud}, identificando limitaciones del enfoque cliente--servidor y justificando la agregación intermedia en Fog como mecanismo para mejorar escalabilidad y gestión regional.  
\textit{Evidencia:} Capítulo de Estado del Arte (~\ref{chap:estado-del-arte}); Síntesis de motivación y decisiones tecnológicas (\textit{Sprint 1~\ref{sect:Sp1}--Sprint 4~\ref{sect:Sp4}}).

\item \textbf{Profundizar en el manejo de tecnologías involucradas (Flower/TFF, Docker/Compose, MQTT).}  
\textbf{Cumplimiento: completo.} Se seleccionó \textit{Flower} como base del flujo federado y se integró con una capa de mensajería MQTT para coordinar la ejecución distribuida. Además, se containerizó parte del sistema y se dejó un stack reproducible para el despliegue y la ejecución.  
\textit{Evidencia:}Automatización y orquestación por configuración en Sprint 10 (HU-070) y stack containerizado de observabilidad (HU-071).

\item \textbf{Diseñar e implementar una arquitectura modular jerárquica con agregación intermedia en Fog.}  
\textbf{Cumplimiento: completo.} Se implementó una capa Fog con agregación parcial parametrizable ($K$) y un \textit{bridge} Fog--Server, habilitando topologías jerárquicas reproducibles y separando responsabilidades (Edge entrena, Fog agrega regionalmente, Server agrega globalmente).  
\textit{Evidencia:} Sprint 7 \ref{sect:Sp7}y consolidación en Sprint 10 (HU-066) y validación de topología avanzada (Fig.~\ref{fig:arch_swell_10runs}).

\item \textbf{Desarrollar un prototipo funcional soportando ML clásico y redes neuronales, con experimentos reproducibles bajo distintos esquemas de partición.}  
\textbf{Cumplimiento: parcial--alto.} Se consolidó un prototipo federado reproducible y funcional, con pipeline configurable y soporte de modelos, incluyendo un modelo neuronal específico para SWELL (\texttt{SwellMLP}). Queda como extensión natural ampliar el catálogo de modelos (p.\,ej., familias CNN/Transformers, \textit{transfer learning}) y soportar selección 100\% declarativa por configuración.  
\textit{Evidencia:} Sprint 9 y 10, Secciones de datasets/particionado y pipeline experimental; Sprint 10 (HU-068 para modelo; HU-070 para launcher config-driven).

\item \textbf{Integrar mecanismos de registro y análisis de métricas experimentales.}  
\textbf{Cumplimiento: completo.} Se integraron métricas federadas exportadas a Prometheus y visualizadas en Grafana, junto con trazas distribuidas en Jaeger para validar el flujo Edge-Fog-Server por ronda y operación. Esto aporta trazabilidad causal (\textit{spans}) y visibilidad operativa (series temporales) del comportamiento del sistema.  
\textit{Evidencia:} Sprint 10, apartado de Observabilidad (HU-006, HU-069, HU-071, HU-003, HU-004); Fig.~\ref{fig:monitoring_stack}, Fig.~\ref{fig:grafana_dashboard_sprint10} y Fig.~\ref{fig:jaeger_traces_sprint10}.  
% (OJO: aquí estaba el fallo de HU: antes ponía HU-068/HU-070 y no correspondía.)

\item \textbf{Evaluar la escalabilidad, eficiencia y privacidad del sistema con pruebas controladas.}  
\textbf{Cumplimiento: parcial.} Se realizaron ejecuciones end-to-end variando topologías, número de clientes y umbral $K$, verificando estabilidad y consistencia operativa. Sin embargo, queda pendiente una evaluación cuantitativa exhaustiva (coste de comunicación, latencias agregadas, eficiencia por ronda, análisis estadístico) y la incorporación de mecanismos avanzados de privacidad (p.\,ej., \textit{secure aggregation} o privacidad diferencial).  
\textit{Evidencia:} Validación experimental en Sprint 10 (HU-062) soportada por instrumentación de métricas/trazas (HU-006, HU-069, HU-071) y despliegue reproducible (HU-070).

\item \textbf{Evaluar el impacto de nodos Fog (latencia, escalabilidad, convergencia, robustez y complejidad).}  
\textbf{Cumplimiento: parcial.} El sistema quedó preparado e instrumentado para medir el impacto Fog (latencias por tramo, tiempos de ronda, actividad regional y trazas causales), y se validó su funcionamiento con distintas topologías. La cuantificación completa (comparativas sistemáticas \emph{con Fog vs.\ sin Fog}, y análisis de convergencia/robustez) se plantea como continuación directa.  
\textit{Evidencia:} Sprint 10 (HU-062 para topologías; HU-006/HU-069/HU-071 para medición). Evidencias visuales en dashboards y trazas (Fig.~\ref{fig:grafana_dashboard_sprint10}, Fig.~\ref{fig:jaeger_traces_sprint10}).

\item \textbf{Aplicar y evaluar el sistema en un caso de uso representativo (detección de estrés).}  
\textbf{Cumplimiento: completo.} Se aplicó el sistema al caso SWELL, integrando dataset y pipeline de entrenamiento federado/centralizado, y añadiendo el modelo \texttt{SwellMLP} para clasificación binaria de estrés. Se ejecutaron experimentos federados reproducibles con trazabilidad completa, reforzando la idoneidad del enfoque para datos sensibles.  
\textit{Evidencia:} Secciones de Caso de uso/Dataset SWELL (\textit{Sección~\ref{sec:swell}}); Sprint 10 (HU-066, HU-068, HU-062) y figuras asociadas (Fig.~\ref{fig:swellmlp_architecture}, Fig.~\ref{fig:arch_swell_10runs}).

\end{enumerate}

\subsection{Revisión de requisitos no funcionales}
\label{sec:revision_rnf}

Por último, respecto a los requisitos no funcionales (Sección~\ref{RequisitosNofuncionales}), el grado de cumplimiento es el siguiente, atendiendo a sus criterios principales:

\begin{itemize}\setlength\itemsep{0.35em}

  \item \textbf{[RNF-001] Seguridad de comunicaciones y secretos.} \textbf{No cumplido.}  

  \item \textbf{[RNF-002] Privacidad y minimización de datos (GDPR/LOPD).} \textbf{Cumplido.}  
  La privacidad se mantiene por diseño: (i) el entrenamiento se realiza en Edge y no se exportan datos crudos, (ii) el dataset SWELL ya está anonimizado, y (iii) el sistema referencia sujetos mediante \textbf{IDs pseudonimizados} (clientes enumerados/asignados por configuración), evitando atributos identificables en manifiestos y resultados. Este enfoque es consistente con los principios de minimización y protección de datos exigidos por GDPR/LOPD \cite{gdpr2016,lopdgdd2018}.

  \item \textbf{[RNF-003] Calidad del modelo (rendimiento y no-regresión).} \textbf{Cumplido.}  
  El modelo específico (\texttt{SwellMLP}) obtiene un rendimiento alto y estable en los experimentos federados (véase HU-068 en Sprint 10~\ref{sect:Sp10}), cumpliendo los criterios de: (i) degradación acotada frente al mejor local y (ii) variación limitada entre ejecuciones bajo la misma configuración/\textit{seed}. La validez se apoya en la ejecución reproducible y la instrumentación asociada a los experimentos (HU-070 y HU-006/HU-069/HU-071).

  \item \textbf{[RNF-004] Reproducibilidad experimental (configuración declarativa).} \textbf{Cumplido (a nivel estructural).}  
  El sistema permite reconstruir determinísticamente la \textbf{arquitectura} desde configuración y manifiesto (HU-070, Sprint 10), incluyendo topología, roles, asignaciones y parámetros efectivos.  
  \textit{Matiz:} el orden temporal exacto de eventos puede variar por condiciones de carrera asociadas a $K$ (disparo de agregación), por lo que la reproducibilidad garantizada es la \textbf{estructural} (arquitectura/plan), no el \textit{timing} distribuido bit-a-bit.

  \item \textbf{[RNF-005] Observabilidad nativa (métricas, logs y trazas).} \textbf{Cumplido.}  
  Se dispone de telemetría extremo-a-extremo: (i) métricas exportadas y visualizadas (Prometheus/Grafana), (ii) trazas distribuidas/correlación por ronda (Jaeger) y (iii) un esquema de instrumentación consistente basado en OpenTelemetry. La evidencia visual se apoya en el stack y dashboards del Sprint~\ref{sect:Sp10} (Fig.~\ref{fig:monitoring_stack}, Fig.~\ref{fig:grafana_dashboard_sprint10} y Fig.~\ref{fig:jaeger_traces_sprint10}).

  \item \textbf{[RNF-006] Fiabilidad y tolerancia a desconexiones.} \textbf{Parcial.}  
  El sistema es operativo en escenarios normales, pero ante pérdidas de clientes/nodos puede quedar bloqueado esperando alcanzar $K$ y acabar dependiendo de \textit{timeouts} del servidor, actualmente no parametrizados como parte de la configuración experimental. Se cumple la estabilidad básica, pero quedan por cerrar: (i) detección/gestión configurable de caídas, (ii) abortos controlados por umbral y (iii) recuperación más explícita de estado. La propia naturaleza asíncrona del patrón de mensajería MQTT hace especialmente relevante formalizar políticas de \textit{timeout} y \textit{retry}.

  \item \textbf{[RNF-007] Rendimiento (latencia y coste operativo).} \textbf{Cumplido en entorno de laboratorio.}  
  Las rondas completan en tiempos razonables y la instrumentación no impide la ejecución práctica; además, el envío de parámetros se mantiene dentro de tamaños manejables para los modelos ensayados. 

  \item \textbf{[RNF-008] Mantenibilidad y calidad del código.} \textbf{Cumplido.}  
  La modularidad por componentes y la separación de responsabilidades facilitan extensión; además, se ha reforzado la calidad con (i) una suite extensa de tests (HU-067), (ii) automatización en CI/CD para linting/formato y (iii) un diseño coherente con principios SOLID \cite{martin2008clean}. En particular, se preserva consistencia de estilo conforme a PEP~8 \cite{pep8}, reforzando legibilidad y mantenibilidad.

  \item \textbf{[RNF-009] Portabilidad y despliegue reproducible.} \textbf{Cumplido.}  
  El despliegue se apoya en contenedores y \textit{compose} (incluyendo el stack de observabilidad del Sprint~\ref{sect:Sp10}), y el entorno de desarrollo queda estabilizado mediante dependencias aisladas, reduciendo variabilidad entre máquinas y facilitando relanzar escenarios.

  \item \textbf{[RNF-010] Interoperabilidad (APIs y mensajería).} \textbf{No aplicable / fuera de alcance.}  
  Aunque la mensajería interna está definida (MQTT y esquemas de mensajes), no se ha priorizado una capa de APIs públicas orientadas a integración con terceros ni exportación estándar más allá de los artefactos de experimento (configuración/manifiestos/resultados).

  \item \textbf{[RNF-011] Usabilidad en CLI y operabilidad.} \textbf{Parcial--alto.}  
  La experiencia de uso mejora sustancialmente con la ejecución \textit{config-driven} (scripts/launcher y presets), pero sigue siendo un entorno orientado a perfiles técnicos. Como trabajo futuro, podría ofrecerse una capa de alto nivel (p.\,ej., ``experimentos como servicio'') para abstraer configuración y ejecución, manteniendo la trazabilidad aportada por la instrumentación \cite{opentelemetry}.

  \item \textbf{[RNF-012] Idioma y documentación.} \textbf{Cumplido.}  
  La base del repositorio, nombres de métricas/logs y la documentación técnica se mantienen en inglés, facilitando replicación y revisión externa.

\end{itemize}

\vspace{0.35em}
\section{Valoración personal y aprendizaje}
\label{sec:valoracion_personal}

A nivel personal, el proyecto ha supuesto un reto completo de ingeniería de software: diseño de una arquitectura distribuida, integración de mensajería, modularización por componentes y validación sistemática del comportamiento del sistema. En particular, diría que el salto cualitativo ha consistido en \textbf{convertir el desarrollo en una plataforma arquitectónicamente configurable}: un proyecto que sienta las bases de un \textbf{framework de experimentación} (no solo ``un prototipo que funciona''), en la que la topología \textit{Edge-Fog-Server}, el reparto regional, y parámetros críticos como el umbral de agregación $K$ se definen de forma declarativa y pueden reproducirse sin modificar el código. Desde mi punto de vista, esto sitúa el proyecto claramente por encima de una demostración básica de aprendizaje federado con un único nodo fog, al aportar un software con el que \textbf{experimentar de manera real} con distintas topologías, cargas y escenarios de despliegue.

Además, el trabajo me ha permitido aplicar y consolidar competencias del máster vinculadas a sistemas distribuidos, contenedores, automatización de despliegue, telemetría y prácticas de calidad software (testing e instrumentación), integrándolas en un caso de uso biomédico realista con restricciones explícitas de privacidad y trazabilidad experimental. En este sentido, la instrumentación del sistema y su observabilidad extremo a extremo quedan reflejadas en Fig.~\ref{fig:monitoring_stack}, Fig.~\ref{fig:grafana_dashboard_sprint10} y Fig.~\ref{fig:jaeger_traces_sprint10}, donde se evidencia tanto la monitorización de métricas como el seguimiento detallado de trazas durante la ejecución federada.



\section{Trabajos futuros}
\label{sec:trabajo_futuro}

A pesar de haber completado el ciclo funcional del prototipo a lo largo de varios sprints, el sistema queda en una posición idónea para evolucionar desde \textit{framework experimental} hacia una plataforma más robusta, segura y evaluada cuantitativamente. En particular, las siguientes líneas conectan directamente con las historias no completadas y que aún se encuentran en el \textit{backlog} (HU-002, Tabla~\ref{tab:hu002}; HU-060, Tabla~\ref{tab:hu060}; HU-011, Tabla~\ref{tab:hu011}; HU-012, Tabla~\ref{tab:hu012}; HU-005, Tabla~\ref{tab:hu005}; HU-010, Tabla~\ref{tab:hu010}; HU-013, Tabla~\ref{tab:hu013}; HU-014, Tabla~\ref{tab:hu014}; HU-064, Tabla~\ref{tab:hu063}) y con los requisitos no funcionales revisados en la Sección~\ref{sec:revision_rnf}.


\begin{enumerate}\setlength\itemsep{0.45em}

  \item \textbf{Seguridad extremo-a-extremo: incorporación segura de nodos y TLS.}  
  Cerrar el gap principal identificado en [RNF-001] mediante: (i) registro y autenticación de nodos Edge/Fog con certificados o tokens (HU-002, Tabla~\ref{tab:hu002}), y (ii) cifrado de comunicaciones con TLS (HU-060, Tabla~\ref{tab:hu060}), como base para despliegues en entornos reales y multi-institución. Esto habilita posteriormente políticas de rotación de secretos y control de acceso coherentes con TLS 1.3 \cite{rfc8446_tls13}.


  \item \textbf{Tolerancia a fallos y robustez distribuida configurable.}  
  Consolidar [RNF-006] abordando explícitamente la recuperación ante desconexiones: detección de fallos, reconexión y reequilibrado de rondas sin pérdida significativa de clientes (HU-012, Tabla~\ref{tab:hu012}). En el contexto asíncrono de MQTT , esta línea incluye parametrizar \textit{timeouts}, reintentos y políticas de degradación controlada cuando no se alcance el umbral $K$.

  \item \textbf{Selección inteligente de clientes y planificación de rondas.}  
  Mejorar eficiencia y estabilidad del entrenamiento incorporando políticas automáticas de selección de clientes por disponibilidad, reputación, ancho de banda y heterogeneidad (HU-005, Tabla~\ref{tab:hu005}). Esta evolución permitiría comparar estrategias de muestreo/selección bajo distintas topologías Fog y medir su impacto con la instrumentación ya integrada (Fig.~\ref{fig:grafana_dashboard_sprint10} y Fig.~\ref{fig:jaeger_traces_sprint10}).

  \item \textbf{Explicabilidad y auditoría del modelo federado.}  
  Incorporar un módulo de interpretabilidad con generación de reportes automáticos (Grad-CAM, SHAP, LIME) para auditoría de sesgos y calidad del modelo final (HU-010, Tabla~\ref{tab:hu010}). Esta línea encaja de forma natural con la observabilidad ya existente (OpenTelemetry/Prometheus/Grafana/Jaeger), complementando métricas operativas con evidencia explicativa del comportamiento del modelo.

  \item \textbf{Medición de contribución y mecanismos de incentivos.}  
  Completar el \textit{ranking} de contribución de clientes mediante valores de Shapley u otros estimadores (HU-014, Tabla~\ref{tab:hu014}), generando rankings persistentes para análisis y auditoría. Además de aportar interpretabilidad a nivel de participantes, puede integrarse con políticas de selección inteligente (HU-005) para diseñar estrategias basadas en utilidad histórica.

  \item \textbf{Entrenamiento híbrido y topologías avanzadas de aprendizaje.}  
  Extender el pipeline hacia escenarios más complejos combinando aprendizaje federado horizontal con \textit{transfer learning} distribuido (HU-013, Tabla~\ref{tab:hu013}), especialmente relevante cuando existan dominios heterogéneos o datos con distinta calidad por región/nodo.

  \item \textbf{Generalización del pipeline por configuración: modelos y arquitecturas de red declarativas.}  
  Evolucionar el enfoque \textit{config-driven} (HU-070) para que tanto el modelo como su arquitectura (p.\,ej., MLP/CNN/\textit{trees}) puedan instanciarse desde configuración mediante una clase/factoría genérica. El objetivo es eliminar acoplamientos específicos (p.\,ej., dependencia directa de \texttt{SwellMLP}, HU-068) y permitir que el sistema cargue dinámicamente el modelo definido en YAML/JSON, manteniendo trazabilidad y reproducibilidad estructural ([RNF-004]).

  \item \textbf{Retomar el enfoque multi-dataset: SWEET y evaluación de viabilidad.}  
  Ampliar el alcance experimental más allá de SWELL retomando el análisis preliminar del dataset SWEET (HU-064, Tabla~\ref{tab:hu063}), caracterizando estructura por sujetos/sesiones, requisitos computacionales y viabilidad de particionado federado. Esta línea permitiría validar la generalidad del framework y comparar comportamiento del sistema bajo datasets con distinta granularidad y distribución.

  \item \textbf{Campaña cuantitativa completa: escalabilidad, coste de comunicación y comparación con/sin Fog.}  
  Aprovechando que el sistema ya se validó end-to-end con topologías múltiples (HU-062) y dispone de telemetría completa (HU-006/HU-069/HU-071; Fig.~\ref{fig:monitoring_stack}, Fig.~\ref{fig:grafana_dashboard_sprint10}, Fig.~\ref{fig:jaeger_traces_sprint10}), se propone una evaluación sistemática: coste de comunicación, latencias por tramo Edge--Fog--Server, throughput, tiempo por ronda, y análisis estadístico comparativo \textit{Fog vs.\ sin Fog} para cuantificar beneficios y trade-offs del nivel intermedio.

\end{enumerate}




\printglossary[title=Glosario] %glosario de términos


%%%%%%%%%%%%%%%%%%%%%%%%%%%%%%%%%%%%%%%%%%%%%



% Citas y referencias utilizadas en el documento. Esto ya por sí solo añade el título "Bibliografía"
\printbibliography


\end{document}
